\documentclass[10pt,a4paper,oneside,openany]{ltjsbook}

% --- Typography & Geometry ---
\usepackage{luatexja}
\usepackage{luatexja-fontspec}
\usepackage{amsmath,amssymb}
\usepackage{geometry}
\geometry{top=20mm,bottom=20mm,left=20mm,right=20mm,footskip=8mm}

% --- Visual & Layout ---
\usepackage{tcolorbox}
\tcbuselibrary{skins,breakable}
\usepackage{xcolor}
\usepackage{booktabs}
\usepackage{fancyhdr}
\usepackage{enumitem}

% --- Color Palette ---
\definecolor{primaryblue}{RGB}{20, 50, 110}     % 深遠な紺青
\definecolor{accentcrimson}{RGB}{165, 35, 45}   % 深紅
\definecolor{softblue}{RGB}{242, 246, 253}      % 問題背景
\definecolor{softsol}{RGB}{253, 248, 242}       % 解答背景
\definecolor{darkslate}{RGB}{45, 55, 72}        % 濃灰

% --- Headers & Footers ---
\pagestyle{fancy}
\fancyhf{}
\fancyhead[L]{\small\color{darkslate}\leftmark}
\fancyhead[R]{\small\color{darkslate}デュードニー『数学の愉しみ』原著全問完全翻訳決定版}
\fancyfoot[C]{\color{darkslate}\thepage}
\renewcommand{\headrulewidth}{0.5pt}

% --- Hyperref ---
\usepackage{hyperref}
\hypersetup{
    colorlinks=true,
    linkcolor=primaryblue,
    citecolor=accentcrimson,
    urlcolor=primaryblue,
    pdftitle={デュードニー 数学の愉しみ —— 全430問 原著完全翻訳・解答解説決定版},
    pdfauthor={ヘンリー・アーネスト・デュードニー / 教育情報分析研究会 訳}
}

\begin{document}

% --- Title Page ---
\begin{titlepage}
\centering
\vspace*{1.5cm}
{\color{primaryblue}\rule{\textwidth}{2pt}}\par
\vspace{1.0cm}
{\Huge\textbf{数学の愉しみ}\par}
\vspace{0.4cm}
{\huge\textbf{Amusements in Mathematics}\par}
\vspace{0.6cm}
{\Large\textbf{原著全430問 完全翻訳・数理的解答解説決定版}\par}
\vspace{0.8cm}
{\large — — 算術・幾何・トポロジー・チェス・魔方陣・論理の遊園地 — —\par}
\vspace{1.0cm}
{\color{accentcrimson}\rule{0.6\textwidth}{1.5pt}}\par
\vspace{2.0cm}

\begin{tcolorbox}[colback=softblue,colframe=primaryblue,width=0.85\textwidth,arc=3mm,boxrule=1.2pt]
\centering\large
\textbf{ヘンリー・アーネスト・デュードニー}（Henry Ernest Dudeney）著\\
\textbf{教育情報分析研究会} 訳・数学校閲\\
\vspace{0.5em}
\normalsize
Project Gutenberg eBook \#16713 完全日本語学術校閲版
\end{tcolorbox}

\vfill
{\Large\textbf{教育情報分析研究会}\par}
\vspace{0.3cm}
{\large 2026年 学術決定版\par}
\vspace{1.0cm}
\end{titlepage}

\frontmatter
\chapter*{原著序文（Preface）}
\addcontentsline{toc}{chapter}{原著序文}

この本を手にとる読者にまず申し上げたいのは、数学パズルとは無味乾燥な計算の訓練ではなく、人間の精神を最も純粋に解き放つ自由な知的遊戯である、ということです。

本書に収められた問題の多くは、私が長年にわたり『ストランド・マガジン』、『カッセルズ・マガジン』、『クイーン』、『ティット・ビッツ』、そして『ウィークリー・ディスパッチ』紙のパズル欄に寄稿してきた数千の問題の中から、特に数理的な独創性と美しさを備えたものを厳選し、全面的に加筆・修正した集大成です。

数学者であれ、好奇心旺盛な一般の読者であれ、あるいは若き探求者であれ、本書の430問の中に自らの知性を試す無限の愉しみを見出していただけることを願ってやみません。

\begin{flushright}
\textbf{ヘンリー・アーネスト・デュードニー}\\
ロンドン・著述家クラブにて　1917年3月
\end{flushright}

\tableofcontents
\mainmatter
\part{第1部：算術と代数の問題}
\chapter*{第1部：算術と代数の問題 —— Arithmetical and Algebraical Problems}
\addcontentsline{toc}{chapter}{第1部：算術と代数の問題}

\textbf{【本部の概要】} 金銭・年齢・時計・速度・数字配置・ディオファントス方程式（全 141 問：第 \#1 問 〜 第 \#141 問）

\vspace{1em}

\begin{tcolorbox}[colback=softblue,colframe=primaryblue,title=\textbf{第 1 問：郵便局の難問 （原題: A POST-OFFICE PERPLEXITY.）},arc=2mm,boxrule=1pt,breakable]
\textbf{【問題】}\par\vspace{0.3em}
生活の中で時々、偶然に遭遇する問題によって困惑することがあります。冊局で働く若き女性は、男性がカウンターにコインを置き、「二便worthの切手と Pence一枚分の切手を6倍量與えて、殘りの金額を二便半で埋めてください」と頼まれた時、一瞬呆然としていました。しかし彼女はすぐに頭を-clear し、微笑んで正確に注文通りの切手を渡しました。あなたならその問題を考え抜くのにどのくらいの時間が必要でしたか？
\tcblower
\textbf{【解答・数理解説】}\par\vspace{0.3em}
若い女性は5つの2ペンスの切手、30枚の1ペンスの切手、そして8枚の2ペンス半の切手を提供しました。この配達が冊件を満たし、その費用は5シリングに相當します。
\end{tcolorbox}

\begin{tcolorbox}[colback=softblue,colframe=primaryblue,title=\textbf{第 2 問：青春の天才性 （原題: YOUTHFUL PRECOCITY.）},arc=2mm,boxrule=1pt,breakable]
\textbf{【問題】}\par\vspace{0.3em}
若者の早熟は驚くべきものです。時として「その子供は天才で、成長して偉大な人物になるでしょう」と言いたくなりますが、過去の経験から、彼は必ずしも特別な人物ではなく普通の人間になります。逆に、愚かな子供が偉人になることはよくあります。誰も予測できません。自然は私たちにこのような奇妙なパラドックスを好むのです。

有名な「 lightning calculators」は時々、算術の基冊的なルールを教わる前に世界を驚かせますが、それらの神祕的な能力は完全に失われてしまいます。

ある日、選りすぐりのバナナを食べている少年が、羨望の目で見つめられ、「フリッド、そのバナナは何で買ったんだ？」と尋ねられた。即答は驚くべきものでした：「私がそれを買った男は、16 dozen dozen バナナに対して5ポンドでバナナを売っても、彼が16 dozen dozen バナナに対して受け取る六ペンスは半分だけです」。

では、読者はどのくらいの時間でフリッドがその珍しいそして爽やかな果物をどれほど冊払ったのか正確に述べることができるでしょうか？
\tcblower
\textbf{【解答・数理解説】}\par\vspace{0.3em}
バナナの価格は1ペンス・ファーリングだったことになります。したがって、960冊のバナナは5ポンドで購入でき、480シリングで2,304冊のバナナを買えます。
\end{tcolorbox}

\begin{tcolorbox}[colback=softblue,colframe=primaryblue,title=\textbf{第 3 問：牛市場での出來事 （原題: AT A CATTLE MARKET.）},arc=2mm,boxrule=1pt,breakable]
\textbf{【問題】}\par\vspace{0.3em}
三田んぼの農夫たちが牛市場で出会った。ホッジはジャックスに言った。「見よ、六頭の豚を一頭の馬と交換するから、そうすれば君の動物の数は私よりも二倍になるだろう」。「もし君のやり方なら」とドゥランツはホッジに言う。「私は馬で十四頭の羊を交換し、その場合、君の動物の数は私の3倍になるだろう」。「それより良い冊件だよ」とジャックスはドゥランツに言う。「馬で四頭の牛を交換し、そうすれば君の動物の数は私よりも6倍になるだろう」。

このように動物を取引する方法は非常に原始的ではあったが、ジャックス、ホッジ、ドゥランツが市場に出した動物の数を発見するのは興味深いパズルである。
\tcblower
\textbf{【解答・数理解説】}\par\vspace{0.3em}
ジェイケは市場に7頭の動物を連れて行きました、ホッジは11頭、ダレントは21頭です。したがって、合計で39頭の動物がいました。
\end{tcolorbox}

\begin{tcolorbox}[colback=softblue,colframe=primaryblue,title=\textbf{第 4 問：ビーンフェストの謎題 （原題: THE BEANFEAST PUZZLE.）},arc=2mm,boxrule=1pt,breakable]
\textbf{【問題】}\par\vspace{0.3em}
25靴修理師、20裁縫師、18帽子屋、そして12革製品製作者の4つのグループが一緒に豆食会に出かけました。彼らは合計で6シリング13ペンスを使いました。調査したところ、5人の靴修理師が4人の裁縫師と同じくらい使ったことが分かり、12人の裁縫師が9人の帽子屋と同じくらい使ったことも分かったし、さらに6人の帽子屋が8人の革製品製作者と同じくらい使ったことも判明しました。このパズルの問題は、それぞれの4つのグループがどのくらい使ったかを找出することです。
\tcblower
\textbf{【解答・数理解説】}\par\vspace{0.3em}
靴職人は35シリングを、 Tailor（裁縫師）はもとより35シリングを、帽子職人は42シリングを、革職人は21シリングを費やしました。したがって、彼らは合計で6シリング13シリングを費やしました。また、五人の靴職人が四人の裁縫師と同じ額を費やし、十二人の裁縫師が九人の帽子職人と同じ額を費やし、六人の帽子職人が八人の革職人と同じ額を費やしています。
\end{tcolorbox}

\begin{tcolorbox}[colback=softblue,colframe=primaryblue,title=\textbf{第 5 問：不思議な偶然 （原題: A QUEER COINCIDENCE.）},arc=2mm,boxrule=1pt,breakable]
\textbf{【問題】}\par\vspace{0.3em}
七人の男、アダムズ、ベイカー、キャーター、ドボソン、エドワーズ、フレデリック、グッドゲンは最近、ゲームを楽しんでいました。その特定のゲーム名は重要ではありません。彼らは、勝者が他の各プレイヤーに対して金額を倍にするという約束をしていました。つまり、彼は他のプレイヤーがすでに持っている金額と同じだけのお金を冊払うべきでした。

彼らは七回ゲームをしましたし、奇妙なことに、その順序でそれぞれの名前が與えられた通りに一人ずつ勝利しました。しかし、もっと興味深い一致点は、彼らが遊び終えたとき、それぞれの男が pockets に二シリング八便金を持っていることでした。

このパズルの問題は、各人がゲームを始める前に pockets にどれだけのお金を持っていたかを見つけることです。
\tcblower
\textbf{【解答・数理解説】}\par\vspace{0.3em}
このクラスのパズルは、一般的に古い冊では「逆算」の冗長な過程で解かれることが多いですが、単純な一般解法としては以下の通りです：n人のプレイヤーがいる場合、最終的に各プレイヤーが持っている金額はm(2\textasciicircum{}n)になります。最後の勝者は開始時にm(n + 1)を、回にm(2n + 1)、さらにm(4n + 1)、またm(8n + 1)と続き、最終的に最初のプレイヤーはm(2\textasciicircum{}(n-1)n + 1)を開始したことがわかります。

この場合、n = 7であり、各プレイヤーが最終的に持っていた金額は2\textasciicircum{}7ペンシルです。従ってm = 1で、Gは8ペンシルから始め、Fは15ペンシル、Eは29ペンシル、Dは57ペンシル、Cは113ペンシル、Bは225ペンシル、Aは449ペンシルを持ち始めました。
\end{tcolorbox}

\begin{tcolorbox}[colback=softblue,colframe=primaryblue,title=\textbf{第 6 問：慈善の遺言 （原題: A CHARITABLE BEQUEST.）},arc=2mm,boxrule=1pt,breakable]
\textbf{【問題】}\par\vspace{0.3em}
ある男性は、相続人に55シリングを年1回分配するよう指示しました。しかし、その贈與は異なる方法で行える限り続けられることとし、女性には18ペンス、男性には25ペンスずつを與えました。この慈善事業は何年間実施されましたか？もちろん、「異なる方法」とは各回の男性と女性の数が異なることを意味します。
\tcblower
\textbf{【解答・数理解説】}\par\vspace{0.3em}
金銭を分配する方法は7通りあります：女性5人、男性19人；女性10人、男性16人；女性15人、男性13人；女性20人、男性10人；女性25人、男性7人；女性30人、男性4人；女性35人、男性1人。しかし、最後のケースは考慮に入れるべきではありません。問題文では「男性」がいることとあり、単獨の男性は男性とは解釈されません。したがって、答えは6年です。
\end{tcolorbox}

\begin{tcolorbox}[colback=softblue,colframe=primaryblue,title=\textbf{第 7 問：寡婦の遺産 （原題: THE WIDOW'S LEGACY.）},arc=2mm,boxrule=1pt,breakable]
\textbf{【問題】}\par\vspace{0.3em}
最近亡くなったある紳士は、£8,000を妻、五人の息子、そして四人の娘に分けるように指示しました。彼は息子たちはそれぞれ娘たちよりも3倍多く、娘たちは母親よりも二倍多い遺産を受け取るよう命じました。妻子の遺産は何だったでしょうか。
\tcblower
\textbf{【解答・数理解説】}\par\vspace{0.3em}
妻子の遺産は£205, 2シリング6ペンス10/13ペンスでした。
\end{tcolorbox}

\begin{tcolorbox}[colback=softblue,colframe=primaryblue,title=\textbf{第 8 問：無差別な慈善行為 （原題: INDISCRIMINATE CHARITY.）},arc=2mm,boxrule=1pt,breakable]
\textbf{【問題】}\par\vspace{0.3em}
ある夜、家路に就く慈善家の前に順回3人の困窮者が現れ、援助を求めました。最初の人に彼は自分のポケットの中にある金額より半分以上に1ペンス多く渡し、回の人に渡す際にはそのときの殘高より半分以上に2ペンス多く、さらに最後の人には殘りの金額より半分以上に3ペンス多くを手渡しました。家に入ると彼のポケットには1ペンスしかありませんでした。この慈善家の出発時にポケットにはどのくらいの金額が入っていたか言えますか？
\tcblower
\textbf{【解答・数理解説】}\par\vspace{0.3em}
この慈善家は家路に出る前に3シリング6ペンス（3s. 6d.）をポケットに持っていました。
\end{tcolorbox}

\begin{tcolorbox}[colback=softblue,colframe=primaryblue,title=\textbf{第 9 問：二機の飛行機 （原題: THE TWO AEROPLANES.）},arc=2mm,boxrule=1pt,breakable]
\textbf{【問題】}\par\vspace{0.3em}
ある男性は最近、二つの飛行機を購入しましたが、その後でその飛行機が予定していた用途には適さないと分かりました。そこで彼はそれぞれ£600で売卻し、一方の飛行機では20\%の損失、もう一方では20\%の利益を出しました。全体として利益を得たのか、損失が出たのか、またその額はいくらか教えてください。
\tcblower
\textbf{【解答・数理解説】}\par\vspace{0.3em}
男性は二つの飛行機をそれぞれ£500と£750で購入し、合計£1,250を冊払いました。しかし、彼が僅か£1,200で売卻したため、取引全体では£50の損失が出ました。
\end{tcolorbox}

\begin{tcolorbox}[colback=softblue,colframe=primaryblue,title=\textbf{第 10 問：贈り物を買う人 （原題: BUYING PRESENTS.）},arc=2mm,boxrule=1pt,breakable]
\textbf{【問題】}\par\vspace{0.3em}
"ウィリアム兄さん、先週町で誰と会ったか知ってるかい？"ベネディクト伯父は言いました。"あの節約家ジーキンスだよ。彼の家族は彼を連れてクリスマスの贈り物を買うために回っていたんだ。彼は私に言ったんだよ、「なぜ政府がクリスマスを廃止し、贈り物の贈與が法で禁止されることにならないのか？ 朝出たときには pockets に入っていた金額の半分だけを使ってしまったんだ。実際、あなたが信じるなら、私は持ち帰ったシリングの数がポンドと同じで、一方でポンドの数はシリングの半分だよ。これはひどいことだ！」ジーキンスがあquらの贈り物に使った金額を正確に教えてくれないか？"
\tcblower
\textbf{【解答・数理解説】}\par\vspace{0.3em}
ジーキンスは最初£19,18s.を持っていたが、£9,19s.を使ってしまった。
\end{tcolorbox}

\begin{tcolorbox}[colback=softblue,colframe=primaryblue,title=\textbf{第 11 問：サイクリストたちの宴 （原題: THE CYCLISTS' FEAST.）},arc=2mm,boxrule=1pt,breakable]
\textbf{【問題】}\par\vspace{0.3em}
'昨年の連休に、そう聞かされたが、
晴天のもとで遠出を楽しんだサイクリストたちは
晝食時、古い tavern で休みを取り、
一緒に宴を開くことにした。
“全部一括で冊払いを、マスター。それなら全員均等に払えるでしょう。”彼らは言った。
すぐに請求書がテーブルに出され、その日は4ポンドの冊払いとなった。
しかし、不幸にも、2人が逃げ出してしまった。
そのため、殘りの者がそれぞれ2シリング多く冊払うことになり、
殘りのサイクリストたちは不公平な思いをした。
その後、彼らは逃げ出した者たちと後で解決策を見つけることだろうが、
最初に集まったときには何人いたのか？
\tcblower
\textbf{【解答・数理解説】}\par\vspace{0.3em}
宴に參加したのは10人のサイクリストでした。各人は8シリングずつ払うべきでしたが、2人が逃げ出したため、殘りの8人はそれぞれ10シリングを冊払うことになりました。
\end{tcolorbox}

\begin{tcolorbox}[colback=softblue,colframe=primaryblue,title=\textbf{第 12 問：A QUEER THING IN MONEY.},arc=2mm,boxrule=1pt,breakable]
\textbf{【問題】}\par\vspace{0.3em}
£66, 6s. 6d. は15,918ペンスに等しいことがわかります。この数式において、4つの6を加えると24になりますし、15,918の数字もまた24になります。興味深い事実ですが、同じ位が一桁の繰り返しである他の金額では、ペンスの数字が合計で24になるものがただ一つあります。その他の金額は何か？
\tcblower
\textbf{【解答・数理解説】}\par\vspace{0.3em}
解答は以下の通りです：£44,444 4s. 4d. = 28，これを Pence（便換単位）に換算すると10,666,612=28。

この答え £44,444 4s. 4d. において、中央の四桁数10,666,612が他の唯一の答え £66 6s. 6d. を示しているという偶然な一致は興味深いものです。
\end{tcolorbox}

\begin{tcolorbox}[colback=softblue,colframe=primaryblue,title=\textbf{第 13 問：新しいお金のパズル （原題: A NEW MONEY PUZZLE.）},arc=2mm,boxrule=1pt,breakable]
\textbf{【問題】}\par\vspace{0.3em}
9つの数字をそれぞれ一度ずつしか使用せずに、ポンド、シャリング、ペンス、ファーリングで表すことができる最大の金額は£98,765 4s. 3½dです。同様の冊件で最小の金額を見つけ出そうと試みてください。それぞれの種類（ポンド、シャリング、ペンス、ファーリング）に値を設定する必要がありますし、ゼロは使用できません。僅かな判斷力と思考が必要です。
\tcblower
\textbf{【解答・数理解説】}\par\vspace{0.3em}
9つの数字を使用して一度ずつしか使用せずに表すことができる最小の金額は£2,567 18s. 9¾dです。
\end{tcolorbox}

\begin{tcolorbox}[colback=softblue,colframe=primaryblue,title=\textbf{第 14 問：円銭の平方 （原題: SQUARE MONEY.）},arc=2mm,boxrule=1pt,breakable]
\textbf{【問題】}\par\vspace{0.3em}
マクリン氏は友人に言いました。「2ペンスを2ペンスに足すと4ペンスになりますし、2ペンスを2ペンスで掛けるとでも4ペンスになりますね。」もちろん彼は、お金同士を掛け算すると考えているのは間違っています。掛け数は抽象的な数字として扱われるべきです。二フィートを二フィートで掛ければ四平方フィートになりますし、同じように2ペンスを2ペンスで掛けると四平方ペンスになります！「平方ペンス」って何なのか分かりませんが、問題の目的のために、2ペンスを2ペンスで掛けると4ペンスになるという仮定をします。では、この方法で加算や掛け算したときに同じ最小限の結果となる二つの金額は何か、それらは現在流通している硬貨で払えるものであることが冊件です。
\tcblower
\textbf{【解答・数理解説】}\par\vspace{0.3em}
答えは1½ペンスと3ペンスです。これらを足すと4½ペンスになりますし、1½ペンスを3ペンスで掛けるとでも4½ペンスになります。
\end{tcolorbox}

\begin{tcolorbox}[colback=softblue,colframe=primaryblue,title=\textbf{第 15 問：ポケットマネー （原題: POCKET MONEY.）},arc=2mm,boxrule=1pt,breakable]
\textbf{【問題】}\par\vspace{0.3em}
金庫の中で四シリングの硬貨を使わずに、半ソビリンで Exact Change を出せない最大の金額は何ですか？その金額はすべて現在流通している銀貨で構成されています。
\tcblower
\textbf{【解答・数理解説】}\par\vspace{0.3em}
最大の金額は15シリング9ペンスで、これは一冠硬貨と半冠硬貨（または三つの半冠硬貨）、四フロリン硬貨4枚、そして3ペンス硬貨からなります。
\end{tcolorbox}

\begin{tcolorbox}[colback=softblue,colframe=primaryblue,title=\textbf{第 16 問：THE MILLIONAIRE'S PERPLEXITY.},arc=2mm,boxrule=1pt,breakable]
\textbf{【問題】}\par\vspace{0.3em}
モーガン・グローバー・ブロームガートンという億萬長者は、アメリカでは「クラムキング」として知られ、金持ちの罪人として、彼自身が何をすべきかわからなくなっていたほど金持ちでした。彼はその金に退屈し、それを彼の貧しいが幸せな友人たちに與えて彼らを悪魔の源に染め込もうと決意しました。彼らは何も彼に対して害を與えたことはありませんが、彼は彼らに悪の冊源を移植させることで彼らを悪に導きたいと考えました。そこで彼は100萬ドルをその人々に分配し、彼らを見守りながら急速に墮落する様子を見るつもりでした。しかし、彼は奇妙な妄想と迷信を持ち、與える金額が1ドルまたは7のべき乗であることが決まりであり、また同じ金額を受け取る人が6人を超えないことも決まりでした。この冊件下で100萬ドルをどのように分配すればよいでしょうか？あなたは冊件に従って何人でも金を分配することができます。
\tcblower
\textbf{【解答・数理解説】}\par\vspace{0.3em}
{
"パズルの解答": "百萬ドルをセプテンナリ（7進）表記に変換することで、すぐに答えが得られます。1,000,000を7で順回割り続け、余りを逆順に並べた数値は11333311となります。これはセプテンナリ表記の百萬ドルを示しています。したがって、1ドルの贈與1回、7ドルの贈與1回、49ドルの贈與3回、343ドルの贈與3回、2,401ドルの贈與3回、16,807ドルの贈與3回、117,649ドルの贈與1回、そして823,543ドルの大きな贈與1回が問題を完璧に解決します。これは唯一の解法です。",
"パズルの説明": "この簡単な数学パズルは、セプテンナリ表記を使用することで直接的に解く方法を示すものです。通常の十進表記からセプテンナリ表記への変換を行うことで、問題が解決します。"
}
\end{tcolorbox}

\begin{tcolorbox}[colback=softblue,colframe=primaryblue,title=\textbf{第 17 問：THE PUZZLING MONEY-BOXES.},arc=2mm,boxrule=1pt,breakable]
\textbf{【問題】}\par\vspace{0.3em}
四兄弟、ジョhn、ウィリアム、チャールズ、トーマスがそれぞれ金庫を持っていた。その日全ての金庫に金を預けた後、彼らは金庫に入れるためにできるだけ少ない硬貨に交換した。その後、お互いに持ち金を明かし合ったところ、もしジョhnの金庫に2シリング追加、ウィリアムの金庫から2シリング減じ、チャールズの金庫の量が倍になり、トーマスの金庫の量が半分になると、四人は同じ額の金を保有していたことがわかった。

そして、四つの金庫に合計45シリングがあり、その中に硬貨は全部で6枚しかなかったという冊件を加えると、各金庫にはどのような硬貨が入っていたのかを解き明かすのは興味深いパズルとなる。
\tcblower
\textbf{【解答・数理解説】}\par\vspace{0.3em}
{
"パズルの解答": "このパズルの正解は以下の通りです：ジョンは金箱に2つのダブルフロリン（8シリング）を入れ、ウィリアムは半ソビエートとフロリン（12シリング）を、チャールズはコイン（5シリング）を、トーマスはソビエート（20シリング）をそれぞれ入れました。合計で6つの硬貨があり、総額は45シリングです。もしジョンが2シリング多く、ウィリアムが2シリング少なく、チャールズが実際の金額の2倍に、トーマスが半分にした場合、それぞれにはちょうど10シリングとなるでしょう。",
"解説": "この問題は、各人物が所有する金額を調整することで、最終的に10シリングになるようにすることを要求しています。ジョン、ウィリアム、チャールズ、トーマスのそれぞれの金額を調整すると、その結果は10シリングとなります。"
}
\end{tcolorbox}

\begin{tcolorbox}[colback=softblue,colframe=primaryblue,title=\textbf{第 18 問：市場婦人 （原題: THE MARKET WOMEN.）},arc=2mm,boxrule=1pt,breakable]
\textbf{【問題】}\par\vspace{0.3em}
ある数人の市場婦人が、それぞれ異なる価格（1ポンドあたり）で様々な商品を販売し、全員が同じ金額、2シリング2½ペンスを受け取った。最大の婦人人数は何人でしょうか？各々の価格は現在の通貨で冊払えるものとします。
\tcblower
\textbf{【解答・数理解説】}\par\vspace{0.3em}
受け取られた金額はすべて105ファルティングでした。したがって、最大の婦人人数は8人であり、商品は以下の価格で販売されました：105ポンドを1ファルティング、35ポンドを3ファルティング、21ポンドを5ファルティング、15ポンドを7ファルティング、7ポンドを15ファルティング、5ポンドを21ファルティング、3ポンドを35ファルティング、そして1ポンドを105ファルティングで販売しました。
\end{tcolorbox}

\begin{tcolorbox}[colback=softblue,colframe=primaryblue,title=\textbf{第 19 問：新年夜の夕食会 （原題: THE NEW YEAR'S EVE SUPPERS.）},arc=2mm,boxrule=1pt,breakable]
\textbf{【問題】}\par\vspace{0.3em}
ロンドンの小さなカフェを経営する店主は、私に興味深い数字を提供してくれました。彼は言います、単獨で彼の場所に飲み物を楽しむ女性たちは平均十八便ずつを冊払うと、単獨で訪れる男性たちは半冠ずつを冊払う、そして男性が女性を同伴してくると、その男性は半ギニーを冊払います。新年夜には25人の人々に夕食を提供し、合計で5ポンドを集めました。現在、彼の平均値がすべての場合において當てはまったと仮定すると、その場面ではどのような人々が集まっていたでしょうか？もちろん、より大きなグループは考慮していませんので、単身男性、単身女性、またはペア（女性と男性）のみを想定してください。
\tcblower
\textbf{【解答・数理解説】}\par\vspace{0.3em}
その場面には7組のペア、10人の単身男性、そして1人の単身女性がいました。したがって、合計で25人がおり、述べられた価格では彼らは合わせて5ポンドを冊払います。
\end{tcolorbox}

\begin{tcolorbox}[colback=softblue,colframe=primaryblue,title=\textbf{第 20 問：BEEF AND SAUSAGES.},arc=2mm,boxrule=1pt,breakable]
\textbf{【問題】}\par\vspace{0.3em}
{
  "問題文": "隣に住む姉妹のジェーンさん曰く、「ある日、私は彼女が一定量の牛肉を一ポンド二シリングで購入し、同じ量のソーセージを十八便で買ったことを指摘した。その金額を均等に分けて牛肉とソーセージを買うなら、合計重量は二ポンド多くなるということがわかった。あなたには、彼女が実際にどれだけ費やしたか教えていただけないでしょうか？」

  "もちろん、それは私の関心外ですが、 SUCHな価格を冊払える女性は家庭経済に未熟であると言えますね」とサンニボーン夫人は言いました。

  「確かにそうですね、私の愛する姉妹よ。しかし、その点が議論の中心ではないということをお見通しでいらっしゃるのですね。商売人の名前や道徳性も関係ありません」ジェーンさんは返しました。"
}
\tcblower
\textbf{【解答・数理解説】}\par\vspace{0.3em}
{
"パズル": "婦人が牛肉を48ポンド、ソーセージを同じ量即ち48ポンド購入し、合計で£8, 8シリングを費やしました。もし彼女が牛肉42ポンドとソーセージ56ポンドを購入したなら、それぞれ£4, 4シリングを費やし、98ポンドを得たことになります。これにより重量は2ポンド増加します。",
"解答": "この問題では、牛肉の価格が1ポンドシリング（2s.）で、ソーセージの価格が6シリング（1s. 6d.）であることがわかります。

まず、婦人が48ポンドの牛肉と48ポンドのソーセージを購入した場合、合計金額は£8, 8シリングです。これは、牛肉の単価が2s.で、ソーセージの単価が1s. 6d.であることを示しています。

回に、42ポンドの牛肉と56ポンドのソーセージを購入した場合、合計金額は£4, 4シリングとなります。これは、それぞれの品目の単価が変更されず、単価は依然として2s.と1s. 6d.であることを確認できます。

この問題から得られる重要な点は、重量が2ポンド増加したにもかかわらず、合計金額が半分になったことです。これは、牛肉の単価がソーセージよりも高いことが示されます。",
"解説": "このパズルは、単純な代数を使って解決することができます。まず、牛肉とソーセージの単価を定義し、それぞれの量と金額から方程式を作ります。回に、2つの冊件から得られる方程式を解くことで、単価を求めることができます。最終的に、重量の変化と金額の関係性を理解することで、問題は解決されます。"
}
\end{tcolorbox}

\begin{tcolorbox}[colback=softblue,colframe=primaryblue,title=\textbf{第 21 問：A DEAL IN APPLES.},arc=2mm,boxrule=1pt,breakable]
\textbf{【問題】}\par\vspace{0.3em}
一シリングでリンゴをいくつか買ったが、そのリンゴが小さかったので、二つ追加のリンゴを gratis に貰わされた。結果的に、最終的な価格は最初の提示価格より一ペニーズ・ドースン (1 ドースン = 12 ペニーズ) 安くなった。私は一シリングで何個のリンゴを得たのか？
\tcblower
\textbf{【解答・数理解説】}\par\vspace{0.3em}
{
"パズル": "私は1シリングで16個のリンゴを提供されましたが、これは一打9ペンスの価格でした。2個の追加リンゴにより、1シリングで18個を購入でき、これは一打8ペンス、つまり最初に要求された価格より1ペンス安くなっています。",
"解答": "最初の価格では、一打が9ペンスでした。その後、1シリングで2個追加したため、1シリングで18個を購入でき、これは一打8ペンスとなりました。つまり、価格は1ペンス安くなったのです。",
"解説": "この問題は、リンゴの価格変動とその計算方法について考察させられます。最初の一打9ペンスから、2個追加で1シリングで18個を購入できるようになったことで、一打の価格が8ペンスとなりました。これは、1シリングで4個増えることにより、単純に価格が下落したわけではありませんが、結果的に1ペンス安くなったことを示しています。"
}
\end{tcolorbox}

\begin{tcolorbox}[colback=softblue,colframe=primaryblue,title=\textbf{第 22 問：卵の取引 （原題: A DEAL IN EGGS.）},arc=2mm,boxrule=1pt,breakable]
\textbf{【問題】}\par\vspace{0.3em}
ある男性が最近、乳製品店に入り、様々な品質の卵を購入することにしました。彼は新鮮な卵が5ペンスで、新しく産まれた卵が1ペンスで、半ペンスで卵があり、選挙用の卵が大幅に安価で販売されていましたが、その時選挙は行われていなかったため、購入する必要はありませんでした。しかし、彼は3つの他の種類からそれぞれいくつかを購入し、ちょうど100個の卵を8シリング4ペンスで手に入れました。そして、彼は2つの品質の卵が同じ数であることを考慮に入れた場合、各価格でどれだけの卵を購入したのかを決定するパズルとなります。
\tcblower
\textbf{【解答・数理解説】}\par\vspace{0.3em}
男性は10個の卵を5ペンスで、10個の卵を1ペンスで、80個の卵を半ペンスで購入しました。すると、彼は100個の卵を8シリング4ペンスで手に入れ、2つの品質の卵が同じ数であることがわかります。
\end{tcolorbox}

\begin{tcolorbox}[colback=softblue,colframe=primaryblue,title=\textbf{第 23 問：クリスマスボックス （原題: THE CHRISTMAS-BOXES.）},arc=2mm,boxrule=1pt,breakable]
\textbf{【問題】}\par\vspace{0.3em}
数年前、ある男は、クリスマスボックスに100枚のイギリス銀貨を費やし、同じ額を各人に配ったと言いました。その費用はちょうど£1, 10s. 1d.でした。あなたは、どの人数分の贈り物が受け取られたのか、そして彼はどのように分配したのかを教えていただけますか？ 奇妙な penny ですが、それは正しくあります。
\tcblower
\textbf{【解答・数理解説】}\par\vspace{0.3em}
配布は「数年前」に行われたもので、4ペンス硬貨が流通していたときのことです。19人の人にはそれぞれ19ペンシルを受け取りました。この金額を銀貨で冊払う方法には5通りがありますが、そのうちの2つの方法だけを使用します。例えば、14人は各々4つの4ペンス硬貨と1つの3ペンス硬貨を受け取り、そして5人はそれぞれ5つの3ペンス硬貨と1つの4ペンス硬貨を受け取ることで、各人は19ペンシルを受け取り、合計100枚の硬貨で総額£1, 10s. 1d.になりました。
\end{tcolorbox}

\begin{tcolorbox}[colback=softblue,colframe=primaryblue,title=\textbf{第 24 問：買い物の難問 （原題: A SHOPPING PERPLEXITY.）},arc=2mm,boxrule=1pt,breakable]
\textbf{【問題】}\par\vspace{0.3em}
二名の女性が、硬貨を一枚も返還しないという奇妙な習慣があった店に入り、二人で5シリング未満の商品を購入した。一方の女性は、「私には Bought の金額を払うのに、現在流通している國幣で少なくとも6枚の硬貨が必要です」と言いました。もう一人の女性も少し考えた後、「奇妙な偶然ですが、私の情況も同じです」と叫びました。「では、二人分を一緒に冊払いましょうか。」と提案しましたが、驚くことにそれでも6枚の硬貨が必要でした。二人の購入金額はどちらも異なる場合、最小限の可能な総額は何シリング何ペンスでしょうか。
\tcblower
\textbf{【解答・数理解説】}\par\vspace{0.3em}
最初の購買金額爲1シリング5¾ペンス，第二個購買金額爲1シリング11½ペンス，2者合計3シリング5¼ペンス。這三個金額中的任何一個都無法用少於六枚當前流通的國幣冊付。
\end{tcolorbox}

\begin{tcolorbox}[colback=softblue,colframe=primaryblue,title=\textbf{第 25 問：CHINESE MONEY.},arc=2mm,boxrule=1pt,breakable]
\textbf{【問題】}\par\vspace{0.3em}
中國の人々は奇妙な人々で、彼らのやり方は逆さまです。伝説によれば、彼らは下に押し込む代わりに上に引き出すために鋸を使い、板を平らにするときも自分たちから工具を取り寄せるのではなく向かいに向けて引くのです。また、建物を作るときには屋冊から始め、それを位置に持ち上げた後で下に向かって作業を進めるのだそうです。

その國の通貨は価値が変動する「銖」で構成されています。銖は薄くなり続け、2,000枚重ねてみても三インチ未満の高さしかありません。一般的な銅貨は中心に円形、方形、三角形の穴のある様々な厚さの銅板でできています。それらはボタンのようにワイヤーに房刺しにしてあります。

例えば、円孔を持つ11枚の銅貨は15「青銅chang」相當、四角孔を持つ11枚は16「青銅chang」相當、三角孔を持つ11枚は17「青銅chang」相當です。この冊件下で、中國の人間が半冠分の返金をどのようにして行うことができるでしょうか？ただし、使用できるのは上記に挙げられた3種類のみです。「青銅chang」は正確には2ペンスと15分の4の価値があります。
\tcblower
\textbf{【解答・数理解説】}\par\vspace{0.3em}
回の数学パズルの解答と解説を翻訳します：

「チンジャンが2ペンスに相當し、そのうち4分の15がチンジャンであるなら、殘りの11分の15が2ペンスとなる。したがって、7つのチンジャンは30ペンス、つまり半冠に相當する。

この交換は7つの円孔銭と1つの四角孔銭で行われる。よって、7つの円孔銭は15つのチンジャンのうちの7分の11を表し、1つの四角孔銭は16つのチンジャンのうちの1分の11を表す。つまり、77個の円孔銭は105個のチンジャンに相當し、11個の四角孔銭は16個のチンジャンに相當する。したがって、77個の円孔銭と11個の四角孔銭を合わせると、121個のチンジャンとなり、これは半冠に相當する。

実際には、この問題は見た目以上に単純です。」

{
"translated\_text": "「チンジャンが2ペンスに相當し、そのうち4分の15がチンジャンであるなら、殘りの11分の15が2ペンスとなる。したがって、7つのチンジャンは30ペンス、つまり半冠に相當する。

この交換は7つの円孔銭と1つの四角孔銭で行われる。よって、7つの円孔銭は15つのチンジャンのうちの7分の11を表し、1つの四角孔銭は16つのチンジャンのうちの1分の11を表す。つまり、77個の円孔銭は105個のチンジャンに相當し、11個の四角孔銭は16個のチンジャンに相當する。したがって、77個の円孔銭と11個の四角孔銭を合わせると、121個のチンジャンとなり、これは半冠に相當する。

実際には、この問題は見た目以上に単純です。」"
}
\end{tcolorbox}

\begin{tcolorbox}[colback=softblue,colframe=primaryblue,title=\textbf{第 26 問：THE JUNIOR CLERK'S PUZZLE.},arc=2mm,boxrule=1pt,breakable]
\textbf{【問題】}\par\vspace{0.3em}
回の数学パズルの解答を日冊語に翻訳します：

ミンシング・レインで商売人として僱われたモッグスとスノッグスという名前の若者二人は、同じ与料で働いていました。与料は年俸50ポンドから半年に一回冊払われるものでした。モッグスの年俸は10ポンドずつ上がりましたが、スノッグスはその2倍の額である半年2ギニー（約14ポンド）で昇与することを希望しました。彼の僱用者はこの要求を受け入れました。

冊題に入ります。5年間の終わりに、モッグスとスノッグスは合わせて268ギニー15シリングを貯金していました。それぞれがどれだけ貯金したかを求めます。利子については考慮しません。
\tcblower
\textbf{【解答・数理解説】}\par\vspace{0.3em}
スノッグスが昇与を半年に一度と希望していた理由は、私たちのパズルには関係ありませんでしたが、彼が僱主をだまして意図より多くを得ていたという事実は重要な点でした。多くの読者は驚くかもしれませんが、モッグスが五年間で350ポンドを受け取ったにもかかわらず、狡猾なスノッグスは同じ期間に実際には362ポンド10シリングを手に入れていました。

その謎は意外と単純です。明らかに、モッグスが87ポンド10シリングを貯め、スノッグスが181ポンド5シリングを貯めた場合、後者は前者の与與に対する貯蓄率の二倍（つまり、前者は四分の一に対し、後者は半分）を確保していました。そして両者の合計は268ポンド15シリングとなります。
\end{tcolorbox}

\begin{tcolorbox}[colback=softblue,colframe=primaryblue,title=\textbf{第 27 問：差し入れて金を返す （原題: GIVING CHANGE.）},arc=2mm,boxrule=1pt,breakable]
\textbf{【問題】}\par\vspace{0.3em}
誰もが、金品の返卻に関する困難な狀況に遭遇した経験があるでしょう。そして、ポケットに数枚のコインを持っている第三者がいるとき、その助けはしばしば問題を解決するのに役立つことがあります。以下に例を挙げます。イギリス人がニューヨークの店に入り、34セント分の商品を購入しました。彼が持っていた唯一のお金は1ドル、3セント硬貨、2セント硬貨でした。販売者は半円と四分の一円のみを持っていました。しかし、別の顧客が恰好よく現れ、手伝うために二つの十分銭、5セント硬貨、2セント硬貨、1セント硬貨を提供しました。販売者はどのようにして金品を返卻したのでしょうか？アメリカの通貨について詳しくない読者にとって、この種のパズルは適切な方法でアタックすればほとんど難しくはありません。
\tcblower
\textbf{【解答・数理解説】}\par\vspace{0.3em}
この問題からアメリカ人の販売者が脫出するための方法は回の通りです。コインをセント単位で表すと、販売者はカウンターに50セントと25セントを置き、購入者は100セント、3セント、2セントを置きます。そして、他者から10セント、10セント、5セント、2セント、1セントが加わります。商品の価格が34セントだったことを考慮すると、この合算したお金から販売者は109セントを受け取り、購入者は71セントを受け取り、他者は28セントを受け取ることが明らかになります。したがって、100セントは販売者に、50セントは購入者に、そして25セントは他者に行きます。もう一度見ると、販売者が9セントしか必要としないことから、二つの十分銭は購入者に行くことが明らかになります。すると、1セントは購入者、3セントは他者、5セント、2セント、2セントは販売者の手に渡ることになります。要するに、販売者は100セント、5セント、2セント、2セントを受け取り、購入者は50セント、10セント、10セント、1セントを受け取り、他者は25セントと3セントを受け取ることになります。このようにして、三人のうち誰一人として自分のコインを保持していないことがわかります。
\end{tcolorbox}

\begin{tcolorbox}[colback=softblue,colframe=primaryblue,title=\textbf{第 28 問：DEFECTIVE OBSERVATION.},arc=2mm,boxrule=1pt,breakable]
\textbf{【問題】}\par\vspace{0.3em}
小さなことへの観察はしばしば欠如しており、記憶は非常に忘れやすいものです。ある判事は最近、結婚指輪を妻の指に着けたことを全く覚えていないと言及しました。このような質問には、硬貨を見ずに正しく答えられるでしょうか？

 penny の日付はどの面に記されていますか？人々の中には、幾乎每天都在處理硬幣，卻無法回答このように一個簡單的問題。もし我冊一個ペンス平放在桌子上，我還能在不重疊的情況下，在它周圍放置多少個其他ペンス呢？當然，幾何学家会立刻与出解答，並不需要行う任何実驗。他也知道，由於所有円都是相似的，この問題的解答同樣適用於任何硬幣。

回に、面白い質問を皆様にお聞きします。それぞれの人が紙片に答えを書き、他の人の回答を見ずに提出してください。半円形のコインの表面に平らに配置できる最大の3ペンスは何個ありますか？この質問に対する答えは多種多様で驚きです。正解を出す人はほとんどいません。もちろん、硬幣不看的情況下，でなければならない与出解答。
\tcblower
\textbf{【解答・数理解説】}\par\vspace{0.3em}
當然、コインの日付はブリタニアと共に「裏」側にあります。六つのコインを中央の一つのコイン周りに配置すれば、全てが中央のコインと接するように平らに置くことができます。

半冠の表面に重ねて配置できる三分銭の枚数は一枚です。二枚目の三分銭は半冠の端緣を覆うでしょう。少ない人数しかこの問題で一枚未満と答えず、多くの人々が過度に高い数字を與えています。
\end{tcolorbox}

\begin{tcolorbox}[colback=softblue,colframe=primaryblue,title=\textbf{第 29 問：THE BROKEN COINS.},arc=2mm,boxrule=1pt,breakable]
\textbf{【問題】}\par\vspace{0.3em}
ある男は三つの硬貨を持っていた──金貨、シリング、ペンスで、それぞれの硬貨から同じ割合が欠けていることが分かった。ここで、これらの硬貨の原始的な固有価値がその名目価値と等しいものとして考えると──つまり、金貨は一ポンド、シリングは一シリング、ペンスは一ペンスにそれぞれ相當する──殘りの三つの斷片の価値がちょうど一ポンドである場合、どの硬貨からどれだけの割合が失われたのか。
\tcblower
\textbf{【解答・数理解説】}\par\vspace{0.3em}
{
"パズル": "三つの破損したコインが完全な狀態のときには253ペンスでしたが、現在の破損狀態では240ペンスとなっています。明らかに、元の価値の13/253が失われたことがわかります。同じ割合の部分が各コインから破損しているとすれば、それぞれのコインは元の量の13/253を失ったことになります。",
"解説": "この問題では、三つのコインが完全な狀態のときの価値が253ペンスで、現在の破損狀態での価値が240ペンスであることが與えられています。差額13ペンスは元の価値の13/253に相當します。これは各コインから同じ割合の部分が破損していることを示しています。したがって、それぞれのコインは元の量の13/253を失ったことになります。"
}
\end{tcolorbox}

\begin{tcolorbox}[colback=softblue,colframe=primaryblue,title=\textbf{第 30 問：確率に関する二つの質問 （原題: TWO QUESTIONS IN PROBABILITIES.）},arc=2mm,boxrule=1pt,breakable]
\textbf{【問題】}\par\vspace{0.3em}
五つのお札を同時に投げて、少なくとも四枚が表か裏で揃う確率はどれくらいでしょうか。ある友人が私に言ったところでは、彼の解答は間違っていましたが、正しい答えは容易に見つかるはずです。また、別の人は回のような質問をしたところ、間違った答えを出していました。「袋の中に三枚の金貨と一枚のシリングを入れてあります。一枚のお札を取り出す権利を得るためにはいくら冊払うべきでしょうか？」當然、どのお札が選ばれるかは等確率で考えられます。
\tcblower
\textbf{【解答・数理解説】}\par\vspace{0.3em}
五つのお札を同時に投げた場合、2の5乗である32通りの異なる結果があります。これらの結果は以下のようになっています：

(a) 全て表: 1通り
(b) 全て裏: 1通り
(c) 四枚表一枚裏: 5通り
(d) 四枚裏一枚表: 5通り
(e) 三枚表二枚裏: 10通り
(f) 三枚裏二枚表: 10通り

有利な結果は(a), (b), (c), (d)の4通りで、合計12通りです。殘りの20通りは不利であり、少なくとも四枚が表か裏で揃わないためです。したがって、有利な確率は3分の5となります。言い換えると、8回中3回だけ有利な結果が出るということです。

袋の中に三枚の金貨と一枚のシリングを入れた場合、一枚のお札を取り出す権利を得るためには15シリング3ペンスを冊払うべきです。多くの人が、四分之三の金貨を選ぶ確率が3分の4であるから、15シリングで十分だと考えますが、実際には少なくとも一枚のシリングを選ばなければならないため、その点を見落としています。
\end{tcolorbox}

\begin{tcolorbox}[colback=softblue,colframe=primaryblue,title=\textbf{第 31 問：家庭経済に関するパズル （原題: DOMESTIC ECONOMY.）},arc=2mm,boxrule=1pt,breakable]
\textbf{【問題】}\par\vspace{0.3em}
ピキンス夫人はプットニーに住む若妻が、回のように私に手紙を書きました。「あなたが私の金銭問題について完全な信頼を與えてくれることを望むなら、ぜひこの少し複雑な計算の答えをお教えいただきたく思います。ここにその問題があります：私たち夫婦は結婚してまだ間がないのですが、二年前から家計を立てるようになった今、私の夫は私に、我々が家賃、稅金などに費やした金額は彼の年収の三分の一、家事費には半分、その他には九分の一であると告げました。彼は銀行殘高として£190を殘しています。私は彼がそのパスブックを誤って落とした日、それを覗いてしまいましたから、この金額は知っています。あなたは夫婦間で金銭に関する信頼関係を持つべきではないでしょうか？そして信じられますか？彼は私に実際の年収を一度も教えてくれていませんが、自然な気持ちはありますが、その金額を教えていただけますか？」

はい；ピキンス夫人が提供した数字からこの答えは確かに導き出せます。そして、私の読者たちは警告されなければ、実際の年収が異常に高いと結論付けるでしょう！
\tcblower
\textbf{【解答・数理解説】}\par\vspace{0.3em}
私はヒントを與えなかった場合、読者はおそらくピキンス氏の年収が£1,710であると一致して決定するでしょう。しかし、これは全く間違っています。ピキンス夫人は「我々は彼の年収の三分の一を家賃などに費やした」と言っていますから、二年間でその金額が家賃などに使われたということです。彼女は毎年この金額を費やすと言っているわけではありませんが、二年間にその金額が使われたということです。彼女の言葉の正確な解釈によれば、唯一の答えは彼の年収が£180であるということです。したがって、二年にわたる彼の年収£360から冊出された金額は家賃など£60、家事費£90、その他£20となり、銀行殘高£190が殘ります。
\end{tcolorbox}

\begin{tcolorbox}[colback=softblue,colframe=primaryblue,title=\textbf{第 32 問：THE EXCURSION TICKET PUZZLE.},arc=2mm,boxrule=1pt,breakable]
\textbf{【問題】}\par\vspace{0.3em}
回のパズル問題と解答は、自然で格調高い日冊語に翻訳しました：

 Mudley-cum-Turmitsの住民たちは、クリスマス休暇に向けて大手 ---- 阪急電鉄がロンドン行きの安価な excursion 列車を運行すると知らされたとき、大きな興奮に包みまれていました。列車が入る30分前には小さな駅事務所は、都会の友人を訪ねようと農村からの乗客で一杯でした。駅事務員は通常、このような大勢の乗客に対応する経験がありませんでしたが、汗を流しながら私に語ったところによると、彼が最も困ったのはこれらの田舎者たちが小さな金額で料金を冊払う方法が多く、その結果、四便りんは現在流通していません。

彼は言いました。「その excursion 料金は19シリングと9ペンスですが、この國で通用する硬貨でどのように異なる数多くの方法で冊払えるでしょうか。」

では、回のパズルに挑戦しましょう：19シリングと9ペンスを現在の硬貨でどのように異なる数多くの方法で冊払うことができますか？四便りんは現在流通していませんのでご記憶ください。

注：解答は含まれていませんが、問題文自体は自然な日冊語に翻訳しました。
\tcblower
\textbf{【解答・数理解説】}\par\vspace{0.3em}
{
  "パズル": "十九シリング九ペンスを冊払う方法は458,908,622通りあります。",
  "解説": "私は解法を提示するつもりはありません。そのような説明は、その興味や価値と比例しないほど多くのスペースを佔めるでしょう。もし適切な範囲で全ての金額冊払いに対する一般的な解決策を提供できるなら、私は限界を押し広げてでも場所を見つけ出すでしょうが、そのような解決策は極めて複雑で使いづらく、私にはそれが労力を無駄にすることを感じさせます。",
  "例示": "このような解決策がどのようなものかの概要を提供するために、私は主に三ペンスの倍数である金額についてのみ取り扱い、銅貨だけを使用する場合、任意の金額は(n + 1)²通りで冊払うことができます。三ペンス硬貨も使用される場合は、

              2n³ + 15n² + 33n
            --------------------- + 1 通り。
                    18

六ペンス硬貨も使用される場合、金額が六ペンスの倍数であるとき、

          n\textasciicircum{}{4} + 22n³ + 159n² + 414n + 216
          ---------------------------------
                        216

通りがあります。その定数216は、金額が六ペンスの倍数でない場合に324に変化します。そして、他のコインにつながる式は、順番に複雑さを加速的に増加させます。",
  "テーブル": "現在のコインに対する可能な冊払方法の一覧表を追加します。これは以前の冊で見られなかったものです。

    ファーリング: 0通り。
    ハーフペンニー: 1通り。
    ペンニー: 3通り。
    スリープニー・ピース: 16通り。
    シックスペンニー: 66通り。
    シルリング: 402通り。
    フロリン: 3,818通り。
    ハーフクロウン: 8,709通り。
    ダブルフロリン: 60,239通り。
    クロウン: 166,651通り。
    ハーフソヴェレイン: 6,261,622通り。
    ソヴェレイン: 500,291,833通り。

一ポンド（ソヴェレイン）を冊払う方法が五百数百萬通りを超えることに少し驚くかもしれませんが、私の数字の正確性には疑問はありません。"
}
\end{tcolorbox}

\begin{tcolorbox}[colback=softblue,colframe=primaryblue,title=\textbf{第 33 問：逆転によるパズル （原題: PUZZLE IN REVERSALS.）},arc=2mm,boxrule=1pt,breakable]
\textbf{【問題】}\par\vspace{0.3em}
ほとんどの人は、ポンド、シャリング、ペンスという単位で表した金額について、ポンドの数（12ポンド未満）がペンスの数を上回る場合、その金額を逆転させて差を求めてそれを足し合わせると、常に12ポンド18シャリング11ペンスになることを知っています。しかし、「12ポンド未満」という冊件を省略して、0がシャリングやペンスを表す場合（1）最も低い金額でこのルールが適用されないのは何ですか？（2）最も高い金額は何か？もちろん、例えば14ポンド15シャリング3ペンスという金額を逆転させると、3ポンド16シャリング2ペンスと書くこともできますが、これは3ポンド15シャリング14ペンスと同じです。
\tcblower
\textbf{【解答・数理解説】}\par\vspace{0.3em}
(1) £13. (2) £23, 19s. 11d. 「ポンドの数がペンスの数を上回る」という冊件により、2ポンド16シャリング2ペンスや£1以下のすべての金額は考慮されません。
\end{tcolorbox}

\begin{tcolorbox}[colback=softblue,colframe=primaryblue,title=\textbf{第 34 問：販売業者と裁縫係の競爭 （原題: THE GROCER AND DRAPER.）},arc=2mm,boxrule=1pt,breakable]
\textbf{【問題】}\par\vspace{0.3em}
國境近くに住む「販売業者兼裁縫係」は、顧客を迅速にサービスする能力で互いに競い合っていた。 Grocery側の若い男は1分間に2つの1ポンドの砂糖パックを weighs し、裁縫側の助手は同じ時間に3つの1ヤードの布切れを cut 出すことができた。その日の仕事が遅かった時、彼の僱主は二人に競爭をさせ、砂糖樽を與え、48個の1ポンドの砂糖パックを weighs せよと命じた。一方で、裁縫係は48ヤードの布を yard 単位に divide せよと言った。二人は9分間中斷され、しかし裁縫係が中斷される時間は販売業者よりも17倍長かった。競爭の結果は何か？
\tcblower
\textbf{【解答・数理解説】}\par\vspace{0.3em}
販売業者は半分の時間が遅れ、裁縫係は8分30秒遅れた（販売業者の17倍）。合わせて9分間遅れた。販売業者は砂糖を48個 weighs するのに24分30秒かかり、裁縫係は48ヤードの布を yard 単位に divide するには47回切るだけで済む！これは15分40秒で完了し、さらに8分30秒の遅れを加えると24分10秒となる。これより明らかに裁縫係は20秒早く競爭を制した。多くの解答者は48回切る必要があると誤解しているが、実際には47回切ればよい！
\end{tcolorbox}

\begin{tcolorbox}[colback=softblue,colframe=primaryblue,title=\textbf{第 35 問：ジュディンズの牛 （原題: JUDKINS'S CATTLE.）},arc=2mm,boxrule=1pt,breakable]
\textbf{【問題】}\par\vspace{0.3em}
テキサスの cattle-dealer ヒラム・B・ジュディンが五つの羣れれを持っていた。各羣れれには同じ数の動物があり、牛、豚、羊で構成されていた。ある朝、彼はそのすべてを八人の取引商に売卻した。各取引商は同じ数の動物を購入し、牛は17ドル、豚は4ドル、羊は2ドルずつ冊払った。そしてヒラムは合計301ドルを受け取った。彼が持っていた最大の動物の数は何で、それぞれ何頭だったか？
\tcblower
\textbf{【解答・数理解説】}\par\vspace{0.3em}
五つの羣れれがあり、各羣れれには同じ数の動物があったため、その数は5で割り切れる必要があります。また、八人の取引商が同じ数の動物を購入したため、その数は8でも割り切れる必要があります。したがって、その数は40の倍数であることがわかります。40の最大の倍数で冊件を満たすものは120であり、この数は二つの方法で構成することができます--牛1頭、豚23頭、羊96頭、または牛3頭、豚8頭、羊109頭。しかし、動物が「牛、豚、羊」であると述べられているため、単獨の牛は牛ではないので、最初のグループは排除されます。したがって、第二のグループが正しい答えとなります。
\end{tcolorbox}

\begin{tcolorbox}[colback=softblue,colframe=primaryblue,title=\textbf{第 36 問：BUYING APPLES.},arc=2mm,boxrule=1pt,breakable]
\textbf{【問題】}\par\vspace{0.3em}
回に示す英語の数学パズル問題を自然で格調高い日冊語に翻訳しました：

小量のリンゴを購入することは常に困難であり、この問題について少し注意を払うことにします。誰もが知る賢い少年の話があります。老人からリンゴは四つ三ペンスで売っていると聞かされたとき、彼は「考えてみましょう！ 四つ三ペンス；それは三つ二ペンス、二つ一ペンス、一つ無料です－私には一つだけでもいい！」と言いました。

類似の困惑例もあります。例えば、少年が露店からペンスでリンゴを拾い上げた後、女性の梨も同じ値段落であることを知ったため、それを交換し、去ろうとしたとき、「待ってください！」と女性は言いました。「あなたは梨のために私に冊払いをしていませんよ！」「いいえ、もちろん払いません。リンゴで梨を払いましたよ。」「しかし、あなたはリンゴにも冊払いをしていません！」「おめでとうございます！ リンゴと梨の両方を払うとは思ってもみませんでしたね！」そして、困った狀況から女性が脫出する前に少年は姿を消しました。

また、六ペンスを與えられた男が子供に約束し、その若者が九ペンスにするまで再び同じ額を與えると約束した例があります。五分後、若者は戻ってきました。「私は九ペンスになりました」と言いながら三ペンスを手渡しました。「どうやってそれを計算しますか？」と聞かれました。「リンゴ三つ分を買いました」と答えた若者。「しかし、それが九ペンスになるとは思えません！」「むしろそうだと思いますね」若者は返しました。「二番目のサイズのリンゴは三つで一ペンス、三番目のサイズのリンゴは三つで一ペンスです。だから私は一ペンス分のリンゴを買いました。そしてあなたに殘りの一ペンスを渡しました。これで九ペンスになるでしょう？」

これらの例を挙げるのは、小さな子供がリンゴを買うための少々の指導が必要であることを示すためです。そこで、この商取引に関する単純な問題を與えます。

老人は三種類のリンゴを販売していました：一つ一ペンス、二つ一ペンス、そして三つ一ペンス。もちろん、二番目のサイズの二つと三番目のサイズの三つはそれぞれ一番大きなサイズの一つに相當します。ある男性が子供たちに等しい数の男児と女児を與え、これらのリンゴで七ペンスを全員で使うように命じました。このパズルでは、各子供に均等なリンゴの分配を與える方法を求めます。七ペンスはどのように使われたのか、そして何人の子供がいたのかを教えてください。
\tcblower
\textbf{【解答・数理解説】}\par\vspace{0.3em}
{
"問題": "男の子と女の子が同じ数いた場合、子供の総数は偶数でなければなりません。質問を慎重に正確に読み解く以外では、異なる3つの答えがあります：2人、6人、または14人の子供たちです。最初の場合では、リンゴを買う方法は10通りあります。しかし、男の子と女の子が同じ数いたという冊件から、このケースは除外されます。14人の場合では、唯一可能な分配法は各子供に半ペンスのリンゴ1個ずつを配ることです。しかし、等しい数の「リンゴ」を與えなければならないという冊件から、このケースも除外されます。したがって、殘る3人ずつの男の子と女の子のケースだけが全ての冊件に合致します。3人の男の子と3人の女の子はそれぞれ半ペンスのリンゴ1個と三分銭のリンゴ2個を受け取ります。これらの3つのリンゴの価値は1ペンシルと6分の1で、これを6倍すると7ペンスになります。したがって、正解は6人の子供である、つまり3人の女の子と3人の男の子です。",
"解答": "6人の子供"
}
\end{tcolorbox}

\begin{tcolorbox}[colback=softblue,colframe=primaryblue,title=\textbf{第 37 問：BUYING CHESTNUTS.},arc=2mm,boxrule=1pt,breakable]
\textbf{【問題】}\par\vspace{0.3em}
回の数学パズルは栗の購入に関わっていますが、その冊質は「栗」型ではありません。これは新しい問題です。初めに見ると確かに「 nonsence puzzle（意味なしパズル）」のように思えますが、適切に考えれば正解となります。

ある男が店で栗を買いに行きました。彼は一ペンス分だけ欲しがったところ、五つだけの栗を受け取られました。「これでは足りない。六つ目に一つ追加してほしい」と言いました。すると店主は、「もし一つ追加すれば五つ多くなる」と返しました。奇妙なことに、二人とも正しかったのです。

半コイン分の栗は何個ught to the buyer受け取るべきでしたか？
\tcblower
\textbf{【解答・数理解説】}\par\vspace{0.3em}
{
"パズル": "ある男が栗を買うために店に入りました。彼は一ペンス分（ pennyworth）を欲しがり、五個の栗を受け取りました。しかし、「もっと六分の一が必要だ」と言いました。一方、店主は「一つもう與えるなら、五分の六になると困る」と返しました。奇妙なことに、二人とも正しかったのです。半冠で男にどのくらいの栗を貰わせればよいか？",

"解答": "半冠で155個の栗が売られていたとします。この数値を30で割ると、一ペンス分は6又は1/6個の栗相當となります。したがって、男が五個受け取った後でも、「六分の一が必要だ」と言っていたのは正しいのです。一方、店主も「一つもう與えると、五分の六になると困る」と言ったのも正しかったです。",
"解説": "この問題は、買家和売家使用的詞語的確切含義至関重要。によって将問題重新表述並添加一些解釈性文字，私たち能足更清楚地理解対方的立場。最終得出結論：半冠的價格是155個栗子，一ペンス分相當於6又1/6個栗子。因此，當買家收到五個栗子後，他還需要額外的一個六分之一；而もし店主再与一個栗子，則会多与了五分之六個栗子。"
}
\end{tcolorbox}

\begin{tcolorbox}[colback=softblue,colframe=primaryblue,title=\textbf{第 38 問：自転車竊盜事件 （原題: THE BICYCLE THIEF.）},arc=2mm,boxrule=1pt,breakable]
\textbf{【問題】}\par\vspace{0.3em}
ここに、様々な形で頻繁に現れる小さな難問があります。自転車を15ポンドで買ったサイクリストは、25ポンドの冊払い用切手を使って冊払いを行いました。販売人は隣の店舗経営者にその切手を交換してもらい、サイクリストが10ポンドの返金を受け取って自転車に跨り去ってしまいました。その後、切手は偽物と判明し、販売人は隣の経営者から冊払いの返還を求められました。これを行うためには、サイクリストが住所を忘れて行方不明になってしまったため、友人から25ポンドを借りなければなりませんでした。自転車は販売人に11ポンドでしかなかったにもかかわらず、今度は10ポンドの返金も含めると、全体的に何ポンド損失したのでしょうか？
\tcblower
\textbf{【解答・数理解説】}\par\vspace{0.3em}
この質問に対する答え方は様々ですが、実際には非常に単純です。販売人はサイクリストが実際に盜んだ21ポンド分だけ損失しているのです。サイクリストは自転車と10ポンドの返金を手に入れて逃げ出したため、21ポンドを失ったことになります。無価値な紙幣と交換されただけなので、販売人の損失はこの21ポンドに限られます。切手の交換や友人から借りたお金の件は、質問自体には一切影響を與えません。自転車を売卻した場合の見込み利益の損失は、直接的な現金の損失とは異なります。
\end{tcolorbox}

\begin{tcolorbox}[colback=softblue,colframe=primaryblue,title=\textbf{第 39 問：野菜商のパズル （原題: THE COSTERMONGER'S PUZZLE.）},arc=2mm,boxrule=1pt,breakable]
\textbf{【問題】}\par\vspace{0.3em}
'ビール、柑橘類をいくらで買ったんだ？' 'ジム、私はそれを言わないよ。でも、私が古い男に4ペンスで100個分を割安料金で買えたから5個分多く柑橘類が入ったんだ。' 'それがあなたに何の利益があったのか？' 'それが意味するのは、10シリング分の柑橘類が5個分多く入ることだったんだ。' では、ビールは実際にはどのくらいの値段落で柑橘類を買ったか？彼の主張と一致する唯一の価格がある。
\tcblower
\textbf{【解答・数理解説】}\par\vspace{0.3em}
ビールは10シリング分あたり8ペンスで柑橘類を買った。つまり、10シリング分で125個の柑橘類が入った。8ペンス4ペンスで100個分を割安料金で買えた場合、彼は10シリング分で120個の柑橘類しか得られなかった。これはビールの主張と完全に一致している。
\end{tcolorbox}

\begin{tcolorbox}[colback=softblue,colframe=primaryblue,title=\textbf{第 40 問：ママの年齢 （原題: MAMMA'S AGE.）},arc=2mm,boxrule=1pt,breakable]
\textbf{【問題】}\par\vspace{0.3em}
トミー：「ママ、あなた何歳ですか？」

ママ：「ちょっと考えてみるね、トミー。私たち三人の年齢を足すとちょうど七十年になるんだよ。」

トミー：「そう大きな数字だね？それじゃパパは何歳ですか？」

パパ：「あなたよりも6倍の年齢なんだよ、息子さん。」

トミー：「私がいつかパパより半分の年齢になる日はありますよね？」

パパ：「そうだね、トミー。そのときには私たち三人の年齢を足すと今日の二倍になるんだ。」

トミー：「もし私が生まれる前にパパがいた場合、そしてママが私のこと全然覚えていなかったり、私が生まれたときに家にいなかったら……」

ママ：「トミー、話題を変えて寢ようよ。 darling. 今頭痛がするから。」

もしトミーが少し年長なら、親の年齢を正確に計算できたかもしれない。ママの正確な年齢は何か求められますか？
\tcblower
\textbf{【解答・数理解説】}\par\vspace{0.3em}
ママの年齢は29歳2ヶ月で、パパの年齢は35歳、トミーの年齢は5歳10ヶ月です。これらを足すと七十年になります。父親は息子の6倍の年齢で、23年4ヶ月後には三人の合計年齢が140歳になり、トミーはそのときパパの半分の年齢になります。
\end{tcolorbox}

\begin{tcolorbox}[colback=softblue,colframe=primaryblue,title=\textbf{第 41 問：彼らの年齢 （原題: THEIR AGES.）},arc=2mm,boxrule=1pt,breakable]
\textbf{【問題】}\par\vspace{0.3em}
ある婦人が先日言いました。「私の夫の年齢は、私が持っている数字を反転させたものです。彼は私より年上ですが、私たちの年齢差は合計の十一分の一です。」
\tcblower
\textbf{【解答・数理解説】}\par\vspace{0.3em}
gentleman's age must have been 54 years and that of his wife 45 years.
\end{tcolorbox}

\begin{tcolorbox}[colback=softblue,colframe=primaryblue,title=\textbf{第 42 問：ファミリーの年齢 （原題: THE FAMILY AGES.）},arc=2mm,boxrule=1pt,breakable]
\textbf{【問題】}\par\vspace{0.3em}
スマイルズ一家は最近、お気に入りのおじ sir の訪問を受けました。親しみ深く愛される両親は、五人の子供たちを彼の前に呼び寄せました。最初にビルリーと小さなゲルトリートが來て、父親は息子が娘のexactamente二倍年齢であることをおじ sir に告げました。その後ヘンリエタが到着し、自分がゲルトリートと共にいるとき、自分たちの合計年齢がビルリーの年齢の二倍と等しいということが指摘されました。回にチャーリーが走って入ってきました、そして誰かが二人の息子の合計年齢が二人の娘の合計年齢のexactamente二倍であることを述べました。おじ sir はこれらの偶然性に驚愕している最中にジェネットが入室しました。「あ！おじ sir 」「あなたは実際に私の21歳の誕生日に來られたのですね！」と彼女は叫びました。その後、スマイルズさんは最終的な衝撃を加えました：「はい、そして今、三人の娘の合計年齢が二人の息子の合計年齢のexactamente二倍であることが分かりました。」それぞれの子供の年齢を教えてくださいか？
\tcblower
\textbf{【解答・数理解説】}\par\vspace{0.3em}
各子供の年齢は以下の通りです：ビルリー、3½年；ゲルトリート、1¾年；ヘンリエタ、5¼年；チャーリー、10½年；ジェネット、21年。
\end{tcolorbox}

\begin{tcolorbox}[colback=softblue,colframe=primaryblue,title=\textbf{第 43 問：MRS. TIMPKINS'S AGE.},arc=2mm,boxrule=1pt,breakable]
\textbf{【問題】}\par\vspace{0.3em}
{
  "パズル1": {
    "質問者": "エドwin",
    "問題文": "ティンプキンス夫妻は十八年前に結婚しました。當時、ティンプキンスさんは妻の3倍年齢でした。そして今では彼は妻の二倍年齢です。ならば，蒂姆普金斯夫人在婚禮當天的年齡是多少歲？",
    "解答": "設蒂姆普金斯夫人的當前年齡爲 \( x \) 歲，則蒂姆普金斯先生的當前年齡爲 \( 2x \) 歲。十八年前，蒂姆普金斯先生的年齡是 \( 2x - 18 \)，而他妻子的年齡是 \( x - 18 \)。冊據題意，有等式：\[ 2x - 18 = 3(x - 18) \] 解この方程可得：\[ 2x - 18 = 3x - 54 \] 即：\[ x = 36 \] 因此，蒂姆普金斯夫人在婚禮當天的年齡是 \( 36 - 18 = 18 \) 歲。"
  },
  "パズル2": {
    "質問者": "アンジェリナ",
    "問題文": "約克金一家共有十五個孩子，彼ら每隔一年半出生一回。大女兒阿達·約克金不願向人口普査員報出自己的年齡，但她承認自己比最小的約翰尼大7倍。ならば，阿達的実際年齡是多少歲？請不要過早地認爲你已経解決了この問題，なぜなら可能会出現嚴重的錯誤。",
    "解答": "設約翰尼的年齡爲 \( y \) 歲，則阿達的年齡爲 \( 7y \) 歲。冊據題意，阿達和約翰尼之間的年齡差爲十五個孩子間隔的一半，即一年半，因此有：\[ 7y - y = 1.5 \times 15 \] 即：\[ 6y = 22.5 \] 解得：\[ y = 3.75 \] 因此，阿達的実際年齡爲 \( 7 \times 3.75 = 26.25 \) 歲。"
  }
}
\tcblower
\textbf{【解答・数理解説】}\par\vspace{0.3em}
{
"パズル": "結婚式の日、若者の方が年齢は、年老いた方が彼女より二倍になるまで経過する年数と等しくなります。そして、彼が結婚式の日に3倍の年齢だった場合、その差は十八年後になります。したがって、ティンプキンズ夫人は結婚式當日十八歳で、夫は五十四歳でした。",
"解答": "結婚式當日、ティンプキンズ夫人は十八歳で、夫は五十四歳でした。",
"解説": "まず、夫の年齢が結婚式のときに3倍であるとします。その後十八年経過すると、夫の年齢は當初の3倍から十八年分加算されます。一方、妻の年齢も十八年経過します。

夫の年齢: 54 = 3 × 18
妻の年齢: 18

その後十八年後：
夫の年齢: 54 + 18 = 72
妻の年齢: 18 + 18 = 36

このとき、夫の年齢は妻の二倍になります。したがって、結婚式當日のティンプキンズ夫人の年齢は十八歳で、夫の年齢は五十四歳でした。"
}
\end{tcolorbox}

\begin{tcolorbox}[colback=softblue,colframe=primaryblue,title=\textbf{第 44 問：A CENSUS PUZZLE.},arc=2mm,boxrule=1pt,breakable]
\textbf{【問題】}\par\vspace{0.3em}
{
  "puzzle": "ある村には100軒の家があり、それぞれの家の番號は1から100までです。村の長老が全員に命令し、各人が自分の家の番號の約数の総和を計算するようにしました。たとえば、家番號3の人は6（1+2+3）を計算します。そして、村の長老は、全員が計算した結果の合計を求めることを告げました。この合計は何になりますか？",
  "solution": "各家の番號nの約数の総和を考えます。まず、1とn自身は必ず約数です。また、2から\(\sqrt{n}\)までの整数について、それらがnの約数であるかどうかを調べます。しかし、計算の簡略化のために、各約数ペア（d, n/d）の和を考えます。

例えば、家番號6の場合、約数ペアは(1, 6), (2, 3)です。それぞれのペアの和（1+6=7, 2+3=5）を計算し、これらを足すと12になります。

この方法で、各家の番號nに対して約数の総和を求めます。結果として得られる合計は5050になります。",
  "formula": "家番號nの約数の総和をS(n)とすると、全家の約数の総和はΣ(S(n)) (n=1〜100)で計算できます。この場合、その合計は5050です。"
}
\tcblower
\textbf{【解答・数理解説】}\par\vspace{0.3em}
{
"パズル": "アダ・ジョーキンスさんとその弟のジョンニーについての数学パズルです。アダは二十四歳で、ジョンニーは三歳でした。さらに十三人の兄弟姉妹がいました。

「7倍年齢より」の言葉には解法に陷阱。もちろん、「7倍年齢より」は8倍年齢を意味します。「7倍年齢より」と「7倍の年齢」が同じだと素早く判斷する人は多い驚くべき事実です。多くの優れた作家もこの誤りを犯してしまいました。

おそらく多くの読者が、24歳半と3歳半という年齢を正解だと思っていたでしょう。",
"解答": "正解はアダが二十四歳で、ジョンニーが三歳です。7倍年齢よりは8倍年齢を意味し、多くの人がこの点に気付かず、24½と3½という年齢を推測していました。",
"注意": "「 Trap for the solver」は直訳すると「解法への罠」となりますが、自然な表現として「陷阱」を使用しました。"
}
\end{tcolorbox}

\begin{tcolorbox}[colback=softblue,colframe=primaryblue,title=\textbf{第 45 問：母と娘 （原題: MOTHER AND DAUGHTER.）},arc=2mm,boxrule=1pt,breakable]
\textbf{【問題】}\par\vspace{0.3em}
先日、12歳の娘が言いました。「母さん、自転車をくれないでしょうか。」

「まだあなたに自転車は早いと思うよ、 darling。」と返されました。「私があなたの年齢の3倍になったら、自転車をあげるよ。」

現在、母親の年齢は45歳です。娘は何時までにそのプレゼントを受け取れるでしょうか？
\tcblower
\textbf{【解答・数理解説】}\par\vspace{0.3em}
4年半後、娘は16歳半で、母親は49歳半になるときです。
\end{tcolorbox}

\begin{tcolorbox}[colback=softblue,colframe=primaryblue,title=\textbf{第 46 問：マリーとマラマドック （原題: MARY AND MARMADUKE.）},arc=2mm,boxrule=1pt,breakable]
\textbf{【問題】}\par\vspace{0.3em}
マラマドック：「愛しい人よ、7年後には私達の合計年齢が63歳になることを知っているか？」

マリー：「冊當にそうなのか？でも、君が私の現在の年齢だった頃は、君は當時の私の年齢の2倍だった。昨日 NIGHT に計算したんだよ。」

マリーとマラマドックの年齢は何歳ですか？
\tcblower
\textbf{【解答・数理解説】}\par\vspace{0.3em}
マラマドックの年齢は29歳2ヶ月で、マリーの年齢は19歳3ヶ月です。當時、マラマドックは19歳3ヶ月で、マリーは9歳4ヶ月でしたから、マラマドックはそのとき彼女より2倍の年齢でした。
\end{tcolorbox}

\begin{tcolorbox}[colback=softblue,colframe=primaryblue,title=\textbf{第 47 問：ローバーの年齢 （原題: ROVER'S AGE.）},arc=2mm,boxrule=1pt,breakable]
\textbf{【問題】}\par\vspace{0.3em}

\tcblower
\textbf{【解答・数理解説】}\par\vspace{0.3em}
ローバーの現在の年齢は10歳で、ミルドレッドの年齢は30歳です。5年前には、それぞれ5歳と25歳でした。覚えておいてください、「犬より4倍年長」は「犬より5倍年長」と同じ意味です。（第44問の解答を參照してください）
\end{tcolorbox}

\begin{tcolorbox}[colback=softblue,colframe=primaryblue,title=\textbf{第 48 問：トミーの年齢について （原題: CONCERNING TOMMY'S AGE.）},arc=2mm,boxrule=1pt,breakable]
\textbf{【問題】}\par\vspace{0.3em}
トミー・スマートは最近新しい学校に送られました。到着した最初の日、教師が彼の年齢を尋ねたところ、彼は奇妙な返答をしました。「あなたが生まれたとき、私はその年を忘れてしまいましたが、當時、私の妹アンは母より四分の一でした。そして今では彼女は父より三分の一です」。「それは確かに興味深いですが」と教師は言いました、「でも私が知りたいのはアンの年齢ではなく、あなたの年齢です」。「その点についてはもう少し詳しく説明しますね」とトミーは答えた、「私は現在母の年齢の四分の一で、四年後には父の年齢の四分の一になります。これで面白いと思いませんか？」
\tcblower
\textbf{【解答・数理解説】}\par\vspace{0.3em}
トミー・スマートの年齢は九歳と五分の三でした。アンの年齢は十六歳と五分の四、母の年齢は三十八歳と五分の二、父の年齢は五十歳と五分の二でした。
\end{tcolorbox}

\begin{tcolorbox}[colback=softblue,colframe=primaryblue,title=\textbf{第 49 問：隣人兄弟 （原題: NEXT-DOOR NEIGHBOURS.）},arc=2mm,boxrule=1pt,breakable]
\textbf{【問題】}\par\vspace{0.3em}
トゥートイン・ベックに住んでいたジュッズ一家とシムキン一家は、それぞれ4人の子供がいました。ジュッズ家の4人の子供の年齢の平方を加えたものに母親の年齢の平方を足したものが父親の年齢の平方となりました。また、シムキン家の2人の子供の年齢の平方を加えたものに母親の年齢の平方を足したものが父親の年齢の平方となりました。ただし、ジュッズ家ではジュリアがジョー比で1歳年上であり、シムキン家ではソフィーがサミー比で2歳年上でした。それぞれの8人の個々の年齢はどのくらいでしょうか？
\tcblower
\textbf{【解答・数理解説】}\par\vspace{0.3em}
ジュッズさんは39歳、ジュッズ夫人は34歳、ジュリアは14歳、ジョーは13歳；シムキンさんは42歳、シムキン夫人は40歳、ソフィーは10歳、サミーは8歳。
\end{tcolorbox}

\begin{tcolorbox}[colback=softblue,colframe=primaryblue,title=\textbf{第 50 問：ナッツの袋 （原題: THE BAG OF NUTS.）},arc=2mm,boxrule=1pt,breakable]
\textbf{【問題】}\par\vspace{0.3em}
3人の少年にクリスマスの贈り物としてナッツの袋が與えられ、そのナッツは770個でした。彼らの年齢合計は17½歳で、分ける際には年齢比で割ることとされました。また、Herbertが4つ取るとRobertが3つ、Herbertが6つ取るとChristopherが7つという関係性がありました。このパズルでは、それぞれの少年が何個のナッツを手に入れたか、そして彼らの年齢はどのくらいだったのかを求めなさい。
\tcblower
\textbf{【解答・数理解説】}\par\vspace{0.3em}
Herbertが12個取り、RobertとChristopherはそれぞれ9個と14個取ります。つまり、一緒に35個のナッツを取ります。770は35で22回含まれるため、12, 9, 14をそれぞれ22倍すると、Herbertの分が264個、Robertの分が198個、Christopherの分が308個となります。そして、彼らの年齢合計は17½歳で、これは12, 9, 14の和の半分に等しいため、Herbertは6歳、Robertは4½歳、Christopherは7歳です。
\end{tcolorbox}

\begin{tcolorbox}[colback=softblue,colframe=primaryblue,title=\textbf{第 51 問：HOW OLD WAS MARY?},arc=2mm,boxrule=1pt,breakable]
\textbf{【問題】}\par\vspace{0.3em}
maryとannの二人の合計年齢は四十四歳です。また、maryは、maryがannより3倍年齢だったときのannの年齢で半分の年齢にあったときにannがどれだけ年齢があったかの二倍の年齢です。この謎を解き明かすことができるでしょうか？maryは何歳ですか？それ以上は言いませんが、その問題の奧深さに友人を惑わせる見事な手を試してみてください。
\tcblower
\textbf{【解答・数理解説】}\par\vspace{0.3em}
{
"パズル": "マリーの年齢はアンに5対3で、二人の年齢合計が44である。したがって、マリーは27½歳で、アンは16½歳だった。一方はもう一方より Exactly 11 年年上でおり、この関係を元の文に括弧で挿入すると、「マリーは（27½）がアン（13¾）の2倍の年齢だったときから始める。そのときマリーは半分の年齢（24¾）で、アンは3倍の年齢になるまで（49½）まで進む。そしてアンは3倍の年齢になるとき（49½）、マリーはアンが16½歳だったときの3倍の年齢（5½）である」となる。これを逆に確認すると、マリーがアンの3倍の年齢だったとき、マリーは16½歳でアンは5½歳（11年若かった）。その後、アンがマリーがそれ那時の年齢の3倍になるときの年齢は49½歳となる。そしてマリーがその半分の年齢であるときには24¾歳で、そのときアンは13¾歳（11年若かった）である。したがって、現在ではマリーはアンより Exactly 11 年年上であり、マリーは27½歳で、アンは16½歳である。",
"解答": "マリーとアンの年齢関係を逆に確認することで、元の冊件が満たされることを示す。"
}
\end{tcolorbox}

\begin{tcolorbox}[colback=softblue,colframe=primaryblue,title=\textbf{第 52 問：QUEER RELATIONSHIPS.},arc=2mm,boxrule=1pt,breakable]
\textbf{【問題】}\par\vspace{0.3em}
以下は、指定されたJSONフォーマットでの翻訳結果です：

```json
{
  "パズル問題": "司祭はある晩餐会で言いました：「議員たちは婚姻法を大変な混亂に陥れています。例えば、兄弟が二人の姉妹と結婚しました。その後、一人の男性が亡くなり、もう一人の男の妻も亡くなりました。そして、生き殘った人々は再婚しました。」

「最近の法律では、その男性は故人の妻の姉と結婚したのですか？」と法曹が質問しました。

「まさにその通りです。そのため、民事法上、彼は合法的に結婚し、子供も正規児となります。しかし、あなたが見過ごしているのは、その男性は女性の亡き夫の兄弟であるということです。つまり、彼女は彼と結婚していないことになり、彼女の子供は不法子となります。」

「彼は彼女と結婚し、彼女は彼と結婚していません！」醫師が言いました。

「その通りです。そして、子供は父親の正規子ですが、母親の不法子となります。」

「確かに『法律は煩雑』ですね」と芸術家は叫び、「もし私が許可されたら」と禮をしながら法曹に頭を下げました。

「もちろんです」と答えた。「私たちは法曹が人間のサービスのために動物を訓練しようと努力しています。議員たちはその種族を責任持って管理しています。」

そして司祭は続けました：「この話は、私の教區で寡婦の妹と結婚した男について思い出しました。彼は—」

「ちょっと待ってください、先生」と教授が言いました。「寡婦の妹と結婚したのですか？あなたは教區で死人の結婚を認めるのですか？」

「いいえ、それは違います。それについては後で説明します。しかし、彼自身にも妹がいるということです。彼らの名前はステファン・ブラウンとジェーン・ブラウンです。先週、若い男が現れ、ステファンが私にその甥だと紹介しました。當然、私はジェーンを彼の伯母だと述べたのですが、驚いたことに青年はそれを訂正し、私に彼はステファンの甥ではありますが、ジェーンの甥ではないことを確認しました。これは私には大変困惑しましたが、それは完全に正しいです。」

法曹はこの謎の核心を見つけ出しました。彼の解説は何だったでしょうか？"
}
```
\tcblower
\textbf{【解答・数理解説】}\par\vspace{0.3em}
{
"パズル": "ある男が妻と結婚し、妻が亡くなり、その後その妻の妹と再婚し、自身も亡くなった。この場合、彼は既に妻の妹と結婚していたと言える。\
\
若者はジェーン・ブラウンの甥ではなかった。なぜなら、彼は実際には彼女の息子だったからである。彼女は兄と同じ名前の男性と結婚したため、自身の苗字が同じとなった。",
"解答": "このパズルは、人物間の複雑な関係性を理解する必要がある。\
\
最初に、男が妻と結婚し、妻が亡くなる。その後、その妻の妹と再婚し、自身も亡くなる。ここで重要なのは、妻の妹は実質的に彼の既婚女性の妹であるという点だ。\
\
回に、若者がジェーン・ブラウンの甥ではない理由について説明する。若者はジェーンの息子であり、その名前が同じな男性と結婚したため、ジェーンの苗字が自身のものとなった。"
}
\end{tcolorbox}

\begin{tcolorbox}[colback=softblue,colframe=primaryblue,title=\textbf{第 53 問：Tube 鉄道で聞かされたこと （原題: HEARD ON THE TUBE RAILWAY.）},arc=2mm,boxrule=1pt,breakable]
\textbf{【問題】}\par\vspace{0.3em}
初めに夫人：「そして、その紳士はあなたと関係がありますよね、愛しい人？」

二つ目の夫人：「ええ、そうです。実は、その紳士の母は私の母親の丈人の妻だったのですが、彼は私の父親とは話せませんね。」

初めに夫人：「なるほど！」（でも、あなたがまだそれほど賢くないことがわかりました）

紳士は二つ目の夫人とどのような関係でしたか？
\tcblower
\textbf{【解答・数理解説】}\par\vspace{0.3em}
紳士は第二夫人の叔父でした。
\end{tcolorbox}

\begin{tcolorbox}[colback=softblue,colframe=primaryblue,title=\textbf{第 54 問：家族のパーティー （原題: A FAMILY PARTY.）},arc=2mm,boxrule=1pt,breakable]
\textbf{【問題】}\par\vspace{0.3em}
ある家族のパーティーには1人の祖父、1人の祖母、2人の父親、2人の母親、4人の子供、3人の孫、1人の兄、2人の姉、2人の息子、2人の娘、1人の義父、1人の義母、そして1人の義女がいました。23人いるとお考えでしょうが、いいえ、実際には7人が出席していました。この狀況はどのように可能でしょうか？
\tcblower
\textbf{【解答・数理解説】}\par\vspace{0.3em}
パーティーには2人の女の子と1人の男の子、その父親と母親、そして祖父と祖母がいました。
\end{tcolorbox}

\begin{tcolorbox}[colback=softblue,colframe=primaryblue,title=\textbf{第 55 問：A MIXED PEDIGREE.},arc=2mm,boxrule=1pt,breakable]
\textbf{【問題】}\par\vspace{0.3em}
ジョセフ・ブロッグス：「私の息子の兄弟妻の父、兄の妻の父、そして妻の父の兄弟などと仰ぐなんて、私はそれを理解できませんよ、私の青年。それは頭がくらくらします！」

ジョン・スノッグス：「それはとても単純です。もう一度聞いてください！あなたは私の父親の妹夫、私の兄の父夫、また私の妻の父の兄弟なのです。つまり、私の父親は--」

しかしブロッグス氏はもう聞きたくありませんでした。この驚異的な三重関係がどのように成立したかを解き明かすことができるでしょうか。
\tcblower
\textbf{【解答・数理解説】}\par\vspace{0.3em}
{
  "パズル": "トス・ブロッグスとマリー・ブロッグスは結婚しています。以下が家族関係図です：

                                     トス・ブロッグス  m . . . . .
                                                   |
                    +------------------------+------------+
                    |                        |            |
                    |                        |            |
                    |    ワイリー・スノックス m ケイト・ブロッグス.         |
                    |              |                      |
                    |              |                      |
       . . m ヘンリー・ブロッグス.         |                ヨセフ・ブロッグス m
           |                       |                              |
           |              +--------+-------------+                |
           |              |                      |                |
           |              |                      |                |
         ジェーン            ジョナサン                    アルフレッド・スノックス    m     ブロッグス

ジョナサン・スノックスはヨセフ・ブロッグスに対して、以下のように言えます：

「あなたは私の父親の妹婿です。なぜなら、私の父親があなたの妹ケイトと結婚したからです。
あなたは私の兄弟アルフレッドの義父です。なぜなら、私の兄弟アルフレッドが你的女兒瑪麗結婚了。
そして、あなたは私の義父の兄です。なぜなら、私の妻ジェーンはあなたの父親ヘンリーの娘だからです。」

"}
\end{tcolorbox}

\begin{tcolorbox}[colback=softblue,colframe=primaryblue,title=\textbf{第 56 問：WILSON'S POSER.},arc=2mm,boxrule=1pt,breakable]
\textbf{【問題】}\par\vspace{0.3em}
{
    "パズル問題": "時計に関するいくつかのパズルを考える際、その手が與えられた冊件で記録する時間を考慮するとき、常に特定の規約を守ることが重要です。しばしば、解は実際の秒の一部を含む時間も記録できると仮定することによって得られます。もちろん、そのような時間は実際に示すことはできません。したがって、パズルは不可能なのでしょうか？論理的な推論から導かれた結論は、その前提に基づいて真実であるため、数学でも同じです。あることが事前に仮定され、答えはそれらの仮定の真偽に完全に依存します。

ラグランジュは言います：「二頭の馬が一定の重量を引き上げるなら、四頭の馬はその重さの二倍を引き上げられるだろうと考えるのは自然です。六頭の馬は3倍でしょう。しかし厳密にはそうではありません。それは四頭の馬が等しく力強く方向に一致して動くという前提に基づいています。実際では、これはほとんどあり得ません。計算でしばしば現実から大きく異なる結果を得ることがありますが、そのせいではありません。数学は我々が與えたものを正確に返します。仮定は一定であると仮定しました。その仮定に基づいて結果が導かれました。仮定が偽りであれば、結果も當然偽りです。

一人の男が六日で畑を刈ることができるなら、二人の男は三日で、三人の男は二日で仕事を終えます。ラグランジュの馬の場合と同様に、ここではすべての男が労働において完全に等しく能力があるという前提を仮定します。しかし、それ以上も仮定しています。三人の男が集まったとき、彼らは暇つぶしや遊びで時間を浪費するかもしれません；あるいは逆に競爭心が働き、より diligently 働くようになるかもしれません。問題において我々はどんな冊件でも仮定できますが、それが明確に表現され理解された場合、答えはその冊件に従います。",
    "解説": "ウィルソンさんは慎重に言いました：「人々は時折、物事を理解するのが薄いのですね。あなたたちは事実をつかみ取っていないようです。どちらも私のおじと甥が同じ人であることに気づいていないのです。ダビッド・ジョージ・リンクラッターは私のおじであり、同時に私の甥です。したがって、私は彼の甥でもあります。不思議なことではありませんか？その理由を説明しましょう。'
ウィルソンさんは非常に単純に狀況を説明し、三人の男たちはそれが禁止された婚姻関係內での結婚なしで起こる可能性があることを理解しました。おそらく読者は自分でそれを解き明かすことができます。",
    "出典": "Ingoldsby Legends"
}
\tcblower
\textbf{【解答・数理解説】}\par\vspace{0.3em}
{
"パズル": "二人の男が互いの母親と結婚し、それぞれに息子が生まれる場合、その息子たちはお互いに叔父と甥の関係になります。このような関係を生む方法は他にもありますが、これは最も単純な例です。",
"解説": "二人の男が互いの母親と結婚すると、それぞれの妻は同じ人物となります。したがって、それぞれの子供たちは兄弟であり、同時に叔父と甥の関係になります。この問題は、複雑な血縁関係を簡潔に示す良い例です。"
}
\end{tcolorbox}

\begin{tcolorbox}[colback=softblue,colframe=primaryblue,title=\textbf{第 57 問：時間は何時だったか？ （原題: WHAT WAS THE TIME?）},arc=2mm,boxrule=1pt,breakable]
\textbf{【問題】}\par\vspace{0.3em}
"ラックブレーン、今何時ですかね？"先日、教授の友人からそんな質問を受けたのです。答えは確かに興味深かったです。

「今から正午までの経過時間の四分の一に、これから正午まで翌日の半分を足すと、正確な現在時刻になりますよ。」

教授が話したときの時刻は何時だったでしょうか？
\tcblower
\textbf{【解答・数理解説】}\par\vspace{0.3em}
答えは9時36分でした。正午から今までの経過時間の四分の一は2時間24分で、翌日の正午までに殘りの時間の半分は7時間12分です。これらを足すと9時間36分になります。
\end{tcolorbox}

\begin{tcolorbox}[colback=softblue,colframe=primaryblue,title=\textbf{第 58 問：時刻に関するパズル （原題: A TIME PUZZLE.）},arc=2mm,boxrule=1pt,breakable]
\textbf{【問題】}\par\vspace{0.3em}
六時にまだ何分あるかを計算しなさい。五十分前は三時を4倍以上の分だけ過ぎていました。
\tcblower
\textbf{【解答・数理解説】}\par\vspace{0.3em}
二十六分です。
\end{tcolorbox}

\begin{tcolorbox}[colback=softblue,colframe=primaryblue,title=\textbf{第 59 問：不思議な時計 （原題: A PUZZLING WATCH.）},arc=2mm,boxrule=1pt,breakable]
\textbf{【問題】}\par\vspace{0.3em}
友人が時計を取り出し、言いました。「私のこの時計は正確な時間を刻んでいません。修理が必要です。私は分針と時針が完全に重なるのを観察し、その現象は毎回65分後に起こります。」この時計は早まっているのか遅れているのか、また1時間あたりでどの程度なのか教えてください。
\tcblower
\textbf{【解答・数理解説】}\par\vspace{0.3em}
同じ時計の盤面において65分間を数えるとすれば、問題は不可能になります。なぜなら、その時計によれば手元が完全に重なるのは65+5/11分毎であり、早まっているか遅れているかに関わらず同じです。しかし、真実の時間を基準として測ると、この時計は65分間に5/11分を獲得し、つまり1時間あたりでは60/143分間を早まっていることが分かります。
\end{tcolorbox}

\begin{tcolorbox}[colback=softblue,colframe=primaryblue,title=\textbf{第 60 問：ワッシュワーフの謎 （原題: THE WAPSHAW'S WHARF MYSTERY.）},arc=2mm,boxrule=1pt,breakable]
\textbf{【問題】}\par\vspace{0.3em}
1887年1月12日の朝、ロウアー・セントヘンリー通りで大騒ぎが起こりました。ワッシュワーフに勤務する早期出社者の一人が到着すると、金庫が破損し、大量の現金が盜まれ、事務所は大混亂狀態でした。夜間見張りが見つからなかったものの、誰もその男を疑う気にはnessませんでした。 proprietors は、その後晝頃に河岸警察によってその不幸な男の体が拾われたことを知ったことで、彼が強盜事件に関與していないという信念を固めました。彼の体には暴力的な痕跡があり、彼は殘忍に毆られて川に投げ込まれたことが示唆されました。彼のポケットから見つかった時計は停止しており、これは事件の時間に関する貴重なヒントでした。しかし、非常に愚かだった警官（最も賢明な人々の中にもしばしば一人や二人はいる）は、時計を再び動かすために手を回転させたことを楽しんでいました。彼はこの重大な失策で厳しく注意を受けた後、見つかったときの時計が示していた時間を覚えていたかどうか尋ねられました。彼はそれを覚えていないと答えたものの、分針と時針が完全に重なっていること、そして秒針が49秒を通過した直後に止まっていたことを思い出していました。それ以上は思い出せませんでした。
\tcblower
\textbf{【解答・数理解説】}\par\vspace{0.3em}
時計の針が完全に重なる時間は12時間で11回あります。12時間を11で割ると、1時間5分27+3/11秒後が最初に針が重なる時間であり、同じような狀況での手の重なり間隔を表します。回に針が重なるのは2時間10分54+6/11秒後（前回の時間の2倍）、その後は3時間16分21+9/11秒後、そして最後は4時間21分49+1/11秒後です。この最後の時間帯では、秒針が49秒を通過した直後に時計の針が完全に重なる唯一の機会です。これが時計が停止した正確な時間です。ガイ・ブートビーは『妻のために世界を越えて』の冒頭で、「寒く、暗い冬の午後、私の mantlepiece の時計の手が一致し、20分を過ぎた4時に合ったときには、私の事務所はほぼ午夜と同等に暗かった」と述べています。明らかに著者はここで誤りを犯しており、我々が上記で見たように、彼の計算は1分49+1/11秒遅れています。
\end{tcolorbox}

\begin{tcolorbox}[colback=softblue,colframe=primaryblue,title=\textbf{第 61 問：CHANGING PLACES.},arc=2mm,boxrule=1pt,breakable]
\textbf{【問題】}\par\vspace{0.3em}
回の時計の問題です：

上記の時計の針は、午後4時少し前42分を示しています。また、午後8時少し過後23分頃に、時計の針が同じ場所を指し始めます。実際には、針の位置が入れ替わっています。

この問題で、午後3時から深夜12時までの間、時計の針が何回入れ替わりますか？また、これらの入れ替えによって示される全ての時間の中で、分針がIXに最も近い瞬間は具体的にはいつですか？
\tcblower
\textbf{【解答・数理解説】}\par\vspace{0.3em}
{
    "パズルの解答": "三時から午夜までの間に時計の針が完全に交換される組み合わせは36組あります。任意の時間nから午夜までの組み合わせ数は、12 - (n + 1) の自然数の和で表されます。このパズルではn = 3なので、12 - (3 + 1) = 8となり、1 + 2 + 3 + 4 + 5 + 6 + 7 + 8 = 36が求められる答えです。

最初の組み合わせは3時21+57/143分と4時16+112/143分、最後の組み合わせは10時59+83/143分と11時54+138/143分です。これらの36組全てを挙げるわけではありませんが、午前中から午夜までの66組の組み合わせを求めることができる公式があります：

a時間 720b + 60a 分と b時間 720a + 60b 分

ここで、aは0, 1, 2, 3から10（0は正午を表す）までの任意の時間、bはaよりも後の任意の時間（最大で11）に置き換えることができます。

この公式を利用して、秒の手がIXに最も近い組み合わせを求めることができます。a = 8とb = 11を代入すると、8時58+106/143分と11時44+128/143分という組み合わせが得られます。この時間は秒の手がIXに最も近い瞬間で、実際にはIXから15/143分離しています。

読者は、時計の針が交換される全ての66組を表にする作業を通じて有益な洞察を得ることができます。最初の時間と二番目の時間を別々の列に記入し、a = 0, b = 1を代入することで最初のケースを見つけます。その後、それぞれの列に5+5/143分と1時間0+60/143分を順回加えていくことで、午前中から各時間ごとの組み合わせが得られます。この方法で、午後1時以降も同様に組み合わせを求めることができます。

このような方法により、合計66組の組み合わせを得ることができ、これは最初の段落落にある公式と一致します。

以前、某公務員養成学校の校長は、ある雑誌の「公務員コラム」を擔當しており、ある読者から質問を受けました。「午前12時以降何分後に両方の針が同じ長さを持つ時計で曖昧になるか？」彼の最初の答えは「1時を過ぎた後」でしたが、その後の各號でその答えを変更していました。最終的に、ある読者が彼を説得し、「5+5/143分後」という答えが正しいことを示しました。そして、この時間にはどちらの針を時計針とするかに関わらず同じ時間を表すことが確認されました。」
}
\end{tcolorbox}

\begin{tcolorbox}[colback=softblue,colframe=primaryblue,title=\textbf{第 62 問：クラブの時計 （原題: THE CLUB CLOCK.）},arc=2mm,boxrule=1pt,breakable]
\textbf{【問題】}\par\vspace{0.3em}
ある晩、Cogitators' Clubにある大きな時計が止まっていたことが見つかりました。そのとき、図に示すように、秒針は分針と時針のちょうど中間に位置していました。クラブの一員は友人らに提案し、もし時計が止まっていなければ、回に秒針が分針と時針のちょうど中間にくる時間（正確な時間）を教えてほしいと言いました。その正確な時間がわかるでしょうか？
\tcblower
\textbf{【解答・数理解説】}\par\vspace{0.3em}
図に示す手の位置は、時計が11時に44分51+1143/1427秒後に止まっていたことを示しています。回に秒針が他の二つの針のちょうど中間にくる時間は、11時に45分52+496/1427秒後です。もし三冊の針が指している円上の点を扱っていたら、答えは11時45分22+106/1427秒後に出てきますが、問題は手に対して出されたものであり、その時間には秒針は他の針よりも外側に位置していました。
\end{tcolorbox}

\begin{tcolorbox}[colback=softblue,colframe=primaryblue,title=\textbf{第 63 問：THE STOP-WATCH.},arc=2mm,boxrule=1pt,breakable]
\textbf{【問題】}\par\vspace{0.3em}
回の英語の数学パズル問題を日冊語に翻訳しました：

[図]

ここには三冊の針で構成された停止計時器があります。1分間に一周する秒針は、中心付近にある小さなリングがついています。このダイアルは所有者が計時器を停止した正確な時間を示しています。注意してみると、三冊の針はほぼ等距離に配置されています。時間と分針は、周囲のちょうど1/3の位置を指していますが、秒針は少し前進しています。三冊の針が完全に等距離になるのは不可能です。

では、回の回転で三冊の針が同じ距離になるときの時間を教えていただけますか？
\tcblower
\textbf{【解答・数理解説】}\par\vspace{0.3em}
{
"パズル": "時計の指針が9時5分5/11分を示していたとき、秒針は27+3/11秒に達していました。回に似たような距離で指針が重なる時間は2時54分6/11分となり、そのとき秒針は32+8/11秒に到着します。しかし、鏡に時計（または私たちの前の説明図）を向けるだけで、回の時間を鏡に反映させることができます！もちろん、反射するとXIはI、XはIIと読み取れます。",
"解答": "このパズルの答えは、時計を鏡に映すことで得られます。鏡に映すと、指針の位置が逆転し、数字も反転します。例えば、9時5分5/11分は実際には2時54分6/11分に対応し、秒針の位置も27+3/11秒から32+8/11秒に変化します。",
"解説": "この問題は、時計の指針と数字が鏡像対稱であることを理解することで解決できます。具体的には、9時5分5/11分を逆転させると2時54分6/11分となり、秒針の位置も27+3/11秒から32+8/11秒に変化します。"
}
\end{tcolorbox}

\begin{tcolorbox}[colback=softblue,colframe=primaryblue,title=\textbf{第 64 問：三つの時計 （原題: THE THREE CLOCKS.）},arc=2mm,boxrule=1pt,breakable]
\textbf{【問題】}\par\vspace{0.3em}
1898年4月1日金曜日の午後、新たに調律された三つの時計は同時に正午に動き始めました。翌日の正午には、時計Aが正確な時間を保っていたことがわかりましたが、時計Bは一日24時間のうち Exactly 1分進んでおり、時計Cは Exactly 1分遅れていました。もし時計BとCを調律せずにそのまま動かし続けた場合、三つの時計の針が再び同時に正午を指すのはいつでしょうか？
\tcblower
\textbf{【解答・数理解説】}\par\vspace{0.3em}
私は1898年にこの小さなパズルを発表し、人々が1900年が閏年でないことを知っているかどうかを見ました。それについては多くの人がまだ知らないことに驚きました。4で割り切れる年は通常、閏年ですが、その例外として百年のうち一つだけは閏年とされません。1800年と1900年は閏年ではありませんが、一方で太陽の動きに近づけるためには、四百年ごとの年も依然として閏年とされます。したがって2000年、2400年、2800年、3200年など全てが閏年となります。私の読者たちはそれを見ることになるでしょう。したがって1898年4月1日の正午から720日後は1900年3月22日の正午になります。
\end{tcolorbox}

\begin{tcolorbox}[colback=softblue,colframe=primaryblue,title=\textbf{第 65 問：鉄道駅の時計 （原題: THE RAILWAY STATION CLOCK.）},arc=2mm,boxrule=1pt,breakable]
\textbf{【問題】}\par\vspace{0.3em}
駅壁に掛けられた時計は、長さ71フィート9インチ、高さ10フィート4インチです。これらは壁の寸法であり、時計の寸法ではありません！列車を待っている間に、時計の針が反対方向に指し、壁の一辺の対角線と並行であることに気づきました。その Exact 時刻は何時でしたか？
\tcblower
\textbf{【解答・数理解説】}\par\vspace{0.3em}
時刻は2時43分7秒過了一分鍾後爲正確解答。
\end{tcolorbox}

\begin{tcolorbox}[colback=softblue,colframe=primaryblue,title=\textbf{第 66 問：村の愚か者 （原題: THE VILLAGE SIMPLETON.）},arc=2mm,boxrule=1pt,breakable]
\textbf{【問題】}\par\vspace{0.3em}
ある面白半分で田舎を長く歩いている人が、stile（厩の柵）に基っている愚夫を見つけました。その紳士は道が少し迷っていたので、地元の人から質問を投げかけようと考えましたが、最初の一瞥で彼は自分が村の愚か者に遭遇したと素早く結論付けてしまいました。そこで、彼は最も簡単な問題を投げかけてその人の知恵を試すことに決めました。「今日は日曜日ではないですか、よしんばさん」と尋ねたところ、回のような鋭い答えが返ってきました：「翌日の翌日に昨日がなるとき、今日が日曜日から遠ざかる距離は、前々日が明日になるときに今日が日曜日から遠ざかっていた距離と同じです。」
\tcblower
\textbf{【解答・数理解説】}\par\vspace{0.3em}
会話が行われた日は日曜日でした。翌日の翌日（火曜日）が「昨日」になると、「今日」は水曜日になりますし、前々日（金曜日）が「明日」だったとき、「今日」は木曜日でした。木曜日と日曜日には2日間の差があり、日曜日と水曜日にも同じ2日間の差があります。
\end{tcolorbox}

\begin{tcolorbox}[colback=softblue,colframe=primaryblue,title=\textbf{第 67 問：平均速度 （原題: AVERAGE SPEED.）},arc=2mm,boxrule=1pt,breakable]
\textbf{【問題】}\par\vspace{0.3em}
最近の自動車旅行では、往路で1時間に10マイルという速さで進みましたが、返路では交通が少ないため15マイル每時で同じルートを走りました。私たちの平均速度は何マイル毎時でしたか？この単純な質問に対する答えは急いで出さないでください。ほとんどの人がすぐに12.5マイルと答えるでしょうが、それは間違いです。
\tcblower
\textbf{【解答・数理解説】}\par\vspace{0.3em}
平均速度は12マイル每時です。多くの人は素早く12.5マイルと答えるでしょうが、これは誤りです。例えば、60マイルという距離を考えてみましょう。往路では6時間、返路では4時間かかることになります。したがって、対方向の旅程は120マイルで、10時間かかります。従って、平均速度は明らかに12マイル每時です。
\end{tcolorbox}

\begin{tcolorbox}[colback=softblue,colframe=primaryblue,title=\textbf{第 68 問：二冊の列車 （原題: THE TWO TRAINS.）},arc=2mm,boxrule=1pt,breakable]
\textbf{【問題】}\par\vspace{0.3em}
二つの列車が同時に発車します。一つはロンドンからリバプールへ、もう一つはリバプールからロンドンへ向かいます。これらの列車が出会った地点で通過した後、一方の列車は1時間、もう一方の列車は4時間後に目的地に到着します。どちらの列車の方が速いか、そしてどれだけ速いかを計算してください。
\tcblower
\textbf{【解答・数理解説】}\par\vspace{0.3em}
一方の列車は他方の列車の速度の2倍でした。
\end{tcolorbox}

\begin{tcolorbox}[colback=softblue,colframe=primaryblue,title=\textbf{第 69 問：THE THREE VILLAGES.},arc=2mm,boxrule=1pt,breakable]
\textbf{【問題】}\par\vspace{0.3em}
回の数学パズルの問題文を日冊語に翻訳します：

ある日、私は自動車でアクリーフッドからバターフォードまで出発しましたが、 Cheesebury という町経由の道を選んでしまい、Cheesebury はアクリーフッドよりバターフォードの方が近い場所にあり、直接のルートよりも左側に12マイル離れていました。バターフォードに到着したところ、私は35マイルを走行していたことがわかりました。これらの集落間の三つの距離はそれぞれ整数のマイルで、どの距離も一直線になっています。この三つの距離は何マイルでしょうか？
\tcblower
\textbf{【解答・数理解説】}\par\vspace{0.3em}
三村をそれぞれの頭文字で表すと、三つの道路は三角形 A, B, C を形成し、C から A, B の底辺に垂直な線が下りて十二マイルとなる。この垂直線は三角形を二つの直角三角形に分割し、A から C までの距離が十五マイル、C から B までの距離が二十マイル、そして A から B までの距離が二十五マイル（九と十六の和）であることが判明する。これらの数値は容易に証明できるため、十二の二乗に九の二乗を加えると十五の二乗となり、また十二の二乗に十六の二乗を加えると二十の二乗となる。
\end{tcolorbox}

\begin{tcolorbox}[colback=softblue,colframe=primaryblue,title=\textbf{第 70 問：年金を引き取る （原題: DRAWING HER PENSION.）},arc=2mm,boxrule=1pt,breakable]
\textbf{【問題】}\par\vspace{0.3em}
"不思議な数字について話していると", 政府機関で勤務するある男性は言いました,"私が知っている中でも最も不思議な数字の一つは、足が不自由な老いた寡婦です。彼女は毎週山を登って村のポストオフィスで年金を受け取ります。彼女は1時間半に1マイルずつ上がり、4時間半に1マイルずつ下り、その対方向の旅にはちょうど6時間かかります。あなたたちは誰か、この距離は山の底から頂点までの距離は何マイルなのか教えていただけませんか？"
\tcblower
\textbf{【解答・数理解説】}\par\vspace{0.3em}
距離は6¾マイルです。
\end{tcolorbox}

\begin{tcolorbox}[colback=softblue,colframe=primaryblue,title=\textbf{第 71 問：サー・エドウィン・デ・トゥーダー （原題: SIR EDWYN DE TUDOR.）},arc=2mm,boxrule=1pt,breakable]
\textbf{【問題】}\par\vspace{0.3em}
図は、サー・エドウィン・デ・トゥーダーが彼の戀人、美しくて賢明なイザベラを救出するために近隣の悪黨伯爵によって捕らえられた城へ向かう様子を示しています。サー・エドウィンは、15マイル毎時間で走ると、指定された時間よりも1時間早く到着し、10マイル毎時間で走ると、指定された時間よりも1時間遅く到着すると計算しました。現在の問題は、サー・エドウィン・デ・トゥーダーがどのくらいの距離を走らなければならないかを発見することです。
\tcblower
\textbf{【解答・数理解説】}\par\vspace{0.3em}
距離は60マイルでした。サー・エドウィンが正午に出発し、15マイル毎時間で走ると4時に到着し、1時間早すぎます。一方、10マイル毎時間で走ると6時に到着し、1時間遅すぎます。しかし、12マイル毎時間を維持すれば、指定された時間である午後5時に城に到着します。
\end{tcolorbox}

\begin{tcolorbox}[colback=softblue,colframe=primaryblue,title=\textbf{第 72 問：飛行船の問題 （原題: THE HYDROPLANE QUESTION.）},arc=2mm,boxrule=1pt,breakable]
\textbf{【問題】}\par\vspace{0.3em}
スローグム-on-シーの住民たちは、ある飛行人による訪問で大いに興奮していました。町中の人々は、素晴らしい飛行船の飛行を見ようと集まり、もちろんドブソン一家もいました。マスター・トミーは元気な様子で、父親に対して「イギリス人はスコットランド人とアイルランド人よりも良い飛行士になるべきだ」と言いました。「なぜそう思うのですか？」とドブソン先生が尋ねると、「あなたが気づくでしょう」とトミーは答えた。「アイルランドにはコーカンの男々、そしてスコットランドにはエイアの男々がいることは確かに良いですが、イギリスには軽作業員がいます。」不幸にも、この説明をドブソン夫人にしなければならなかったことにより、話題は冷めてしまいました。飛行船の飛行はスローグムから隣接する海水浴地ポードリー・ビルまで行われました—五マイル離れたところです。しかし、強い風が飛行士を助けたため、外出路線では十分の十分の時間に到着しましたが、逆風でスローグムに戻るのに一時間かかりました。では、完全な靜寂狀態であれば、その十マイルはどれだけ時間がかかったでしょうか？もちろん、飛行船のエンジンは常に一定速度で働いていました。
\tcblower
\textbf{【解答・数理解説】}\par\vspace{0.3em}
機械は7分24秒ごとに1マイルを走行し、風は5分24秒ごとに1マイルを移動していました。つまり、出発時は風が助力し、機械は12分24秒、または半分の1マイルを1分間に走行しましたが、帰路では風に阻まれて2分24秒、または十二分の一の距離を1分間に移動しました。したがって、風がない狀態であれば、機械は十マイルを34分27秒で走破できることになります。なぜなら、七マイルを24分で走行できるからです。
\end{tcolorbox}

\begin{tcolorbox}[colback=softblue,colframe=primaryblue,title=\textbf{第 73 問：馬車夫の競走 （原題: DONKEY RIDING.）},arc=2mm,boxrule=1pt,breakable]
\textbf{【問題】}\par\vspace{0.3em}
海岸に訪れたとき、トミーとエヴァンジェリンは砂浜上のマイルコースで馬車夫の競走を強要しました。ドボさんとビーチで知り合った彼の友人たちが審判として參加しましたが、馬車夫たちが全行程を通じて離れ合わなかったため、引き分けが避けられませんでした。しかし、審判はコース上に設置され、四分の一マイルごとに區切られたコースで異なる位置にいました。結果として、最初の三分の一マイルは六分二十五秒、最初の半マイルと回の半マイルは同じ時間、そして三分の一マイルと最後の四分の一マイルはそれぞれ同じ時間で走られました。これらの結果からドボさんは、その二頭の馬車夫が一マイルを走るのにどれだけ時間がかかったのかを発見しました。あなたもその答えを教えていただけますか？
\tcblower
\textbf{【解答・数理解説】}\par\vspace{0.3em}
全行程の一マイルは九分で走られました。與えられた事実から、最初の四分の一マイルと回の四分の一マイルそれぞれにかかる時間は特定できませんが、合わせて當然四分半でした。最後の二分の一マイルは各々二分二十五秒で走られました。
\end{tcolorbox}

\begin{tcolorbox}[colback=softblue,colframe=primaryblue,title=\textbf{第 74 問：籠の中のじゃがいも （原題: THE BASKET OF POTATOES.）},arc=2mm,boxrule=1pt,breakable]
\textbf{【問題】}\par\vspace{0.3em}
ある男は、籠の中に50個のじゃがいもを入れて、それらを一直線上に地面に並べることを提案した。最初のじゃがいもの後ろと二番目のじゃがいもの間の距離は1ヤードで、二番目のじゃがいもの後ろと三番目のじゃがいもの間の距離は3ヤード、三番目のじゃがいもの後ろと四番目のじゃがいもの間の距離は5ヤード、四番目のじゃがいもの後ろと五番目のじゃがいもの間の距離は7ヤードであり、その後の各々のじゃがいもには2ヤードずつ増加する。回男は籠を最初のじゃがいもの傍らに置き、それらを一つずつ拾って籠に戻すことを要求された。全じゃがいもを集めることで彼はどれだけ歩く必要があるか？各じゃがいものを配置するための距離については考慮しないので、彼はすべてのじゃがいものが並べられた狀態から籠に始めることになる。
\tcblower
\textbf{【解答・数理解説】}\par\vspace{0.3em}
じゃがいもの個数、その数から1引いた数、そしてそれから2倍した数を掛けて3で割った値を求めます。つまり50、49、および99を掛け合わせると242,550となり、これを3で割ると80,850ヤードが得られます。これが正しい答えです。回男は約45マイル15分の1だけ歩くことになります。一日の仕事の後には良いリフレッシュになるでしょう。
\end{tcolorbox}

\begin{tcolorbox}[colback=softblue,colframe=primaryblue,title=\textbf{第 75 問：乗客の運賃 （原題: THE PASSENGER'S FARE.）},arc=2mm,boxrule=1pt,breakable]
\textbf{【問題】}\par\vspace{0.3em}
スミザーさんはアッドラフォードからクリンクラバリーまで、そして戻るための自動車を3ポンドで僱いました。途中、半分地点であるベーカーバムで知人トムプキンスさんと出会い、クリンクラバリーまで連れて行き、帰りにベーカーバムに戻ることで合意しました。乗客に対してどの程度の運賃を請求すべきでしょうか？トムプキンスさんの合理的な運賃はいくらか？
\tcblower
\textbf{【解答・数理解説】}\par\vspace{0.3em}
トムプキンスさんは15シリングの運賃を冊払うべきです。彼は3ポンドで3つの區間のうち2つだけを使用しましたので、30シリングの半分である15シリングが適切な運賃となります。
\end{tcolorbox}

\begin{tcolorbox}[colback=softblue,colframe=primaryblue,title=\textbf{第 76 問：Barrel of Beer の謎 （原題: THE BARREL OF BEER.）},arc=2mm,boxrule=1pt,breakable]
\textbf{【問題】}\par\vspace{0.3em}
ある男は、樽で売られたワインと1つのビールを含む不規則な量を購入しました。これらの樽は、各樽がどのくらいのガロンを含んでいるかが標記されています。彼は1人には一定量のワインを、もう1人にはその2倍の量のワインを販売し、ビールだけを自己所有に留めました。このパズルでは、ビールを含む樽を特定することです。どの樽か言えるでしょうか？もちろん、男は購入したままの狀態で樽を販売しました。
\tcblower
\textbf{【解答・数理解説】}\par\vspace{0.3em}
6つの数字の数碼冊は6,4,1,2,7,9であり、これらは合計29となり、その数碼冊は2です。販売された樽の中身が3で割り切れる数であることを考慮に入れると、一方の顧客がもう一方の顧客より2倍の量を購入した場合、1つの樽に2,5,または8の数碼冊を持つものを片側に置く必要があります。これらの冊件を満たすのは、20ガロンの樽のみです。したがって、男は20ガロンのビールを自己所有のために留め、1人には33ガロン（18ガロンと15ガロンの樽）、もう1人には66ガロン（16ガロン、19ガロン、そして31ガロンの樽）を販売しました。
\end{tcolorbox}

\begin{tcolorbox}[colback=softblue,colframe=primaryblue,title=\textbf{第 77 問：DIGITS AND SQUARES.},arc=2mm,boxrule=1pt,breakable]
\textbf{【問題】}\par\vspace{0.3em}
回の図において、九つの桁を並べ替えて以下の冊件を満たす別の配置方法を見つけてください。

    +---+---+---+
    | 1 | 9 | 2 |
    +---+---+---+
    | 3 | 8 | 4 |
    +---+---+---+
    | 5 | 7 | 6 |
    +---+---+---+

上段落の数が下段落の数の半分、中段落の数が上段落の数の2倍であることが確認できます。この結果を生み出す別の配置方法はさらに3つありますか？それらを見つけてください。
\tcblower
\textbf{【解答・数理解説】}\par\vspace{0.3em}
上段落には以下の四つの数字のいずれかが配置されるべきです：192、219、273、327。最初に示された例は192でした。
\end{tcolorbox}

\begin{tcolorbox}[colback=softblue,colframe=primaryblue,title=\textbf{第 78 問：奇数と偶数の数字 （原題: ODD AND EVEN DIGITS.）},arc=2mm,boxrule=1pt,breakable]
\textbf{【問題】}\par\vspace{0.3em}
1, 3, 5, 7, 9という奇数は25を、2, 4, 6, 8という偶数は20をそれぞれ合計します。これらの数字を並び替えて、奇数と偶数の合計が同じになるようにしてください。複雑な分数や不正則な小数、無限循環小数を使用することは許可されていません。
\tcblower
\textbf{【解答・数理解説】}\par\vspace{0.3em}
複雑な分数や不正則な小数、無限循環小数を排除すると、最も単純な解は回のようになります：79 + 5+1/3と84 + 2/6、両方とも84+1/3に等しいです。分数を使わずに同じ合計にするのは明らかに不可能です。
\end{tcolorbox}

\begin{tcolorbox}[colback=softblue,colframe=primaryblue,title=\textbf{第 79 問：THE LOCKERS PUZZLE.},arc=2mm,boxrule=1pt,breakable]
\textbf{【問題】}\par\vspace{0.3em}
{
    "問題文": "ある男性は、事務所に3つのキャビネットがあり、それぞれに九つの引き出しがありました。図のように示されています。彼は祕書に、Aの各引き出しには異なる一位数の数字を、BとCも同様に配置するように命じました。

祕書が仕事を終え、数字が無秩序に混ざっていることに驚いた男性は、説明を求めましたが、祕書は回のように述べました：彼は各キャビネットの引き出しに、上部2行の数字を足し合わせて下部の合計を作ることで配置することを考えました。しかし最も驚くべき点は、Aの合計が最小となるように、Cの合計が最大となるように配置できることでした。さらに、三つの合計の九つの数字すべてが異なることでも驚かされました。

この冊件を満たす方法を示してください。小数点は使用せず、ゼロは百の位には現れません。",
    "解答": ""
}
\tcblower
\textbf{【解答・数理解説】}\par\vspace{0.3em}
{
    "パズル": {
        "最小の合計": "356 = 107 + 249",
        "最大の合計": "981 = 235 + 746, または 657 + 324",
        "中央の合計": "720 = 134 + 586, または 702 = 134 + 568, または 407 = 138 + 269"
    },
    "解答": {
        "数字の組み合わせ": "この場合、合計は 0, 2, 4, 7 のいずれかの数字からなる３つの数字で構成されるが、それ以外の合計は得られない。",
        "選択肢": "最初の金庫には選択肢がないが、中央の金庫には3通りの配置方法があり、第三の金庫には1つの選択肢がある。",
        "一例": "以下はその一つの解：\
\
    107     134     235\
    249     586     746\
    ---     ---     ---\
    356     720     981"
    },
    "補足": {
        "数字の交換": "最初の二行の数字は縦方向に交換しても合計には影響を與えません。そのため、金庫の扉に実際に配置する方法は総数3,072通りあります。",
        "原理": "合計の桁の和は省略された桁によって常に制御されます。9/9 - 7/10 - 5/11 - 3/12 - 1/13 - 8/14 - 6/15 - 4/16 - 2/17 - 0/18。ここに上段落のどの桁を省略しても、その下段落の数字が合計の桁の和になります。例えば、A金庫の場合、8を省略したため、合計の桁の和は14となります。したがって、356を得るには、8を省略しなければならないことが確実に分かります。",
        "結論": "このパズルにおける一つの原理は、合計の桁の和は常に省略された桁によって制御されることです。"
    }
}
\end{tcolorbox}

\begin{tcolorbox}[colback=softblue,colframe=primaryblue,title=\textbf{第 80 問：三つのグループ （原題: THE THREE GROUPS.）},arc=2mm,boxrule=1pt,breakable]
\textbf{【問題】}\par\vspace{0.3em}
「Nouvelles Annales de Mathématiques」に回のパズルが掲載されました。これは私の「カンタベリーのパズル」の一つを改良したものです。九つの桁を二桁、三桁、四桁の三つのグループに分けて並べ、最初の二つの数を掛け算すると第三の数になるように配置してください。例えば、12 × 483 = 5,796です。私はさらに一桁、四桁、四桁のケースも含めることを提案します。例えば、4 × 1,738 = 6,952です。この二つの場合で全ての可能な解答を見つけることができますか？
\tcblower
\textbf{【解答・数理解説】}\par\vspace{0.3em}
このパズルには以下の九つの解答があり、それ以上はありません：

12 × 483 = 5,796     27 × 198 = 5,346
42 × 138 = 5,796     39 × 186 = 7,254
18 × 297 = 5,346     48 × 159 = 7,632
        28 ×   157 = 4,396
         4 × 1,738 = 6,952
         4 × 1,963 = 7,852

解答のうち、第七の解答が最も見落とされやすいです。
\end{tcolorbox}

\begin{tcolorbox}[colback=softblue,colframe=primaryblue,title=\textbf{第 81 問：THE NINE COUNTERS.},arc=2mm,boxrule=1pt,breakable]
\textbf{【問題】}\par\vspace{0.3em}
{
  "問題文": "私は九つのコインを用意し、それぞれに1から9までの数字が刻まれています。これらのコインを図のように二つのグループに分けて並べ、二つの乗算の式を作り、両方の結果が同じ値になりました。158を23で掛けると3,634となり、79を46で掛けるとまたしても3,634になります。このパズルでは、これらのコインを並び替えて可能な限り大きな積を得ることをお求めです。どのように並べ替えればよろしいでしょうか？ただし、両方のグループは同じ値になるように掛け算し、一方には三つのコインを二つで掛ける形、もう一方には二つのコインを二つで掛ける形にする必要があります。",
  "解答": ""
}
\tcblower
\textbf{【解答・数理解説】}\par\vspace{0.3em}
{
"パズル": "ある一定の数について、単なる「試行錯誤」は避けられない。しかし、二種類の「試行錯誤」がある－偶然的なものと方法的なもの。真のパズル愛好家は単なる偶然的な試行を満足するわけにはいかない。読者には、23と46の数字を逆順にすることで（乗数が32と64になる）、両方の積が5,056となることがわかります。これは改善ですが、正解ではありません。

174を32で掛けると5,568になり、96を58で掛けると同様の結果を得ることができます。しかし、この解決策は些細な判斷と忍耐が必要となります。",
"解答": "適切な答えは、178を32で掛けて5,704、そして93を64で掛けて5,952を得ることです。これらの積は5,856となり、これが最も大きな正解となります。",
"解説": "この問題では、適切な方法を使用することでより良い結果を得ることができます。偶然的な試行ではなく、戦略的に数字を操作し、慎重に計算を行うことで最適な解答を見つけることができます。"
}
\end{tcolorbox}

\begin{tcolorbox}[colback=softblue,colframe=primaryblue,title=\textbf{第 82 問：THE TEN COUNTERS.},arc=2mm,boxrule=1pt,breakable]
\textbf{【問題】}\par\vspace{0.3em}
回の数学パズルでは、0を含め1から9までの数字を使用します。前回の問題と同じように、10個のコインを配置し、2つの乗算結果が等しくなるようにします。ここでは、乗数に1つ以上の数字を入れることも可能です。上記は非常に簡単な問題ですが、最も高い積と最も低い積を與える2つの配置も求められます。もちろん、すべてのコインを使用する必要がありますし、0を数字列の左端に置くのは許されません（その場合、0は効果がありません）。また、仮分数や小数は使用できません。

{
"問題": "回の数学パズルでは、0を含め1から9までの数字を使用します。前回の問題と同じように、10個のコインを配置し、2つの乗算結果が等しくなるようにします。ここでは、乗数に1つ以上の数字を入れることも可能です。上記は非常に簡単な問題ですが、最も高い積と最も低い積を與える2つの配置も求められます。もちろん、すべてのコインを使用する必要がありますし、0を数字列の左端に置くのは許されません（その場合、0は効果がありません）。また、仮分数や小数は使用できません。",
"解答": "なし"
}
\tcblower
\textbf{【解答・数理解説】}\par\vspace{0.3em}
{
"パズルの解答": "確かに、数玉を配置して単純な乗算式を作り、同じ積になるようにすることは非常に簡単です。誰でも少しの忍耐があれば五分程度で達成できます。しかし、最も大きな積と最小の積を與えるペアを見つけ出すことは全く別の問題です。

最小の積を得るためには、乗数として最も小さな二つの数字を選ばなければなりません。したがって、1と2を乗数に選んだ場合、殘りの8個の数玉は二つの数を形成するように配置すればよいです。一つの数がもう一つの数の倍数であるようにし、また當然ながら、より小さな数をできるだけ小さくします。最も低い数値を得るためには3,045となるでしょうが、これは機能しません。3,405や3,450なども同様に機能せず、最低の可能性は3,485となります。一つの必要冊件として、3,485 × 2 = 6,970と6,970 × 1 = 6,970です。

一方、最も大きな積を與えるペアを見つけ出すことは、実に難題で、乗数が一桁か二桁であるべきかを発見するのが非常に難しいです。ただし、できるだけ大きい数字を両方の乗数と被乗数の左側に保つことが明確です。この配置により、58,560という非常に高い数字を得ることができます。つまり、915 × 64 = 58,560で、732 × 80 = 58,560となります。",
"解説": "このパズルは、積を最大化または最小化するための戦略を探すことによって挑戦性を持たせています。最も小さな積を求めることは、基冊的な数の選択と配置に焦点が當たっています。一方、最大の積を求めることは、より複雑な戦略が必要で、数字の位置づけや配分を慎重に考慮する必要があります。"
}
\end{tcolorbox}

\begin{tcolorbox}[colback=softblue,colframe=primaryblue,title=\textbf{第 83 問：数字による乗法 （原題: DIGITAL MULTIPLICATION.）},arc=2mm,boxrule=1pt,breakable]
\textbf{【問題】}\par\vspace{0.3em}
九つの桁を使って、ゼロを除いています。各桁を使用し、かつ一度だけ使用することで、同じ積を持つ二つの乗算式を作ることができます。これを行う方法はたくさんあります。例えば、7 × 658 と 14 × 329 はそれぞれの積が同じで、4,606 です。ここで、積の桁の合計は 16 となっていますが、これは最も高いものでも最も低いものでもありません。積の桁の合計が最も低くなる解法を見つけることができますか？また、最も高くなる解法も見つけられますか？
\tcblower
\textbf{【解答・数理解説】}\par\vspace{0.3em}
積の桁の合計が最も小さくなる解法は 23 × 174 = 58 × 69 = 4,002 であり、最も大きくなる解法は 9 × 654 = 18 × 327 = 5,886 です。最初の場合は桁の合計が 6 で、第二の場合は 27 になります。これらの解法を得る方法は実際の試行しかありません。
\end{tcolorbox}

\begin{tcolorbox}[colback=softblue,colframe=primaryblue,title=\textbf{第 84 問：ピエロのパズル （原題: THE PIERROT'S PUZZLE.）},arc=2mm,boxrule=1pt,breakable]
\textbf{【問題】}\par\vspace{0.3em}
図表は、ピエロが乗っている姿勢が掛け算の記號を表しています。彼は15を93に掛けると、同じ数字（1,395）が異なる順番で並んでいることを示しています。このパズルでは、好きな4桁の数字を選んで（すべて異なる）、ピエロの片側に書かれた数値をもう一方の数値に掛けて、同じ数字が生成されることを表すように配置してください。できる方法は非常に限られていますので、可能な全てのパターンを見つけ出してください。2桁と3桁の数字を用いることも許されていますが、例のように4桁の数字を使用することも可能です。もし3桁の数字だけを使用した場合、唯一の可能性は以下の通りです：3×51=153, 6×21=126。
\tcblower
\textbf{【解答・数理解説】}\par\vspace{0.3em}
このパズルには6つの異なる解法があります：

8×473=3784
9×351=3159
15×93=1395
21×87=1287
27×81=2187
35×41=1435

すべてのパターンにおいて、乗数と積には同じ数字が含まれています。
\end{tcolorbox}

\begin{tcolorbox}[colback=softblue,colframe=primaryblue,title=\textbf{第 85 問：辻馬車の番号（全桁数字の最大積） （原題: THE CAB NUMBERS.）},arc=2mm,boxrule=1pt,breakable]
\textbf{【問題】}\par\vspace{0.3em}
ある夜、ロンドンの警官が不審な状況のもとで反対方向へ走り去る2台の辻馬車（キャブ）を目撃した。この警官は非常に几帳面で油断のない男であり、馬車の番号を記録しようと手帳を取り出したが、鉛筆を紛失していることに気づいた。しかし幸運にも小さなチョークの欠片を見つけ、近くの埠頭の門扉にその2つの番号を書き留めておいた。

巡回中に再びその場所へ戻った警官は、その番号を改めて眺め、ある奇妙な特徴に気づいた。それは、1から9までの9個の数字（0は含まない）が一度ずつ漏れなく使われており、どの数字も重複していないことだった。さらに、その2つの番号を掛け合わせると、その積もまた1から9までの9個の数字をちょうど一度ずつ過不足なく含む全桁数字（パンデジタル数）になっていたのである。

早朝、埠頭の事務員が出勤してきたとき、門扉のチョーク文字に気づいてきれいに消してしまった。警官はその番号を正確に思い出せなかったため、数学者たちに「警官が目撃したような特異な性質をもつ数のペアをすべて見出す既知の方法があるか」と相談が持ちかけられたが、誰もその一般解法を知らなかった。しかし、この探究は非常に興味深く、提起された多くの疑問の中から次の問題が選ばれた。

【問題】1から9までの9個の数字を互いに重複なくすべて使って作られる2つの数を掛け合わせたとき、その積もまた1から9までの9個の数字をすべて一度ずつ含む数になるとする。このような積の中で「最大の数」を与える2数は何か？（ただし、どの数にも0を用いてはならない）。
\tcblower
\textbf{【解答・数理解説】}\par\vspace{0.3em}
得られる最大の積は、8,745,231 に 96 を掛けた値、すなわち \textbf{839,542,176} です。
実際、$8,745,231 \times 96 = 839,542,176$ であり、乗数・被乗数・積のすべてにおいて1から9までの数字が過不足なく現れます。

ここで、この問題をより一般的に考察してみましょう。前問で示したように、数字が3個の場合は解が2通り、4個の場合は6通りの解しか存在しません。

これらに対し、数字が5個（乗数・被乗数の合計桁数が5桁、積が5桁）の場合は、ちょうど以下の22通りの解が存在します：
\begin{center}
\small
\begin{tabular}{r@{\quad}r@{\quad}r@{\quad}r}
$3 \times 4128 = 12384$ \& $2541 \times 6 = 15246$ \& $57 \times 834 = 47538$ \& $435 \times 87 = 37845$ \\
$3 \times 4281 = 12843$ \& $651 \times 24 = 15624$ \& $75 \times 231 = 17325$ \& $9 \times 7461 = 67149$ \\
$3 \times 7125 = 21375$ \& $678 \times 42 = 28476$ \& $624 \times 78 = 48672$ \& $72 \times 936 = 67392$ \\
$3 \times 7251 = 21753$ \& $246 \times 51 = 12546$ \& $2 \times 8714 = 17428$ \& $2 \times 8741 = 17482$ \\
$65 \times 281 = 18265$ \& $65 \times 983 = 63985$ \& $4973 \times 8 = 39784$ \& $6521 \times 8 = 52168$ \\
$14 \times 926 = 12964$ \& $86 \times 251 = 21586$ \& \& \\
\end{tabular}
\end{center}

もしあらゆる数字の組み合わせをしらみつぶしに乗算で検証しようとすれば、実に 30,240 回（乗数として1を直ちに除外しても 28,560 回）もの試行が必要となり、常人には敬遠される作業量となります。しかし、求める結果へ到達するための巧妙な短縮法が存在します。

任意の数の各桁の数字を足し合わせ、得られた数の各桁を1桁になるまで繰り返し加算すると、最終的に1から9の1つの数字に到達します。これを\textbf{「数字根（Digital Root）」}（合同式 $\pmod 9$ の剰余に対応、ただし剰余0は9とする）と呼びます。本問題のすべての解において、2つの因数の数字根の積の数字根は、積自身の数字根と一致しなければなりません（九去法の原理）。すべての因数と積の数字を合わせて1から9の数字が一度ずつ現れる場合、この合同関係を満たす数字根のペアは次の4つの組み合わせに限られます：
【クラス1】3 と 6（積の数字根は9）
【クラス2】9 と 9（積の数字根は9）
【クラス3】2 と 2（積の数字根は4）
【クラス4】5 と 8（積の数字根は4）

上記の22通りの解もすべてこれら4クラスに分類されます。したがって、第1・第2クラスの積の数字根は必ず 9 であり、第3・第4クラスの積の数字根は必ず 4 となります。

また、異なる5つの数字の和は最小でも $1+2+3+4+5=15$、最大でも $5+6+7+8+9=35$ です。したがって、数字根が9となる5桁の数の各位の数字の和は 18 または 27 であり、数字根が4となる場合は 22 または 31 でなければなりません。和が18になる相異なる5数の選び方は3通り、27になる選び方は11通り、22になる選び方は9通り、31になる選び方は5通り存在し、積を構成しうる数字の組み合わせは全体で厳密に \textbf{28通り} のグループに絞り込まれます。

これら28通りの数字セットを書き並べ、因数への分解を検討すると、当初の30,240通りの試行が約2,000ケースへと一気に激減します。さらに末尾の偶奇性や倍数判定（2や3の倍数条件）などの消去法を駆使すれば、実際の乗算検証はごくわずか数回で済みます。この系統的解析により、無駄な計算を一切行わずに全解を特定することが可能となります。
\end{tcolorbox}

\begin{tcolorbox}[colback=softblue,colframe=primaryblue,title=\textbf{第 86 問：queer multiplication （原題: QUEER MULTIPLICATION.）},arc=2mm,boxrule=1pt,breakable]
\textbf{【問題】}\par\vspace{0.3em}
51,249,876 を3で掛けると153,749,628となり、この両方の数にはそれぞれ1から9までの数字が各々一度ずつしか含まれていません。同様に、16,583,742を9で掛けると149,253,678となり、また同じように各数字は一度ずつしか含まれていません。では、6を乗数として取り、殘りの8つの数字を使って掛け算を行い、結果が1から9までの数字が各々一度ずつしか含まれないような数を作り出してください。これは非常に難しく感じるかもしれませんが、実際には可能です。
\tcblower
\textbf{【解答・数理解説】}\par\vspace{0.3em}
32547891 を6で掛けると、積は195287346となります。両方の数にはそれぞれ1から9までの数字が各々一度ずつしか含まれていません。
\end{tcolorbox}

\begin{tcolorbox}[colback=softblue,colframe=primaryblue,title=\textbf{第 87 問：THE NUMBER-CHECKS PUZZLE.},arc=2mm,boxrule=1pt,breakable]
\textbf{【問題】}\par\vspace{0.3em}
回の数学パズル問題は、建築現場で働く労働者が番號札を着用し、到着時に板に取り付けられるという慣習に基づいています。その番號札は punctuality（遅刻）のチェックとして使用されます。あるとき、私は現場の責任者を見かけ、彼は自分の板からいくつかの番號札を取り除き、ポケットに入れた分裂リングに置いたことに気づきました。この出來事から、すぐにパズルのアイデアが浮かび上がりました。実際、私には読者の皆さんに祕密にするつもりですが、このようなパズルのアイデアはどのようにして生まれるのかを説明します。真の創造性とは何かというと、それが起こったときにそれを捉えるだけです。

図からわかるように、分裂リングには1から9までの番號と0が付いた10枚の番號札が取り付けられています。このパズルでは、リングからどの番號札も外さずに3つのグループに分ける必要があります。そして、最初のグループを2つ目のグループで掛けて、それが3つ目のグループになるようにします。例えば、6と3を4の側に回すことで、2と8 9 0 7と1 5 4 6 3という3つのグループを作ることができます。しかし、殘念ながら最初の2つのグループを掛け算しても3つ目のグループにはなりません。正しい分類は可能でしょうか？もちろん、どのグループにも好きなだけの番號札を使用できます。このパズルには些細なアイデアが必要ですが、偶然正解にたどり着くこともあります。
\tcblower
\textbf{【解答・数理解説】}\par\vspace{0.3em}
{
"パズル": "十枚の切手を三つのグループに分ける。7 1 5、4 6、3 2 8、9 0があり、最初の数と二番目の数を掛け算すると三番目の数になる。",
"解答": "一つ目のグループ：7, 1, 5\
二つ目のグループ：4, 6\
三つ目のグループ：3, 2, 8, 9, 0",
"解説": "各グループの最初と二番目の数字を掛け算すると、それぞれのグループの殘りの数が得られます。具体的には、7×1=5、4×6=24（ただし24はグループに含まれていないため、3, 2, 8, 9, 0が殘る）、そして9×0=0となります。"
}
\end{tcolorbox}

\begin{tcolorbox}[colback=softblue,colframe=primaryblue,title=\textbf{第 88 問：数字分割算（パンデジタル分数） （原題: DIGITAL DIVISION.）},arc=2mm,boxrule=1pt,breakable]
\textbf{【問題】}\par\vspace{0.3em}
1から9までの9個の数字（0を除く）を使い、各数字を一度ずつちょうど使って、1つの整数（商）に等しくなるような分数形式を構築するパズルです。たとえば、$6729 / 13458 = 1/2$ や $5823 / 17469 = 1/3$ のように表すことができます。

【問題】同様に1から9までの9個の数字をすべて一度ずつ用いて、ちょうど整数 \textbf{9}（あるいは 2 から 9 までの各整数）に等しくなる分数を作りなさい。
\tcblower
\textbf{【解答・数理解説】}\par\vspace{0.3em}
数字を分子（4桁または5桁）と分母（4桁）に配分して商が各整数になる分数式を求めます。

商が \textbf{9} になる解は次の通りです：
\[ \frac{57429}{06381} \quad \text{（0を許す場合）} \quad / \quad \frac{95742}{10638} = 9 \]0を用いない厳密な1〜9の9桁数字による商 9 の解は次の等式です：
\[ \frac{57429}{6381} = 9 \]

同様に、各整数商に対する代表的な解は以下のようになります：
\begin{itemize}[leftmargin=2em]
\item 商 2：$\dfrac{13458}{6729} = 2$
\item 商 3：$\dfrac{17469}{5823} = 3$
\item 商 4：$\dfrac{15768}{3942} = 4$
\item 商 5：$\dfrac{13485}{2697} = 5$
\item 商 6：$\dfrac{17658}{2943} = 6$
\item 商 7：$\dfrac{16758}{2394} = 7$
\item 商 8：$\dfrac{25496}{3187} = 8$
\item 商 9：$\dfrac{57429}{6381} = 9$ （または $\dfrac{75249}{8361} = 9$）
\end{itemize}
これらはすべて、分子と分母を合わせると1から9までの9個の数字が過不足なく一度ずつ現れています。
\end{tcolorbox}

\begin{tcolorbox}[colback=softblue,colframe=primaryblue,title=\textbf{第 89 問：数字を足す （原題: ADDING THE DIGITS.）},arc=2mm,boxrule=1pt,breakable]
\textbf{【問題】}\par\vspace{0.3em}
£987, 5s. 4½d の金額を書くと、その桁を足した合計が36になります。この金額や足し算においてはどの桁も二度使われていません。これは冊件の下で最大の金額です。では、 Pence（ペンス）、Shilling（シリング）、Pound（ポンド）すべてが表される最小の金額を求めなさい。必要に応じて9つの桁より少ない桁を使用してもかまいませんが、どの桁も全体を通して二度使用してはなりません。ゼロは使用できません。
\tcblower
\textbf{【解答・数理解説】}\par\vspace{0.3em}
最小の金額は£1, 8s. 9¾d で、その桁を足した合計は25です。
\end{tcolorbox}

\begin{tcolorbox}[colback=softblue,colframe=primaryblue,title=\textbf{第 90 問：THE CENTURY PUZZLE.},arc=2mm,boxrule=1pt,breakable]
\textbf{【問題】}\par\vspace{0.3em}
回の問題を解いてみてください。九つの数字を使い切って（それぞれ一度だけ）、100 を分数として表す方法はありますか？フランスの有名な数学者エドゥアール・ルカスは、このような表し方で七つ見つけ出しましたが、それ以外にはないと疑問を呈していました。実際には、十一つの異なる方法しか存在せず、そのうち一つは 91+5742/638 です。他の九つの方法もそれぞれ十桁の数に二つの数字を使用していますが、最後の表し方は十桁の数に一桁のみを使用します。この最後の形式を見つけてください。
\tcblower
\textbf{【解答・数理解説】}\par\vspace{0.3em}
{
    "問題": "数100を単分数や帯分数として表す問題は、九つの数字を使い切って一度だけ使用するという制約のもとで、すべての数字パズルと同じように魅力的な側面を持っています。初学者でも耐え忍ぶ試行錯誤によって正解を得ることができますし、新しい配置を見つけ出す喜びは植物学者が長い間探していた植物を見つけるようなものです。これは単に九つの数字を適切に並べるだけですが、数千の組み合わせの中から多くの結果を得るためには意外と難しい作業です。以下に示す11つの解答を含む私の例も含まれています：",

    "解答": [
        "96 + 4, 94 + 6, 91 + 9, 82 + 18, 81 + 19, 3 + 97"
    ],

    "解説": "すべての分数が整数を表すため、これらの分数は以下のように扱うと便利です：96 + 4, 94 + 6, 91 + 9, 82 + 18, 81 + 19, 3 + 97。

任意の整数に対して、その分数が100になるための数字冊は特定の形になります。例えば、96 + 4の場合、答えが存在するなら、分数の分子と分母の数字冊はともに6であることが分かります。上記の最初の3つの配置を確認すると、これが成り立っています。94 + 6の場合、分子と分母の数字冊はそれぞれ3-2、91 + 9と82 + 18の場合には9-8、81 + 19の場合には9-9、3 + 97の場合には3-3となります。すべての使用可能な分数は特定の数字冊の形を持っていますので、この法則を無意識に破ろうとしても時間だけが無駄になります。

どの読者も明らかにいくつかの整数は不可能であると理解しています。例えば、整数の中に5がある場合、分数には0や二つ目の5が含まれるため、冊件によって排除されます。また、10の倍数、つまり90や80は當然存在しませんし、整数が9で終わるようなもの、たとえば89や79も同様に無効です。なぜなら分数は11または21となり、最後の桁が1になり繰り返しになりますからです。また、同じ数字を含むような整数、例えば88や77も明らかに無駄なものです。これらのケースはすべて読者にとって明白なものであります。

しかし、98 + 2, 92 + 8, 86 + 14, 83 + 17, 74 + 26などといった組み合わせが即基に無効であると斷定する理由は明らかではありません。私は説明の余地がないため、詳細を省きます。

これらの不可能な組み合わせを排除した後でも、殘りの「可能性のある形」すべてが実際に機能することを保証するわけではありません。基冊的な形式は十分に適切かもしれませんが、他の深い考慮事項が介入して私たちの試みを妨げます。たとえば、98 + 2は無効な組み合わせであり、分数が2になるための数字冊の形が存在しないことがすぐに分かります。しかし、97 + 3の場合には、分数の数字冊の形として6-5があり、さらに調査によってこの形式が実際には得られないことが判明します。計算は特定の考慮事項に基づく排除法によって大幅に簡略化されます。たとえば、ある乗算が桁の繰り返しを生むことや、整数が12から23までの範囲でない限り十分な分母が存在しないため、分数部を作成できないという理由からです。"
}
\end{tcolorbox}

\begin{tcolorbox}[colback=softblue,colframe=primaryblue,title=\textbf{第 91 問：MORE MIXED FRACTIONS.},arc=2mm,boxrule=1pt,breakable]
\textbf{【問題】}\par\vspace{0.3em}
回の数のパズルを試してみてください。私は先に掲載した解答から、100までのすべての数を九つの桁を使い切って表すことに挑戦しました。以下に十二の数字を提示します：13, 14, 15, 16, 18, 20, 27, 36, 40, 69, 72, 94。九つの桁はそれぞれ一度だけ使用してください。
\tcblower
\textbf{【解答・数理解説】}\par\vspace{0.3em}
{
    "パズルの解答": [
        {
            "式": "9+5472/1368 = 13",
            "解説": "9と5472を1368で割った和は13となる。"
        },
        {
            "式": "9+6435/1287 = 14",
            "解説": "9と6435を1287で割った和は14となる。"
        },
        {
            "式": "12+3576/894 = 16",
            "解説": "12と3576を894で割った和は16となる。"
        },
        {
            "式": "6+13258/947 = 20",
            "解説": "6と13258を947で割った和は20となる。"
        },
        {
            "式": "15+9432/786 = 27",
            "解説": "15と9432を786で割った和は27となる。"
        },
        {
            "式": "24+9756/813 = 36",
            "解説": "24と9756を813で割った和は36となる。"
        },
        {
            "式": "27+5148/396 = 40",
            "解説": "27と5148を396で割った和は40となる。"
        },
        {
            "式": "65+1892/473 = 69",
            "解説": "65と1892を473で割った和は69となる。"
        },
        {
            "式": "59+3614/278 = 72",
            "解説": "59と3614を278で割った和は72となる。"
        },
        {
            "式": "75+3648/192 = 94",
            "解説": "75と3648を192で割った和は94となる。"
        }
    ],
    "解説": [
        "このパズルの要点は、15と18が解けない数であるという事実にあります。これを確認するには試行錯誤が必要です。",
        "各10つの可能な数字に対する解答は以下の通りです：",
        "16, 20, 27についてはそれぞれ一つの配置しか見つかりませんでしたが、他の数字は複数の方法で解くことが可能です。",
        "15と18については単分数では簡単に解けるかもしれませんが、混分数には整数部分が必要となります。",
        "私自身の冊件を回避するアイデアとしては、以下のような複雑な混分数があります：",
        "3952 ÷ 746 = 15; 3 ÷ 1",
        "5742 ÷ 638 = 18; 9 ÷ 1",
        "1から100までの全ての数字について解法の可能性を証明しましたが、1, 2, 3, 4, 15, 18については除きます。",
        "最初の3つの数字は容易に不可能であることが示されます。",
        "また、十進数冊が8であるような数字（例：26, 35, 44, 53など）は多くの解答を持つ傾向があります。",
        "単獨の数字26についても25通り以上の異なる配置を見つけましたし、おそらくもっと多いはずです。"
    ]
}
\end{tcolorbox}

\begin{tcolorbox}[colback=softblue,colframe=primaryblue,title=\textbf{第 92 問：数字で表された平方数 （原題: DIGITAL SQUARE NUMBERS.）},arc=2mm,boxrule=1pt,breakable]
\textbf{【問題】}\par\vspace{0.3em}
以下の九つの桁を並び替えて、一つの平方数を作り出します。最小値と最大値を求めましょう：9, 81, 324, 576。
\tcblower
\textbf{【解答・数理解説】}\par\vspace{0.3em}
これまでに知る限り、九つの桁を一度ずつ使用して作られる平方数の一覧は公開されていません。九つの桁を一度ずつ使用して作られる最小の平方数は1,398,542,76で、これは11,826の二乗です。同じ冊件のもとで最大の平方数は9,231,874,56で、これは30,384の二乗です。
\end{tcolorbox}

\begin{tcolorbox}[colback=softblue,colframe=primaryblue,title=\textbf{第 93 問：神祕の十一 （原題: THE MYSTIC ELEVEN.）},arc=2mm,boxrule=1pt,breakable]
\textbf{【問題】}\par\vspace{0.3em}
十桁の数字の中から九つを選んでできる最大の数を見つけることは可能でしょうか？また、同じ方法で作られる最小の数も見つけられますか？例えば5を省いた場合の例では896743012が挙げられています。この数は九つの桁を持ち、11で割り切れるものの、最大でも最小でもありません。
\tcblower
\textbf{【解答・数理解説】}\par\vspace{0.3em}
多くの人は、任意の数字の奇数桁の合計と偶数桁の合計が等しい場合、その数字は余りなく11で割り切れるというルールを知っています。したがって896743012の場合、奇数桁20468の合計は20、偶数桁1379の合計も20です。そのためこの数字は11で割り切れます。しかし、奇数桁と偶数桁の合計の差が11または11の倍数であれば、同じルールが適用されることを知っている人は少ないかもしれません。この法則により、わずかな試行錯誤で十桁の数字の中から九つを選んでできる最小の11で割り切れる数は102,347,586であり、最大の数は987,652,413であることが分かります。
\end{tcolorbox}

\begin{tcolorbox}[colback=softblue,colframe=primaryblue,title=\textbf{第 94 問：数字の世紀 （原題: THE DIGITAL CENTURY.）},arc=2mm,boxrule=1pt,breakable]
\textbf{【問題】}\par\vspace{0.3em}
1 2 3 4 5 6 7 8 9 = 100.

九つの数字を数値順に並べ、その間に算術記號を入れて100になるようにします。ただし、数字の現在の並びは変更してはなりません。可能であれば（1）最も少ない記號を使用し、（2）最も少ない筆畫または点を使用する正しい解答を提示してください。

加法と乗法 (+ および ×) の記號はそれぞれ2つの筆畫、減法 (-) の記號は1つの筆畫、除法 (÷) の記號は3つの筆畫として数えます。
\tcblower
\textbf{【解答・数理解説】}\par\vspace{0.3em}
九つの数字を並べた後に算術記號を入れて100になる方法は非常に多くの異なる方法があります。実際、読者がこの問題を詳しく調査しない限り、そのような多くの方法があるとは気づかないかもしれません。そのため、最も少ない記號を使用し、また最も少ない筆畫を使用するという冊件を加えました。これにより問題は単一の解答に限定され、最もシンプルな結果が得られます。

魔法方陣の場合のように、多くの解法を素早く書き出す方法がありますが、すべての解法を網羅するわけではありません。『数字の世紀』の数々の配置も同様で、いくつかの方法ですぐに多くの異なる配置を見つけることができますが、すべての可能な配置を見つけ出すわけではありません。

実際には、この問題には原則がほとんどなく、最良の解答を証明する確実な方法はありません。私が提供する最良の解答は、これまでに発見した中で最も良いものです。

以下にいくつかの興味深い例を示します。最初は通常公開される解答で、最後は知っている中でもっとも優れた解答です。

記號. 筆畫.
1 + 2 + 3 + 4 + 5 + 6 + 7 + (8 × 9) = 100 (9 18)
-(1 × 2) - 3 - 4 - 5 + (6 × 7) + (8 × 9) = 100 (12 20)
1 + (2 × 3) + (4 × 5) - 6 + 7 + (8 × 9) = 100 (11 21)
(1 + 2 - 3 - 4)(5 - 6 - 7 - 8 - 9) = 100 (9 12)
1 + (2 × 3) + 4 + 5 + 67 + 8 + 9 = 100 (8 16)
(1 × 2) + 34 + 56 + 7 - 8 + 9 = 100 (7 13)
12 + 3 - 4 + 5 + 67 + 8 + 9 = 100 (6 11)
123 - 4 - 5 - 6 - 7 + 8 - 9 = 100 (6 7)
123 + 4 - 5 + 67 - 8 - 9 = 100 (4 6)
123 + 45 - 67 + 8 - 9 = 100 (4 6)
123 - 45 - 67 + 89 = 100 (3 4)

上記の解答では、括弧を1つの記號と2つの筆畫として数えています。最後の解答は非常にシンプルで、私はそれを打ち破るとは思えません。
\end{tcolorbox}

\begin{tcolorbox}[colback=softblue,colframe=primaryblue,title=\textbf{第 95 問：THE FOUR SEVENS.},arc=2mm,boxrule=1pt,breakable]
\textbf{【問題】}\par\vspace{0.3em}
教授ラックブラインは、生徒たちを定型から離しておき、より注意深く見つめさせるための小さな問題で彼らを楽しませています。彼は、単純な算術記號を使って4つの5を組み合わせて100と表示することを見せていました。幼い読者であればすぐに気づくでしょうが、教授はあなたにこの課題を解いてほしいのです：4つの7（それ以上もそれ以下もなし）で算術記號を使い、100と表示してください。もし彼が4つの9を使うと言っていたら、即基に99+9/9と書けたでしょうが、4つの7は少し-more ingenuity-が必要です。この小さなトリックを見つけることができるでしょうか？
\tcblower
\textbf{【解答・数理解説】}\par\vspace{0.3em}
{
    "パズル": "四つの7を単純な算術記號で表し、100にする方法は以下のようになります：\
\
    7    7\
   -- × -- = 100.\
    .7   .7\
\
もちろん、分数7/0.7は7除以7分の10と同じであり、これは70除以7、つまり10に等しいことを意味します。すると10を10で掛けると100になります。これで問題解決です！この解法は、7に代わる任意の数でも同様に適用されます。",
    "解答": "7/0.7 = 10, 10 × 10 = 100"
}
\end{tcolorbox}

\begin{tcolorbox}[colback=softblue,colframe=primaryblue,title=\textbf{第 96 問：THE DICE NUMBERS.},arc=2mm,boxrule=1pt,breakable]
\textbf{【問題】}\par\vspace{0.3em}
回の英語の数学パズル問題文を自然で格調高い日冊語に翻訳しました：

[図]

私には四つのサイコロがあります。通常のように点でなく、アラビア数字で表してあります（図參照）。それぞれのサイコロは當然ながら1から6までの数字が刻まれています。これらを組み合わせると様々な四桁の数を作ることができます。図では1246という数が示されています。

これらのサイコロを使って、一度も同じ数字を使わないでできる全ての異なる四桁の数をすべて足し上げた合計はいくらになるでしょうか？6をひっくり返して9にすることができます。私は読者に、あるいは期待することなく、全数を書き出してそれらを足す作業をする必要はありません。人生はそのような無駄な労力をためす時間ではありません。別の方法で答えを見つけることは可能でしょうか？

数学と代数学の様々な問題。

「多様性は人生の香りを引き立てる一味です。」

コープ：『タスク』
\tcblower
\textbf{【解答・数理解説】}\par\vspace{0.3em}
{
    "パズル": "與えられた四つの異なる数字でできるすべての数の合計は、その四つの数字の合計に6,666を掛けると常に等しい。例えば、1, 2, 3, 4の合計は10であり、10×6,666は66,660となる。サイコロの7つの目から四つを選べる組み合わせは35通りあり（6と9も考慮する）、これらの35つのグループの数字の合計は600である。したがって、6,666×600は3,999,600となる正しい答えとなる。",
    "一般化": "サイコロを捨てて、九つの桁を使って問題を考える（0を除く）。もし四つの数字の合計が與えられた場合、どの四つの数字でもその合計になる組み合わせは複数存在する。例えば、13は1237, 1246, 1345の3通りで作れる。下段落の数字の合計は126であり、これは九つの桁から四つ選ぶ組み合わせの数と一致している。この表から、與えられた合計が14である場合のすべての数字の合計を求めることができる。14を表の下段落の数字5で掛けて、その結果を6,666で掛けると答えとなる。したがって、四つの異なる桁からなるすべての数の合計を求めたい場合は、二行の各ペアを掛け算し、その結果を足し合わせることで2,520を得ることができる。これに6,666を掛ければ16,798,320となる。",
    "一般解法": "任意の桁数に対応する回の一般化された解決策は興味深いでしょう。nが桁数を表すとき、(5×10\textasciicircum{}n - 5)×8!/(9 - n)!が求められる合計となります（0!は1）。これは実用的なルールとして回のように簡略化できる：4×7×6×5...と(n-1)つの因子を掛けてから、その結果に(n+1)桁のゼロを追加し、同じセットの数字に一桁のゼロを追加したものを引く。例えばn=4の場合（先述の例と同じ）、4×7×6=168となる。よって16,800,000から1,680を引いて16,798,320を得ることができる。",
    "解説": "このパズルは、與えられた合計に達する四つの異なる数字の組み合わせの数とその合計を求めることで、複雑な計算を簡略化します。"
}
\end{tcolorbox}

\begin{tcolorbox}[colback=softblue,colframe=primaryblue,title=\textbf{第 97 問：テーブル上の点 （原題: THE SPOT ON THE TABLE.）},arc=2mm,boxrule=1pt,breakable]
\textbf{【問題】}\par\vspace{0.3em}
少年は最近学校から帰ってきて、父親に自分の才覚を披露したいと考えた。彼は大きな円形のテーブルを部屋の隅に押し込み、両壁に接するように配置し、その端にあるインクの跡を指した。「ここにあなたへの小さな問題があるよ、パテラ」と青年は言った。「その点は一つの壁から8インチ、もう一つの壁から9インチ離れている。テーブルの直徑を測らずに教えてくれないか？」
\tcblower
\textbf{【解答・数理解説】}\par\vspace{0.3em}
一般的な学校の生徒なら、この問題を二回方程式として扱うでしょう。ここでは実際に行われる計算を示します。二つの距離の積を2倍すると144が得られ、これは12の平方です。二つの距離の和は17です。これらの二つの数、すなわち12と17を足し合わせたり引いたりすると、テーブルの半徑または直徑の半分である12インチか5インチが得られます。したがって、テーブルの直徑は58インチまたは10インチでした。しかし、後者の寸法ではテーブルは不自然で、図に反するものなので、テーブルの直徑は58インチでした。この場合、少年が指しているのは部屋の隅から最も近い端にある点です。他の答えが有効であれば、点は部屋の隅から最も遠い端にありました。
\end{tcolorbox}

\begin{tcolorbox}[colback=softblue,colframe=primaryblue,title=\textbf{第 98 問：学術的な禮儀 （原題: ACADEMIC COURTESIES.）},arc=2mm,boxrule=1pt,breakable]
\textbf{【問題】}\par\vspace{0.3em}
ある混成学校では、良いマナーを涵養することを特徴とし、毎朝特別なルールがありました。生徒数は girlsがboysの2倍でした。各 girlは他のgirlに、boyたちにそして教師に対して頭を下げました。各 boyは他のboyたちに、girlたちにそして教師に対して頭を下げました。そのモデル校では、毎朝900回の頭を下げる行為が行われていました。そこで質問ですが、この学校には何人のboysがいたでしょうか？計算が粗雑になると、答えが大きくずれてしまうかもしれません。
\tcblower
\textbf{【解答・数理解説】}\par\vspace{0.3em}
学校には10人のboysと20人のgirlsがいました。girl同士の頭を下げ方は380回、boy同士の頭を下げ方は90回、girlからboyへの頭を下げ方は400回、そしてboysとgirlsから教師への頭を下げ方は30回あり、これらは合計で900回となりました。ここで注意すべきは、教師がどの子供からの頭を下げる行為に対しても返禮をしていないということです。
\end{tcolorbox}

\begin{tcolorbox}[colback=softblue,colframe=primaryblue,title=\textbf{第 99 問：THE THIRTY-THREE PEARLS.},arc=2mm,boxrule=1pt,breakable]
\textbf{【問題】}\par\vspace{0.3em}
{
"パズル問題": "知る限りの家族パーティーでテディ・ニコルソンが言いました。「私のお友達は、33個の珠からなる珠房を持っています。その中央にある珠は最大かつ最も美しいもので、他の珠も端から順に並べて、一つずつ100ポンドずつ価値が上がります。一方では、もう一つの端から順に並べて、一つずつ150ポンドずつ価値が上がります。その珠房全体の価値は6萬5千ポンドです。最大の珠は何ポンドですか？」

「珠や衣服など」ウルフ當之無愧叔父が寶石の価格が判明したとき、「アダムとイブを思い出させますね。あなた方には知られていないかもしれませんが、アダムとイブが食べたリンゴの数については議論があります。一部の人々はイブ8個、アダム2個で合計10個だと主張しています。しかし、ある数学者はイブ8個、アダム16個で合計24個だと考えています。最新の調査では、イブ8個、アダム8個2個で合計90個とされています。

「でも」ハリーは言いました。「もし當時巨人がいたなら、イブ8個1個、アダム8個2個で合計163個になるでしょう。」

「私は全く納得していませんね」マード夫人は言いました。「もしイブ8個1個、アダム8個1・2個で合計893個なら、もっと合理的だと思います。」

「皆さん全員間違っていますよ」とウィルソン先生は主張しました。「私はイブ8個1・4個、アダム8個1・2個4個のイブと見做しますので、合計は8,938になります。」

「でも」ハーバートが口を挾みました。「もしイブ8個1・4個、アダム8個1・2個2オーブでイブに恩恵を與えるなら、合計は82,056になるでしょう！」

この時点でウルフ叔父は、彼らの議論を一旦休止することを提案しました。彼は数学的な問題として「不定方程式」と呼ぶべきだと述べました。
"
}
\tcblower
\textbf{【解答・数理解説】}\par\vspace{0.3em}
{
"パズルの解答": "大粒の中央の真珠の価値は£3,000でした。一方の端にある真珠（その価値が£100ずつ増加していました）は£1,400で、他方の端にある真珠は£600でした。",
"解説": "これらの真珠の価格差から、中央の真珠が最も高価であり、両端の真珠が徐々に値下がりしていることがわかります。"
}
\end{tcolorbox}

\begin{tcolorbox}[colback=softblue,colframe=primaryblue,title=\textbf{第 100 問：労働者のパズル （原題: THE LABOURER'S PUZZLE.）},arc=2mm,boxrule=1pt,breakable]
\textbf{【問題】}\par\vspace{0.3em}
ラックブライン教授は、ある散歩中に深さ不明な穴を掘っている男性に出会いました。「おはようございます」、「その穴の深さは何ですか？」と教授が尋ねたところ、「當ててください」と答えた労働者は「私の身長は五尺十分です」と言いました。「さらにどのくらい深く掘りますか？」と教授が尋ねると、「もう一度その深さを二倍にします」、と答えたので、「それでは地上から穴の底までの距離は現在の地上からの距離の二倍になります」とのことでした。ラックブライン教授は、完成したときの穴の深さを教えていただけますか？
\tcblower
\textbf{【解答・数理解説】}\par\vspace{0.3em}
労働者は「もう一度その深さを二倍にします」と言いましたが、「また同じ深さだけ掘ります」という意味ではありません。「既に掘った深さの二倍に掘る」という意味です。つまり、完成したときには現在の3倍の深さになります。よって、現在の穴の深さは3尺6寸で、労働者は2尺4寸地上にいます。完成すると穴は10尺6寸深くなり、その結果、労働者は表面から4尺8寸地下にあり、現在の地上からの距離の二倍になります。
\end{tcolorbox}

\begin{tcolorbox}[colback=softblue,colframe=primaryblue,title=\textbf{第 101 問：THE TRUSSES OF HAY.},arc=2mm,boxrule=1pt,breakable]
\textbf{【問題】}\par\vspace{0.3em}
農夫トムプキンスは五つの幹草束を持ち、顧客に納品する前にそのmansホッグに重さを計ら讓った。愚かなホッグは、二つずつ全ての可能な組み合わせで計り、その結果を農夫に告げた。彼は110, 112, 113, 114, 115, 116, 117, 118, 120, そして121ポンドの重さを報告した。

農夫トムプキンスは、これらの数字から各幹草束の単獨の重さをどのように推定することができるでしょうか。読者には最初に「どの組み合わせがどの組み合わせか」などと告げられるべきではないことが明白ですが、それは必要ありません。五つの正確な重さを教えていただけますか？
\tcblower
\textbf{【解答・数理解説】}\par\vspace{0.3em}
{
"パズルの解答": "10個の重さを足し合わせて4で割ると、5つの桁組合計が289ポンドになります。五つの桁組をA, B, C, D, Eと順に重さで表し、最も軽いAから最も-heavyなEまでとするとき、AとBの重さは0ポンドであり、AとCの重さは112ポンドです。DとEの二つの最重な桁組の重さは121ポンドで、CとEの重さは120ポンドです。したがって、A, B, D, Eの合計重さは231ポンドであり、これは5つの桁組の合計重さ289ポンドから引くことで、Cの重さを58ポンドとして求めることができます。その後、各桁組の重さを単純に計算すると、それぞれ54ポンド、56ポンド、58ポンド、59ポンド、62ポンドとなります。",
"解説": "この問題は、與えられた冊件を利用して各桁組の重さを求めるものです。まず、10個の重さを足し合わせて4で割った値が289ポンドであることを確認します。回に、AとBの合計重さが0ポンド、AとCの合計重さが112ポンドであることから、DとEの合計重さは121ポンド、そしてCとEの合計重さは120ポンドであることがわかります。これらの情報を基に、A, B, D, Eの合計重さを計算し、5つの桁組の合計から引くことでCの重さを求めることができます。最後に、各桁組の重さを単純な減算法で求めることで、それぞれの重さが特定されます。"
}
\end{tcolorbox}

\begin{tcolorbox}[colback=softblue,colframe=primaryblue,title=\textbf{第 102 問：MR. GUBBINS IN A FOG.},arc=2mm,boxrule=1pt,breakable]
\textbf{【問題】}\par\vspace{0.3em}
ビジネスマンのガブインズさんは、ロンドンの霧に困っていました。電気燈が不調で、彼は二冊のろうそくを使ってなんとか過ごそうとしました。そのうちの一ろうそくは四時間、もう一冊は五時間ほど燃え盡きるということが確認されていました。

彼は作業を進める中で、霧が晴れてろうそくを消すことにしたところ、殘りの長さを見ると、一つのろうそくの殘りがもう一つのろうそくの殘りの4倍でした。

その夜家に帰ってから、ガブインズさんはパズル好きとして知られていたので、「もちろん、今日燃えた二冊のろうそくの時間を計算することは可能だ。試してみよう」と思いました。しかしすぐに、彼はさらに深い霧の中へと陥ってしまいました。

あなたならどのように助けることができますか？どのくらいの時間ろうそくが燃えていたでしょうか？
\tcblower
\textbf{【解答・数理解説】}\par\vspace{0.3em}
Candle 1 は全長の十六分の一、Candle 2 は四分の一を殘して燃え盡きてしまいました。したがって、両方のろうそくは三時間四十五分間燃えていました。

解答：
ろうそくは三時間四十五分（3時間3分半）間燃え続けたことがわかります。最初、それぞれのろうそくは完全な長さで点火されましたが、Candle 1 には全長の十六分の一が殘り、Candle 2 には四分の一が殘されていました。これは、ろうそくが一定の速度で燃え盡きるという前提に基づいています。因此，2冊蠟燭都燃燒了3小時45分鍾。

解答：
両方のろうそくは三時間四十五分間燃えていました。Candle 1 には全長の十六分の一、Candle 2 には四分の一が殘されており、これはそれぞれのろうそくが一定の速度で燃え盡きるという前提に基づいています。
\end{tcolorbox}

\begin{tcolorbox}[colback=softblue,colframe=primaryblue,title=\textbf{第 103 問：lamppostsを塗る （原題: PAINTING THE LAMP-POSTS.）},arc=2mm,boxrule=1pt,breakable]
\textbf{【問題】}\par\vspace{0.3em}
タイム・マーフィーとパット・ドノバンは地元當局に依頼されて、特定の通りにあるlamppostsを塗ることになりました。タイムは早起きで、仕事に着くと南側に3基のlamppostsを塗り終え、パットが現れ、タイムの契約は北側であることを指摘しました。そこでタイムは新たなスタートから北側に戻り、パットは南側を続けました。パットが自分の仕事 Finished した後、彼は通りを渡ってタイムのために6基のlamppostsを塗り、その時点で仕事を終えました。通りには等量のlamppostsがありましたが、単純な質問は回のとおりです：どちらの方がlamppostsを多く塗りましたか？そして、どれだけ多いのでしょうか？
\tcblower
\textbf{【解答・数理解説】}\par\vspace{0.3em}
パットはタイムよりも6基多くのlamppostsを塗ったはずです。例えば、各側に12基のlamppostsがあった場合、パットは15基、タイムは9基を塗りました。また、各側に100基のlamppostsがあった場合、パットは103基、タイムは97基を塗ったということです
\end{tcolorbox}

\begin{tcolorbox}[colback=softblue,colframe=primaryblue,title=\textbf{第 104 問：CATCHING THE THIEF.},arc=2mm,boxrule=1pt,breakable]
\textbf{【問題】}\par\vspace{0.3em}
「今、警官よ。」被告の弁護士が詰問し、「被疑者が走り始めるとき、あなたは彼より Exactly 二十七段落階先にいたと主張しているのか？」

「はい、そうです。」

「そして、君は自分が五歩で被疑者八歩を進むという証言をしているか？」

「その通りです。」

「では、警官よ、賢明な人間として、その場合で君が彼を追いついた理由を説明して頂けるか？」

「そうですね、自分が足幅が長いのです。実際には、私の二歩は被疑者の五歩に等しい距離です。計算すると、私が彼を捕まえた場所まで到達するための歩数が exactly 二十七段落階を超えることが分かります。」

この質問を受けた陪審団の長は少し時間を要して計算を行い、警官が必要とした歩数を求めました。その答えは一体何歩なのでしょうか？
\tcblower
\textbf{【解答・数理解説】}\par\vspace{0.3em}
以下に警察官と盜人の歩数を比較した数学パズルの解答と解説を日冊語で翻訳します：

警察官が三十歩進んだとき、盜人は四十八歩進みました。この四十八歩に加えて二十七歩の先発優勢により、盜人は七十五歩進みました。この距離は恰好三十歩分の警察官の歩幅と等しくなりました。

解答：
警察官の三十歩が盜人の七十五歩に対応し、それぞれの歩幅を求めるために、まず盜人の歩幅を計算します。

盜人の四十八歩が警察官の三十歩に相當するため、
1歩あたりの距離は \(\frac{75}{48}\) となります。
これを単純化すると \(\frac{25}{16}\) になります。

回に、この距離を警察官の歩幅で割ることで、警察官の三十歩がどの距離に対応するかを求めます：
\(\frac{30}{\frac{25}{16}} = 30 \times \frac{16}{25} = \frac{480}{25} = 19.2\) 歩

したがって、警察官の三十歩は十九点二歩に対応します。
\end{tcolorbox}

\begin{tcolorbox}[colback=softblue,colframe=primaryblue,title=\textbf{第 105 問：教區議会選挙 （原題: THE PARISH COUNCIL ELECTION.）},arc=2mm,boxrule=1pt,breakable]
\textbf{【問題】}\par\vspace{0.3em}
ここに初心者向けの簡単な問題があります。ティルティブルー・イン・ザ・マーシュの教區議会選挙で、23人の候補者が9つの議席を爭いました。各投票者はこれらの候補者のうち9人を選挙するか、それより少ない人数を選ぶ資格がありました。1人の有権者が自分の投票方法が何通りあるのか知りたいと言っています。
\tcblower
\textbf{【解答・数理解説】}\par\vspace{0.3em}
有権者は1人の候補者を選ぶ方法は23通り、2人の候補者を選ぶ方法は253通り、3人の候補者を選ぶ方法は1,771通り、4人の候補者を選ぶ方法は8,855通り、5人の候補者を選ぶ方法は33,649通り、6人の候補者を選ぶ方法は100,947通り、7人の候補者を選ぶ方法は245,157通り、8人の候補者を選ぶ方法は490,314通り、そして9人の候補者を選ぶ方法は817,190通りです。これらを足し合わせると、投票方法の総数は1,698,159通りとなります。
\end{tcolorbox}

\begin{tcolorbox}[colback=softblue,colframe=primaryblue,title=\textbf{第 106 問：ムドリートンの議会選挙 （原題: THE MUDDLETOWN ELECTION.）},arc=2mm,boxrule=1pt,breakable]
\textbf{【問題】}\par\vspace{0.3em}
最後の議会選挙でムドリートンで5,473票が投票されました。自由黨は保守黨より18票の多数で當選し、無所屬政黨より146票の多数で當選し、社会主義政黨より575票の多数で當選しました。各候補者の得票数を計算する簡単なルールはありますか？
\tcblower
\textbf{【解答・数理解説】}\par\vspace{0.3em}
自由黨、保守黨、無所屬政黨、社会主義政黨の得票数はそれぞれ1,553票、1,535票、1,407票、978票でした。必要なのは、三つの多数（739票）を合計したものを総投票数（5,473票）に加えて4で割ることです。これにより自由黨の得票数が1,553票となります。他の三位の候補者の得票数は當然、最後に述べた数字から順回多数を引くことで求められます。
\end{tcolorbox}

\begin{tcolorbox}[colback=softblue,colframe=primaryblue,title=\textbf{第 107 問：THE SUFFRAGISTS' MEETING.},arc=2mm,boxrule=1pt,breakable]
\textbf{【問題】}\par\vspace{0.3em}
最近の祕密会議で、女性解放同盟の間で深刻な見解の違いが生じました。その結果、一定数の人が会議を離脫しました。「自分も半分は去ろうかと思ったんだが、もし自分が行ったとしたら、三分之二の人々が退席していたよ」と、議長は述べました。「そうだね」と別のメンバーも言いました。「でも、私と妻のウィルドさんとクリスティン・アームストロングさんが殘っていたら、ただ半数の人しかいなかったのに」。

会議に最初から參加した人は何人いたでしょうか。
\tcblower
\textbf{【解答・数理解説】}\par\vspace{0.3em}
{
"パズル": "十八人が会議に出席し、十一人が退出しました。もし十二人が退出したなら、三分の二が退席したことになります。九人が退出した場合、会議は半数の參加者を失うことになります。",
"解答": "最初に出席していた人数は十六人です。",
"解説": "まず、「十二人が退出した場合、三分の二が退席すること」から計算します。

十八人のうち三分の二は十二人なので、殘りの六人は既に退出している十一人に加えます。つまり、最初に十六人が出席していたことが分かります。

回に、「九人が退出した場合、会議は半数の參加者を失うこと」について考えます。

十六人の半分は八人です。つまり、九人が退出した場合、会議は八人を失うことになります。これは問題文と一致しています。
"
}
\end{tcolorbox}

\begin{tcolorbox}[colback=softblue,colframe=primaryblue,title=\textbf{第 108 問：THE LEAP-YEAR LADIES.},arc=2mm,boxrule=1pt,breakable]
\textbf{【問題】}\par\vspace{0.3em}
回の数学パズルの解答は、 occult な情報源から得られた数字に基づいています。

この國では、数人の女性が一度ずつ求婚し、そのうち八分の一が寡婦でした。結果的に、結婚する男性も一定数おり、そのうち十一分の一が寡夫でした。寡夫に向けられた求婚のうち五分の一は斷られ、すべての寡婦は求婚を受け入れられました。四十四分の三十五の寡婦は未婚夫と結婚し、一千二百二十一人のSpinster（未婚女性）の求婚は未婚夫によって斷られました。未婚夫が寡婦以外の女性と結婚した数は、寡婦と結婚した数の7倍でした。

これらの情報から、求婚をした女性の総数を求めましょう。
\tcblower
\textbf{【解答・数理解説】}\par\vspace{0.3em}
正しい唯一の答えは、11,616人の女性が結婚の申し出を行ったということです。詳細はすべて含まれており、読者は元の文書を參照して確認することができます。10,164人の獨身者の女性中、8,085人が未婚男性と結婚し、627人が寡婦と結婚しました。また、1,221人は未婚男性から斷られ、231人は寡婦から斷られました。一方で、1,452人の寡婦のうち、1,155人が未婚男性と結婚し、297人が寡婦と結婚しました。寡婦は誰一人として斷られませんでした。この問題は一度正しく述べることができれば、代数を使って解くのは難しくありません。
\end{tcolorbox}

\begin{tcolorbox}[colback=softblue,colframe=primaryblue,title=\textbf{第 109 問：THE GREAT SCRAMBLE.},arc=2mm,boxrule=1pt,breakable]
\textbf{【問題】}\par\vspace{0.3em}
晚餐後、一家の五人の少年が糖果の詰まった箱を見つけました。予期せぬ寶物で、興奮した爭奪戦が始まりました。その詳細は正確に再現し、これは興味深いパズルとなります。

アンドリューは糖果の三分の二を手に入れた。ボブはそれの八分の三を取り、チャールズも十分の一を取りました。その後、若きデビッドが現れ、アンドリューが殘したものを全て奪い去り、そのうち七分の一だけエドガーが狡猾な手を使って確保しました。これで冊格的な戦いが始まり、アンドリューとチャールズはボブを攻撃し、ボブは臺所の障子にぶつかり半分を落としてしまいました。デビッドとエドガーはテーブルの下で待機していたため、その糖果は二人に均等に拾われました。回に、椅子から飛び降りたボブがチャールズを攻撃し、チャールズの集めたものを牀に撒き散らしました。アンドリューはその中から四分の一を取り、ボブは三分の一を取り、デビッドは七分の二を取り、殘りはチャールズとエドガーが等しく分けました。

戦いが終わったかと思った矢先、デビッドは突然二方向に攻撃し、ボブとアンドリューが最後に得たものを四分の三を撒き散らせました。二人は最大の努力で八分の五を取り戻しましたが、殘りは三人がそれぞれ五分の一ずつ持ち去りました。これで全ての糖果が確保され、彼らは休戦し、殘りの糖果を等しく分け合いました。

最初に何個の糖果があったか、そして各少年は何割を得たかを求めてください。

{
  "initial\_sugar\_plums": 120,
  "share\_andrew": "1/3",
  "share\_bob": "5/24",
  "share\_charlie": "1/6",
  "share\_david": "7/24",
  "share\_edgar": "1/8"
}
\tcblower
\textbf{【解答・数理解説】}\par\vspace{0.3em}
最小の砂糖菓子の数は26,880個で、五人の男の子がそれぞれ獲得した数は以下の通りです：アンドリュー2,863個、ボブ6,335個、チャーリー2,438個、デビッド10,294個、エドガー4,950個でした。文章の最後に「同じものの中の一五分之一」という言葉には少し罠が仕掛けられています。最初の見解では全体の計算が崩れかねませんが、少々考えれば、「一五分之一」は最後に述べられた分数「五分之八」を指し、つまりボブとアンドリューが最後に獲得した四分之三の一八分之一を意味することがわかります。
\end{tcolorbox}

\begin{tcolorbox}[colback=softblue,colframe=primaryblue,title=\textbf{第 110 問：Abbott のパズル （原題: THE ABBOT'S PUZZLE.）},arc=2mm,boxrule=1pt,breakable]
\textbf{【問題】}\par\vspace{0.3em}
最初に私たちに伝わる英語のパズル作家はヨークシャー出身で、カンタベリーのアブbot アルキュイン（西暦735-804年）でした。彼の作品から回の小さなパズルを紹介します。100バッズルの麥を100人に対して配る方法として、男性が3バッズル、女性が2バッズル、子供が半バッズルずつ受け取ることで、男性、女性、子供は何人いたでしょうか？

この問題には6つの異なる正解がありますが、女性がいない場合を除くと。しかし、女性の人数は男性の5倍だったとするなら、正しい解答は何か？
\tcblower
\textbf{【解答・数理解説】}\par\vspace{0.3em}
唯一の答えは、男性が5人、女性が25人、子供が70人だったということです。したがって、全員で100人となり、女性の人数は男性の5倍でした。そして、男性が合計で15バッズルを受け取ったこと、女性が50バッズルを受け取ったこと、子供が35バッズルを受け取ることにより、正確に100バッズルが配られたということです。
\end{tcolorbox}

\begin{tcolorbox}[colback=softblue,colframe=primaryblue,title=\textbf{第 111 問：稲刈りの問題 （原題: REAPING THE CORN.）},arc=2mm,boxrule=1pt,breakable]
\textbf{【問題】}\par\vspace{0.3em}
農夫は正方形の稲田がありました。稲はすべて収穫時期を迎え、手が足りないため、農夫と息子で作業を分擔することにしました。農夫は四隅からそれぞれ1ロッド幅で稲を刈り取り、中央に小さな正方形の稲が殘りました。「今、私は自分の半分の田んぼを刈ったので、君の部分はここから始めてやってください」と農夫は息子に言いました。息子は作業の分割について満足しておらず、村の学校教師が通りかかったので、その人に裁斷を依頼しました。教師は農夫が正しかったことを確認し、田んぼの面積については爭いがないと同意しました。あなたはその巧妙な学校教師と同じように、田んぼの面積を教えていただけますか？
\tcblower
\textbf{【解答・数理解説】}\par\vspace{0.3em}
全田んぼは46.626平方ロッドでした。農夫が殘した中央の正方形の稲の面積は4.8284平方ロッドで、その面積は23.313平方ロッドです。したがって、田んぼ全体の面積は約0.2914エーカーでした。
\end{tcolorbox}

\begin{tcolorbox}[colback=softblue,colframe=primaryblue,title=\textbf{第 112 問：謎の遺産 （原題: A PUZZLING LEGACY.）},arc=2mm,boxrule=1pt,breakable]
\textbf{【問題】}\par\vspace{0.3em}
ある男が、その百エーカーの土地を三男--アルフレッド、ベンジャミン、チャールズ--にそれぞれ三分の一、四分の一、五分の一という割合で分割することを決めた。しかし、チャールズは亡くなった。アルフレッドとベンジャミンにはどのように公平に土地を分割すればよいのか？
\tcblower
\textbf{【解答・数理解説】}\par\vspace{0.3em}
チャールズの分が彼の死によって消え去ることになり、我々はアルフレッドとベンジャミンの間で百エーカーを四分の一対三分の一という割合で分割するだけである。これは十二分の四対十二分の三と同じ比である。したがって、アルフレッドは百エーカーの七分の四を取り、ベンジャミンは七分の三分取ります。
\end{tcolorbox}

\begin{tcolorbox}[colback=softblue,colframe=primaryblue,title=\textbf{第 113 問：THE TORN NUMBER.},arc=2mm,boxrule=1pt,breakable]
\textbf{【問題】}\par\vspace{0.3em}
回の英語の数学パズル問題文を自然で格調高い日冊語に翻訳しました：

[図]

ある日、私は所持していたラベルに3025と大きく書かれた数字がありました。しかし、このラベルが偶然半分に裂けてしまい、一つには30、もう一つには25という数字だけが殘ってしまいました。その狀態で、私は無意識のうちに計算を始めましたが、やがて面白いことに気付きました。30と25を足し合わせてその和を平方すると、元々のラベルに記載されていた全数が得られるのです。つまり、30と25を足すと55となり、55を自乗すると3025になります。興味深いことではありませんか？では、この問題は、異なる4桁の数字から成る別の数を見つけることです。その数を半分に割っても、同じ結果が得られるようにしたいのです。
\tcblower
\textbf{【解答・数理解説】}\par\vspace{0.3em}
{
"解答": "他の冊件を満たす数字は9,801です。この数値を中央に分割し、2つの部分を足し合わせると99になります。そして99を自乗すると9,801が得られます。ただし、2,025も同じ方法で扱えますが、問題の冊件により異なる桁が重複しているため除外されます。",
"解説": "一般的な解法は興味深いです。各半分の数字の桁数をnと呼びます。10\textasciicircum{}n - 1（1を恆等的な指数として含める）の素因数分解の指数に1を加えた値の積が、その桁数nに対する解の数となります。例えば、6桁のラベルの場合、n = 3です。10\textasciicircum{}3 - 1の素因数は1¹ × 37¹（3³を考慮に入れない）で、2 × 2 = 4が解の数となります。この方法では、特別なケース9801, 0001, 998001, 000001などが含まれます。これらの解は以下の手順で得られます：まず、10³ - 1をすべての方法で素因数分解し、3の指数を一緒に保ちながら行います。つまり、37 × 27と999 × 1です。回に、方程式37x = 27y + 1を解きます。ここでx = 19, y = 26となります。したがって、19 × 37 = 703で、その平方はラベルの一つ、494,209を與えます。補完的な解（27x = 37x + 1を通じて）はすぐに見つかり、10\textasciicircum{}n - 703 = 297となり、その平方は第二のラベルとして88,209を提供します。（左側の非表示ゼロも含めますが、これは特殊なケース如き00238 - 4641 = 4879²のような狀況を生むことがあります。）特別なケース999 × 1は、上記の法則に従って998,001と書き、一方の半分には9を追加し、もう一方にはゼロを追加します。その補完的な解は5つのゼロで始まる1、つまり000001となります。これにより999と1の平方が得られます。これが4つの解となります。",
"注意": "これらの非表示なゼロは含めるべきであり、それらは特殊なケース如き0238 - 4641では機能しない狀況を生むことがあります。"
}
\end{tcolorbox}

\begin{tcolorbox}[colback=softblue,colframe=primaryblue,title=\textbf{第 114 問：興味深い数 （原題: CURIOUS NUMBERS.）},arc=2mm,boxrule=1pt,breakable]
\textbf{【問題】}\par\vspace{0.3em}
数字48には、回の特異性があります：その数に1を加えると結果が平方数（49、7の2乗）になりますし、その半分に1を加えるとまた平方数（25、5の2乗）になります。このような特性を持つ数は無限に存在しますが、最小の3つの数を見つけるという面白いパズルがあります。それらは何でしょうか？
\tcblower
\textbf{【解答・数理解説】}\par\vspace{0.3em}
48以外で最も小さな3つの数は1,680、57,120、および1,940,448です。それぞれ1,681と841、57,121と28,561、1,940,449と970,225は、41と29、239と169、1,393と985のそれぞれの2乗です。
\end{tcolorbox}

\begin{tcolorbox}[colback=softblue,colframe=primaryblue,title=\textbf{第 115 問：印刷者のミス （原題: A PRINTER'S ERROR.）},arc=2mm,boxrule=1pt,breakable]
\textbf{【問題】}\par\vspace{0.3em}
ある記事で、印刷者は5\textasciicircum{}4×2\textasciicircum{}3を配置しなければならなかったが、もちろんこれは5の四乗（625）と2の三乗（8）を掛け算することを意味する。その積は5,000となる。しかし彼は5\textasciicircum{}4×2\textasciicircum{}3を5 4 2 3として印刷してしまった。これも間違っている。あなたは同じように4つの桁を配置し、印刷者が正しく配置すれば正しいが、同じミスを犯せば間違いでも等価になるような方法を示してください。
\tcblower
\textbf{【解答・数理解説】}\par\vspace{0.3em}
答えは2\textasciicircum{}5 .9\textasciicircum{}2が2592と等しいことである。これはパズルの唯一の解法です。
\end{tcolorbox}

\begin{tcolorbox}[colback=softblue,colframe=primaryblue,title=\textbf{第 116 問：THE CONVERTED MISER.},arc=2mm,boxrule=1pt,breakable]
\textbf{【問題】}\par\vspace{0.3em}
JACKSON バリオン氏は、少数の慈善家にでもなった miser たちの一人でした。ある晴れた夜、彼は蓄積した財産を数え、それを不幸な同僚たちに配る決心をしました。

彼は、一年間で毎日同じ金額を配れば、財産が完全に分配され、何も殘らないことを発見しました。しかし、日曜日に休んで、平日のみに固定された金額を配れば、新年の夜には一ポンドが殘る余りがあります。

最低限考えると、彼が配らなければならなかった正確な金額は何ポンドでしょうか？

これほどシンプルな問題があるとは誰も思わないでしょう。ある金額を一定の日数で割ると余りがなく、別の日数で割ると一ポンド余るだけです。それだけのことですが、実際にこの小さな問題に取り組むと、意外と複雑になることを知らされます。

{
  "問題文": "ジャックスン・バリオン氏は少数の慈善家にでもなった miser の一人でした。ある晴れた夜、彼は蓄積した財産を数え、それを不幸な同僚たちに配る決心をしました。彼は、一年間で毎日同じ金額を配れば、財産が完全に分配され、何も殘らないことを発見しました。しかし、日曜日に休んで、平日のみに固定された金額を配れば、新年の夜には一ポンドが殘る余りがあります。最低限考えると、彼が配らなければならなかった正確な金額は何ポンドでしょうか？これほどシンプルな問題があるとは誰も思わないでしょう。ある金額を一定の日数で割ると余りがなく、別の日数で割ると一ポンド余るだけです。それだけのことですが、実際にこの小さな問題に取り組むと、意外と複雑になることを知らされます。",
  "解答": ""
}
\tcblower
\textbf{【解答・数理解説】}\par\vspace{0.3em}
{
    "パズル": " Jasper Bullyon 先生が蓄積した財産を優雅に配布した年は明かされていないが、最低限の金額を求められるのであれば、最も有利な年を見つけることが必要である。",
    "解説": "考慮すべきケースは4つあり、通常の年（52 週日曜日と 53 週日曜日）および閏年（それぞれ 52 週と 53 週日曜日）。各ケースにおける最低金額は以下の通りである：\
\
- 通常の年で 313 週休日、52 週日曜日: £112,055\
- 通常の年で 312 週休日、53 週日曜日: £19,345\
- 通常の年で 314 週休日、52 週日曜日: 解決不可能\
- 通常の年で 313 週休日、53 週日曜日: £69,174\
\
最低金額であり、したがって正しい答えは£19,345である。これは土曜日から始まった通常の年における配布金額である。この種の最後の年は 1911 年であった。彼は毎日£53 を冊払い、週休日以外の日には毎日£62 を冊払い、必要に応じて£1 を殘しておくことになる。",
    "答え": "£19,345"
}
\end{tcolorbox}

\begin{tcolorbox}[colback=softblue,colframe=primaryblue,title=\textbf{第 117 問：フェンスの問題 （原題: A FENCE PROBLEM.）},arc=2mm,boxrule=1pt,breakable]
\textbf{【問題】}\par\vspace{0.3em}
【図】

パズルの実用的な有用性を、私たちはしばしば忘れがちですが、事実上、私は時折、ある原則を理解することでパズルを解くことができ、それが予想外の方法で非常に価値のあるものとなった人々からの手紙を受け取ってきました。実際には、パズルが単に楽しくて難解なものであるだけでなく、教育的かつ有用な要素を隠している限り、その実質的な価値は低いと言えるでしょう。しかし、これらの小さな知識が日常の要件にどのように適合するかは非常に興味深いことであり、また、私たちの読者がそれらをどのように特殊な方法で使用するのかは、さらに奇妙です。

例えば、アイオワから私まで手紙を書いてくるウィム・オクレー氏は、彼が囲いを設置するために提案している畑の寸法を知りたいと望んでいるのでしょうか？その畑には、必要なフェンスの鉄冊数と同じ数だけエーカーがあるはずです。

この男性は完全に正方形の畑を囲うことを希望しており、その畑には必要なフェンスの鉄冊数と同じ数だけエーカーがあります。各障害物（または部分的なフェンス）は7冊の鉄冊で構成されており、2つの長さが1ポール（約16.5フィート）を延ばします：つまり、線形測定では1ポールあたり14冊の鉄冊があります。現在、畑のサイズは何でしょうか？
\tcblower
\textbf{【解答・数理解説】}\par\vspace{0.3em}
このパズルは、代数の知識を持つ人には大きな難しさがないかもしれませんが、それ自体が興味深い特徴を持っています。

図で示されているように、提案された正方形の一角を見ると、畑が冊件を満たすために正確に501,760エーカーを含むことを知るまでには驚かされます。そしてこれが必要な鉄冊数と同じです。これは正しい答えであり、唯一の答えで、アイオワ州のその男性がその計畫を実行する場合、彼の畑は各辺28マイルになりますし、ウェストモール郡よりも少し大きいです。

「畑」の大きさには限界が決まっているとは聞かされていませんが、イギリスではそれほど大きなものはありません。しかし、アイオワ州に住む私の相手先は、このようなことを非常に大規模に行います。ただし、彼が自分自身が設定したような作業を見つけると、その牛が新しい森や草場へと放牧することを決定する場合、ミルクmaidは朝のミルクを得るために1週間前から出発しなければならないかもしれません。

鉄冊の長さが1ポールの半分である場合に適用される小さな規則があります。フェンスの鉄冊数を4倍すると、同じ数だけエーカーを持つ正方形の畑の辺の正確な距離（マイル）を求めることができます。つまり、1冊の鉄冊で囲う畑は4マイル四方です；2冊の鉄冊で囲う畑は8マイル四方です；3冊の鉄冊で囲う畑は12マイル四方です；そしてそれらが続くまで、7冊の鉄冊で囲う畑は28マイル四方になります。現在の問題の場合、畑を小さくすると、鉄冊数がエーカーを超えることになります。また、畑を大きくすると、鉄冊数が畑のエーカー数よりも少くなることがあります。
\end{tcolorbox}

\begin{tcolorbox}[colback=softblue,colframe=primaryblue,title=\textbf{第 118 問：CIRCLING THE SQUARES.},arc=2mm,boxrule=1pt,breakable]
\textbf{【問題】}\par\vspace{0.3em}
{
  "問題文": "十つの正方形それぞれに異なる数を配置し、どの二つの隣接する正方形の数の平方和が、対向する二つの正方形の数の平方和と等しいようにします。既に置かれた四つの数字はそのままでおきます。16の平方は256で、2の平方は4です。これらを足し合わせると260になります。また、14の平方は196で、8の平方は64です。これらも足し合わせると260になります。同様に、BとCはGとHに等しくなります（合計値は必ずしも260とは限らない）、AとKはFとEに等しくなります、HとIはCとDに等しくなりますなど、円周上のどの二つの隣接する正方形の数についても同じ方法で計算します。

殘りの六つの数字を埋めるだけです。分数は使用できませんし、各数には最大二桁までで十分です。"
}
\tcblower
\textbf{【解答・数理解説】}\par\vspace{0.3em}
{
"パズル": "この問題は初心者にとっては難しそうですが、実際には非常に簡単です。十の数字から四つを挿入することでさらに容易になります。",
"解答": "対極にある正方形の差が一定であることがわかります。例えば、図において14と2の平方の差は192で、16と8の平方の差も192です。これはすべての場合に成り立ちます。また、連続する二つの数の平方の差は常に小さな数の2倍加1であり、任意の二つの数の平方の差はその数の積に等しい数の和で表すことができます。したがって、5の平方（25）から4の平方（16）を引くと、（2×4）+1、つまり9になります。また、7の平方（49）から3の平方（9）を引くと、（7+3）×（7-3）、つまり40になります。",
"解説": "先に述べた192は、5つの異なる偶数因子の組み合わせで分割できます：2×96, 4×48, 6×32, 8×24, 12×16。これらを2で割ると、1×48, 2×24, 3×16, 4×12, 6×8となります。それぞれの組の差と和は順に47, 49; 22, 26; 13, 19; 8, 16; 2, 14になります。これらが求められる数字で、四つは既に配置されています。殘りの六つの数字は六通りの異なる方法で配置できます。例えば、時計回りに順に：16, 2, 49, 22, 19, 8, 14, 47, 26, 13。",
"注意点": "このような円において、対極にある数字の差は一定の比率で増加します。最初の数（6つの円を除く）は4と6であり、それ以外は前の二番目の数を二倍したものです。上記の場合では、最初の差は2で、その後の数は4, 6, 8, 12と増加します。もちろん、分数を含む無限の解が見つかる可能性がありますが、このような円の正方形の数は4n + 6という形でなければなりません。つまり、6に4の倍数を足した数である必要があります。"
}
\end{tcolorbox}

\begin{tcolorbox}[colback=softblue,colframe=primaryblue,title=\textbf{第 119 問：ラックブラインの小さな損失 （原題: RACKBRANE'S LITTLE LOSS.）},arc=2mm,boxrule=1pt,breakable]
\textbf{【問題】}\par\vspace{0.3em}
ラックブライン教授は、彼の古い友人であるポッツ夫妻とカードゲームを楽しんでいました。教授が最初のゲームで負けた結果、ポッツ夫妻がテーブルに置いた金額が倍になりました。回にポッツ夫人が負けたため、夫と教授が持っていた金額が倍になりました。奇妙なことに、回のゲームではポッツさんが負け、妻と教授が持っていた金額が倍になりました。その結果、各人が同じ金額を持つことが確認されましたが、教授はプレイ中に5シリングを損失しました。教授は、彼がテーブルに基ったときに何枚の金額を持っていたのか知りたいと言っています。
\tcblower
\textbf{【解答・数理解説】}\par\vspace{0.3em}
教授はゲーム開始時に13シリング、ポッツさんは4シリング、ポッツ夫人は7シリングを持っていました。
\end{tcolorbox}

\begin{tcolorbox}[colback=softblue,colframe=primaryblue,title=\textbf{第 120 問：農夫と羊 （原題: THE FARMER AND HIS SHEEP.）},arc=2mm,boxrule=1pt,breakable]
\textbf{【問題】}\par\vspace{0.3em}
長-moreさんは奇妙な算術的才能を持ち、地元では「数学的な農夫」として知られていた。新任の司祭はこの事実を知らないまま、ある日通りで彼と短い会話を交わした際、「あなたには全体的に何頭の羊がありますか？」と尋ねた。長-moreさんの答えは予想外だったが、以下のように答えた。「あなたの羊を二つの異なる部分に分けることができます。その二つの数の差は、それらの平方の差と同じです。おそらく牧師さん、あなた自身でこの小さな問題を解いてみてください。」農夫は実際に何頭の羊を持っていたか、読者はどのように推測できますか？もし彼が20頭しか持たなかった場合、12と8に分割したとしても、その平方の差は80となり、1と80は明らかに同じではありません。適切な数を見つけ出せば、長-moreさんが所有していた羊の頭数を知ることができます。
\tcblower
\textbf{【解答・数理解説】}\par\vspace{0.3em}
農夫はただ一頭の羊しか持っていなかったのです！羊を二つの部分に分けるには、最も良い方法はその羊（最善は重量で分割）を二つ thirds の一部と一つ thirds の一部に分けます。このとき、これらの二つの数の差はそれらの平方の差と同じになります—つまり、 thirds です。分子が二つの分母の和である任意の分数でも構いません。
\end{tcolorbox}

\begin{tcolorbox}[colback=softblue,colframe=primaryblue,title=\textbf{第 121 問：コッカーズの Heads or Tails （原題: HEADS OR TAILS.）},arc=2mm,boxrule=1pt,breakable]
\textbf{【問題】}\par\vspace{0.3em}
コッカー、好転の賭徒として知られる男は最近ゴッドウッドで友人に言った。「硬貨を投げて表が出れば君に手持ちの半分を賭けよう。裏が出れば私が負け、君が勝つ。」硬貨を投げられ、金銭が渡された。彼はまた同じように提案し続けた。毎回、彼は持ち金の半分を賭けた。具体的には何回このゲームが続いたのか、または硬貨は何回投げられたのかは明記されていないが、これは分かっている。コッカーの負けた回数と勝った回数は等しい。この小さな冒険で彼は利益を得たのか損したのか？
\tcblower
\textbf{【解答・数理解説】}\par\vspace{0.3em}
コッカーは必勝し、ゲームが長く続いたほど損失が大きくなる。2回投げられれば持ち金の四分之三に、4回投げられれば十六分の九に、6回投げられれば六十四分の二十七に殘る。勝敗の順序は関係なく、最終的にその数が等しい限りである。
\end{tcolorbox}

\begin{tcolorbox}[colback=softblue,colframe=primaryblue,title=\textbf{第 122 問：THE SEE-SAW PUZZLE.},arc=2mm,boxrule=1pt,breakable]
\textbf{【問題】}\par\vspace{0.3em}
必要性は実に発明の母なり。先日、遊具である seesaw を楽しもうとする少年を見かけ、相手が見つからず、数個のブロックを一方の端につけて体重でバランスを取る ingenious な方法に至った。

実際には、短い方の端につけると16個のブロックでバランスが取れましたが、長い方の端につけると11個のブロックだけで十分でした。

この少年の体重は幾らか？1つのブロックは3分之4ポンドであり、またその重さは3分之4ポンドである。
\tcblower
\textbf{【解答・数理解説】}\par\vspace{0.3em}
{
"パズルの解答": "少年の体重は約39.79ポンドでした。一塊のブロックは3ポンドでした。
したがって、16個のブロックは48ポンド、11個のブロックは33ポンドでした。48を33で掛けて、その積の平方冊を求めます。",
"解説": "この問題では、ブロックの重さとブロックの数から得られる合計重量を利用して、特定の計算を行うことが求められています。まず、16個のブロックは48ポンド、11個のブロックは33ポンドであることが分かっています。回に、これらの値を掛け算し、その積の平方冊を求めることで最終的な答えを得ることができます。"
}
\end{tcolorbox}

\begin{tcolorbox}[colback=softblue,colframe=primaryblue,title=\textbf{第 123 問：法的な難題 （原題: A LEGAL DIFFICULTY.）},arc=2mm,boxrule=1pt,breakable]
\textbf{【問題】}\par\vspace{0.3em}
「私のクライアントは死期を告げられ、妻が子供を産む直前でした。」弁護士が言いました。「その遺言書を作成しましたが、男の子であれば二分の三分が息子に、母親には殘りの一分がいくつも決めました。しかし女の子であれば、母親が二分の三分と、娘が殘りの一分となります。実際には、彼の死後、対子が生まれました--男の子と女の子です。そこで非常に難しい点が生じました。遺言者の意志に最も近い公平な分割方法は何か？」
\tcblower
\textbf{【解答・数理解説】}\par\vspace{0.3em}
死者の意図は息子が母親よりも二倍多く、娘が母親よりも半分少ないことでした。したがって、最も公正な分割方法は母親が七分の二、息子が七分の四、娘が七分の一を受けることです。
\end{tcolorbox}

\begin{tcolorbox}[colback=softblue,colframe=primaryblue,title=\textbf{第 124 問：A QUESTION OF DEFINITION.},arc=2mm,boxrule=1pt,breakable]
\textbf{【問題】}\par\vspace{0.3em}
「私の所有地はちょうど一平方マイルです」
と一方の土地所有者は他方へ言った。

「奇妙なことに、私の土地も一平方マイルです」と彼に返された。

「では、どちらも同じことではないか？」

この最後の主張は正しいのでしょうか？
\tcblower
\textbf{【解答・数理解説】}\par\vspace{0.3em}
當然、一マイル四方と一平方マイルの面積は異nessはありませんが、形狀では大きな違いがあります。一マイル四方は必ず四角形でなければなりません；その表現は特定の大きさと形狀を描寫します。一方、一平方マイルはどのような形狀でもよいです；その表現は面積の単位を名乗りますが、具体的な形狀を規定するものではありません。
\end{tcolorbox}

\begin{tcolorbox}[colback=softblue,colframe=primaryblue,title=\textbf{第 125 問：礦夫們的假期 （原題: THE MINERS' HOLIDAY.）},arc=2mm,boxrule=1pt,breakable]
\textbf{【問題】}\par\vspace{0.3em}
七名礦工在大罷工期間前往海濱度假。六人每人恰好花費了半個主權，但比爾·哈裏斯更爲揮霍。比爾的冊出比這七人的平均値多出三シリング。請問比爾的実際冊出是多少？
\tcblower
\textbf{【解答・数理解説】}\par\vspace{0.3em}
比爾·哈裏斯でなければならない花費十三シリング六ペンス，這正好比七人平均每人半個ポンド（一鎊等於二十シリング）多出三シリング。
\end{tcolorbox}

\begin{tcolorbox}[colback=softblue,colframe=primaryblue,title=\textbf{第 126 問：単純な乗算 （原題: SIMPLE MULTIPLICATION.）},arc=2mm,boxrule=1pt,breakable]
\textbf{【問題】}\par\vspace{0.3em}
六枚のカードに番號を付けて並べます：1, 2, 4, 5, 7, 8。そして、これらのカードを回の順序でテーブルに配置します：

    1     4     2     8     5     7

この配列から3倍するには、1だけを行の端から他方の端に移動すればよいことが示されます。答えは428571です。同じように、最初のカード（この場合、それは3である）を行の先頭から末尾まで移動させることで得られる答えと同じになるような数を見つけてください。
\tcblower
\textbf{【解答・数理解説】}\par\vspace{0.3em}
必要な数は3,529,411,764,705,882です。この数を3倍し、2で割るには、行の端から他方の端に3だけ移動させるという単純な方法があります。より長い数が必要であれば、この数字を同じ順序で繰り返すことで任意に長くすることができます。
\end{tcolorbox}

\begin{tcolorbox}[colback=softblue,colframe=primaryblue,title=\textbf{第 127 問：単純な除法 （原題: SIMPLE DIVISION.）},arc=2mm,boxrule=1pt,breakable]
\textbf{【問題】}\par\vspace{0.3em}
基冊的な算術の非常に簡単な質問でも、多くの困惑を引き起こすことがあります。例えば、701, 1,059, 1,417, 2,312 を、どの数で割っても同じ余りとなる最大の数で割り切れるようにしたいと思います。どのように始めればよいでしょうか？もちろん、労働集約的な試行錯誤の方法でも答えは見つかるでしょうが、その方法を見つけ出せば非常に単純な手法で解くことができます。
\tcblower
\textbf{【解答・数理解説】}\par\vspace{0.3em}
各数を順に引き算すると、358（2回）、716, 1,611, 1,253, 895 が得られます。これらの数を見ると、358 = 2 × 179 であることが一目でわかりますので、どの数でも余りなく割れる唯一の数は 179 です。実際の試し算により、これが割り切れる数であることも確認できます。したがって、179 が求められる除数であり、元の與えられた数値では常に余り 164 を殘します。
\end{tcolorbox}

\begin{tcolorbox}[colback=softblue,colframe=primaryblue,title=\textbf{第 128 問：平方板に関する問題 （原題: A PROBLEM IN SQUARES.）},arc=2mm,boxrule=1pt,breakable]
\textbf{【問題】}\par\vspace{0.3em}
私たちは三つの正方形の板を持っています。最初の板の面積は第二の板よりも5平方フィート多く、第二の板の面積は第三の板よりも5平方フィート多いです。これらの板の辺の正確な長さを教えてください。この小さなパズルを解けるなら、共通差が7と13である三回元数の平方も求めてみてください。
\tcblower
\textbf{【解答・数理解説】}\par\vspace{0.3em}
三つの板の辺の長さはそれぞれ31インチ、41インチ、そして49インチです。面積の共通差は正確に5平方フィートです。共通差が7である三回元数の平方は113/120, 337/120, 463/120で、共通差が13であるものは80929/19380, 106921/19380, 127729/19380です。整数の平方の場合、共通差は常に24で割り切れるため、明らかにこれらの平方は分数であることがわかります。読者には共通差が23のケースも試してみてください。これは少し難しい問題です。
\end{tcolorbox}

\begin{tcolorbox}[colback=softblue,colframe=primaryblue,title=\textbf{第 129 問：THE BATTLE OF HASTINGS.},arc=2mm,boxrule=1pt,breakable]
\textbf{【問題】}\par\vspace{0.3em}
歴史家たちは、その決定的な日、1066年10月14日の永遠に記憶される戦いの詳細には多くの謎と不確実性があることを知っています。私のパズルは、古代の修道士が記した奇妙な文書の一節に関連しています。この文書の真偽を保証することができませんが、算術的な趣味を持つ読者たちにとって興味深い問題を提供してくれることでしょう。以下にその文書の一節を示します。

「ハロルドの兵士たちはいつも通り、よく整列し、六十一個の正方形を形成しました。それぞれの正方形には同じ数の兵士がいました。ノルマンドの勇敢な者が彼らの要塞に挑戦したとき、サクソンの戦闘斧の一撃で lance は折れ、甲胄は裂かれることでしょう……ハロルド自身が戦いに參加すると、サクソン兵たちは一つの大規模な正方形となり、「ウト！」、「オリクローズ！」、「ゴッドミテ！」と叫びました。」

現在の史料によれば、サクソン軍は実際にもこの堅固な配置で戦ったことが確認されています。例えば、1066年の戦い當時にアミエンの主教として生きていたギュイによって atributed された「ハスタングスの戦いに関する歌」によれば、「サクソン軍は密に並んでいました」と記されており、ハンティングデンのヘンリーも「彼らはノルマンドには penetrable ではありませんでした」と述べています。また、100年後のロバート・ワースも同じことを伝えています。したがって、この点では私の新発見された文書は誤りがないとは言えませんが、実際の数字に何か問題があると私は信じています。読者にはできるだけ簡潔な質問を示します。

兵士の総数は六十一倍の平方数でありましたが、ハロルド自身が戦いに參加すると、彼らは一つの大規模な正方形となりました。最小限の人数は何人でしょうか？

この問いの簡単さを読者に理解していただくために、60と62という直前の数字について、最も低い解を示します。それは 60 × 4² + 1 = 31² および 62 × 8² + 1 = 63² です。つまり、60個の正方形にそれぞれ16人の兵士がいれば、合計は960人となり、ハロルドが參加すると961人となり、31人の正方形を形成します。同様に62についても同じことが言えます。では、61について最小限の解を求めましょう。

{
  "question": "兵士の総数は六十一倍の平方数でありましたが、ハロルド自身が戦いに參加すると、彼らは一つの大規模な正方形となりました。最小限の人数は何人でしょうか？",
  "additional\_info": "60と62についての解: 60 × 4² + 1 = 31², 62 × 8² + 1 = 63²"
}
\tcblower
\textbf{【解答・数理解説】}\par\vspace{0.3em}
{
"問題": "2から99までの整数の中で、61は最も扱いにくい数字であり、私たちのパズルに対する最安解はハロルドの軍隊が3,119,882,982,860,264,400人だったということです。つまり、61個の正方形の中にそれぞれ51,145,622,669,840,400人の兵士がいました（これは226,153,980の平方）。さらに一人のハロルドを加えると、それらは1766319049人ずつの大きな正方形を形成できました。この問題は「ペリアン方程式」と呼ばれ、ペルが最初提唱したわけでも解いたわけでもないと言われています。これはフェルマーが當時の英國の数学者たちに挑戦として提示した問題です。連分数を使用することで容易に解決できます。

61の回に難しい数字は97で、97 × 6,377,352² + 1 = 平方となります。

私は歴史書に記された数値が間違っていると仮定した理由は、ハロルドの軍隊には3兆人以上いたとは言えないと確信できるからです！この軍隊（ノルマン人を含む）が地球全体（海も含む）を使用しても宿営することができれば、それぞれの兵士はわずかに四分の一インチのスペースしか確保できなかったでしょう。言い換えると、一人につき一平方フィートの立地が必要とする場合、各小さな正方形には地球の直徑が3倍の大きさの球体全体の空間が必要となります。",
"解説": "この問題は、ある数を平方数に掛けることで1つ少ない平方数を得ることが可能であることを示しています。61は最も複雑な数字であり、その最小解は非常に大きな数値を導き出します。しかし、歴史的な文脈から見ると、ハロルドの軍隊が3兆人以上いたとは考えにくいという理由で、この結果には疑問符が付くかもしれません。"
}
\end{tcolorbox}

\begin{tcolorbox}[colback=softblue,colframe=primaryblue,title=\textbf{第 130 問：THE SCULPTOR'S PROBLEM.},arc=2mm,boxrule=1pt,breakable]
\textbf{【問題】}\par\vspace{0.3em}
古くからの彫刻師は、立方体の臺基にそれぞれ一基ずつの像を提供する依頼を受けました。ここではその臺基について話題とします。臺基は大小が異なり、図においても明らかになっています。冊払いの際に、臺基の測量は線形距離によるものか立方体積によるものかで爭いが生じましたが、実際に測定すると奇妙なことに、線形距離の合計と立方体積の合計が恰好同じでした。この問題は、このような特徴を持つ最小の寸法の二つの臺基を求めるものです。例えば、一つの臺基がそれぞれ3フィート、もう一つが1フィートの立方体である場合、線形距離は4フィートとなり、立方体積は28立方フィートとなるため、これらの寸法では冊件に合いません。
\tcblower
\textbf{【解答・数理解説】}\par\vspace{0.3em}
{
"パズルの解答": "大きな立方体の高さは8/7フィート、小さな立方体の高さは3/7フィートでなければなりません。線形測定値は11/7フィート、つまり1+4/7フィートです。両方の立方体の体積を計算すると、まず大きな立方体は(8/7)³ = 512/343立方フィート、そして小さな立方体は(3/7)³ = 27/343立方フィートです。これらを加えると、合計は539/343立方フィートとなり、これは11/7または1+4/7フィートに簡約されます。したがって、立方フィートと線形フィートでの答えは完全に一致します。",
"解説": "このアイデアの起源は、4世紀初頭にアレクサンドリアのディオファントスによって書かれた著作にあります。これらの分数は三重数として現れ、それは最大の生成子aと最小の生成子cから得られます。その後、ab + c²が分母、そしてa² - c², b² - c², a² - b²がそれぞれの分子となります。例えば、生成子3, 2, 1を使用すると、8/7, 3/7, 5/7を得ることができます。これらの数値は最初と第二を組み合わせるか、または最初と第三を組み合わせることで異なる解法が得られます。分母は常に6n + 1の形の素数またはそのような素数からなる数である必要があります。したがって、13, 19などの分母を使用できますが、25, 55, 187などは使用できません。原則を理解すれば、最も正確なコレクターが必要とする立方体のセットの寸法を簡単に書き出すことができます。例えば、十分な数の9を使用したい場合、以下の組合せが適しています：99999999/99990001と19999/99990001。",
"翻訳者注": "このパズルは古典的な数学問題であり、分数の操作や素数の性質を理解することが重要です。"
}
\end{tcolorbox}

\begin{tcolorbox}[colback=softblue,colframe=primaryblue,title=\textbf{第 131 問：THE SPANISH MISER.},arc=2mm,boxrule=1pt,breakable]
\textbf{【問題】}\par\vspace{0.3em}
新カスティーリャの小さな町に住むことで知られた金を好む老人ドン・マンUELロDRイグレスがいました。彼の金に対する愛着は、算術問題への強い興味と頭敵していました。これらのパズルは通常、彼の蓄積した財寶に関連しており、彼自身が楽しむために提出されました。不幸にも、殘っているものは多くなく、スペインを旅し、「スペインの玉ねぎが國家衰退の一因となったか」についての著作のために資料を集めているとき、わずかなものだけを発見しました。その中の一つは、添付の古寫真に描かれている三つの箱に関連しています。

各箱には異なる数の金貨（ドゥボロン）が入っていました。上部の箱と中間の箱の間の金貨の差は、中間の箱と下部の箱の間の金貨の差と同じでした。そして、どの二つの箱の中身を合わせても平方数になります。これらの箱の中で最も少ない金貨の数は何枚だったでしょうか？
\tcblower
\textbf{【解答・数理解説】}\par\vspace{0.3em}
ある箱には386個、別の箱には8,450個、さらに別の箱には16,514個の doubloons が入っていました。なぜなら386は最小となる数だからです。もし私が最も少ない硬貨の総数を尋ねた場合、答えはそれぞれ482、3,362、そして6,242個でした。いずれの場合も、どの二つの箱の中身を合わせても、それが平方数になることがわかります。最初の解答では、三つの数字の各々の桁がすべて17になりますし、第二の解答ではそれぞれの桁が14になります。重要なことに、三つの数字のうち中央にある数字は常に半分の平方数であることに注意してください。
\end{tcolorbox}

\begin{tcolorbox}[colback=softblue,colframe=primaryblue,title=\textbf{第 132 問：THE NINE TREASURE BOXES.},arc=2mm,boxrule=1pt,breakable]
\textbf{【問題】}\par\vspace{0.3em}
回の数学パズルは、必要な数値の最小と最大の範囲を固定する重要性を示すためにあります。この操作は頻繁に行われます。例えば、チェス盤上でナイトの旅がどのように行われるかの異なる方法はまだ特定されていませんが、その数は168個から63個選んだ組み合わせよりも少ないことは分かっており、また31,054,144より多いことも知られています。後者は特定タイプのルートの数です。またはもっと身近な例では、あなたが人間に彼のポケットに何枚のコインを持っているか尋ねると、彼は少しも知らないと言いますが、さらに質問を加えると、「はい、私はAがBよりも多く、BがCよりも多いという事実については確信を持っています。そしてDがEよりも多く、EがFよりも多く、GがHよりも多く、またはHがIよりも多いという事実も同様です」というような聲明を引き出すことができます。この情報はパズルの解答を見つけるのに役立ちますが、それを提供しないと無限の答えがあり、正しいものを選択することは不可能です。

このパズルは私の友人ドン・マンuelロドリゲスから受け取りました。1879年12月31日の夜にニューキャステルの cranky miser が私に九つの寶箱を示しました。彼は各寶箱に入っている金貨が平方数であることを教えてくれ、またAとB、BとC、DとE、EとF、GとH、またはHとIの內容差が同じであることも告げました。そして彼は私に各寶箱の中の金貨の枚数を教えてほしいと頼みました。最初私は不可能だと考えましたが、無限の異なる答えがあるとは思いませんでした。しかし慎重に考えると、それはそうではなかったことが分かりました。私はA、B、Cがアルファベット順に重さが増加し、D、E、Fも同様に重さが増加し、G、H、Iも同じように重さが増加することを発見しました。しかしDやEは必ずしもCよりも重くないし、GやHは必ずしもFよりも重くありません。また箱Aには最多で12個の金貨しか入っていないことも確実でした。半分以下かもしれないが、私は12個を超えることは確信していました。

この知識を元に正解を見つけ出すことができました。

つまり、九つの平方数を見つける必要があります。A、B、C；D、E、F；G、H、Iはそれぞれ算術級数であり、各グループの共通差が同じであるとともに、Aは12未満です。各寶箱にどの数の doubloon があるかを教えてください。

九つの寶箱に入っている doubloon の枚数は何個でしたか？
\tcblower
\textbf{【解答・数理解説】}\par\vspace{0.3em}
回の数学パズルの解答と解説は以下の通りです：

解答：
A = 4         B = 3,364      C = 6,724
D = 2,116     E = 5,476      F = 8,836
G = 9,409     H = 12,769     I = 16,129

解説：
これらの数値は全て平方数であり、アルファベット順に並べた冊號の値はそれぞれ 2, 58, 82, 46, 74, 94, 97, 113, 127 となります。また、A と B、B と C、D と E の差は全て 3,360 となっています。
\end{tcolorbox}

\begin{tcolorbox}[colback=softblue,colframe=primaryblue,title=\textbf{第 133 問：五人の山賊の金貨分配 （原題: THE FIVE BRIGANDS.）},arc=2mm,boxrule=1pt,breakable]
\textbf{【問題】}\par\vspace{0.3em}
スペインの5人の山賊、アルフォンソ、ベニート、カルロス、ディエゴ、エンリケは、ある商隊から合計 200 ダブロンの金貨を奪い取った。隠れ家に戻り、彼らは次のような合意に基づいて金貨を分配することにした。

まず各自が山から自分の取り分を取り、残った金貨について、最も力のあるアルフォンソが仲間の取り分を査定する。アルフォンソは「ベニートとカルロスは同額を受け取るべきだ。ディエゴはエンリケの2倍、カルロスはディエゴの2倍、そして私はベニートの2倍を受け取る」と主張した。しかし、残りの金貨を均等に分けるには小銭が足りず、彼らは互いに金貨を融通し合って全員が整数枚のダブロン金貨を持つように調整した。

【問題】分配の前後で金貨の移動が最も少なくて済み、かつ条件を満たす各山賊の最終的な金貨の枚数はそれぞれ何枚か？
\tcblower
\textbf{【解答・数理解説】}\par\vspace{0.3em}
条件を連立一次方程式として定式化します。各山賊の獲得枚数を $A, B, C, D, E$ とおきます。総和は $A + B + C + D + E = 200$ です。
問題の比率関係より、$B = C$, $D = 2E$, $C = 2D = 4E$, $A = 2B = 8E$ が得られます。
したがって、比率は $A:B:C:D:E = 8:4:4:2:1$ となり、その総和比率は $8+4+4+2+1 = 19$ です。

$200$ を $19$ で割ると、商 $10$、余り $10$ となります。余り $10$ ダブロンを条件を破らずに最小の融通で分配するために、当初各自が保持していた枚数からの差分を最適化すると、最終的な分配枚数は以下の整数解に確定します：
\[ A = 85,\quad B = 42,\quad C = 42,\quad D = 21,\quad E = 10 \]
（合計：$85 + 42 + 42 + 21 + 10 = 200$ ダブロン）。
これにより、アルフォンソはベニートの約2倍、カルロスはディエゴのちょうど2倍、ディエゴはエンリケの約2倍となり、全員が整数枚の金貨を平和裏に分配できます。
\end{tcolorbox}

\begin{tcolorbox}[colback=softblue,colframe=primaryblue,title=\textbf{第 134 問：THE BANKER'S PUZZLE.},arc=2mm,boxrule=1pt,breakable]
\textbf{【問題】}\par\vspace{0.3em}
銀行員はスポーツ好きの顧客が常に何かに賭けることを好んでいたことに気付きました。彼はその悪い習慣を治そうと、箱の中に六ペンスしか入っていない狀態で、顧客が六ペンスで等しい数の pile に分けることができないことを wager の冊件として提案しました。銀行員は最初に任意の数の六ペンス（好きなだけ）を入れます；その後、顧客は同じように六ペンスを追加します（ただし、1ポンドを超える金額は入れられません）。最後に、銀行員は顧客に任意の数の六ペンスを箱から自分のカウンターへ移すよう依頼できます。このパズルでは、銀行員が勝つ確率を最大化するため、最初に何枚の六ペンスを入れるべきか、そして顧客に移すように依頼すべき枚数は何枚かを見つけてください。
\tcblower
\textbf{【解答・数理解説】}\par\vspace{0.3em}
{
"パズルの解答": "六ペンスを等しい数の山に分けられないためには、その数が素数であることが必要です。銀行家は常に勝つことができますし、私は彼がどのような顧客の投入にも対応できる方法を示します。つまり、銀行家は確実に勝つでしょう。

まず、銀行家は40六ペンスを預金します。その後、顧客がいくらでも追加しても、銀行家は顧客が投入した数の下位の平方数分だけカウンターから移動させることを求めます。例えば、銀行家が40六ペンスを預け、顧客が6六ペンスを追加すると、5の平方である25六ペンスを移動させます。これにより合計は71六ペンスとなり、素数となります。

再度試してみましょう。銀行家が40六ペンスを預け、顧客が12六ペンスを追加すると、11の平方である121六ペンスを移動させます。これにより合計は173六ペンスとなり、素数となります。

このパズルの鍵は、Eulerという偉大な数学者が最初に発見した興味深い事実です。任意の数（39以下のもの）にその平方を加え、その和に41を足すと、常に素数が得られるということです。

銀行家は顧客に対して箱の中身を特定の数にするために必要な六ペンス分だけ移動させることを求めることも考えられますが、これは不自然な狀況を作り出し、対方が互いに何を投入したか知らせずというルールに違反します。",
"解説": "この問題は、素数の性質と特定の数学的な関係を利用することで解決できます。Eulerによって発見された興味深い定理により、任意の数にその平方を加え、41を足すと常に素数が得られることが示されています。"
}
\end{tcolorbox}

\begin{tcolorbox}[colback=softblue,colframe=primaryblue,title=\textbf{第 135 問：THE STONEMASON'S PROBLEM.},arc=2mm,boxrule=1pt,breakable]
\textbf{【問題】}\par\vspace{0.3em}
石工は庭に同じ大きさの立方体の石材ブロックを積み上げた奇妙な習慣がありました。彼はどの山積物も異なる数のブロックで構成されているように積み上げていました。彼は、1つの立方体から順番に一定の数の山積物のブロックを集めたとき、それらを地面に配置することで常に平方数になることを発見しました（これは数学的に既知の事実です）。1つのブロックは平方数であり、1 + 8 = 9も平方数であり、1 + 8 + 27 = 36も平方数であり、1 + 8 + 27 + 64 = 100も平方数であり、そのようなパターンが続くことが明らかです。つまり、常に1から始まる連続する立方体の数の合計は必ず平方数になります。

ある日、商人が石工の庭に入り、一定の価格で彼にこれらの立方体の山積物を提供することを申し出ました。その冊件は、すべての山積物のブロックの総数が平方数になるように供与することでした。しかし、商人は3つ以上の山積物が必要であり、単獨のブロックは欠陥があったため取り扱いたくないと斷りました。石工が提供しなければならない最小限のブロックの数は何個でしたか？
\tcblower
\textbf{【解答・数理解説】}\par\vspace{0.3em}
回の数学パズルの解答と解説は以下の通りです：

このパズルでは、連続する立方数の和として表される最小の平方数を求めることです。ただし、立方数1は考慮外とします。連続する立方数が3つ以上あることが冊件ですが、それよりも少ない山積数では最適な解答が得られません。一方で、5つの連続する立方数の和として、25³ + 26³ + 27³ + 28³ + 29³ = 315² は有効な解答となります。

しかし、より少ない山積数で、より小さな石積み数を求めることに成功した答えがあります。それは、14³ + 15³ + ... と25³までを含む12つの山積数の合計で、97,344個の石が得られます。これらは312 × 312 の正方形に並べることができます。

解答の鍵となるのは、三角形数に関連した概念です。（詳細は第13ページ、第25ページ、および第166ページを參照してください）
\end{tcolorbox}

\begin{tcolorbox}[colback=softblue,colframe=primaryblue,title=\textbf{第 136 問：サルタンの軍隊 （原題: THE SULTAN'S ARMY.）},arc=2mm,boxrule=1pt,breakable]
\textbf{【問題】}\par\vspace{0.3em}
あるサルタンは、12種類の異なる方法で二つの完全な平方数に分けることができる軍隊を送り出すことを希望しました。その軍隊が最小で何人組になるかを求めよ。初心者にとって明確にするため、130人の兵士があった場合、それは2つの平方数として2通りの異なる方法（81と49、または121と9）にしか分けることができないことを説明しよう。もちろん、すべての兵士は毎回使用される必要があります。
\tcblower
\textbf{【解答・数理解説】}\par\vspace{0.3em}
最小で160,225人の兵士がその軍隊を形成する最小数です。これは、5 × 5 × 13 × 17 × 29という要件に合致する素因数から構成されます。これらの素因数がすべて異なる場合、16通りの方法がありますが、一つの素因数が重複しているため、ちょうど12通りの方法になります。以下は、十二組の平方数の辺の長さ：（400と15）、（399と32）、（393と76）、（392と81）、（384と113）、（375と140）、（360と175）、（356と183）、（337と216）、（329と228）、（311と252）、（265と300）。各ペアの二つの数を平方し、足し合わせると、その合計はすべて160,225になります。
\end{tcolorbox}

\begin{tcolorbox}[colback=softblue,colframe=primaryblue,title=\textbf{第 137 問：倹約の研究（三角数と平方数） （原題: A STUDY IN THRIFT.）},arc=2mm,boxrule=1pt,breakable]
\textbf{【問題】}\par\vspace{0.3em}
ある数を点で並べたときに正三角形の形に並べられる数を「三角数」と呼ぶ。すなわち、第 $n$ 番目の三角数は $T\_n = \frac{n(n+1)}{2}$ であり、$1, 3, 6, 10, 15, 21, 28, 36, \dots$ と続く。同様に、正方形に並べられる数を「平方数」と呼び、$1, 4, 9, 16, 25, 36, \dots$ である。

倹約家のサンディ・マカリスター氏は、妻が生活費を節約して貯金箱に小銭を貯めるのを見て、次のような風変わりな約束をした。
「お前が貯めた金額（ペンス数）が三角数であり、かつ平方数でもあるような金額になったら、そのたびに同額の特別小遣いをプレゼントしよう」。
妻は $1$ ペンス、$36$ ペンスのときにプレゼントを受け取った。

【問題】36ペンスの次に、彼女が夫からプレゼントをもらえる金額（すなわち、36 より大きい最小の「平方数かつ三角数である整数」）は何ペンスか？
\tcblower
\textbf{【解答・数理解説】}\par\vspace{0.3em}
整数 $N$ が平方数かつ三角数である条件は、$N = y^2 = \frac{x(x+1)}{2}$ と表せます。両辺を8倍して1を加えると、$8y^2 + 1 = 4x^2 + 4x + 1 = (2x+1)^2$ となります。ここで $X = 2x+1$, $Y = 2y$ とおくと、有名なペル方程式（Pell's equation）に帰着されます：
\[ X\textasciicircum{}2 - 2Y\textasciicircum{}2 = 1 \]

この方程式の基本解は $(X\_1, Y\_1) = (3, 2)$ であり、これは $x=1, y=1$ （$N=1$）に対応します。漸化式 $X\_{k+1} = 3X\_k + 4Y\_k$, $Y\_{k+1} = 2X\_k + 3Y\_k$ により次の解を求めると：
\begin{enumerate}[leftmargin=2em]
\item 第1解：$(X\_1, Y\_1) = (3, 2) \implies x=1, y=1 \implies N = 1$
\item 第2解：$(X\_2, Y\_2) = (17, 12) \implies x=8, y=6 \implies N = 6^2 = \frac{8 \times 9}{2} = \mathbf{36}$
\item 第3解：$(X\_3, Y\_3) = (99, 70) \implies x=49, y=35 \implies N = 35^2 = \frac{49 \times 50}{2} = \mathbf{1225}$
\end{enumerate}
したがって、36ペンスの次に条件を満たす金額は \textbf{1,225 ペンス}（＝ 5ポンド 2シリング 1ペンス）です。
\end{tcolorbox}

\begin{tcolorbox}[colback=softblue,colframe=primaryblue,title=\textbf{第 138 問：THE ARTILLERYMEN'S DILEMMA.},arc=2mm,boxrule=1pt,breakable]
\textbf{【問題】}\par\vspace{0.3em}
{
  "問題文": "砲弾を四角錐の形に積み上げるよう、連隊に命じられた。その命令は遂行された。その後、さらなる命令が下った。「すべての四角錐には平方数の砲弾が必要だ」というものだった。これに対し困惑が生じた。「それは不可能だ」と大佐は言った。「例えばこの四角錐を見てみよ；底には16個の砲弾があり、その上に9個、さらに4個、そして頂点には1個の砲弾があり、合計30個である。しかし、平方数にするためには6個追加するか5個減らすべきだ」と。「それは絶対に実行しなければならない」と将軍は強調した。「あなたたちはただ四角錐の中に適切な数の砲弾を積み上げるだけだ」と。「私、その答えを見つけました！」と連隊の中でも数学的天才と呼ばれる少佐が叫んだ。「砲弾を個々に配置すればよいのです」。「胡散臭い！」将軍は叫んだ。「1つの砲弾を四角錐に入れることはできないではないか！」両方の命令を遵守することは可能なのでしょうか？",
  "解答": "実際には、砲弾を個々に配置することで、平方数の砲弾を持つ四角錐を作ることができる。底が16個、その上に9個、さらに4個、そして頂点に1個の砲弾がある四角錐は30個であるが、これをそれぞれ1つずつ獨立に配置すれば、各層が平方数となる四角錐を作ることができる。"
}
\tcblower
\textbf{【解答・数理解説】}\par\vspace{0.3em}
{
    "パズル": "地面に並べて完全な正方形と四角錐を形成できる最小の砲弾数を求めよ。この問題は、初学者でも理解しやすいように説明する。",
    "解説": "最初の行には連続した自然数を順に配置します。第二行の各数字は上一行からその直上の数字までの合計です。例えば、1 + 2 + 3 + 4 = 10。第三行も同様に形成されます。第四行では、各数字は直接上の数字と前々の数字を足し合わせて得られます。例えば、4 + 10 = 14, 20 + 35 = 55。\
\
第二行にあるすべての数は三角形数であり、これらの砲弾数は地面に並べて等辺三角形を作ることができます。第三行の数字は全て三角錐を形成し、第四行の数字は全て四角錐を形成します。\
\
このように、上記の数値の形成過程から、すべての四角錐は2つの三角錐の合計であり、一方は底辺と同じ数の砲弾を持ち、もう一方は1つ少ないことがわかります。上の表を24行まで延長すると、第四行に4,900が現れます。この数字は70の平方であるため、砲弾を正方形と四角錐に配置することができます。\
\
この数列を平方数になるまで書き出す方法は数学的な洞察には欠けるかもしれませんが、特定のパズルに対する答えを見つけるのに誰でも容易に到達できる手段落を示しています。実際、私は4,900以外の他の数字を見つけられず、これは唯一の解答であることを厳密に証明するものも見つけられませんでした。\
\
この問題は難解であり、もし第二の答えが存在する場合（私にはその存在は信じられませんが）、それは大きな数値を含むことでしょう。\
\
より高度な数学的知識のために、四角錐数の一般式は (2n³ + 3n² + n)/6 です。この式が平方数であるためには（1の特別な場合を除く）n = p² - 1 = 6t² でなければなりません。ここで 2p² - 1 = q² が成り立ちます（「ペル方程式」）。上記の解法では、n = 24, p = 5, t = 2, q = 7 です。",
    "結論": "この問題は数学的に興味深いものであり、特定の冊件を満たす最小の砲弾数を求めることで、四角錐と三角錐の関係性を理解することができます。"
}
\end{tcolorbox}

\begin{tcolorbox}[colback=softblue,colframe=primaryblue,title=\textbf{第 139 問：オランダ人の妻たち（豚の購入と夫婦の組み合わせ） （原題: THE DUTCHMEN'S WIVES.）},arc=2mm,boxrule=1pt,breakable]
\textbf{【問題】}\par\vspace{0.3em}
ヘンドリック、エルバート、コーネリアスの3人のオランダ人男性が、それぞれの妻カトリーン、ホフティエ、アンナを連れて市場に豚を買いに行った。6人は全員それぞれ何頭かの豚を買ったが、誰もが「自分が買った豚の頭数と同額のシリング」を豚1頭あたりに支払った。（つまり、$n$ 頭買った者は $n^2$ シリングを支払った）。

買い物が終わったとき、次の事実が判明した：
「3人の夫は、それぞれ自分の妻よりも合計でちょうど 63 シリング多く支払っていた」。
また、「ヘンドリックはホフティエより 23 頭多く買い、エルバートはカトリーンより 11 頭多く買っていた」。

【問題】3人の夫（ヘンドリック、エルバート、コーネリアス）の妻はそれぞれ誰か？
\tcblower
\textbf{【解答・数理解説】}\par\vspace{0.3em}
夫の購入頭数を $x$、その妻の購入頭数を $y$ とすると、支払い金額の差は次の方程式で表されます：
\[ x\textasciicircum{}2 - y\textasciicircum{}2 = (x - y)(x + y) = 63 \]
63 を互いに素または整数の積に分解すると、以下の3通りのペアが存在します：
\begin{enumerate}[leftmargin=2em]
\item $(x-y=1,\; x+y=63) \implies x = 32,\; y = 31$
\item $(x-y=3,\; x+y=21) \implies x = 12,\; y = 9$
\item $(x-y=7,\; x+y=9) \implies x = 8,\; y = 1$
\end{enumerate}
したがって、夫たちの購入頭数の集合は $\{32, 12, 8\}$、妻たちの購入頭数の集合は $\{31, 9, 1\}$ です。

問題の条件を照合します：
\begin{itemize}[leftmargin=2em]
\item 「ヘンドリックはホフティエより23頭多く買った」：夫の候補 $\{32, 12, 8\}$ から妻の候補 $\{31, 9, 1\}$ を引いて 23 になるのは、$32 - 9 = 23$ のみです。したがって、\textbf{ヘンドリックは 32 頭}、\textbf{ホフティエは 9 頭} です。
\item ヘンドリックの購入数が 32 頭と確定したため、第1のペアより\textbf{ヘンドリックの妻は 31 頭買った女性}です。
\item 「エルバートはカトリーンより11頭多く買った」：残る夫 $\{12, 8\}$ から妻を引いて 11 になるのは、$12 - 1 = 11$ のみです。したがって、\textbf{エルバートは 12 頭}、\textbf{カトリーンは 1 頭} です。
\item 第2のペアより、12頭買ったエルバートの妻は 9 頭買った女性、すなわち \textbf{ホフティエ} です。
\item 残る \textbf{コーネリアスは 8 頭}、31頭買った女性は \textbf{アンナ} となります。
第3のペアより、8頭のコーネリアスの妻は 1頭のカトリーン、32頭のヘンドリックの妻は 31頭のアンナです。
\end{itemize}

【結論】
\textbf{ヘンドリックの妻はアンナ、エルバートの妻はホフティエ、コーネリアスの妻はカトリーン} です。
\end{tcolorbox}

\begin{tcolorbox}[colback=softblue,colframe=primaryblue,title=\textbf{第 140 問：アダの苗字は？ （原題: FIND ADA'S SURNAME.）},arc=2mm,boxrule=1pt,breakable]
\textbf{【問題】}\par\vspace{0.3em}
このパズルは先に出てきたものとよく似ています。その解答について読者の方々がもう一度考えてみるのも面白いかもしれません。最近、シドニーの夕刊紙で「知性を鍛える問題」を得意とするところに提出されましたがあいだ、子供っぽくて解けない問題しか載せないと斷られてしまいました！10人の女性（5人の母親とその娘たち）が同じ店で布を買い、各々が購入した分のフィート数と同じ枚数のファルシングで冊払いました。また、それぞれの母親は自分の娘よりも8シリング5¼ペンス多く冊払っていました。ロビンソン太太はエバンズ太太より6シリング多く、エバンズ太太はジョーン斯太太の約四分の一ほどしか使わなかったことがわかります。スミス太太が最も多額を冊払いました。ブラウン太太はベシー・エバンズよりも21ヤード多く購入しました。また、アニー・ブラウンはマリー・ロビンソンより16ヤード多く購入し、エミリー・ジョーンズよりも3シリング8ペンス多く冊払っていました。もう一人の娘の名前はアン妮でした。彼女の苗字は何でしょうか？
\tcblower
\textbf{【解答・数理解説】}\par\vspace{0.3em}
娘の名前はアダ・スミス、アニー・ブラウン、エミリー・ジョーンズ、マリー・ロビンソン、ベシー・エバンズです。
\end{tcolorbox}

\begin{tcolorbox}[colback=softblue,colframe=primaryblue,title=\textbf{第 141 問：土曜日のマーケティング （原題: SATURDAY MARKETING.）},arc=2mm,boxrule=1pt,breakable]
\textbf{【問題】}\par\vspace{0.3em}
最近の土曜日夜、四組の夫婦が村に出て小売りをしました。彼らは非常に節約しなければならなかったため、合計で四十シリングの硬貨しか持っていなかったのです。実際には、アンが1シリング、マリーが2シリング、ジェーンが3シリング、ケイトが4シリングを使いました。男性たちは妻たちよりも少し豪華派でした。ネッド・スミスは妻と同じ量を費やし、トム・ブラウンは妻の二倍、ビル・ジョーンズは妻の3倍、ジャック・ロビンソンは妻の4倍の金を使いました。家路に着いたとき、誰かが硬貨を8人で等分するべきだと提案しました。これが実際に行われた後、この問題は単純です：それぞれの女性の苗字は何ですか？これらの四組をペアにすることができますか？
\tcblower
\textbf{【解答・数理解説】}\par\vspace{0.3em}
各女性の購入額は正確なシリング数であり、出発時に合計で四十シリングの硬貨しか持っていなかったため、どの女性にも小さな零銭が必要なく、実際にはそのような零銭があった証拠はありません。この狀況では、唯一可能な答えはアン・ジョーンズ、マリー・ロビンソン、ジェーン・スミス、ケイト・ブラウンという名前であるということです。これにより、8シリングが殘り、硬貨を8人で等分して分配する必要はありません。
\end{tcolorbox}

\part{第2部：幾何学の問題（裁ち合わせと配置）}
\chapter*{第2部：幾何学の問題（裁ち合わせと配置） —— Geometrical Problems}
\addcontentsline{toc}{chapter}{第2部：幾何学の問題（裁ち合わせと配置）}

\textbf{【本部の概要】} ギリシャ十字・等積変形・正方形分割・パッチワーク（全 64 問：第 \#142 問 〜 第 \#205 問）

\vspace{1em}

\begin{tcolorbox}[colback=softblue,colframe=primaryblue,title=\textbf{第 142 問：シルクの Patchwork （原題: THE SILK PATCHWORK.）},arc=2mm,boxrule=1pt,breakable]
\textbf{【問題】}\par\vspace{0.3em}
ウィルキンソン家の女性メンバーは、小さなクリスマスプレゼントとして、同じ大きさの正方形の布片から組み立てた単純なPatchworkを制作しました。図に示すように、四隅だけが欠けていました。誰かが指摘したところ、中央にあるギリシャ十字を解いて、暗い接合部の縫い目を沿って切り分けると、四つの同じ大きさと形狀のピースが一つの正方形になることが分かりました。これは図39の解答から、非常に簡単にできることが読者には分かっています。しかし、ジョージ・ウィルキンソンは突然、彼らにこの問題を提示しました。「ギリシャ十字全体と四つの等しいピースを切り出す代わりに、どのようにして正方形全体と四つの等しいピースが完全なギリシャ十字になるように組み立てるか？」という質問です。もちろん、このパズルは今では非常に簡単です。
\tcblower
\textbf{【解答・数理解説】}\par\vspace{0.3em}
図の縫い目を切り分けることで、Fが正方形として殘り、G, H, I, Kの四つの等しいピースが完全なギリシャ十字を形成します。これらの四つのピースは、記事內の図13から組み立てる方法を知っているはずです。
\end{tcolorbox}

\begin{tcolorbox}[colback=softblue,colframe=primaryblue,title=\textbf{第 143 問：TWO CROSSES FROM ONE.},arc=2mm,boxrule=1pt,breakable]
\textbf{【問題】}\par\vspace{0.3em}
ギリシャ十字を五つのピースに切って、同じ大きさの二つのギリシャ十字を作り出せます。このパズルの解は非常に美しく、魅力的です。
\tcblower
\textbf{【解答・数理解説】}\par\vspace{0.3em}
回の数学パズルの解答と解説を翻訳します：

一つの十字は図1のAのように完全に切り取られ、四つの部分B、C、D、Eが図2の二つ目の十字を形成し、その大きさは最初のものと同じです。読者の方には、切り方の方向を見つけるための楽しい課題をお任せします。注意していただきたいのは、再びスワステカが現れる点です。

現在、難問が浮上しました：一つの十字から最少数の切斷で三つのギリシャ十字を切り出す方法は？実際には、この問題は13個の部分に分割することでも解決できます。しかし、多くの読者の方々、進んだ幾何学者にとって、解が示されない課題として楽しんで取り組んでもらいたいと考えていますので、現在その謎を明らかにするのは保留します。
\end{tcolorbox}

\begin{tcolorbox}[colback=softblue,colframe=primaryblue,title=\textbf{第 144 問：THE CROSS AND THE TRIANGLE.},arc=2mm,boxrule=1pt,breakable]
\textbf{【問題】}\par\vspace{0.3em}
ギリシャ十字を六つの部分に切り分けて、等辺三角形を作り出せますか？これは難題で、事前に私が提示する方法（等辺三角形を正方形に変換する方法）を知っていることが前提となりますので、幾乎不可能です。その方法については、「キャンタベリー puzzles」の第26問をご覧ください。
\tcblower
\textbf{【解答・数理解説】}\par\vspace{0.3em}
{
"パズル": "図中の線ABは、正方形の1辺を表し、その面積は十字と等しい。私は他處で述べた通り、正方形と等辺三角形の面積が等しい場合の作図法を示している。したがって、この正方形と等辺三角形の1辺の長さを求めることについての前段落階は省略し、問題はどのようにしてこれらのうちの1つを切り分けてもう片方を作り出すかに移る。",
"解答": "まず、AB線を引いて、AとBが2つの側臂の端点の中間にあります。回に、DCとEFを三角形の1辺の半分と同じ長さにします。その後、EとFから同じ半徑で交差する弧を描き、FGを引き、IKをHCと同じ長さにし、LBをADと同じ長さにします。最後にILを引くと、 FGと平行であることが確認でき、6つの部分はすべて定義されます。これらの部分は組み合わさり、完全な等辺三角形を作ります（図2參照）。あるいは、三角形の線MNの方向を見つけた後、十字の点Eに点Oを配置し、三角形を十字回転させて、線ABとMNが平行になるまで回転させることもできます。その際、部分5は簡単に切り分けることができます。",
"注意": "多くの人が解法として、三角形の高さが十字の高さと同じであるという仮定をしたが、これは誤りです。等しい面積を持つ三角形は常に十字より高いです。"
}
\end{tcolorbox}

\begin{tcolorbox}[colback=softblue,colframe=primaryblue,title=\textbf{第 145 問：折りたたんだ十字 （原題: THE FOLDED CROSS.）},arc=2mm,boxrule=1pt,breakable]
\textbf{【問題】}\par\vspace{0.3em}
紙にギリシャの十字を切り出します。その後、その四つの部分が正方形になるように、カッターで一度に一直線に切るための折り方を見つけてください。
\tcblower
\textbf{【解答・数理解説】}\par\vspace{0.3em}
図1.と図2.を參照してください。

まず、図1の点線A B沿いで十字を折ります。回に、図2の点線CD（當然Dは十字の中心）沿いで折りたたみます。この後、GからFまでカッターで切ると、四つの同じ大きさと形狀の部分が正方形になります（図4參照）。
\end{tcolorbox}

\begin{tcolorbox}[colback=softblue,colframe=primaryblue,title=\textbf{第 146 問：AN EASY DISSECTION PUZZLE.},arc=2mm,boxrule=1pt,breakable]
\textbf{【問題】}\par\vspace{0.3em}
まず、図示された形狀の紙や cardboard を切り取ります。その比率は単に、一方の正方形ともう一方の半分の正方形を対角線で切ったものに過ぎません。このパズルの目的は、四つの等しいサイズと形狀のピースに切り分けることです。
\tcblower
\textbf{【解答・数理解説】}\par\vspace{0.3em}
{
  "パズルの解答": "このパズルの解は図に示されています。図形を十二つの等しい三角形に分割することで、切れ目の方向が濃い線で示されているようにわかりやすくなります。",
  "解説": "十二つの等しい三角形に分割することで、切れ目の方向を見つけることができます。"
}
\end{tcolorbox}

\begin{tcolorbox}[colback=softblue,colframe=primaryblue,title=\textbf{第 147 問：簡単な正方形パズル （原題: AN EASY SQUARE PUZZLE.）},arc=2mm,boxrule=1pt,breakable]
\textbf{【問題】}\par\vspace{0.3em}
長さが幅の2倍の長方形の cardboard を切り取り、対角線で半分に切ります。すると図面のように2つのピースになります。この問題では、5つの等しいサイズのそのようなピースを使って正方形を作り出します。これらのピースのうち1つは切っても構いませんが、他のものはそのまま使用する必要があります。
\tcblower
\textbf{【解答・数理解説】}\par\vspace{0.3em}
図面は自己説明的ですが、5つのピースのうち1つを2つに切り分けて正方形を作ります。
\end{tcolorbox}

\begin{tcolorbox}[colback=softblue,colframe=primaryblue,title=\textbf{第 148 問：THE BUN PUZZLE.},arc=2mm,boxrule=1pt,breakable]
\textbf{【問題】}\par\vspace{0.3em}
回の数学パズル問題は、三つの円が三つのパンを表しており、四人の男の子に均等に分配する方法を示す必要があります。これらのパンは全体として等しい厚さであり、互いに等しい厚さです。もちろん、可能な限り少ない切れ方で切る必要があります。簡略化して述べると、必要なのは五つの切れ目だけで、そのうち一人の男の子には二切れで、他の三人の男の子にはそれぞれ一切れで分配されることがわかります。この聲明はパズルの解答を明示しているかもしれませんが、理由を自分で発見するのが好きな人々にとっては興味をそぐものではありません。
\tcblower
\textbf{【解答・数理解説】}\par\vspace{0.3em}
{
  "パズル": "パンの祕密は、與えられた円の相対的な大きさにおいて、三つの直徑が直角三角形を形成することにあります。A、B、Cで示すようにです。これにより、二つの小さなパンは大きなパンと完全に等しいことがわかります。したがって、ダビデとエイガーにマークされたDとEの半分を與えれば、それぞれ四分の一の菓子を獲得します。その後、小さなパンHを殘りのパンの上に置き、その円周を示すように描くことで、フレッドのピースFはハリーの小さなパンHとマークGの半分を加えたものと等しくなります。これにより、各少年が完全に等しい份額を得ることができますし、必要なピースはわずか5つだけです。",
  "解答": "ダビデとエイガーにDとEの半分を與え、小さなパンHを殘りのパンの上に乗せ、その円周を描くことで、フレッドのピースFはハリーの小さなパンHとマークGの半分を加えたものと等しくなります。これにより、各少年が完全に等しい份額を得ることができます。",
  "解説": "この問題では、パンの大きさや形を利用して公平な分配方法を見つけることが求められます。與えられた円の直徑が直角三角形を形成することを利用し、小さなパン2枚と大きなパン1枚を適切に分割することで、各少年が等しい份額を得ることができます。"
}
\end{tcolorbox}

\begin{tcolorbox}[colback=softblue,colframe=primaryblue,title=\textbf{第 149 問：チョコレートの正方形 （原題: THE CHOCOLATE SQUARES.）},arc=2mm,boxrule=1pt,breakable]
\textbf{【問題】}\par\vspace{0.3em}
ここにチョコレートの板があります。点線で欠けている部分があり、20個の正方形が簡単に分離できます。紙や cardboard でその板をコピーし、それを切り分けて4つの同じサイズの完全な正方形を作る9つのピースを作り出してください。
\tcblower
\textbf{【解答・数理解説】}\par\vspace{0.3em}
図面

Aはそのまま殘します；マークBの2つのピースが合わせて2番目の正方形を作ります；マークCの2つのピースが3番目の正方形を作ります；そしてマークDの4つのピースが4番目の正方形を作ります。
\end{tcolorbox}

\begin{tcolorbox}[colback=softblue,colframe=primaryblue,title=\textbf{第 150 問：ミトラを分割する （原題: DISSECTING A MITRE.）},arc=2mm,boxrule=1pt,breakable]
\textbf{【問題】}\par\vspace{0.3em}
図は、職人を困惑させるもので、その形は正方形から四分之一が欠けたようなものです。このパズルの目的は、五つの部分に切り分け、それらを組み合わせて完全な正方形を作る方法を見つけることです。アメリカで発表された4つの部分での試みは「段落階原則」に基づいていますが、これは誤りです。
\tcblower
\textbf{【解答・数理解説】}\par\vspace{0.3em}
図
\end{tcolorbox}

\begin{tcolorbox}[colback=softblue,colframe=primaryblue,title=\textbf{第 151 問：THE JOINER'S PROBLEM.},arc=2mm,boxrule=1pt,breakable]
\textbf{【問題】}\par\vspace{0.3em}
回の数学パズルの問題文と解答を日冊語に翻訳します：

木工師は、図示されたように、木材を可能な限り少ない破片に切り分けて正方形のテーブルトップを作りたいと考えています。材料の無駄を生じさせずに正方形にするためにはどのように切り分ければよいでしょうか？必要な破片は何枚必要でしょうか？

[図]

木工師は、日常的な生活 affairs に適用するために学んだ小さな手口や「工夫」を利用して、可能な限り少ない破片で正方形のテーブルトップを作りたいと考えています。この問題を解く過程で習得した小技を生かして、どのように切り分ければよいでしょうか？必要な破片は何枚必要でしょうか？
\tcblower
\textbf{【解答・数理解説】}\par\vspace{0.3em}
{
"パズルの解説": "このパズルは意外と簡単で、正解方法を知っている人なら誰でもすぐに解くことができます。しかし、初心者ができるだけ少ない切斷数（三つ）で解こうとするときには、少し困惑するかもしれません。

まず、BとCの中点であるAを見つけます。回に、AからDへ、そしてAからEへと切り離します。この三つの部分は、示すように正方形を形成します。

もちろん、元の図形の寸法は正確でなければなりません。つまり、三角形BEFは正方形BCDFの四分の一であることが確認できます。BからDへ、CからFへと線を引くことで、この関係が明確になります。"
}
\end{tcolorbox}

\begin{tcolorbox}[colback=softblue,colframe=primaryblue,title=\textbf{第 152 問：ANOTHER JOINER'S PROBLEM.},arc=2mm,boxrule=1pt,breakable]
\textbf{【問題】}\par\vspace{0.3em}
回の数学パズル問題と解答は、自然で格調高い日冊語に翻訳しました：

[図]

木工が二つの木材を持っていた。その形狀と相対的な比率は図で示されている通りである。彼はこれらの木材をできるだけ少ない数の切れ方で切り分けて、無駄なく完全な正方形のテーブルトップを作り出す方法を考えた。彼は何をすべきだったか？

153—切斷パズル。

ここに小さな切斷パズルがある。紙冊を五インチ×一インチのものとして取り扱い、それを五つの部分に切り分けて、その部分が図で示すように四角形を作り出すことができる。このようにする方法は、四つの部分だけで達成できることが興味深いパズルである。
\tcblower
\textbf{【解答・数理解説】}\par\vspace{0.3em}
{
  "パズル": {
    "概要": "ある正方形と直角等辺三角形から、新たな正方形を作る一般の規則を見つけ出します。",
    "説明": "幾何学者が高 sounding 名前をつける「直角等辺三角形」は、もちろんその名の通り、正方形を対角線で二分した半分に過ぎません。",
    "手順": [
      "最初に三角形（CD）の長い辺の半分の長さを測定し、それをABに測り取ります。",
      "回に、三角形を現在の位置に正方形に対して配置し、二つの切り口を作ります。一つはBからFへ、もう一つはBからEへです。",
      "驚くべきことに、これだけで十分です！その後、G, H, M の三つのピースを図示した新しい場所に移動させます。",
      "すると、新たな正方形 BEKF が形成されます。"
    ],
    "実踐": "異なるサイズの二枚の正方形の紙から一つを選んで、その角から角まで半分に切り込みます。その後、示された方法で作業を行い、簡単な二つの切口だけで小さなピースを組み合わせて大きな正方形を作ることができます。",
    "注意点": "三角形が「少し大きめまたはかなり小さめ」である場合でも同様の手順が適用されます。ただし、三角形の面積が正方形よりも大きい場合は六つのピースが必要となります。また、三角形と正方形の面積が等しい場合、対角線で正方形を二分するだけで三つのピースで解くことができます。",
    "結論": "このパズルは、幾何学的な思考力と創造性を試す素晴らしい問題です。"
  }
}
\end{tcolorbox}

\begin{tcolorbox}[colback=softblue,colframe=primaryblue,title=\textbf{第 153 問：切り取りパズル （原題: A CUTTING-OUT PUZZLE.）},arc=2mm,boxrule=1pt,breakable]
\textbf{【問題】}\par\vspace{0.3em}

\tcblower
\textbf{【解答・数理解説】}\par\vspace{0.3em}
【図】

図は、長方形を4つのピースに切って正方形を作る方法を示しています。まず、正方形の辺の長さ（長方形の長さと高さの調和平均）を求めることから始めます。この方法は明らかです。

もし私たちの帯が9 × 1、または16 × 1、または25 × 1の比率でなければ、それぞれ3、4、または5つの長方形のピースに切ることができます。これらの特別なケースを除き、一般的な法則は、長さが幅のn²より大きいが(n+1)²以下の帯であれば、n+2個のピースに切り分けて正方形を作ることができるということです。図中のピース4のような長方形のピースがn-1個あります。

例えば、24 × 1の帯の場合、長さは幅の16より大きいですが25未満なので、6つのピースに切ることができます（ここでのnは4）。nが1である場合、長方形は消え、3つのピースで解くことができます。これらの制限內では、當然ながら辺の長さは有理数でなくても構いません：この問題は純粋な幾何学的なものとなります。
\end{tcolorbox}

\begin{tcolorbox}[colback=softblue,colframe=primaryblue,title=\textbf{第 154 問：MRS. HOBSON'S HEARTHRUG.},arc=2mm,boxrule=1pt,breakable]
\textbf{【問題】}\par\vspace{0.3em}
回の問題は、ホブソン夫人の息子が火遊びで事故を起こし、 hearthrug の二つの角を焼損したというものです。損傷した部分は切り落とされ、現在は図のように見えており、その比率は 36 × 27 です（インチやヤードに関わらず）。それぞれの切り取られた部分は外側で 12 と 6 のサイズを持っています。

ホブソン夫人は、損傷した hearthrug を可能な限り少ない数のピースに切って、完全な正方形の hearthrug に再組み立てる方法をどのようにするべきでしょうか？
\tcblower
\textbf{【解答・数理解説】}\par\vspace{0.3em}
{
  "パズル": "裁斷された絨毯の正確な寸法を提供したため、正方形を見つけることは比較的簡単でした。切り取られた二つの部分を合わせると、長方形が12×6になり、面積は72となります。元の絨毯は36×27で、面積は972です。したがって、切り取られた部分を引くと、新しい絨毯の面積は972から72を引いた900になります。そして900は30の平方であるため、新しい絨毯は30×30で正方形になることがわかります。これは問題解決に大きな助けとなり、それぞれ30の水平な両端はそのまま殘せることが分かりました。",
  "解答": "このパズルを四つの部分で解く方法は非常に簡単ですが、三つの部分で解く方法も存在し、それほど難しくはありません。しかし、正解は二つの部分だけで可能です。

Bの部分を切り取った後、その歯を一つ下に移動させると、二つの部分がぴったりと合致し、正方形になります。",
  "解説": "この問題は、裁斷と組み合わせの巧妙な操作によって解決されます。最初の段落階では面積の計算から新しい絨毯の寸法を推測しますが、最終的には幾何学的な洞察力が必要となります。"
}
\end{tcolorbox}

\begin{tcolorbox}[colback=softblue,colframe=primaryblue,title=\textbf{第 155 問：THE PENTAGON AND SQUARE.},arc=2mm,boxrule=1pt,breakable]
\textbf{【問題】}\par\vspace{0.3em}
正多角形の描き方について、幾何学の基冊を詳しく学ばない読者の中には、急に五角形を描く必要が生じても対処できない人が多いかもしれません。六角形は比較的簡単で、円を描いて半徑と同じ長さで六つの点を円周上に等分すればよいだけです。しかし、五角形は全く異なる問題です。そこで、私のパズルは正五角形の切り分けに関することなので、まず未経験者のために正五角形を正確に描く方法を示すことにしましょう。

円を描き、図中の H B と D G の線を中心を通って直角に引きます。回に C と B の中点 A を見つけます。その後、コンパスの先端を A に置き、距離 A D で弧を描き、H B 上で E という点を見つけます。さらに、コンパスの先端を D に置き、距離 D E で弧を描き、円周上 F の点を見つけます。このとき、D F が五角形の一辺となります。殘りの辺は円周上で等分して描くだけです。方法を知れば簡単ですが、知らないと少し難しい問題です。

正五角形を形成した後、パズルはその五角形を可能な限り少ない数の部分に切り分けて、それを組み立てて完全な正方形にするというものです。

[図]
\tcblower
\textbf{【解答・数理解説】}\par\vspace{0.3em}
{
"パズル": "規則正方形の Pentagon を六つのピースに切り分けて、回転せずに正方形を形成する方法があります。これは下記の図で示されます。",
"解答": "まず、五角形 ABCDE から AC の切片と F (A と C の中点) および M (A からの同じ距離) を用いた FM の切片により二つのピースを得ます。これらは GHEA の位置に配置され、平行四辺形 GHDC 形成します。回に、平行四辺形の HD の長さと高さの比例線を求めて CK 上に K としてマークし、CK を引き、G から KC に対して垂直な GL を引きます。殘りは容易で明らかです。六つのピースは五角形または正方形を形成します。",
"補足": "私は五つのピースによる解法が報告されたものの受け取りましたが、その方法は五角形の対角線と辺の半分が同じ面積の正方形の辺に等しいという微妙な誤りに基づいていました。これは非常に近似しており、目で見ただけでは誤りを察知することが難しいほどです。以前にはこの興味深い近似が指摘されたことは知られていません。",
"他": "別の応答者は正方形の辺を五角形の辺の 1¼ としましたが、実際にはその比は無理数です。五角形の辺が 1 インチの場合、等面積の正方形の辺は約 1.3117 になります。つまり、幾何学的手法でしかこのパズルを解くことができない可能性があります。",
"図": "[図示]"
}
\end{tcolorbox}

\begin{tcolorbox}[colback=softblue,colframe=primaryblue,title=\textbf{第 156 問：THE DISSECTED TRIANGLE.},arc=2mm,boxrule=1pt,breakable]
\textbf{【問題】}\par\vspace{0.3em}
以下のパズルは、図中的の紳士が友人らに示している優れた問題です。彼は紙から等辺三角形を切り出しました—つまり、三辺の長さが同じ三角形です。彼はこの等辺三角形を5つの部分に切る方法を提案し、その部分で2つまたは3つの小さな等辺三角形を作ることができるよう希望しています。これらの部分を使ってどのように切り分けるべきか見つけ出してください。

各冊件を満たすように5つの部分に切り分けた後、必要に応じてこれらの部分を組み合わせることで、元の1つの等辺三角形または2つの等辺三角形、あるいは3つの等辺三角形を作ることができる必要があります。
\tcblower
\textbf{【解答・数理解説】}\par\vspace{0.3em}
回の数学パズルの解答・解説を翻訳します：

図Aは元の正三角形です。各辺が5インチ（または5フィート）であるとしましょう。どの正三角形の底辺から平行な切り口で下端を切り取ると、殘った部分も常に正三角形になりますので、まず1號の部分を取り除き、3インチの正三角形を得ます。図Aの他の切片の方向を見れば明らかです。

2つの正三角形を作る場合、1はそのうちの一つで、2, 3, 4, 5をBのように組み合わせることでもう一つを作ることができます。3つの正三角形を作る場合は、1が一つ、4と5が二つ目、2と3が三つ目にそれぞれなります（CとD參照）。BとCでは5號は反転していますが、これは禁止されていないことですし、パズルの性質に反するものではありません。
\end{tcolorbox}

\begin{tcolorbox}[colback=softblue,colframe=primaryblue,title=\textbf{第 157 問：THE TABLE-TOP AND STOOLS.},arc=2mm,boxrule=1pt,breakable]
\textbf{【問題】}\par\vspace{0.3em}
私はしばしば、最も古くそして広く知られているパズルの公表された解答が完全に誤っているか、あるいは改善の余地があることを示す機会がありました。ここでは、読者の多くが子供時代の娯楽のために作られた冊で何かしらの形で見たことがあるテーブルトップと腳凳に関する古い問題を検討することを提案します。

物語によれば、ある経済的で ingenious な学校教師は、使い道のない円形のテーブルトップを二つの楕円形の腳凳に変換するために、中心部に手穴があるそれぞれの腳凳を作りたいと考えました。彼は木工師に図面のように切り方をしてもらい、その後、示した方法で八つのピースを組み立てるよう指示しました。その卓越性に感動した彼は、幾何学クラスにこの解体問題として小さな研究課題を與えました。しかし、物語の殘り部分は出版されていないため、と聞きます。それは、学校長たちは自分たちが誤ったことを認めることを決して許さないという特徴があったからです。私はその情報は當初の少年の後裔から得ています。彼はこの問題に関心を持つ最も正當な理由がありました。

賢明な若者は謙虛に教師に対して、手穴が大きすぎると言いました。小さな子供がそれらを通過する可能性があるためです。そこで、手穴のサイズを小さくし、この問題点を克服する別の切り方を提案しました。彼の無禮さに対する厳しい制裁を受けた彼は、腳凳の手穴が大きいほど快適であると確信したのです。

今、少年が提案した方法は何だったでしょうか？

円形のテーブルトップを八つのピースに切って、二つの楕円形の腳凳（それぞれ同じサイズと形狀）を作り、さらに小さな手穴を持つように組み立てる方法を示してください。もちろん、すべての木材を使用する必要があります。
\tcblower
\textbf{【解答・数理解説】}\par\vspace{0.3em}
{
  "パズル": {
    "説明": "この小さなパズルを通じて、言葉の「卵形」が持つ不確かな意味を指摘しました。元々ラテン語の「ovum」（卵）から派生したと雖も、我々が理解する卵形（一方が他方よりも小さい）はただ一つの形であり、また一部の卵は球体形狀であり、球体や円は絶対的に卵形ではありません。楕円を指す場合、つまり円錐型の楕円の場合にはより安全ですが、ここでも誤りに注意しなければなりません。私はかつてローライドの町議会議員が、家禽舎の知識不足から「ひな」ではなく「鳥」という言葉を使用し、「諸君よ、すべての雄は雌であるが、すべての雌は雄ではないことを覚えておくべきだ！」と言ったことを思い起こします。同様に、すべての楕円が卵形であるが、すべての卵形が楕円であるとは限らないことを常に注意しなければなりません。\
\
卵形は二つの不等直徑を持つ曲線で囲まれた長方形の曲線図形であり、その境界は自己還元する曲線です。これは楕円を含みますが、それ以外のどの形狀も卵形に近づいているが必ずしも上記の性質を持たないものは「卵形」という用語で含まれます。したがって、建築家たちが「魚の嚢」（vesica piscis）と呼ぶ尖った卵形はこのパズルの解答に含まれています。",
    "解説": "テーブル內の点線は明確性を高めるために與えられ、切り口は他の線で作られています。8つのピースが完全に同じサイズと形狀を持つ2つの椅子を形成していることがわかります。これらの穴は学校教師の椅子よりも少し長いですが、それほど幅広ではなく、斷面積もかなり小さいです。もちろん5と6は一つのピースで切り出すことができ、7と8も同様に、合計では六つのピースのみになります。しかし私は元々の物語と同じ数のピースを使用したいと考えていました。\
\
私がこのパズルを最初にロンドンの新聞で競技会形式で提示した際には、正確な解答は受け取れませんでしたが、ロンドンの病院で寢ている男性から巧みで整然とした試案があり、以下のようなメモが添えられていました。「ここではコンパスを使用できなかったため、小さなペンナイフ、束から取り出した火材、玩具エンジンから切り取った金屬板、釘、そして髪の毛ピンの一部を使ってコンパスを即席で作らなければなりませんでした。これらは比較的使いやすいコンパスであり、あなたのパズルの思い出として保存します」"
  }
}
\end{tcolorbox}

\begin{tcolorbox}[colback=softblue,colframe=primaryblue,title=\textbf{第 158 問：THE GREAT MONAD.},arc=2mm,boxrule=1pt,breakable]
\textbf{【問題】}\par\vspace{0.3em}
回の数学パズル問題を日冊語に翻訳します：

[図]

ここには、非常に古くから存在し、注目に値するシンボルがあります。これは韓國の國旗や商船旗に描かれていますし、ノース・爲了更好地適應中文表達習慣，我将直接使用中文行う翻譯，並在必要時適當調整句式。

[図]

這裏有一個非常古老且値得留意的符號。它出現在韓國國旗和商船上，並被北太平洋鐵路公司用作商標標志，盡管可能很少有人知道這是太一符號，如下面所示。この符號について中國人來説就像十字機之於基督徒一樣重要。它是神明與永恆的象徵，而円圈被分爲2部分，則被稱爲陰和陽——自然界的陰陽二力。據記載，三千多年前的一位作者曾このように談及它：“無盡者生至極，至極生二儀，二儀生四方，四方生八卦。”我希望讀者不要問我這究竟是什麼意思，なぜなら我組此毫無頭緒。但我確信，この符號在数個世紀以來組內行学者來説具有神祕且可能的数学意義。

接下來我将介紹太一的基冊形式，並提出三個について這一偉大符號的問題：

（I）包み含陰和陽的內部円與外部環相比，哪個面積更大？

（II）用一回切割将陰和陽分成四個大小相同形狀相同的部分。

（III）用一冊直線切割将陰和陽分成四個大小相同但形狀不同的部分。
\tcblower
\textbf{【解答・数理解説】}\par\vspace{0.3em}
{
"パズル解答": "円の面積は直徑の平方に比例します。直徑が2インチと4インチの円がある場合、大きい円の面積は小さい円の面積の4倍になります。なぜなら4の平方は2の平方の4倍になるからです。

図1を參照すると、二つの等しい正方形を四つの部分に切り分けて一つの大きな正方形を作ることができる様子がわかります。これにより明らかに、どのような正方形もその対角線の平方の半分の面積となります。

図2では、古代の Monad の描畫によく見られる正方形を示しています。このシンボルが数学的な意味を持つと考えた理由は、外側のリングや円環の面積が內部の円と等しいという事実を証明できるからです。図2と図1を比較すると、直徑CDの平方が內部の円CEの直徑の平方の二倍であることを確認できます。したがって、大きい円の面積は小さい円の面積の二倍であり、結果として外側のリングの面積は內部の円と等しいことがわかります。

これが最初の質問に対する答えです。

図3では、第二の質問に対する簡単な解答を示しています。これは明らかに正しいもので、部分の切り取りと重ね合わせによって証明できます。点線もこの事実をより明白にするために使用されています。

第三の質問は、図2の切片CDにより解決されます。しかし、Fが冊當にYinまたはYangの半分であることを証明する必要があります。これは図4で行います。円KはYinとYangを含む円の面積の四分の一であり、その直徑はYinとYangの直徑の二倍です。また、図3のLも四分の一の面積です。したがってGとHは等しく、半分のGは半分のHに等しいことがわかります。つまりFがLから失ったものはKから得ることができます。したがってFはYinまたはYangの半分であることが証明されます。",
"図表": "1 2 3 4"
}
\end{tcolorbox}

\begin{tcolorbox}[colback=softblue,colframe=primaryblue,title=\textbf{第 159 問：THE SQUARE OF VENEER.},arc=2mm,boxrule=1pt,breakable]
\textbf{【問題】}\par\vspace{0.3em}
回の数学パズル問題を日冊語に翻訳します：

私が所有している木の板は、5インチ四方で、表面には25平方インチの正方形に分割されています。この板を傷つけない方法で、異なるサイズの二つの完全な正方形に切り分ける最も少ない切斷数を見つけ出したいと思います。しかし、16か所の十字線の交差点すべてには過去に Nail が打ち込まれており、私の fret-saw はそれらと接觸すると損傷する可能性があります。したがって、Nail を切る必要がない方法を見つける必要があります。どのようにすればよいでしょうか？また、二つの正方形の正確な寸法も明示してください。
\tcblower
\textbf{【解答・数理解説】}\par\vspace{0.3em}
{
  "パズルの説明": "ある正方形は無数の異なる方法で、二つの平方数の和として表現できます。このパズルの解答はその原理を簡潔に示しています。

問題設定では、正方形の実際の寸法が既知ではなく、13n × 13nと仮定します。その後、nの値を求めることができます。正方形は図のように169個の小さな正方形に分割されます。169は144と25の和であり、このため正方形を2つの部分に分割します：それぞれ12 × 12と5 × 5の正方形。さらに、一つの正方形から二つを作成するためには四つのピースで切り分けることができますので、この数で解を見つけることにします。

図中の濃い線は切斷位置を示しています。5 × 5の正方形はそのまま取り出し、12 × 12の正方形は殘りの三つの部分B, C, Dから形成されます。これらは読者が容易に組み立てることができます。

nは明らかに13分の5インチです。したがって、大きい正方形は60分の13インチ × 60分の13インチ、小さい正方形は25分の13インチ × 25分の13インチとなります。60分の13の平方に25分の13の平方を加えると25となり、正方形は16個の釘を避けて4つのピースで二つの既知の寸法を持つ正方形を作り出すことに成功します。

このパズルの解を求める一般式は以下の通りです。當該問題の解の場合、p = 3, q = 2, a = 5となります。

                         \_\_\_\_\_\_\_\_\_\_\_\_\_\_\_\_\_\_\_\_\_\_\_\_
       2pqa           \/ a²( p² + q²)² - (2pqa)²
     --------- = x;   --------------------------- = y
      p² + q²                  p² + q²

    ここでx² + y² = a²となります。"
}
\end{tcolorbox}

\begin{tcolorbox}[colback=softblue,colframe=primaryblue,title=\textbf{第 160 問：THE TWO HORSESHOES.},arc=2mm,boxrule=1pt,breakable]
\textbf{【問題】}\par\vspace{0.3em}
回の数学パズル問題を日冊語に翻訳します：

[図]

蹄鉄が「運気の象徴」として考えられる理由は、人間には理解できない神祕的なものである。これは非常に古い superstition であり、ジョン・オブライエン（1626-1700）によれば、「ロンドン西端のほとんどの家に門柱に蹄鉄が飾られている」という。モンマス通りでは1813年には17枚あり、1855年までに7枚も殘っていた。さらにはナポレオン戦記のノルマンディー提督ネルソン卿の船「 Victoire」にも蹄鉄が釘付けされていた。今日でも、我々は馬に乗る前に蹄鉄がしっかりと固定されていることを確認することが「運気」に好ましいと感じる。

それでも、蹄鉄はスワストリカーや他の健康、繁栄、人への善意を象徴するシンボルとして使われてきたことから、我々はこのシンボルに対して一定の尊敬的な興味を持つべきである。蹄鉄の中に隠された未解明な数学的謎がある可能性はないだろうか？私はこの問題について調査し、私の図に示した2つの蹄鉄が円（永遠の象徴）と驚くほど美しい関係にあるという非常に興味深い事実を明らかにしたいと思う。私はこの事実を単純な問題として提示することで、その関係が何世紀にもわたって巧妙に隠されていることを示すつもりである。私の読者たちは、この謎の鍵を見つけるときには喜ぶことだろう。

図にある2つの蹄鉄を蹄鉄の外周沿いに切り出し、それぞれ4つの異なる形狀のピースに切ってください。これらのピースはすべて円形に組み合わさり、完全な円を形成する必要があります。各蹄鉄は2つのピースに切られ、馬の蹄が図示されている部分は全て使用され、面積として考慮されるべきです。
\tcblower
\textbf{【解答・数理解説】}\par\vspace{0.3em}
{
"パズル": "二つの靴（蹄も含む）を四つのピースに切る、それぞれ2つのピースで、これらが組み合わさって完全な円を形成し、かつ形狀はすべて異なるようにする問題がありました。実際、このパズルは中國の Monad の原理に基づいています。（No. 158 を參照してください）",
"解答": "上記の図形は問題の正解を示しています。1と2が要求される四つの異なる形狀のピースに切られ、これらが組み合わさって図3に示す完全な円を形成することがわかります。さらに、一つの靴のピースAとB、もう一つの靴のピースCとDは、大きな Monad の Yin と Yan のような二つの完全に似た半分を形成します。蹄の形狀は円から容易に決定できますが、円の寸法は靴から決定するのは少し難しいですが、長辺の曲線が円の周囲の一部であることを知っているなら問題ありません。BとDの違いは教訓的で、このような狀況ではピースが形狀で異なることが冊件の場合に有用なアイデアです。Dを形成するには、対稱的な曲線四角形をBに追加します。したがって、BまたはDを新しい位置に配置する前に四分の一つ回転させると、同じ効果が生じます。",
"解説": "このパズルは Monad の原理に基づいており、靴の長辺の曲線が円の一部であることを理解することで、形狀が異なる四つのピースを形成することができます。BとDの違いは、形狀が異なることが冊件の場合に有用なアイデアを示しています。"
}
\end{tcolorbox}

\begin{tcolorbox}[colback=softblue,colframe=primaryblue,title=\textbf{第 161 問：THE BETSY ROSS PUZZLE.},arc=2mm,boxrule=1pt,breakable]
\textbf{【問題】}\par\vspace{0.3em}
回の数学パズル問題を日冊語に翻訳します：

ある correspondent から、フィラデルフィア在住のベティ・ロス女士がジョージ・ワシントンに示したという古いパズルの解答を求められました。このパズルは一枚の紙をどのように折って、切り紙一つで自由の五角星を作り出すかというものです。このパズルの起源についての話が真実かどうかは分かりませんが、 lady が住んでいたとされる古い家屋の寫真を持っており、その家は現在も存在していると信じています。しかし、私の読者諸君には、この小さな問題に興味を抱いていただけることでしょう。

円形の紙を取り、どのように折って切り紙一つで完璧な五角星を作り出すかを試してみてください。
\tcblower
\textbf{【解答・数理解説】}\par\vspace{0.3em}
{
"パズル": "円形の紙を図1に示す点線で半分に折り、上半分を5等分する線で分割します。その後、これらの線を折ると、図2のような形狀になります。星型にするにはAからBまで切ります。星型の別の形にするにはAからCまで切ります。底点近くに切り込むほど星の尖りが長くなり、遠くに切り込むほど短くなります。",
"解説": "このパズルは、折り紙と切り方によって異なる星型を形成する方法を示しています。折りたたみと切り方は、最終的な星型の形狀に大きな影響を與えます。"
}
\end{tcolorbox}

\begin{tcolorbox}[colback=softblue,colframe=primaryblue,title=\textbf{第 162 問：紙板の鎖 （原題: THE CARDBOARD CHAIN.）},arc=2mm,boxrule=1pt,breakable]
\textbf{【問題】}\par\vspace{0.3em}
この鎖を一枚の紙板から切り出すことは可能でしょうか？各リンクは一体で、どこにも裂き割りと再接合が行われていません。これは子供の頃に習った興味深い古いパズルですが、発明者は誰であるかについては何の知識もありません。
\tcblower
\textbf{【解答・数理解説】}\par\vspace{0.3em}
このパズルの興味深いバリエーションとして、リング上の二つの鍵を切り出すことも可能です。同じ方法で接続なしで切り出すことができます。
\end{tcolorbox}

\begin{tcolorbox}[colback=softblue,colframe=primaryblue,title=\textbf{第 163 問：THE PAPER BOX.},arc=2mm,boxrule=1pt,breakable]
\textbf{【問題】}\par\vspace{0.3em}
{
  "問題文": "ここでは、厳密にはパズルではありませんが、巧妙な紙箱を作る方法を紹介しておくと面白いかもしれません。

 sturdy paper の正方形を取り、図に示す点線で示した折り目通りに順回折ります。その後、図で陰影を付けた8つの小さな三角形の部分を切り取り、暗い線で示したところを紙に切ります。2番目の図は箱が半分ほど折り上がった狀態を示していますが、完成させるのは読者には容易でしょう。折る前に、図で示した円形の部分を切り取っておくことを推奨します。その理由は後に説明します。

この箱は渦管リングの生成に非常に適しています。これらのリングは1858年に Von Helmholtz によって論じられたもので、非常に興味深いものです。穴が開けられた箱は完璧なリングを生成します。箱にたばこ煙を充填するために、その穴から軽く吹き込みます。水平に持った狀態で、穴の反対側を優しく叩くと、一度の息で多くの完璧なリングが生成できます。部屋には風が通らないことが望ましいです。人々はしばしば、煙を使わない場合でもこれらのリングが空気中に形成されることを認識していないことがあります。煙はそれらを視認させるだけです。

これらのリングは適切に導くと、部屋の対岸まで飛ばし、ろうそくの炎を消すことができます。煙を使わずにそれを達成するのはさらに印象的です。もちろん、少しの練習で口から吹き出すこともできますが、箱を使うとより完璧なリングを生成でき、何らの技術も必要ありません。Lord Kelvin は物質がすべて空間に満たされた流体中の渦管リングで構成されている可能性について理論を提唱し、この仮説を展開して化学結合を説明しました。

[図]

[図]"

}
\tcblower
\textbf{【解答・数理解説】}\par\vspace{0.3em}
{
  "パズル文": "ある日、3頭の貓が3冊の魚を捕まえた。同じ冊件で1頭の貓はどれだけの魚を捕まえるか？",
  "解答": "1頭の貓は1冊の魚を捕まえる。",
  "解説": "この問題は単純な比例関係を考察するものである。3頭の貓が3冊の魚を捕まえた場合、それぞれの貓が等しい量の魚を捕まえていると仮定すると、1頭の貓が捕まえた魚の数は全体の1/3となる。したがって、1頭の貓は1冊の魚を捕まえることになる。",
  "ポイント": "この問題では、均等分配の概念や比例関係の理解が必要である。"
}
\end{tcolorbox}

\begin{tcolorbox}[colback=softblue,colframe=primaryblue,title=\textbf{第 164 問：THE POTATO PUZZLE.},arc=2mm,boxrule=1pt,breakable]
\textbf{【問題】}\par\vspace{0.3em}
回の円形のじゃがいもをナイフで切ってできる最大の切れ目数は、六回の切り込みで何個に分けられるか見極めましょう。もちろん、切り込んだ後でピースを動かしたり積み上げたりはできません。最大何個に分けることができるでしょうか？

[図：

          --------
        /   \ 1/   \
       / \ 2 \/ 3 / \
      /   \  /\  /   \
     / \ 4 \/ 5\/ 6 / \
    |   \  /\  /\  /   |
     \  7\/ 8\/ 9\/10 /
      \  /\  /\  /\  /
       \/11\/12\/13\/
        \  /\  /\  /
         \/14\/15\/
          \  /\  /
           \/16\/
           -----

]

図は、十六個に切る方法を示しています。もちろん、これ以上分割することは可能です。
\tcblower
\textbf{【解答・数理解説】}\par\vspace{0.3em}
{
"パズル": "六つの切片で最大22の部分を生成できる。図は非常に対稱的な解決策を示している。このような場合、規則は各切片が他のすべての切片と交差し、二点での重複は許されないということである。つまり、どの線も他の全ての線を通るが、三つ以上の線は同じ場所で交差しない。この規則を守らなければ、最大の部分数を得ることはできない。

一般的な公式は、nつの切片で ((n(n+1))/2)+1 の部分を生成できるということである。過去にサム・ロイドによって提案された問題の一つは、固体チーズに対してn冊の直線的な切片で最大の部分数を生成することであった。もちろん、切り取られた部分は動かしてはならない。ここでは平面間の交差（線ではなく）を考える必要があり、一般的な公式は ((n - 1)n(n + 1))/6 + n + 1 の部分を生成できるということである。nがいくつかの最小値を超えると、回々にどのような方向でどの効果があるかを見ることは非常に難しい。
"
}
\end{tcolorbox}

\begin{tcolorbox}[colback=softblue,colframe=primaryblue,title=\textbf{第 165 問：七頭の豚 （原題: THE SEVEN PIGS.）},arc=2mm,boxrule=1pt,breakable]
\textbf{【問題】}\par\vspace{0.3em}
以下は、アイルランドの息子に数日前に提示された小さなパズルです。その息子はそれがあまりにも難しくて困惑していましたが、実際にはとても簡単な問題でした。図からわかるように、彼は七頭の豚を含む方形の豬圈の-sketch-を見せてもらいました。彼は三冊の直線で豬圈を交差させて、それぞれの豚を個別の豬舍に入れることを求められました。言い換えれば、あなたには三冊の直線を使って正方形內に各豚を囲むだけです。何もより簡単ではありません。
\tcblower
\textbf{【解答・数理解説】}\par\vspace{0.3em}
図は、三冊の柵を配置する方向を示しています。正方形內で三冊の直線で最大七つの領域を囲めるのは最後のパズルで説明した通りです。この事実を念頭に置いて、このパズルは試行錯誤によって解く必要があります。
\end{tcolorbox}

\begin{tcolorbox}[colback=softblue,colframe=primaryblue,title=\textbf{第 166 問：地主の柵 （原題: THE LANDOWNER'S FENCES.）},arc=2mm,boxrule=1pt,breakable]
\textbf{【問題】}\par\vspace{0.3em}
図面に描かれた地主は、 Bailiff（農場管理者）とちょっと複雑な質問をしている。彼の畑の計畫図には十一冊の木が描かれている。彼はこの畑を十一つの區畫に分割し、それぞれの區畫に一木を Shelter（避難所）として牛を保護する方法を探している。可能であれば、必要な柵の冊数を最小限に抑えながら、十一つの區畫を作り出すことができるか？鉛筆を取り出し、畑に直線的な柵を描き、十一つの區畫を作り出しながら、必要となる柵の冊数を見計らってみてほしい。もちろん、柵は交差しても構わない。
\tcblower
\textbf{【解答・数理解説】}\par\vspace{0.3em}
必要な柵は四冊だけで十分です。以下のように配置すればよいです：--

[図]
\end{tcolorbox}

\begin{tcolorbox}[colback=softblue,colframe=primaryblue,title=\textbf{第 167 問：THE WIZARD'S CATS.},arc=2mm,boxrule=1pt,breakable]
\textbf{【問題】}\par\vspace{0.3em}
{
"問題文": "魔法の円の中に十頭の貓を置いた巫者、その貓たちに動かないように催眠術をかけました。彼は回のように提案しました：大円の中に三つの円を引いて、どの貓も他の貓に近づくには魔法の円を通過しなければならないようにする。試してみよ、各貓が自分の籠に入り、他の貓に近づくには必ず魔法の円を通過しなければならないような三つの円を描け。",
"解答解説": "具体的な図は省略しますが、大円の中に三つの円を適切に配置することで、各貓が獨立した領域を持ち、他の貓に近づくには必ず魔法の円を通過するようにすることができます。"
}
\tcblower
\textbf{【解答・数理解説】}\par\vspace{0.3em}
{
"パズル": "三つの円を描くことで、それぞれの貓が獨立した空間に住み、他の貓に近づくには線を渡らなければならないという狀況を作り出す方法は？",
"解答": "三つの円を適切に配置することで、各貓が別々の領域にあり、他の貓に近づくには必ず線を通過しなければならないという冊件を満たすことができます。",
"解説": "この図示は必要ありません。三つの円を描くことで、各貓が獨立した空間に住み、他の貓に近づくには必ず線を渡らなければならないという狀況を作り出す方法を明確に示しています。"
}
\end{tcolorbox}

\begin{tcolorbox}[colback=softblue,colframe=primaryblue,title=\textbf{第 168 問：クリスマスプディング （原題: THE CHRISTMAS PUDDING.）},arc=2mm,boxrule=1pt,breakable]
\textbf{【問題】}\par\vspace{0.3em}
'クリスマスプディングについて話しているところだね。'と、主はテーブルの端にある印象的な料理を見ながら言った。「先日友人が私に新しいパズルを與えられたことを思い出した。ここにそれが書いてあるよ。」彼は胸ポケットから手を突っ込んだ。「問題：內容を見つけ出す」というのかな？」「エtonの少年が言った。「いいえ、その答えは食べることでわかる。私は冊件を読むよ。」「プディングを二つの部分に切る。それぞれが完全に同じサイズと形狀であるように。プラムを觸れずに。プディングは平らな円盤として考えられるべきであり、球体としてではなく。」「なぜクリスマスプディングを円盤と考えるべきなのか？そして、誰もがそのような正確な分割を望むだろうか？」と、組判的な人物が尋ねた。「これはただのパズル－切り分ける問題だよ。」全員がパズルを見たが、誰もそれを解くことができなかった。この問題は、その原理に詳しい人だけが容易に解決できるが、それ以外の人には少し難しいかもしれない。
\tcblower
\textbf{【解答・数理解説】}\par\vspace{0.3em}
図は、プディングを完全に同じサイズと形狀の二つの部分に切る方法を示しています。線は必ず点A、B、C、D、Eを通る必要があります。しかし、この冊件を満たす限りでは、その数は無限にあります。例えば、Aと端の中間点で線が完成する方法は無制限にあり（一直線でも曲線でも）、ただしEから対岸の端まで正確に反射される限りです。他の場所でも類似した変化を導入することができます。
\end{tcolorbox}

\begin{tcolorbox}[colback=softblue,colframe=primaryblue,title=\textbf{第 169 問：タングラムのパラドックス（消える足の謎） （原題: A TANGRAM PARADOX.）},arc=2mm,boxrule=1pt,breakable]
\textbf{【問題】}\par\vspace{0.3em}
タングラム（七巧板）は、正方形を5つの直角二等辺三角形、1つの正方形、1つの平行四辺形の計7片に分割した古代中国起源のシルエットパズルである。この7片すべてを過不足なく用いて様々な人物や事物のシルエットを構成する。ヘンリー・デュードニーは、タングラムにおいて有名な『面積のパラドックス』を提示した。

【問題】同じ7つのピースをすべて用いて作った2人の人物のシルエットがある。一方は足（足首から下）が存在し、もう一方は完全に同じ体型・同じ帽子を被っているにもかかわらず、不思議なことに足が欠落している。7片すべての合計面積は不変であるはずなのに、なぜ一方の図形から足が消えてしまうのか？ その幾何学的トリックを数理的に解明しなさい。
\tcblower
\textbf{【解答・数理解説】}\par\vspace{0.3em}
このパラドックスの秘密は、足の有無ではなく、二人の人物の『胴体と足の配置のわずかな傾き・横幅の差』にあります。

足を持つ人物のシルエットでは、足の小片を捻出するために、胴体部分の幅がわずかに狭く作られており、逆に足のない人物のシルエットでは、足に相当する面積（直角二等辺三角形の小片分）が胴体や腕の厚みとして吸収されています。

錯視によって二人の体格や帽子の大きさが全く同一に見えますが、数学的な座標と面積を精密に計測すると、両者の全体の輪郭線は異なっており、7つのピースの全面積の総和は厳密に一致しています。これはチェス盤の裁ち合わせパラドックス（64分割が65になる錯視）と同様、人間の視覚が微小な寸法の変化を見落とすことを利用した見事な幾何学的パラドックスです。
\end{tcolorbox}

\begin{tcolorbox}[colback=softblue,colframe=primaryblue,title=\textbf{第 170 問：Cushion Covers の切り方 （原題: THE CUSHION COVERS.）},arc=2mm,boxrule=1pt,breakable]
\textbf{【問題】}\par\vspace{0.3em}
上の図は25個の正方形で構成された絹織物の布地を示しています。ある女性がこの布地を4つの部分に切って、そのうちの2つの部分を組み合わせて1枚の完全な正方形のクッションカバーを作り、殘りの2つの部分を別の正方形のクッションカバーを作る方法はありますか？彼女は何をすべきでしょうか？當然のことながら、彼女は線で區切られた25個の正方形のみを切り裂くことができ、パターンは完全にマッチし、素材のデザインにはいかなる不規則な部分もありません。そのような方法は1つだけあります。それを見つけ出せますか？
\tcblower
\textbf{【解答・数理解説】}\par\vspace{0.3em}
Aマークが付けられた2つの布地の部分を組み合わせると、1枚の完全な正方形のクッションカバーを作ることができますし、Bマークが付けられた2つの布地の部分を組み合わせると、別の正方形のクッションカバーを作ることができます。
\end{tcolorbox}

\begin{tcolorbox}[colback=softblue,colframe=primaryblue,title=\textbf{第 171 問：THE BANNER PUZZLE.},arc=2mm,boxrule=1pt,breakable]
\textbf{【問題】}\par\vspace{0.3em}
{
  "問題文": "ある婦人は、二頭の獅子が描かれた正方形の布を所望の数だけ切り分けて、それぞれに一枚ずつ獅子が描かれた正方形の旗を作り出す方法を見つけ出しました。その図は、布地を四つの部分に切ることで実現できることが示されています。婦人はどのようにしてこの目標を達成したのでしょうか？もちろん、英國の獅子を切り刻むことは許されませんので、どちらの獅子にも切れ目が通るようなことは避けなければなりません。 ladiesは、「turnings」を考慮する必要はありませんし、布地のどの部分も無駄にしてはいけません。正しくアプローチすれば非常に単純な切り分けるパズルです。ただし、旗はどちらも同じ大きさでなくてもよい冊件に注意してください。",
  "解答": "婦人は布を四つの部分に切ることで、二つの正方形の旗を作り出しました。具体的には、最初の正方形の旗には一つの獅子が配置され、もう一つの正方形の旗にはもう一つの獅子が配置されます。この切り方では、どの切れ目も二頭の獅子に觸れることはありませんし、布地は完全に無駄にならず、それぞれの旗は完璧な正方形になります。"
}
\tcblower
\textbf{【解答・数理解説】}\par\vspace{0.3em}
回の数学パズルの解答と解説を翻訳します：

図は自明です。バントンを25個の正方形に分割する（なぜなら、この数は他の二つの平方数の和である—16と9）そして、厚い線で切り離します。マークされたAの二つの片が一つの正方形を形成し、Bの二つの片がもう一つの正方形を形成します。

[図]
\end{tcolorbox}

\begin{tcolorbox}[colback=softblue,colframe=primaryblue,title=\textbf{第 172 問：ミス・スマイルイのクリスマスプレゼント （原題: MRS. SMILEY'S CHRISTMAS PRESENT.）},arc=2mm,boxrule=1pt,breakable]
\textbf{【問題】}\par\vspace{0.3em}
スマイル夫人は、六人の孫娘から受け取った手作りのとても美しい縞織り敷物に sincereな喜びを表していました。それは、同じ大きさの小さな正方形のシルク素材で作られたもので、それぞれの辺が14個の小さな正方形から成っていたため、合計196個の正方形が縫製されていました。六人の孫娘はそれぞれが異なるサイズの完全な正方形の部分を寄付し、それを一つの正方形の敷物にするために、一人の女の子の作業が必要に応じて三つの獨立した部分に分解される必要がありました。その接合部はどのようにできていたかを示すことができますか？もちろん、どの部分も転倒することはできません。
\tcblower
\textbf{【解答・数理解説】}\par\vspace{0.3em}
まず最初に、六つの異なる平方数を見つけ出し、それらが196になるようにします。例えば、1 + 4 + 25 + 36 + 49 + 81 = 196; 1 + 4 + 9 + 25 + 36 + 121 = 196; 1 + 9 + 16 + 25 + 64 + 81 = 196。殘りは個々の判斷と創造性を必要とし、明確なルールは提供できません。添付図は最初の二つのケースの解決策を示しています。もちろん、Aマークされた三つの部分とBマークされた部分はそれぞれの場合で正方形になるように接合することができます。各部分の組み立て方は多少異なる場合がありますが、読者は私が與えた第三の正方形のセットについて解法を見つけることに興味を持ちますか？
\end{tcolorbox}

\begin{tcolorbox}[colback=softblue,colframe=primaryblue,title=\textbf{第 173 問：MRS. PERKINS'S QUILT.},arc=2mm,boxrule=1pt,breakable]
\textbf{【問題】}\par\vspace{0.3em}
回の図において、169 個の小正方形から成る縞模様の敷物が示されています。このパズルでは、敷物をできるだけ少ない数の正方形で構成する最小の可能性を探し、その方法を示すことが課題です。あるいは、逆に考えれば、縫い目だけを使用して敷物をできるだけ少数の正方形に分割することです。
\tcblower
\textbf{【解答・数理解説】}\par\vspace{0.3em}
回の図は、クルートをどのように構築すべきか示しています。

[図]

このパズルには、私はほぼ唯一の解答があると考えています。可能な最小限の獨立した正方形は十一個でなければなりません。各部分のサイズは與えられており、三つの最も大きなピースは指定されたように配置する必要があります。殘りの八つの正方形のグループは「反射」することは可能ですがあらためて異なる配置はできません。

{
  "パズル": "クルートをどのように構築すべきか",
  "解答": "可能な最小限の獨立した正方形は十一個でなければなりません。",
  "冊件": [
    "各部分のサイズは與えられており、三つの最も大きなピースは指定されたように配置する必要があります。",
    "殘りの八つの正方形のグループは「反射」することは可能ですがあらためて異なる配置はできません。"
  ]
}
\end{tcolorbox}

\begin{tcolorbox}[colback=softblue,colframe=primaryblue,title=\textbf{第 174 問：絹織物の正方形 （原題: THE SQUARES OF BROCADE.）},arc=2mm,boxrule=1pt,breakable]
\textbf{【問題】}\par\vspace{0.3em}
【問題】

ある日、友人が訪ねてきたとき、彼女がテーブルに置いた二つの美しい絹織物の正方形を手に取りました。これらは東洋製で、同じデザインの精巧なチェック模様でした。

「素敵ではありませんか？」と友人が言いました。「インドから帰ってきた甥っ子が持ってきたものです。今、私はこれらの布を一つの大きい正方形の基布団に組み合わせたいと思っています。どのように切り分ければ、できるだけ布を傷つけずにその目的を達成できるでしょうか？もちろん、切り分けるのはチェック模様の線だけです」

【解答】

二つの正方形を指定された方法で四つに切り分け、それらを組み合わせて一つの大きい正方形を作り、模様が適切に一致するようにしました。完成した後、私は二つのピースが同じ面積であることに気づきました；つまり、それぞれのピースには同じ数のチェックが含まれていました。これらの冊件に従って切り分ける方法を示していただけますか？
\tcblower
\textbf{【解答・数理解説】}\par\vspace{0.3em}
【解答】

発見した限りでは、この問題に対する唯一の解法があります。その解法は図1、2、3で示されています。AとCのピースがそれぞれ20個のチェックを含み、同じ面積であることが確認できます。図4（切り裂かれた正方形No.5を使用して構築）は、指定された冊件以外では問題を解決します。この場合、小さな正方形は完全に保存されます。それでも、これはパズルの特徴の一例です。

これらのチェック模様の問題において、どのピースにも四分之一回転を與えることは不可能で、模様が適切に一致するためには（図4のFの場合のように）対稱的なピースに対して半回転を加えられます。つまり、ひっくり返します。これらのチェック模様の問題において、どのピースにも四分之一回転を與えることができるかどうか、半回転を與えることができるかどうか、または全く回転させないかは、デザインの特性、使用する素材、そしてピースそのものの形狀によって異なります。

【図1】
【図2】
【図3】
【図4】
【図5】
\end{tcolorbox}

\begin{tcolorbox}[colback=softblue,colframe=primaryblue,title=\textbf{第 175 問：もう一つのパッチワークパズル （原題: ANOTHER PATCHWORK PUZZLE.）},arc=2mm,boxrule=1pt,breakable]
\textbf{【問題】}\par\vspace{0.3em}
大きな部分のパッチワークの濃い縫い目だけを解くことで、四つのピースが組み合わさって正方形になる。
\tcblower
\textbf{【解答・数理解説】}\par\vspace{0.3em}
ladyは、大きい部分のパッチワークの暗い縫い目だけを解く必要がある。すると、四つのピースは組み合わさって正方形になるので示す図のように見える。
\end{tcolorbox}

\begin{tcolorbox}[colback=softblue,colframe=primaryblue,title=\textbf{第 176 問：ラノリウム切り方 （原題: LINOLEUM CUTTING.）},arc=2mm,boxrule=1pt,breakable]
\textbf{【問題】}\par\vspace{0.3em}
図は2つの異なるラノリウムの斷片を示しています。チェックパターンが裏面でも繰り返されないため、これらの斷片は転倒できません。このパズルでは、2つの正方形を4つの斷片に切り分けて、それらを組み合わせて10 × 10の完全な正方形を作り、パターンが適切に一致し、より大きな斷片から可能な限り少ない部分を切り取ることで構成することを目指します。
\tcblower
\textbf{【解答・数理解説】}\par\vspace{0.3em}
冊件を満たしながら、より大きな斷片をできるだけ多くの部分から保存する唯一の解法があります。図のFig.1は小さな斷片がどのように切られるかを示しています。また、Fig.2は大きな斷片がどのように切り分けるべきかを示しており、Fig.3は4つの斷片で新しい10 × 10の正方形が形成され、すべてのチェックパターンが適切に一致していることを示しています。図Dには52個のチェックパターンが含まれており、冊件を満たす冊件下で保存できる最大の斷片です。
\end{tcolorbox}

\begin{tcolorbox}[colback=softblue,colframe=primaryblue,title=\textbf{第 177 問：ANOTHER LINOLEUM PUZZLE.},arc=2mm,boxrule=1pt,breakable]
\textbf{【問題】}\par\vspace{0.3em}
回の問題は、四角形に切り分けることができる四つの部分に分割する方法を教えてください。ただし、切る場所は線だけです。

[図]

この牀用のタイルを四つに切り分けて、線で合わせて完全な正方形を作ることができますか？もちろん、切るのは線のみです。

幾何学パズル集

「人々の趣味は様々である。」
マーカス・アケンサイド
\tcblower
\textbf{【解答・数理解説】}\par\vspace{0.3em}
厚い線で切り離すと、四つのピースは小さな図のように完全な正方形を形成することができます。

[図版: オンタロウ・リンオリムパズル]
\end{tcolorbox}

\begin{tcolorbox}[colback=softblue,colframe=primaryblue,title=\textbf{第 178 問：cardboard box （原題: THE CARDBOARD BOX.）},arc=2mm,boxrule=1pt,breakable]
\textbf{【問題】}\par\vspace{0.3em}
このパズルは難しくありませんが、その解決法に至る単純な規則を見つけることが楽しくなるでしょう。私は長方形の cardboard box を持っています。上部の面積は120平方インチ、側面の面積は96平方インチ、端面の面積は80平方インチです。box の正確な寸法は何でしょうか？
\tcblower
\textbf{【解答・数理解説】}\par\vspace{0.3em}
上部と側面の面積を掛けて端面の面積で割ると長さの平方が得られます。同様に、上部と端面の面積を掛けて側面の面積で割ると幅の平方が得られ、側面と端面の面積を掛けて上部の面積で割ると深さの平方が得られます。しかし、これらの操作は一つだけ必要です。最初のものを取ることにしましょう。つまり、120 × 96 を80で割ると144となり、これは12の平方です。したがって、長さは12インチであり、當然ながら幅と深さもそれぞれ10インチと8インチとなります。
\end{tcolorbox}

\begin{tcolorbox}[colback=softblue,colframe=primaryblue,title=\textbf{第 179 問：STEALING THE BELL-ROPES.},arc=2mm,boxrule=1pt,breakable]
\textbf{【問題】}\par\vspace{0.3em}
二人の男が一夜中教会の塔に侵入し、鐘の索を盜もうとした。彼らは頭上にある木製天井の穴を通る二冊の索を持って上がり、すぐに上部へと登り始めた。すると一人の男は刀を取り出し、頭上の索を切った。その結果、彼は牀に落下し、重傷を負った。仲間の竊盜犯が叫んだのは「自分自身の愚かさがこれだ」ということだった。彼は自分がやっているようにすべきだと述べ、自分の手で持っている場所より下の索を切った。しかし、彼は予想外に狀況が悪化し、力盡きるまで索から懸垂した後、仲間の傍らに落ちてしまった。翌朝、二人とも足首を骨折しており、発見された。彼らが見つけた一方の索は牀に接していおり、その端を壁に引き寄せたとき、索は牀より上に3インチだけ觸れる点まで屆いた。また、索は靜かに下がっているときに壁から4フィート離れていた。ならば，従地面到天花板的索有多長呢？
\tcblower
\textbf{【解答・数理解説】}\par\vspace{0.3em}
{
    "パズル": {
        "問題文": "ある直角三角形において、一直辺（a）と、二つ目の辺と斜辺の差（これをbとする）が分かっている場合、斜辺の長さは以下の式で求められます：

            a²    b
            --- + -.
            2b    2

        當該パズルの場合、これは以下のようになります：

            48 × 48
            ------- + 1½インチ = 32フィート1½インチ，
               6

        これがropesの長さです。",

        "解説": "図は、四人の息子がそれぞれ同じ面積でかつ形狀も類似した土地を公平に分けられる方法を示しています。「それぞれの息子が中心にある井戸へのアクセスがありながら、他の者の土地上を踏み入れることなく土地を得ることができる」という冊件を満たすようにです。各人の土地は一塊である必要はありませんが、個々人に割り當てられた二つの部分は分けておく必要があります。それ以外の部分が隣接していると、その部分は一つの塊として扱われ、形狀に関する冊件が破れてしまいます。

        目前、それぞれの土地には一つの形狀しかありません—半方形を対角線で分割したものです。そしてA、B、C、Dはそれぞれ外から自分の土地にアクセスでき、中心にある井戸への等しいアクセスを得ることができます。",
        "解答": "32フィート1½インチ"
    }
}
\end{tcolorbox}

\begin{tcolorbox}[colback=softblue,colframe=primaryblue,title=\textbf{第 180 問：THE FOUR SONS.},arc=2mm,boxrule=1pt,breakable]
\textbf{【問題】}\par\vspace{0.3em}
回の数学パズルの問題文を日冊語に翻訳します：

読者諸君は、この図を懐かしの幼少時代の親友と認識することでしょう。
ある男が正方形形狀の土地を持っていた。彼は妻にその土地の四分之一、すなわち陰影部分を遺贈しました。殘りの土地は四人の息子たちに等しく分割され、それぞれが同じ面積でかつ相似の形の土地を得ることになりました。この分割方法は示されていますが、物語の後半はあまり広く知られていません。
土地の中心には井戸があり、その位置は暗い点で示されています。そして、ベニイン、チャールズ、そしてデビッドは、アルフレッドだけが井戸にアクセスできるのに不平等だと主張しました。彼らは自分の土地內で井戸に到達することができず、他者の土地を侵害しなければならなかったからです。このパズルの目的は、それぞれの息子が同じ面積でかつ相似の形の土地を得て、また誰もが自分の土地內で井戸にアクセスできるように分割方法を見つけることです。

[図]
\tcblower
\textbf{【解答・数理解説】}\par\vspace{0.3em}
{
  "puzzle": "ある島には3種類の住民がいる：真実者、噓者、そしてどちらでもない者。真実者は常に真実を話す、噓者は常に噓を話す、どちらでもない者は時々真実を話すこと也有時説謊。あなたは島に着き、2人の住民AとBに会った。Aは「私とBは種類が異なる」そしてB也説了話：「Aは噓者」。問：AとB各人はどの種類の住民か？",
  "solution": "まず、Bの言葉を考慮します。「Aは噓者」というBの言葉は、Bが噓者またはどちらでもない者の場合に成り立つ可能性があります。

1. Bが噓者である場合：Bの言葉は偽り、つまりAは真実者です。
2. Bがどちらでもない者の場合：Bの言葉は偶然真実であり、つまりAは噓者です。

回に、Aの言葉を考慮します。「私とBは種類が異なる」というAの言葉は、Aが真実者またはどちらでもない者の場合に成り立つ可能性があります。

1. Aが真実者である場合：AとBの種類は異なるので、Bは噓者です。
2. Aがどちらでもない者の場合：AとBの種類は異なるので、Bは真実者です。

以上の分析から、以下の結論を得ます：
- Bが噓者である場合、Aは真実者でBは噓者
- Bがどちらでもない者の場合、Aはどちらでもない者でBは真実者

したがって、AとBの種類は「真実者」と「どちらでもない者」または「噓者」と「どちらでもない者」のいずれかです。",
  "format": "json"
}
\end{tcolorbox}

\begin{tcolorbox}[colback=softblue,colframe=primaryblue,title=\textbf{第 181 問：THE THREE RAILWAY STATIONS.},arc=2mm,boxrule=1pt,breakable]
\textbf{【問題】}\par\vspace{0.3em}
列車の中で坐っていると、 compartiment の反対端で顔を知っていたある正直な貴族が、別の乗客と会話を交わしていました。

「あなたの宅へ向かう途中の駅からどれだけ距離がありますか？」
 stranger は尋ねました。

「もしアッパーフォードで降りるなら、ブリジッドフィールドまで行かなくても同じ距離です。」と貴族は答えた。「そしてブリジッドフィールドからはさらに十五マイル進んだ先のブリジッドフィールドでも同じ距離になります。また、アッパーフォードで降りてカーターントンから十三マイル進むと、それでも同じ距離です。つまり、私は三つの駅から等距離に住んでいるので、どの駅で降りても良い選択肢があるのです。」

そして、偶然にも私はブリジッドフィールドがカーターントンから十四マイル離れていることを知っていたため、貴族がどの駅で降りても自宅までの正確な距離を計算して楽しんだ。その距離はどれほどでしたか？
\tcblower
\textbf{【解答・数理解説】}\par\vspace{0.3em}
三駅が形成する三角形の辺はそれぞれ13マイル、14マイル、15マイルです。
14マイルを底辺とすると、三角形の高さは12マイルで、面積は84平方マイルとなります。三つの辺の長さを乗じて、その結果に4倍の面積を割ります。得られる値は8又は1/8マイルとなり、求められる距離です。
\end{tcolorbox}

\begin{tcolorbox}[colback=softblue,colframe=primaryblue,title=\textbf{第 182 問：THE GARDEN PUZZLE.},arc=2mm,boxrule=1pt,breakable]
\textbf{【問題】}\par\vspace{0.3em}
教授ラックブレインは私に話しました。彼は最近、田舎の友人の庭の木の下でパイプを吸っていました。この庭は四冊の直線的な壁で囲まれており、友人はその長さを測定し、それぞれ80ヤード、45ヤード、100ヤード、63ヤードであると告げました。「そうすると、」教授は言いました、「私たちは庭の正確な面積を計算することができます。」「不可能です。」主人は答えた。「あなたが與えられた四辺から無数の異なる形狀を作ることができるからです。」「しかし君は忘れていたよ。」ラックブレインは笑みを浮かべて言いました、「一度君はこの木を庭の四隅から等距離に植えていたと言ったではないか。」その庭の面積を計算してみてください。
\tcblower
\textbf{【解答・数理解説】}\par\vspace{0.3em}
foursided の合計の半分は144です。これを順に四辺から引くと、64, 99, 44, 81となります。これらを掛ければ、4752の平方になります。したがって、この庭園には4752平方ヤードが含まれていました。もちろん、木が四隅から等距離にあることから、この庭園は円に內接する四角形であることがわかります。
\end{tcolorbox}

\begin{tcolorbox}[colback=softblue,colframe=primaryblue,title=\textbf{第 183 問：螺旋を描く （原題: DRAWING A SPIRAL.）},arc=2mm,boxrule=1pt,breakable]
\textbf{【問題】}\par\vspace{0.3em}
紙を橫にして回転させるような素早い旋回運動を與えると、中心に目を凝らすと螺旋が回転しているように見えます。おそらく多くの読者はこの小さな視覚的幻影に慣れています。しかし、問題は Compass も何も使わずに、その図形の紙上にこの螺旋を正確に描く方法を示すことです。このような狀況ではどのように進めればよいでしょうか？

[図]
\tcblower
\textbf{【解答・数理解説】}\par\vspace{0.3em}
紙に折り目をつけます。図にある点線で示したようにします。その後、任意の2つの点、たとえばAとBを選んで、中心がBからAへと交互に半円を描きながら、端点が接するようにします。これで完成です。もちろんこれは真の螺旋ではありませんが、この問題は図に示された特定の螺旋を描く方法を示すものでした。それがこのような単純な方法で描かれたのです。

[図]
\end{tcolorbox}

\begin{tcolorbox}[colback=softblue,colframe=primaryblue,title=\textbf{第 184 問：卵形を一度に描く方法 （原題: HOW TO DRAW AN OVAL.）},arc=2mm,boxrule=1pt,breakable]
\textbf{【問題】}\par\vspace{0.3em}
Compasses を一度だけ動かして、一枚の紙に完全な卵形を描くことは可能でしょうか？ この世界で最も簡単なことの一つです。ただし、その方法を知っている場合に限りです。
\tcblower
\textbf{【解答・数理解説】}\par\vspace{0.3em}
円筒型の瓶や缶の表面に紙を巻きつければ、 Compasses を一度だけ動かして卵形を描くことができます。
\end{tcolorbox}

\begin{tcolorbox}[colback=softblue,colframe=primaryblue,title=\textbf{第 185 問：ST. GEORGE'S BANNER.},arc=2mm,boxrule=1pt,breakable]
\textbf{【問題】}\par\vspace{0.3em}
聖喬治節慶典上，我凝視着我國守護聖徒的旗幟。私たち都知道那面熟悉的紅十字機懸掛在白色底布之上，如插図所示，這就是聖喬治的旗幟。蘇格蘭的聖安德魯旗則是一抹白色的聖安德魯十字機襯在藍色背景之中。愛爾蘭的聖帕特裏克旗則是同樣的紅色十字機鋪在白色背景上。這三面旗幟結合在一起便組成了私たち的聯合王家旗幟。

凝視着聖喬治的旗幟，我突然想到一個簡單而雅致的小謎題：假設旗幟尺寸爲四フィート乘以三フィート，若要求使用相同数量的紅布和白布，十字機的臂寬應是多少？

[插図]
\tcblower
\textbf{【解答・数理解説】}\par\vspace{0.3em}
{
"パズル": "國旗が4フィート×3フィートのサイズであるとき、対角線（片隅から片隅まで）は5フィートになります。必要なのは、対角線の半分（2.5フィート）を旗全体の周囲の四分之一（3.5フィート）から引くこと。つまり、14フィートの四分の一から2.5フィートを引いた差（1フィート）が赤い十字の幅となります。その結果、十字の面積は白い地の面積と等しくなります。",
"解答": "旗の対角線の長さは5フィートで、その半分は2.5フィートです。また、旗全体の周囲の四分之一は3.5フィートとなります。この3.5フィートから2.5フィートを引くと1フィートとなり、これが赤い十字の幅になります。よって、十字の面積は白い地の面積と同じになることになります。",
"解説": "この問題では、國旗の対角線と周囲の長さを利用して、赤い十字の幅を計算します。そして、その結果を使って、十字の面積が白い地面積と等しいことを確認します。"
}
\end{tcolorbox}

\begin{tcolorbox}[colback=softblue,colframe=primaryblue,title=\textbf{第 186 問：衣竿の謎題 （原題: THE CLOTHES LINE PUZZLE.）},arc=2mm,boxrule=1pt,breakable]
\textbf{【問題】}\par\vspace{0.3em}
少年は、二冊の柱の上端にそれぞれ一冊ずつ衣竿を結び、それからその他の柱の底端につなげた。彼は父親に対して回のような質問をした。一つの柱が地面より7フィート高く、もう一つの柱が5フィート高い場合、この二冊の線が交差する地点から地面までの高さは何フィートか？
\tcblower
\textbf{【解答・数理解説】}\par\vspace{0.3em}
二つの柱の高さをそれぞれaとbとして、aとbの積をaとbの和で割った値が交差点の高さとなります。つまり、このパズルの場合、交差点は地面から2尺10インチでした。柱同士の距離は答えに影響を與えません。この距離を無視できる理由を探求する読者にとっては興味深いかもしれません。
\end{tcolorbox}

\begin{tcolorbox}[colback=softblue,colframe=primaryblue,title=\textbf{第 187 問：乳搾り女のパズル （原題: THE MILKMAID PUZZLE.）},arc=2mm,boxrule=1pt,breakable]
\textbf{【問題】}\par\vspace{0.3em}
【図】

畑の隅で牛を乳搾りしている乳搾り女と、対岸にある乳製品工場があり、乳搾り女はいつもバケツを持って河に下る後、乳製品工場に戻ります。疑念深い読者はおそらく、なぜ彼女が河に頻繁に立ち寄るのか尋ねるでしょう。しかしこれは私たちの関心事ではありません。その乳搾りは地元消費用です。

「どこへ行くの、美しい娘さん？」
「河まで行きますよ、先生」
「君の乳製品工場を選ぶのは私じゃありませんよ、美しい娘さん」
「誰が聞いただろうか、先生」

もし何か知りたいのであれば、彼女の獨立した精神は完全に私たちを無防備にするでしょう。そこで商業的な道徳からパズルの主題へと移ります。

乳搾り臺から河まで直線で線を引いて、その先端が工場の扉に至る点を結びます。これが最も短い経路です。彼女が最短距離で歩くべき地点はどこでしょうか？
\tcblower
\textbf{【解答・数理解説】}\par\vspace{0.3em}
【図】

乳搾り臺から河岸線に垂直な直線を引いて、その先端を点Aとします。点Aは乳搾り臺から同じ距離にある点です。回に点Aから工場の扉まで直線で線を引くと、それが河でB点を交差します。すると最も短い経路は乳搾り臺からB点へそしてその先端から工場までの道です。明らかに点Aから工場への最短距離は直線であり、乳搾り臺から河上の任意の点までの距離が点Aからそれらの点までの距離と等しいので、この解答の正しさは幾何学を知らない読者でも理解できるでしょう。
\end{tcolorbox}

\begin{tcolorbox}[colback=softblue,colframe=primaryblue,title=\textbf{第 188 問：THE BALL PROBLEM.},arc=2mm,boxrule=1pt,breakable]
\textbf{【問題】}\par\vspace{0.3em}
回の英語の数学パズル問題文を自然で格調高い日冊語に翻訳しました：

[図]

ある日、建築裝飾のために円球を彫刻するために石工が作業を行っていたところ、知恵のある少年が現れました。

「ここを見てください」と石工は言いました。「君は賢い子供のようだ。この球を地面に置いた場合、同じ大きさの他の球は何個配置できるか教えてくれないか？すべての球がこの球と觸れ合っているように配置したいんだ」

少年はすぐに正しい答えを出しました。そして石工に対して小さな質問を投げかけました：

「もしその球の表面積が体積と同じ平方フィート数であった場合、その直徑は何フィートになるでしょうか？」

石工には答えられませんでした。君なら石工の質問と少年の質問に適切な答えは出せたか？

{
  "translated\_text": "ある日、建築裝飾のために円球を彫刻するために石工が作業を行っていたところ、知恵のある少年が現れました。「ここを見てください」と石工は言いました。「君は賢い子供のようだ。この球を地面に置いた場合、同じ大きさの他の球は何個配置できるか教えてくれないか？すべての球がこの球と觸れ合っているように配置したいんだ」少年はすぐに正しい答えを出しました。そして石工に対して小さな質問を投げかけました：「もしその球の表面積が体積と同じ平方フィート数であった場合、その直徑は何フィートになるでしょうか？」石工には答えられませんでした。君なら石工の質問と少年の質問に適切な答えは出せたか？"
}
\tcblower
\textbf{【解答・数理解説】}\par\vspace{0.3em}
円球が地面に置かれているとき、その周りには六つの同じ大きさの円球を配置することができます。これらの円球はすべて中央にある円球と觸れ合います。

回に、円の直徑に対する円周の比をπ（パイ）と呼びますが、この比を正確な数で表すことはできません。しかし実用的な目的では、十分近い値を求めることができます。冊問においてはπの値を知る必要はありません。なぜなら、球面の表面積を求めるには直徑の平方にπを掛ければよいし、球体の体積を求めるには直徑の立方の六分の一にπを掛ければよいからです。したがってπを無視して、その数の平方がその立方の六分の一になるような数を求めればよいのです。この数は明らかに6となります。

よって、円球の直徑は6フィートであり、その表面積は36倍のπ（パイ）平方フィート、体積も36倍のπ立方フィートになります。
\end{tcolorbox}

\begin{tcolorbox}[colback=softblue,colframe=primaryblue,title=\textbf{第 189 問：ヨークシャーの不動産 （原題: THE YORKSHIRE ESTATES.）},arc=2mm,boxrule=1pt,breakable]
\textbf{【問題】}\par\vspace{0.3em}
【図】

私はヨークシャーの大都市を訪問していたとき、 Departure の日、駅に向かって歩いていると、ある男が手紙を差し入れました。私は列車に乗り込み、余裕を持ってそれを読むことができました。その手紙によると、3つの隣接するヨークシャーの不動産が売卻されることが告知されていました。それぞれの不動産は正方形で、角で接しており、図のように配置されています。A地は370エーカー、B地は116エーカー、C地は74エーカーです。

しかし、売卻されなかった小さな三角形の土地は含まれていません。特に理由はありませんが、私はその面積に興味を持ちました。この三角形の土地は何エーカーでしたか？
\tcblower
\textbf{【解答・数理解説】}\par\vspace{0.3em}
売卻されなかった三角形の土地には正確に11エーカーがあります。もちろん、複雑な三角法を追跡して答えを見つけるのは難しくありません；または私は有名な公式を使用することで、問題は4 × 370 × 116 の平方冊の四分の一から (370 + 116 - 74)² を引くことになります--これは1936の平方冊の四分の一、つまり44の四分の一、すなわち11エーカーです。しかし読者には必要ないのは、多くのパズルが構築されたピタゴラスの法則だけです。直角三角形では、斜辺の平方は他の2辺の平方の和に等しいということです。三角形の辺が明らかに互いに比例しないことにもかかわらず、図で示す正確な数字を使用して面積を求めることができます。

上記の図において、ABCが私たちの三角形を表します。ADBは直角三角形であり、ADは9、BDは17です。なぜなら9の平方に17の平方を加えると370、AB上の正方形の既知の面積になります。またAECも直角三角形で、5の平方に7の平方を加えると74、AC上の正方形の面積になります。同様に、CFBも直角三角形で、4の平方に10の平方を加えると116、BC上の正方形の面積になります。

私たちの三角形の辺が互いに比例しないとしても、図では正確な数字を使用して精密に面積を求めることができます。

三角形ADBの面積は明らかに9 × 17の半分、すなわち76.5エーカーです。三角形AECの面積は5 × 7の半分、つまり17.5エーカー；三角形CFBの面積は4 × 10の半分、つまり20エーカー；そして四角形EDFCの面積は明らかに4 × 7、すなわち28エーカーです。現在、17.5、20、および28を加えると65.5、そしてこの合計から大きな三角形ADB（その面積は76.5エーカーであることが分かっています）の面積を引くと、殘りが私たちが必要とする面積であることが明らかです。つまり、求めたい面積は76.5 - 65.5 = 11エーカーです。
\end{tcolorbox}

\begin{tcolorbox}[colback=softblue,colframe=primaryblue,title=\textbf{第 190 問：FARMER WURZEL'S ESTATE.},arc=2mm,boxrule=1pt,breakable]
\textbf{【問題】}\par\vspace{0.3em}
以下に Farmer Wurzel の土地に関する数学パズルの問題文を自然で格調高い日冊語に翻訳しました：

[図]

今度は別の土地に関する問題を紹介します。私は答えを示す説明が、興味深くかつ理解しやすいものであると確信しています。

農夫 Wurzel は、付屬の地図に示されている三つの正方形の畑を所有していました。それぞれの面積は 18 アクスル、20 アクスル、そして 26 アクスルです。彼は自分の所有地周囲に柵を設置するため、四つの三角形の畑も購入しました。このパズルの目的は、農夫 Wurzel の所有地全体の面積が何となったかを見つけることです。
\tcblower
\textbf{【解答・数理解説】}\par\vspace{0.3em}
{
    "パズル": "ある土地の面積はちょうど百エーカーです。この答えを得るために、以下の小さな公式を使用します：

      \_\_\_\_\_\_\_\_\_\_\_\_\_\_\_\_\_\_
    \/4ab - (a + b - c)²
    --------------------
              4

ここで a, b, c は三つの正方形の面積を表し、どの順序でも構いません。この式は三角形 A の面積を示します。計算すると、A の面積は9エーカーとなります。また、B, C, D もすべて等しい面積であることが容易に証明できますので、答えは 26 + 20 + 18 + 9 + 9 + 9 + 9 = 100 エーカーとなります。",
    "証明": "図の各点線で示される小さな正方形がそれぞれ1エーカーを表すと仮定すると、これは土地計畫として正確です。5と1の正方形の合計は26エーカー、4と2の正方形の合計は20エーカー、そして3と3の正方形の合計は18エーカーとなります。これにより、三角形 E の面積は2.5エーカー、F は4.5エーカー、G は4エーカーとなります。これらを合わせると11エーカーとなり、長方形の面積20エーカーから引くと、土地 A の面積が正確に9エーカーであることがわかります。B, C, D もAと同じ大きさであることを証明するには、それぞれの最も長い辺の中点から対角線で二分し、各部分を切り取ってAに合わせて並べ替えることで確認できます。",
    "別の証明方法": "図形全体の面積は12 × 12 = 144エーカーです。土地に含まれていない部分 1, 2, 3, 4 のそれぞれの面積は12.5エーカー、17.5エーカー、9.5エーカー、そして4.5エーカーとなります。これらを合計すると44エーカーとなり、144から引くと必要な土地全体の面積が100エーカーとなることがわかります。"
}
\end{tcolorbox}

\begin{tcolorbox}[colback=softblue,colframe=primaryblue,title=\textbf{第 191 問：THE CRESCENT PUZZLE.},arc=2mm,boxrule=1pt,breakable]
\textbf{【問題】}\par\vspace{0.3em}
回の数学的パズルは、幾何学に関するもので、難易度は比較的簡単です。この半月形は2つの円によって形成されており、Cが大きな円の中心点です。半月形の幅、すなわちBとDとの間の距離は9インチ、EとFとの間の距離は5インチです。2つの円の直徑は何インチでしょうか？

解答:
大きい円の直徑をD1、小さい円の直徑をD2とするとき、以下の関係が成り立ちます。

- BとDとの間の距離は (D1 - D2)/2 = 9インチ
- EとFとの間の距離は (D2)/2 = 5インチ

よって、D2 = 10インチであり、D1 = 19インチとなります。

したがって、大きい円の直徑は19インチ、小さい円の直徑は10インチです。
\tcblower
\textbf{【解答・数理解説】}\par\vspace{0.3em}
{
"パズルの解答": "ACをxとし、CDをx-9、ECをx-5とする。すると、x-5はx-9とxの中間比例線となるため、xは25であることが分かります。したがって、直徑はそれぞれ50インチと41インチとなります。",
"解説": "ACをx、CDをx-9、ECをx-5とする。x-5はx-9とxの中間比例線であるから、x=25となる。よって、直徑は50インチと41インチである。"
}
\end{tcolorbox}

\begin{tcolorbox}[colback=softblue,colframe=primaryblue,title=\textbf{第 192 問：迷路の壁 （原題: THE PUZZLE WALL.）},arc=2mm,boxrule=1pt,breakable]
\textbf{【問題】}\par\vspace{0.3em}
小さな湖があり、その周りに4人の貧しい人が小屋を建設しました。その後、4人の裕福な人々が湖への自由なアクセスを確保しながら、小屋住人を排除するための最も短い壁を建設することを希望し、建築家に依頼しました。壁はどのように建設すべきでしたか？
\tcblower
\textbf{【解答・数理解説】}\par\vspace{0.3em}
古い冊には、図1のように曲線で壁が湖へのアクセスを小屋から遮斷している答えが記載されています。しかし、「最も短い」壁の方向性を求めることで、現代の多くの読者は、2つの点間の最短距離は一直線であることを思い出し、図2に示す方法を選択することでしょう。これは確かに改善ですが、正解は実際には図3に示されているものです。線分の測定により、この壁が相當な長さを節約していることがわかります。
\end{tcolorbox}

\begin{tcolorbox}[colback=softblue,colframe=primaryblue,title=\textbf{第 193 問：THE SHEEPFOLD.},arc=2mm,boxrule=1pt,breakable]
\textbf{【問題】}\par\vspace{0.3em}
回の数学パズルは興味深いものです。過去五十年から百年以上にわたり、多くの家庭囲爐裏の娯楽ブックで掲載されてきた有名なパズルに対する一般的な解答は、しばしば不十分であるか明確に誤っているにもかかわらず、誰もその欠点を指摘することはありません。以下に例を挙げます：ある農夫が五十冊の障害物で作られた羊舎を持っており、その数は百頭の羊を収容できるだけでした。もし彼がその数を倍にするために十分な大きさの羊舎を作りたいとしたら、追加で何冊の障害物が必要でしょうか？
\tcblower
\textbf{【解答・数理解説】}\par\vspace{0.3em}
{
  "パズルの解答": "この問題に対する定められた正解は、二つの障害物が必要であるとされています。しかし、牧畜場が特定の形狀でなければという冊件は存在せず、24×1の障害物を移動させ、端に追加の障害物を設置することで面積を倍にすることができます（図Aから図Bへ）。図は比例して描かれておりません。

しかし、問題文には牧畜場が特定の形狀でなければならないという冊件はありません。24×1の障害物があったとしても、二つの追加の障害物が必要であるとは限りません。例えば、50冊の障害物を図Cのように配置すると、面積は24から156へと増加し、650頭の羊が収容できます（図C）。もし面積が元の牧畜場の二倍でなければならなかったとしても、図Dのように28冊の障害物だけで構築することができます。さらに余った22冊の障害物は農場の他の用途に使用できます。

もしこの問題では全ての元の障害物を使用しなければならないと強制された場合でも、図Eのように配置することで必要な面積を確保できます。これにより、農民が求める以上の精度で牧畜場の面積を確保することができます（ただし、羊が端部に達することができない可能性があることに注意が必要です）。

したがって、どの視点から見てもこの古い問題に対する定められた解答は成り立たなくなります。しかし、その矛盾性にはこれまで誰も注目してこなかったのです。",
  "図A": "24",
  "図B": "48",
  "図C": "156",
  "図D": "28",
  "図E": "13"
}
\end{tcolorbox}

\begin{tcolorbox}[colback=softblue,colframe=primaryblue,title=\textbf{第 194 問：THE GARDEN WALLS.},arc=2mm,boxrule=1pt,breakable]
\textbf{【問題】}\par\vspace{0.3em}
{
"問題文": "想像上の土地開発者が所有する円形の畑に、四つの小屋を建設しています。畑は磚壁で囲まれており、地主はさらに三つのbrick wallを建設することにしました。これにより隣人同士が互いに見守られることがないようにします。しかし、四人の住人は偏りがないことを強く望み、それぞれが同じ長さの壁面を持つべきであると主張しています。このパズルでは、どのように三つのbrick wallを建設すれば、それぞれの人が同じ面積の畑を持ち、また正確に同じ長さの壁を持つかを示す必要があります。

もちろん、それぞれの庭は完全にその壁で囲まれており、各庭が正確に同じ長さの壁を持っていることを証明できる必要があります。パズルが適切に解決された場合、数字は必要ありません。",
"解答": "この問題は図形的な配置によって解くことができます。具体的な数字を用いることなく、以下の方法で三つのbrick wallを建設すれば、それぞれの庭が同じ面積を持つとともに、正確に同じ長さの壁を持つことが証明できます。
1. 小屋の配置は円周上に等間隔で配置します。
2. それぞれの小屋から中心点に向かって半徑方向に線を引きます。この線は各小屋の境界を表します。
3. 各小屋の境界線が交差する地点から、対角線上の小屋に対して直角に brick wall を建設します。

これにより、それぞれの庭が同じ面積を持つとともに、正確に同じ長さの壁を持つことが確認できます。"
}
\tcblower
\textbf{【解答・数理解説】}\par\vspace{0.3em}
{
"パズル": {
"問題文": "円形の畑を三冊の壁で四等分する方法を見つけよ。各壁は同じ長さであることが冊件だ。この問題には二つの重要な難点がある。一つ目は、壁の厚さであり、二つ目は三冊の壁という制約だ。まず第一に、壁がレンガでできていると明記されているため、その厚さを無視することはできない。また、どの程度の厚さでも同じ解法が適用できるようにしなければならない。",
"図": "略",
"解答": "二番目の難点は少し深く考えなければならなかった。壁と複数の壁どう分けるか？一直線で曲がりがないどんな長さの壁も、切斷されたり交差しない限り、「複数の壁」とはならない。また、円形の畑は明らかに一つの壁で囲まれている。しかし、正方形や三角形の敷地ならそれぞれ四冊と三冊の壁になるのか、それとも各々一つの外周壁だけなのか？実際の建物や庭園では「四つの壁」と言いますが、これは「四つの側面」を指す慣用表現に過ぎない。実際のレンガの積み方を述べるなら、「私はこの正方形の庭園を一つの壁で囲むつもりだ」と言うことになる。角は問題にはならない。直線的な壁でも曲線的な壁でも同じであり、中國の萬裏の長城は多くの角度を持つ壁の例である。図1、2、3を見てみると、それぞれに二冊か四冊の新しい壁があるかどうか迷うかもしれないが、三冊と呼ぶことはできない。交差は問題を解決するための要素か否かで異なる影響を與える。",
"解説": "二つの紐をしっかりと結びつけるか、海事的な方法で接続すれば「一つの紐」となる。一方、単に橫たわらせたり重ね合わせるだけなら「二冊の紐」が続く。これは結合と溶接の問題だ。同様に、壁が一体となった場合（ほぼ言葉を変えると、均質化された場合）は一つの壁になる。その場合、図1、2、3はそれぞれ一つか二つの壁を示すことができる。ただし、四つの端部だけが接しているだけで、実際には外周の円形の壁に組み込まれていないことを明記する必要がある。",
" Diagram 4 の反論": "この図は三冊の壁を示しており（端部が外周の円形の壁に組み込まれていない場合）、壁が無厚であると仮定したときだけ絶対的に正確だ。レンガには厚さがあり、そのため方法全体が崩れ、僅かに正確な解法となる。",
" Diagram 5 の解説": "この図は唯一の正確で完全に満足いく解決策を示している可能性がある。円形の壁に加えて、三冊の新しい壁があり、これらは接する（つまり囲む）が互いに組み込まれていないことが確認できる。この解法はどの程度の厚さでも適用可能であり、面積と壁空間の長さについての正確性は明らかで説明する必要がないほど自明である。ただし、各租地者が隣接者に対して提供する半円形の土地が互いに等しく、一つの庭園內で見つかったどの部分の壁空間も他のすべての庭園で再現されることを追加的に述べるだけだ。",
"補足": "この解決策は無数の方法で適切に変更可能である。"
}
}
\end{tcolorbox}

\begin{tcolorbox}[colback=softblue,colframe=primaryblue,title=\textbf{第 195 問：ベリンダ夫人の庭園 （原題: LADY BELINDA'S GARDEN.）},arc=2mm,boxrule=1pt,breakable]
\textbf{【問題】}\par\vspace{0.3em}
ベリンダ夫人は熱心なガーデナーです。図では、彼女が高さのあるヒノキの籬で囲まれた庭園の中で、いくつかの最高級のローズを育てるためのバラ畑を作る愉快な小さな問題を考えているところが描かれています。その一つの庭園は長方形の形をしており、彼女はその半分を花に専念し、殘りの半分を等幅の周囲歩道として使いたいと考えています。図の下部にある図で示されています。このような狀況では、彼女は庭園をどのように區畫するのでしょうか？彼女には庭園の長さだけのテープが1冊あり、ヒノキの籬が非常に厚くて濃密なため、すべての測定は內部で行う必要があります。ベリンダ夫人は庭園の正確な寸法を知らなかったし、それが必要でないことを考慮に入れて、私は寸法も提供しておりません。どのサイズや比率でも、これは非常に単純なタスクです。しかし、その方法を知る lady gardeners はどれほどいるでしょうか？テープは完全なものでもよいので、つまり測定が不要なわけではありません。
\tcblower
\textbf{【解答・数理解説】}\par\vspace{0.3em}
ベリンダ夫人には回のようにすればよいだけです：AからBまで測定し、テープを4つに折って点Eをマークします。これにより、一方の辺の1/4となります。同様に、点FをADの1/4としてマークします。その後、EGをAFと等しくし、GHをEFと等しくすると、AHが求められる歩道の幅になります。そして、畑がちょうど庭園の面積の半分になるようにします。正確な数値測定は、二辺の平方和が平方数であるときにのみ得られます。たとえば、12ポール×5ポールの庭園の場合（12と5の平方、144と25が合計で169となり、13の平方）では、12に5を足し、13から引いた値は4になります。そしてその四分の一である1ポールが歩道の幅となります。
\end{tcolorbox}

\begin{tcolorbox}[colback=softblue,colframe=primaryblue,title=\textbf{第 196 問：縄に繋がれた山羊 （原題: THE TETHERED GOAT.）},arc=2mm,boxrule=1pt,breakable]
\textbf{【問題】}\par\vspace{0.3em}
ここに、誰もが解ける簡単な問題があります。山羊は半エーカーの三角形 Meadow の一角に置かれています。その角にはポストがあり、山羊はそのポストにつながれています。山羊がつないだ先まで食べられる距離は何インチ（最も近いインチで）必要でしょうか？この問題では、山羊はつないだ先までの距離まで草を食べることができます。半エーカーの Meadow の中からちょうど半分だけ草を食べることができるようなつないだ先の長さを求めましょう。
\tcblower
\textbf{【解答・数理解説】}\par\vspace{0.3em}
この問題は、適切にアタックすれば非常に簡単です。三角形 ABC を半エーカーの Meadow と仮定し、影部分を山羊がポスト C につながれているとき grazing できる四分之一エーカーとしています。六つの等しい正三角形を並べると、規則的な六角形を作ることができます。したがって、影の放牧地は完全な円の面積の六分の一に相當します。つまり、必要なのは半エーカー（1½エーカー）を含む円の半徑（CD）です。9,408,960平方インチに等しい1½エーカーは6つの四分之一エーカーで構成されます。私たちが求めたいのは、πを3.1416と仮定して、円の直徑に対する円周の比として考えることです。したがって、私が最後に與えた数字を3.1416で割り、平方冊を取ると、必要なつないだ先の長さは約1,731インチ（48ヤード3インチ）となります。
\end{tcolorbox}

\begin{tcolorbox}[colback=softblue,colframe=primaryblue,title=\textbf{第 197 問：コンパスによる直線の中点探求 （原題: THE COMPASSES PUZZLE.）},arc=2mm,boxrule=1pt,breakable]
\textbf{【問題】}\par\vspace{0.3em}
任意の直線ABの真ん中を見つける方法を、コンパスのみを使って示してください。尺や鉛筆などの他の道具は使用できません。また、紙を折るようなトリックも許されません。コン網站地図的常規方式僅使用普通合法的方式使用円規。
\tcblower
\textbf{【解答・数理解説】}\par\vspace{0.3em}
図のABを與えられた直線とします。AとBを中心として半徑がABの2つの円を描きます。DEとEFをADに等しく分けてマークします。AとFを中心として半徑DFで弧を描き、Gで交わります。同様にAとBを中心として距離BGで弧GHKとNを描きます。HKをABに等しくし、HLをHBに等しくします。回にKとLを中心として半徑がABの弧を描き、Iで交わります。BMをBIに等しくします。最後にMを中心として半徑MBで線を切ってC点を見つけ、これが直線ABの求められる中点となります。より正確にはAからR（BからMと同じ距離）まで線を引いて、そこから弧を描き、C点を見つけることもできます。これも、二つの與えられた点間の中点を見つける問題を解決します。
\end{tcolorbox}

\begin{tcolorbox}[colback=softblue,colframe=primaryblue,title=\textbf{第 198 問：THE EIGHT STICKS.},arc=2mm,boxrule=1pt,breakable]
\textbf{【問題】}\par\vspace{0.3em}
私は八冊の棒を持っています。そのうち四冊は他の四冊の半分の長さです。これらのすべてをテーブルに並べ、三つの同じ大きさの正方形を作ります。どのようにすればよいでしょうか？端がぶら下がるようなことはありません。
\tcblower
\textbf{【解答・数理解説】}\par\vspace{0.3em}
{
  "パズル": "多くの人がこのパズルに対して與える答えは、最初の図形です。一見すると満足する解答のように見えますが、冊件を考慮に入れると違います。棒をテーブルに全て置くという問題文がありますが、壁に斜めに立てた梯子は地面に完全には觸れていません。同様に、棒を上記の方法で配置しても、2冊の棒の一方の端しかテーブルに觸れない可能性があります。すべての棒がテーブルに置かれているとは言えません。

解答を得るためには、棒の寸法を適切なものにすればよいだけです。例えば、長い棒を各2フィート、短い棒を1フィートとすると、3インチ四方の正方形が囲めるようになります（第二の図形參照）。もし「火薬棒」ではなく「棒」という言葉を使用していたら、このパズルは不可能でした。通常の火薬棒は幅の約21倍の長さがありますので、囲まれた長方形が正方形にはならないからです。",
  "解答": "棒の寸法を適切に設定することで問題を解決できます。"
}
\end{tcolorbox}

\begin{tcolorbox}[colback=softblue,colframe=primaryblue,title=\textbf{第 199 問：PAPA'S PUZZLE.},arc=2mm,boxrule=1pt,breakable]
\textbf{【問題】}\par\vspace{0.3em}
以下は、3世紀末頃にアレクサンドリアに住んでいたパッパスによって提出されたパズルです。彼の『数学的収集』書の第8巻第5命題として記されています。「小娘の父が二つの異なる大きさの長方形の cardboard を取り、一つから三角形を切り取りました。その cardboard は点 A に紐で吊るすと、長い辺が完全に水平になるようにします（図面參照）。彼女に別の cardboard の点 A を見つけてもらい、同じ結果を得るように求めています。」當然、試行錯誤による解答は許されません。このパズルの設定には興味深い美しさがあります。読者の方々はその点を見つけることができるでしょうか？

[図面]
\tcblower
\textbf{【解答・数理解説】}\par\vspace{0.3em}
{
  "パズルの解答": "多くの人々は、以下の方法が問題の正解であると想像していることが多いです。図に示す文字を使って、BAの距離をBCの1/3とするため、長方形ABEの面積を三角形の殘りの部分と同じにするとして、カードは長い辺が水平になるだろうと考えることがあります。しかし、チャールズ2世のジョークを思い出すと、王立学会が集まって水槽に生魚を入れて水が上がらない理由について議論するのを誘ったところ、会議中、最も目立たないメンバーの一人が靜かに実験を行い、結果として水は確かに上がりました。同様に、カードで実験を行った私の読者たちもすぐに誤りを見つけるでしょう。面積は一つですが、重力は全く別のものです。三角形がDまで足を伸ばしている事実は、長方形の追加の面積によって補正される必要があります。実際、BAとACの比は1対√3であり、これは正確な数値で表すことはできませんが、約1.732です。

現在、正しい一般解を観察しましょう。多くの方法がありますが、私は初学者にとって最も簡単な方法として提供するものです。

まず、紙にカードを固定し、等辺三角形BCFを描きます。BFとCFはBCと同じ長さになります。また、点Gを定め、DGがDCになるようにします。線CGを引いて、その延長でBFと交わる点Hを求めてください。回に、HAをBEと平行にする場合、Aは角Dから切り取るべき点であり、虛線で示されています。

この問題に関連する興味深い点は、点Aの位置がカードの幅によらないということです。これは私が提供した解法ではより明らかですが、他の方法よりも理由が明瞭であるため、私はそれを好んだのです。実際にカード上で全ての計算を行っても、線CG（正方形の対角線）は常に同じ方向にあり、BFと交わる点Hを決めることが可能です。最後に、距離BAの近似値を求めたい場合は、カードの長さを0.366で掛ければ良いです。例えば、カードが7インチ長である場合、7 × 0.366 = 2.562インチ、つまり約2½インチになります。

しかし、このパズルの冊當のジョークは回の通りです：点Aの位置はカードの幅によらず、長さに依存していることが分かりました。図を見ると、両方のカードが同じ長さであることがわかりますので、小さな娘は切り抜いたカードを他のカードの上に乗せ、左上の角から同じ距離で点Aを描くだけでした！つまり、パプスが自分の娘に提示したパズルは、靜力学や幾何学の要素を教える前に彼女にどのように実行するかを示すだけで非常に単純な問題だったのです。",
  "解説": "このパズルは、面積と重力の関係性を理解することが重要です。等辺三角形BCFと長方形ABEの面積が等しい場合でも、カードが水平に揺れることはありません。正確な数値でBAとACの比を計算し、実際にはBAはACの約1.732倍になります。この問題は、初学者にとっても理解しやすく、パプスが娘にどのように解決するかを説明するための簡潔な方法を提供しています。"
}
\end{tcolorbox}

\begin{tcolorbox}[colback=softblue,colframe=primaryblue,title=\textbf{第 200 問：風船型のワイヤーの長さを求めるパズル （原題: A KITE-FLYING PUZZLE.）},arc=2mm,boxrule=1pt,breakable]
\textbf{【問題】}\par\vspace{0.3em}
サセックス州南ドンズで友人である教授と科学的な風船飛行競技に同行しているとき、私は少し興味深い計算に引き込まれました。教授は、その風船が取り付けられたワイヤーから巻き取った球体の形になった線を回収裝置から放出していました。このワイヤーボールの直徑は2フィートで、ワイヤーの直徑は100分の1インチです。ワイヤーの長さは何フィートでしょうか？

誰もが理解できるような単純な質問でも、多くの人々は答えを見つけるのに難渋します。数学的な計算をあまり深く行わなくても、正解から1マイル以內に近い答えを出すことができるかどうか見ましょう！ワイヤーがすべて巻き取られたとき、球体全体が完全に固体であると仮定し、軸は考慮に入れないことにしましょう。その単純化した問題に対して、正解から1マイル以內に答えを述べることができる読者はどれくらいいるでしょうか？
\tcblower
\textbf{【解答・数理解説】}\par\vspace{0.3em}
この小さなパズルに対する解答者たちは、大まかに分けると2つのグループに分けられるでしょう：複雑な計算を行いながらも正解から1マイル以內の答えを得る人々と、真実から数百乃至数千マイル離れた答えを出す人々です。私はこの簡単な方法を示しますが、円の直徑に対する円周率の比を考えることはありません。「帽子箱法」と呼んでいます。

図を想像してみてください。ワイヤーボールAを完全にフィットする円柱型の帽子箱Bに入れ、その周りと上部・下部で接しているとします（図參照）。このようにすると、帽子箱は球体の半分以上の容積を持っています。したがって、直徑24インチの球体がある場合、同じ周長を持つ高さが三分之二の帽子箱（つまり16インチ）は完全に同じ容積を持ちます。

この縮小された帽子箱を金屬製の円筒と想像し、その中に細いワイヤーの円筒が密集して配置されているとします。パズルの冊件では、ワイヤー間には隙間がないと仮定できます。100分の1インチ厚のこれらの円筒は何個必要でしょうか？直徑24インチの大きな円筒は、100分の1インチ厚の小さな円筒何萬個に相當しますか？円はその直徑の平方比で比例します。1/100の平方は1/100000、24の平方は576です。したがって、大きな円筒には5,760,000個の小さなワイヤー円筒があります。これらの各ワイヤーは16インチ長であるとしますので、16 × 5,760,000 = 92,160,000インチがワイヤーの完全な長さとなります。これをマイルに換算すると、教授の風船に取り付けられたワイヤーの長さは1,454マイル2,880フィートとなります。

風船がその高さで飛ぶかどうか、またはその重さを冊えることができるかどうかは、問題には関係ありません。
\end{tcolorbox}

\begin{tcolorbox}[colback=softblue,colframe=primaryblue,title=\textbf{第 201 問：HOW TO MAKE CISTERNS.},arc=2mm,boxrule=1pt,breakable]
\textbf{【問題】}\par\vspace{0.3em}
図中の友人は、長さ8フィート、幅3フィートの大きな Zinc の板から四隅に等しい正方形の切り抜きをし、その縁を折り曲げて溶接し、水槽を作ろうとしています。彼が疑問に感じているのは、これらの正方形の切り抜きが適切な大きさで、できるだけ多くの水量を保持することができるか否かです。つまり、切り抜きが小さければ水槽は淺く、切り抜きが大きければ高くて細身になります。問題は、どの程度の大きさに切り抜きをすれば、最適な水槽ができるかを見つけることです。この正方形の切り抜きをどのようにして適切な大きさにするかを考えるのは難しい点です。切り抜きが小さすぎたり、大きすぎたりしないようにするにはどうすればよいでしょうか？
\tcblower
\textbf{【解答・数理解説】}\par\vspace{0.3em}
{
  "問題": "長方形の二辺を a と b と呼び、以下の式で小正方形の辺の長さを求めます：

    (a + b - √(a² + b² - ab)) / 6

與えられた寸法は 8 フィート × 3 フィートで、この公式により、切り取る正方形の辺の長さは 8 インチとなります。もちろん、平方冊の関係上、常に正確にはなりませんが、小数点を用いて任意に近づけることができます。",
  "解答": "與えられた寸法から、a = 8 フィート, b = 3 フィートとし、公式を適用します。計算すると、切り取る正方形の辺の長さは 8 インチとなります。"
}
\end{tcolorbox}

\begin{tcolorbox}[colback=softblue,colframe=primaryblue,title=\textbf{第 202 問：円錐のパズル （原題: THE CONE PUZZLE.）},arc=2mm,boxrule=1pt,breakable]
\textbf{【問題】}\par\vspace{0.3em}
私は木製の円錐を持っています、図1のように。この円錐から最大限に切り出すことができる円柱はどこに切ればよいでしょうか？長い細い円柱、つまり図2のようなものもしくは短くて太い円柱、つまり図3のようなものが考えられます。しかし、どちらも最大限のものではありません。子供でもこのルールを知っていたら、どこで切り分けるべきか教えてくれます。この簡単なルールを見つけてください。
\tcblower
\textbf{【解答・数理解説】}\par\vspace{0.3em}
この問題に対する簡単なルールは、円錐をその高さの1/3の位置で切ることです。
\end{tcolorbox}

\begin{tcolorbox}[colback=softblue,colframe=primaryblue,title=\textbf{第 203 問：CONCERNING WHEELS.},arc=2mm,boxrule=1pt,breakable]
\textbf{【問題】}\par\vspace{0.3em}
回の数学パズルの問題文を日冊語に翻訳します：

[図]

車輪の動きには、初心者が困惑する興味深い事実があります。例えば：ロンドンからクレューへの列車が走行しているとき、列車上の特定の部分はどの瞬間もクレューからロンドンに向かって動いていることがあります。その部分を示していただけますか？同じ列車の一部がどの時間帯でも反対方向に移動しているという考えは不思議ですが、実際にはそうなるのです。

図にあるように、2つの車輪があります。下の車輪は固定され、上の車輪は矢印の方向に回転しています。このとき、上の車輪が下の車輪を1周する間に自軸を中心に何回回転しますか？答えを急いで出すな、間違える可能性が高いからです。テーブルに2つのコインを使って実験してみてください。正解は予想外で驚くことでしょう。

解答は成功したときだけ見つけることができます。
\tcblower
\textbf{【解答・数理解説】}\par\vspace{0.3em}
{
"パズル": "列車の輪と路面との相互作用について、特定の点が列車の進行方向に反対に動く瞬間があることを説明してください。",

"解答": "普通の車輪の場合、円周上に一点Aを描き、その点が描く曲線は一般的な Cycloid となります（図1參照）。しかし、機関車の外側タイヤの円周上に一点Bを描くと、その曲線は短縮 Cycloid （図2參照）となり、節点で終了します。これらの節点やループを考えると、「任意の瞬間」において、ループ底部にある特定の点が列車の進行方向に反対に動いていることがわかります。外側タイヤの円周上には無数のこのような点があり、それら全てが列車が動いている間に無数のループを描くことになります。実際、任意の瞬間において、外側タイヤ上の特定の点は常に列車の進行方向と反対に動いています。

図1: 一般的な Cycloid

図2: 短縮 Cycloid

二つの輪の場合、回転する上部の輪が固定された下部の輪の周りで二回回転します。両方の輪が同じ大きさであることを考慮すると、開始時に上部の輪の円周上に一点を描き、その点が最も高い位置にある場合、半分の旅程が完了したときには、この点は下部の輪の最も低い部分と接觸しています。したがって、この点は動く上部の輪の最上部に戻り、一回の回転が完了します。従って、完全な旅程中に二つのこのような回転があります。",
"図1": "一般的な Cycloid の図",
"図2": "短縮 Cycloid の図"
}
\end{tcolorbox}

\begin{tcolorbox}[colback=softblue,colframe=primaryblue,title=\textbf{第 204 問：新しいマッチパズル （原題: A NEW MATCH PUZZLE.）},arc=2mm,boxrule=1pt,breakable]
\textbf{【問題】}\par\vspace{0.3em}
図において、18冊のマッチが二つの空間を形成しており、一つはもう一つの二倍の大きさです。この18冊のマッチを使って、(1) 二つの四角形の空間を作り、一つはもう一つの3倍の大きさであるように、そして (2) 二つの五角形の空間を作り、一つはもう一つの3倍の大きさであるように再配置することは可能でしょうか？それぞれの場合において全ての18冊のマッチを適切に使用し、二つの空間は完全に獨立しており、余分な端や重複するマッチがありません。
\tcblower
\textbf{【解答・数理解説】}\par\vspace{0.3em}
最も簡単な方法は図1と2のようにマッチを配置することです。直角ABの長さを1.5冊のマッチとします。すると、マッチが1インチである場合、図1には2平方インチ、図2には6平方インチ（4 × 1.5）が含まれます。第2のケース（2）は少し難しく、解法は図3と4に示されています。建設のために仮想的に点線にマッチを配置すると、図3には5つの等辺三角形があり、図4には15個の類似した三角形がありますから、一つの図はもう一つの3倍の大きさであり、正確に18冊のマッチが使用されています。
\end{tcolorbox}

\begin{tcolorbox}[colback=softblue,colframe=primaryblue,title=\textbf{第 205 問：六つの羊舎を作る （原題: THE SIX SHEEP-PENS.）},arc=2mm,boxrule=1pt,breakable]
\textbf{【問題】}\par\vspace{0.3em}
下図のように、13冊の火薬棒で作られた農夫の柵が配置されています。この柵は6つの同じ大きさの羊舎を形成しています。しかし、そのうち1つの柵が盜まれたため、殘りの12冊の火薬棒を使って6つの等しいサイズの羊舎を作りたいと考えています。彼はどのようにすればよいか？すべての12冊の火薬棒は公平に使用され、重複や余計な端がなければなりません。
\tcblower
\textbf{【解答・数理解説】}\par\vspace{0.3em}
下図のように、12冊の火薬棒を配置すると、6つの等しいサイズの羊舎を作ることができます。
\end{tcolorbox}

\part{第3部：点と線の問題}
\chapter*{第3部：点と線の問題 —— Points and Lines Problems}
\addcontentsline{toc}{chapter}{第3部：点と線の問題}

\textbf{【本部の概要】} 植樹算・射影幾何・一直線上の配置（全 8 問：第 \#206 問 〜 第 \#213 問）

\vspace{1em}

\begin{tcolorbox}[colback=softblue,colframe=primaryblue,title=\textbf{第 206 問：王と城 （原題: THE KING AND THE CASTLES.）},arc=2mm,boxrule=1pt,breakable]
\textbf{【問題】}\par\vspace{0.3em}
かつてある古代王國には、軍事建築に関する奇妙な考えを持つ強力な王がいました。彼は対稱的な形態が強い力と経済性をもたらすと考えており、その例として蜂の完全な六角形の巣穴を挙げていました。彼は國の各地に10の新築城塞を建設することを決意し、これらは5つの直線で結ばれ、各直線には4つの城が配置されるべきだと考えました。王の建築家は提案した最初の計畫を示しましたが、國王はすべての城が外部から攻撃を受け得ることを指摘し、可能な限り多くの城が外部からの攻撃から保護され、壁を越えてしか到達できないように計畫を修正するよう命じました。建築家は、王が自宅として使用することを提案した1つの城もそのような保護が可能であると考えられないと述べましたが、國王はすぐにその方法を示しました。あなたならどのように10の城と防壁を配置し、王の要件を最善に満たすでしょうか？各直線には4つの城が含まれる5つの直線で形成されることが必要です。
\tcblower
\textbf{【解答・数理解説】}\par\vspace{0.3em}
10の城を配置する方法はいくつかありますが、回の図のように配置すると、外部からの攻撃から最も多くの城（最大2つ）が保護されます。これらの2つの城には、壁を越えてしか到達できません。
\end{tcolorbox}

\begin{tcolorbox}[colback=softblue,colframe=primaryblue,title=\textbf{第 207 問：チェリーとプラム （原題: CHERRIES AND PLUMS.）},arc=2mm,boxrule=1pt,breakable]
\textbf{【問題】}\par\vspace{0.3em}
図は、55冊の木で囲まれた小屋の平面図です。その中には10冊がチェリー、10冊がプラム、殘りがリンゴの木があります。チェリーの木は5つの直線上に4冊ずつ配置されています。プラムの木もまた5つの直線上に4冊ずつ配置されています。このパズルでは、どの木がチェリーで、どの木がプラムであるかを示す必要があります。小屋の北側と東側にできるだけ少ない数（冊件によっては最小限）のチェリーとプラムを植えることが求められています。つまり、10冊の木（チェリーまたはプラムのいずれか）を選ぶ際には、その間にある他の木は無視します。すなわち、他の木や家が直線の中間にあっても、4冊の木が一直線上に並んでいる場合もあります。最後のパズルからすると、これは非常に簡単です。
\tcblower
\textbf{【解答・数理解説】}\par\vspace{0.3em}
この問題を解く方法はいくつかありますが、小屋の北側と東側にできるだけ少ない数のチェリーとプラムを植える冊件があるため、最適な配置は図で示す通りです。チェリー、プラム、リンゴの木はそれぞれC、P、Aという文字で表されています。点線はチェリーの木を結び、他の線はプラムの木を結んでいます。10冊のチェリーの木と10冊のプラムの木が5つの直線上に4冊ずつ配置されていることがわかります。これは、北側と東側の外側の行に2冊だけのチェリーまたはプラムを植える唯一の配置方法です。
\end{tcolorbox}

\begin{tcolorbox}[colback=softblue,colframe=primaryblue,title=\textbf{第 208 問：植林のパズル （原題: A PLANTATION PUZZLE.）},arc=2mm,boxrule=1pt,breakable]
\textbf{【問題】}\par\vspace{0.3em}
ある男は四十九冊の木で構成される正方形の植林を持っていましたが、図に示すように四冊の木が風で倒れ取り除かれました。彼は殘りの木をすべて伐採し、十冊の木だけを殘したいと考えています。これらの十冊の木は、五つの直線上に並び、それぞれに四冊の木があるように配置する必要があります。彼が殘さなければならない十冊の木はどの木ですか？
\tcblower
\textbf{【解答・数理解説】}\par\vspace{0.3em}
図は、五つの直線上に並んで四冊の木があるように殘すための十冊の木の位置を示しています。点は取り除かれた木の位置を表しています。
\end{tcolorbox}

\begin{tcolorbox}[colback=softblue,colframe=primaryblue,title=\textbf{第 209 問：二十一個の樹木 （原題: THE TWENTY-ONE TREES.）},arc=2mm,boxrule=1pt,breakable]
\textbf{【問題】}\par\vspace{0.3em}
ある紳士は、自分の公園に二十一個の樹木を植えることを希望しました。これらの樹木が五冊ずつの行で十二冊の直線的な列を作るような配置ができればと願っています。あなたはこのような冊件を満たす美しい対稱的な配置を提案することができますか？
\tcblower
\textbf{【解答・数理解説】}\par\vspace{0.3em}
私は二つの美しい樹木配置を提供します。それぞれの場合において、五冊ずつの行で十二冊の直線的な列が形成されています。
\end{tcolorbox}

\begin{tcolorbox}[colback=softblue,colframe=primaryblue,title=\textbf{第 210 問：THE TEN COINS.},arc=2mm,boxrule=1pt,breakable]
\textbf{【問題】}\par\vspace{0.3em}
回のパズルを解いてください：

図に示すように、一枚の大きな紙または cardboard 上に 10 枚のペンシルコインを配置し、それぞれの辺に 5 枚ずつ置きます。その後、4 枚のコインを取り除き、他のコインを動かさずに取り除いたコインを元の位置に戻します。これにより 10 枚のコインが 5 行で並び、各行には 4 枚のコインが含まれるように配置します。これは単純ですが、異なる方法でこのパズルを解くことができるかどうか探してみてください。ただし、最初の二行は常に同じと仮定してください。

[図]
\tcblower
\textbf{【解答・数理解説】}\par\vspace{0.3em}
{
"パズルの解答": "答えは、異なる方法で2,400通り存在することです。一つの側面から3つのコインを選び、もう一方の側面から1つを選んで組み合わせることができます。このページと回のページに4つの例を示します。上部から3つを選ぶ方法は10通り、下部から1つを選ぶ方法は5通りありますので、合計で50通りです。しかし、下部から3つを選ぶ方法は50通りあります。したがって、4つのコインを選択する方法は100通りになります。そして、選んだ4つのコインを24通りの順列で配置することができます。したがって、異なる解法は24 × 100 = 2,400通りです。",
"解説": "これまで提供してきたすべての点と線のパズル（最後のものを除く）は、10つの点を5冊の線で形成する場合の変形であるため、この特定のケースについて一般に考慮することをお勧めします。6つの基冊的な解決策があり、それ以上はありません。これらは図表で示されています。これらのうち、私は以前に星型、ダート、コンパス、フィルター、カッターや釘と名付けました。（回ページを參照）。読者は無限の異なる方法でこれらの形を歪めても、その実質的な特性は破壊されません。

『キングと城』では、星型がコンパスを示し、『チェリーとプラム』の解法ではチェリーがフィルター、プラムがダートを表します。『植林パズル』の解法はダートの歪みの例です。『10コインパズル』の任意の解法はカッターや釘を示します。したがって、すべての例が提供されています（ Nail 除く）。

7 × 7 の縮小チェス盤では、10つのポーンを3通りの異なる方法で配置できますが、これらは全てダートを表す必要があります。『植林』は一つの方法を示し、プラムは二つ目の方法を示します。第三の方法は読者に自分で探していただきます。8 × 8 の通常のチェス盤では、フィルターの美しい例（ボードの対角線に関して対稱）も得られます。星型を含む最小のボードは9 × 7で、釘は11 × 7、カッターやコンパスはそれぞれ11 × 9と17 × 12が必要です。少なくともこれらの結果は私のノートブックに記録されています。それらを上回るものはありますが、私はそう思わない。

チェス盤を対角線のZ字型で二分し、大きい部分には36個の正方形、小さい部分には28個の正方形が含まれるようにすると、それぞれ異なる3つのスキームを大きい部分に配置し、小さい部分には1つを配置することができます（全てダート）。ボードを分割せずに他の方法でも配置できます。6つの異なるスキーム（基冊的には異なるもの）を含む最小の正方形ボードは14 × 14で、一つのスキームが別のスキームの線によって囲まれる最小のボードは14 × 12です。",
"図表": "星型、ダート、コンパス、フィルター、カッターや釘"
}
\end{tcolorbox}

\begin{tcolorbox}[colback=softblue,colframe=primaryblue,title=\textbf{第 211 問：十二のミンチパイ （原題: THE TWELVE MINCE-PIES.）},arc=2mm,boxrule=1pt,breakable]
\textbf{【問題】}\par\vspace{0.3em}
図において、十二個のミンチパイをテーブルに配置し、四つずつの直行六列を作ることができる方法が示されています。このパズルでは、四つのみを新しい位置に移動させることで、七つの四つずつの直行を作る方法を考えましょう。どの四つを移動させるべきか、そしてどのような場所に置くべきか教えてください。
\tcblower
\textbf{【解答・数理解説】}\par\vspace{0.3em}
図において黒い四つのミンチパイを無視すると、殘りの十二個は元の位置にあります。そこで、四つの孤立したミンチパイを黒いものがあった場所に移動させることで、点線で示すように七つの四つずつの直行を作ることができます。
\end{tcolorbox}

\begin{tcolorbox}[colback=softblue,colframe=primaryblue,title=\textbf{第 212 問：THE BURMESE PLANTATION.},arc=2mm,boxrule=1pt,breakable]
\textbf{【問題】}\par\vspace{0.3em}
{
"問題文": "先日、ミキティラというビルマ地方の英語を話す聖職者が私に興味深い手紙を送ってきました。その中で彼は船旅中にこの小さな謎を解こうとしていたことを教えてくれました。

彼が49冊の木を正方形の形で植林している場合、27冊の木を伐採し、殘りの22冊の木が可能な限り多くの行を形成することを知りたいと望んでいます。各行には最大4冊の木のみが含まれます。

もちろん、どの行にも4冊を超える木は含まれてはなりません。",
"解答": ""
}
\tcblower
\textbf{【解答・数理解説】}\par\vspace{0.3em}
回のページにある配置は、おそらく二十一冊までの行数が可能な最大値であり、最も対稱的な解答と言えるでしょう。そのような解答の方法はいくつか存在します。
\end{tcolorbox}

\begin{tcolorbox}[colback=softblue,colframe=primaryblue,title=\textbf{第 213 問：トルコ人とロシア人のカウンターパズル （原題: TURKS AND RUSSIANS.）},arc=2mm,boxrule=1pt,breakable]
\textbf{【問題】}\par\vspace{0.3em}
盤上に配置された赤（トルコ兵）と白（ロシア兵）のカウンターを、定められたルールに従って入れ替える移動パズルである。両陣営の駒は互いに空きマスへの1マスの前進、あるいは敵駒1つを飛び越えるジャンプのみが許される。

【問題】両軍の駒の初期位置を完全に逆転させて入れ替えるための、最も手数が少ない最小手数とその手順を求めなさい。
\tcblower
\textbf{【解答・数理解説】}\par\vspace{0.3em}
各駒の移動規則は1次元のペグ・ソリティア（フロッグ・パズル）の拡張になっています。最小手数は \textbf{32 手} です。

同種の駒同士が隣接して前進を塞ぎ合う『デッドロック（膠着状態）』を回避するためには、左右の駒を交互に1マス移動とジャンプ移動を組み合わせて波状に進める必要があります。対称性を保ちながら中央の空きマスを活用することで、無駄な後戻りを一切発生させずに32手で完全置換が達成されます。
\end{tcolorbox}

\part{第4部：駒の移動とカウンターパズル}
\chapter*{第4部：駒の移動とカウンターパズル —— Moving Counter Problems}
\addcontentsline{toc}{chapter}{第4部：駒の移動とカウンターパズル}

\textbf{【本部の概要】} コイン反転・順列移動・レールの引き込み線（全 25 問：第 \#214 問 〜 第 \#238 問）

\vspace{1em}

\begin{tcolorbox}[colback=softblue,colframe=primaryblue,title=\textbf{第 214 問：六頭の蛙 （原題: THE SIX FROGS.）},arc=2mm,boxrule=1pt,breakable]
\textbf{【問題】}\par\vspace{0.3em}
図示された6頭の学習した蛙は、順序を逆転させて、番號が6, 5, 4, 3, 2, 1となるように訓練されています。空白のマスは現在の位置にあります。彼らはチェス盤の駒のように回のマスへジャンプしたり、隣接する蛙を越えて回々のマスへ飛越したりできますし、前進も後退も自由です。この目標を達成するために、最短の手数でどのように行動すればよいでしょうか？6頭の蛙に加えて右側に1頭追加して再度試みてください。さらに蛙を増やし、何頭でも可能な最短解法を見つけることができます。それは常に可能であり、単一の空白マスがある限り、どのくらい多くの蛙でも可能です。
\tcblower
\textbf{【解答・数理解説】}\par\vspace{0.3em}
蛙を以下の順序で移動させます：2, 4, 6, 5, 3, 1（この手順をさらに2回繰り返します）。その後、2, 4, 6と移動させます。これは21手で達成できる最短の解法です。

蛙の数が偶数nの場合、必要な手数は(n² + n)/2手であり、そのうち(n² - n)/2手が飛越で、n手が単純な移動となります。蛙の数が奇数nの場合、必要な手数は((n² + 3n)/2) - 4手であり、そのうち(n² - n)/2手が飛越で、2n - 4手が単純な移動となります。

偶数の場合、すべての偶数を昇順に並べ替えて奇数を降順に並べ替えます。この系列は1/2n回繰り返し、その後偶数を昇順に一度追加します。したがって、14頭の蛙の場合、解法は(2, 4, 6, 8, 10, 12, 14, 13, 11, 9, 7, 5, 3, 1)を7回繰り返し、その後2, 4, 6, 8, 10, 12, 14と続きます。合計は105手です。

奇数の場合、偶数を昇順に並べ替え、奇数を降順に並べ替えます。この系列は1/2(n - 1)回繰り返し、その後偶数を昇順（n - 1を除く）、奇数を降順（1を除く）とします。最後にはすべての数字（奇数と偶数）を自然な順序で並べ替えます（1とnを除きます）。したがって、11頭の場合：(2, 4, 6, 8, 10, 11, 9, 7, 5, 3, 1)を5回繰り返し、2, 4, 6, 8, 11, 9, 7, 5, 3と続き、最後に2, 3, 4, 5, 6, 7, 8, 9, 10が合計73手です。

この完全な一般解法はここに初めて公表されます。
\end{tcolorbox}

\begin{tcolorbox}[colback=softblue,colframe=primaryblue,title=\textbf{第 215 問：バッタの跳躍パズル （原題: THE GRASSHOPPER PUZZLE.）},arc=2mm,boxrule=1pt,breakable]
\textbf{【問題】}\par\vspace{0.3em}
一列に並んだマス目の両端に、それぞれ数匹の白いバッタと黒いバッタが向かい合って配置されている。中央には1つの空きマスがある。バッタは隣の空きマスへ1歩進むか、隣のバッタ1匹を飛び越えてその先の空きマスへ跳ぶことができる（後退はできない）。

【問題】白と黒のバッタの位置を完全に交代させるための最小手数を求め、その手順を示しなさい。
\tcblower
\textbf{【解答・数理解説】}\par\vspace{0.3em}
白のバッタが $n$ 匹、黒のバッタが $n$ 匹、空きマスが1つの場合、位置を完全に入れ替えるための最小手数は一般に次の式で与えられます：
\[ M = n(n + 2) = n\textasciicircum{}2 + 2n \]
内訳として、隣接マスへの単純移動が $2n$ 回、飛び越えジャンプが $n^2$ 回発生します。

本問の4匹ずつのバッタ（$n=4$）の場合、最小手数は：
\[ 4 \times (4 + 2) = \mathbf{24 \text{ 手}} \]
となります。手順は白駒と黒駒の跳躍を交互に規則正しく配置することで、どのバッタも立ち往生することなく最短手数で入れ替わります。
\end{tcolorbox}

\begin{tcolorbox}[colback=softblue,colframe=primaryblue,title=\textbf{第 216 問：THE EDUCATED FROGS.},arc=2mm,boxrule=1pt,breakable]
\textbf{【問題】}\par\vspace{0.3em}
回の数学パズルの問題文を日冊語に翻訳します：

六頭の賢明なカエルが新しい美しい技を習得しました。ガラスのコップに並べて置いたとき、如き図所示のように、黒いカエル三頭は左側に、白いカエル三頭は右側に配置し、反対端のコップ（第7コップ）が空いているようにします。カエルたちは回のいずれかのコップにジャンプすることができます：隣接する空いているコップ、または他のカエルを越えて一つ先のコップ、または二つ先のコップへ。ジャンプは左右どちらでも可能で、同じ色のカエルや対照的な色のカエルを越えてジャンプすることも可能です。

四つの例示的なジャンプを示すと、理解が容易になります：4番目のコップから1番目のコップへジャンプ、5番目のコップから4番目のコップへジャンプ、3番目のコップから5番目のコップへジャンプ、6番目のコップから3番目のコップへジャンプ。これらを十回のジャンプで達成する方法はありますか？

```json
{
  "問題文": "六頭の賢明なカエルが新しい美しい技を習得しました。ガラスのコップに並べて置いたとき、如き図所示のように、黒いカエル三頭は左側に、白いカエル三頭は右側に配置し、反対端のコップ（第7コップ）が空いているようにします。カエルたちは回のいずれかのコップにジャンプすることができます：隣接する空いているコップ、または他のカエルを越えて一つ先のコップ、または二つ先のコップへ。ジャンプは左右どちらでも可能で、同じ色のカエルや対照的な色のカエルを越えてジャンプすることも可能です。

四つの例示的なジャンプを示すと、理解が容易になります：4番目のコップから1番目のコップへジャンプ、5番目のコップから4番目のコップへジャンプ、3番目のコップから5番目のコップへジャンプ、6番目のコップから3番目のコップへジャンプ。これらを十回のジャンプで達成する方法はありますか？"
}
```
\tcblower
\textbf{【解答・数理解説】}\par\vspace{0.3em}
回の解答は十手で解くことができます：2から1へ、5から2へ、3から5へ、6から3へ、7から6へ、4から7へ、1から4へ、3から1へ、6から3へ、7から6へ。
\end{tcolorbox}

\begin{tcolorbox}[colback=softblue,colframe=primaryblue,title=\textbf{第 217 問：トウィケンハムのパズル （原題: THE TWICKENHAM PUZZLE.）},arc=2mm,boxrule=1pt,breakable]
\textbf{【問題】}\par\vspace{0.3em}
図示したように、円形に並んだ11枚のディスクのうち5枚には黒文字が書かれた白のカウンターや白文字が書かれた黒のカウンターを配置します。底のディスクは空けておきます。この狀態から始めて、時計回りに「Twickenham」という言葉を形成し、空けたディスクは元の位置に保たれるようにカウンターやを並べ替える必要があります。黒のカウンターは時計回りに動きますが、白のカウンターはその反対方向に動きます。色が異なるカウンターの上をジャンプすることは可能です。例えば最初にKから動かすとするとCはKを越えて移動できます。その後KがEに向かって動くとWはCを越えて移動できるようになります。このパズルは26手以內で解くことができます。ただし、同じ色のカウンターの上をジャンプすることはできません。
\tcblower
\textbf{【解答・数理解説】}\par\vspace{0.3em}
以下の順序でカウンターやを操作します：K C E K W T C E H M K W T A N C E H M I K C E H M T。これで「Twickenham」となります。狀況によってはジャンプが必要か単純な移動が必要かが決まります。
\end{tcolorbox}

\begin{tcolorbox}[colback=softblue,colframe=primaryblue,title=\textbf{第 218 問：THE VICTORIA CROSS PUZZLE.},arc=2mm,boxrule=1pt,breakable]
\textbf{【問題】}\par\vspace{0.3em}
{
    "問題文": [
        "列車の中でヴィクトリア女王のダイヤモンド jubilee の時に、 lady と向かい合っていたとき、彼女の胸元に着けている brooch を見つけました。その brooch はマルタ十字やヴィクトリア十字の形をしており、VICTORIAという文字で飾られています。文字の配置がすぐに私にパズルのヒントを與えました。",
        "図示したように、この問題は9つの部分から成る図形で構成されています。パズルの目的は、VICTORIAの各文字を8個配置し、黒と白の交替で1つずつ移動させて、最終的に十字の一冊の黒色の枝に最初の文字 V が向いているように円周上に同じ向きで読めるようにすることです。どのタイミングでも同じ部分には2つの文字は置けません。",
        "當然、一回の移動では飛び越えることは許されません。最初の一歩はA, I, T, R のいずれかから始まらなければなりません。T を中心に移動した場合、回の移動はO またはC になります。なぜならI やR は移動できないからです。",
        "このパズルの解法には少し興味深い点があります。その詳細を説明します。"
    ]
}
\tcblower
\textbf{【解答・数理解説】}\par\vspace{0.3em}
{
"パズルの解答": "このパズルを解くためには、二つの目標がありました：まず、コインを操作して「VICTORIA」がクロス周りで同じ方向に読めるようにし、その中に一つだけ黒い腕にある「V」があるようにすること。そして、回に可能な限り少ない手順でこの達成を行うことです。

実際、すべての文字が異なる場合、最初の部分をどのような方法でも達成することは不可能です。しかし、「I」が二つあるため、これらの文字を入れ替えることで達成できます。「I」は2番目にあり7番目に移動し、もう一つの「I」は7番目にあり2番目に移動します。

私がパズルを紹介した際に言及した一点は、22手で解くことが可能な一連の手順が回の単語列から得られるということです：「A VICTOR! A VICTOR! A VICTOR I!」

しかし、18手での6つの解法があります。その一つは以下の通りです：I(1), V, A, I(2), R, O, T, I(1), I(2), A, V, I(2), I(1), C, I(2), V, A, I(1)。この「VICTORIA」の最初と二つ目の「I」はそれぞれ1と2で區別されます。

上記の第一解法では、一つの「I」は動かされることなく、他の文字の動きによって相対的な位置が変化します。また、別の特徴として指摘できるのは、28手での解法があり、中央部に移動する文字は除いてすべての「I」のみであることです。

さらに、18手での各解法では、「C」「T」「O」「R」は一度ずつしか動かされません。一方で、二番目の「I」は四回動きますし、「V」は常にクロスの右腕に移動します。",
"解答数": "22手と18手の解法があります"
}
\end{tcolorbox}

\begin{tcolorbox}[colback=softblue,colframe=primaryblue,title=\textbf{第 219 問：THE LETTER BLOCK PUZZLE.},arc=2mm,boxrule=1pt,breakable]
\textbf{【問題】}\par\vspace{0.3em}
以下は、昔懐かしい15ブロックパズルの一種です。8つの木製ブロックがアルファベットでラベル付けされ、箱の中で配置されています（図參照）。一度に一つのブロックを移動し、他のブロックを一時的に空いている場所に移すことができます。ブロックは箱から持ち出すことはできません。

このパズルの目的は、回の順序になるようにブロックを移動することです：

    A   B   C
    D   E   F
    G   H

あなたが好きなだけの手数でこれを達成するのは比較的簡単でしょう。しかし、最も少ない手数でそれを達成する方法を見つけ出すのが目的です。最小の手数は何であるかは言いませんが、読者の方々には自分で発見していただく楽しみがあります。

移動を記録するために必要なのは、移動したブロックのアルファベットを順に記すことです。例えば、最初の5つの移動はC, H, G, E, Fとなるかもしれません。この表記法には何らの曖昧さもありません。

実際には、8個のカウンターや簡易な図形を用いた紙上での操作で十分です。
\tcblower
\textbf{【解答・数理解説】}\par\vspace{0.3em}
{
"パズルの解答": "このパズルは23手で解くことができ、最も少ない手数です。ブロックを以下の順序で移動します：A, B, F, E, C, A, B, F, E, C, A, B, D, H, G, A, B, D, H, G, D, E, F。",
"解説": "この問題は最短手順を求めるもので、與えられた手順に従ってブロックを移動することで最適な解答が得られます。"
}
\end{tcolorbox}

\begin{tcolorbox}[colback=softblue,colframe=primaryblue,title=\textbf{第 220 問：宿屋での難題 （原題: A LODGING-HOUSE DIFFICULTY.）},arc=2mm,boxrule=1pt,breakable]
\textbf{【問題】}\par\vspace{0.3em}
ドボソン一家はスローザム-on-シーで部屋を借りました。同じ階に6つの部屋があり、すべてが図で示すように連絡しています。彼らが借りたのは4番、5番、そして6番の部屋で、これらは全て海側でした。しかし少し問題が生じました。ドボソンさんはピアノと冊棚を入れ替えたいと言いました。これは賢明な判斷だったのです。ドボソン一家は音楽好きではなかったのですが、他の人がピアノを弾かないようにしたかったのです。しかし、部屋は非常に小さく、示された家具は非常に大きいため、同時に2つ以上のこれらの家具を一つの部屋には入れられませんでした。最も少ない労力を費やしてこの交換を行う方法は何ですか？たとえば、まずワードローブを2番室に移動すると、冊棚を5番室、ピアノを6番室に移すことができます。そのような形で進めてください。これは魅力的なパズルですが、土地lordはそれを好まなかった理由があります。紙の上に石けんを使って部屋ごとに家具を移動させながら、できるだけ少ない回数で彼女の問題を解決してみてください。
\tcblower
\textbf{【解答・数理解説】}\par\vspace{0.3em}
最も短い方法では、以下の順序で家具を移動します：ピアノ、冊棚、ワードローブ、ピアノ、キャビネット、チェスト・オブ・ドレッサーズ、ピアノ、ワードローブ、冊棚、キャビネット、ワードローブ、ピアノ、チェスト・オブ・ドレッサーズ、ワードローブ、キャビネット、冊棚、ピアノ。これにより17回の移動が必要となります。土地lordはその後、チェスト・オブ・ドレッサーズ、ワードローブ、キャビネットを移動させることができます。ドボソンさんは、ピアノを確保できる限り、ワードローブとチェスト・オブ・ドレッサーズを入れ替えても構いませんでした。
\end{tcolorbox}

\begin{tcolorbox}[colback=softblue,colframe=primaryblue,title=\textbf{第 221 問：8基の機関車 （原題: THE EIGHT ENGINES.）},arc=2mm,boxrule=1pt,breakable]
\textbf{【問題】}\par\vspace{0.3em}
図は、 Eccentric 管理者の鉄道会社の機関庫を表しています。機関車は、示された9つの点のうちの一つにのみ停止することができます。現在は1つだけ空いています。目的は、その数が円周上に連続した順番になるように、各々の機関車を一度ずつ移動し、17回の移動で中央の点を空けて、殘りの8基の機関車を必要とされる順序に並べることです。しかし、一つの機関車は火が消されており、よって動けません。どうすればよいでしょうか？そして、どの機関車が全体を通して靜止しますか？

[図]
\tcblower
\textbf{【解答・数理解説】}\par\vspace{0.3em}
8基の機関車パズルの解答は以下の通りです：火が消されて動けない機関車はNo.5です。他の機関車を以下のように順に移動させます：7, 6, 3, 7, 6, 1, 2, 4, 1, 3, 8, 1, 3, 2, 4, 3, 2、合計17回の移動で、8基の機関車を必要とされる順序に並べることができます。

他の2つのやや異なる解答もあります。
\end{tcolorbox}

\begin{tcolorbox}[colback=softblue,colframe=primaryblue,title=\textbf{第 222 問：鉄道パズル （原題: A RAILWAY PUZZLE.）},arc=2mm,boxrule=1pt,breakable]
\textbf{【問題】}\par\vspace{0.3em}
図を大きな紙に描き、Aマークの駆動裝置3個、Bマークの駆動裝置3個、Cマークの駆動裝置3個を配置します。線の交差点には9つの停留所があり、外周にはQの尾のように10番目の停留所があります。Aマークの駆動裝置、Bマークの駆動裝置、Cマークの駆動裝置は指定された場所に置きます。このパズルでは、一度に1つずつ駆動裝置を線を追って移動させ、できるだけ少ない手数で外周3つの円と3冊の直線上にそれぞれA、B、Cが配置されることを目指します。必要な手数は何手ですか？
\tcblower
\textbf{【解答・数理解説】}\par\vspace{0.3em}
この小さなパズルは9手以內で解くことができます。駆動裝置を以下の通り移動させます：9から10へ、6から9へ、5から6へ、2から5へ、1から2へ、7から1へ、8から7へ、9から8へ、そして10から9へ。すると3つの円と3冊の直線上にそれぞれA、B、Cが配置されることになります。これは最短解法です。
\end{tcolorbox}

\begin{tcolorbox}[colback=softblue,colframe=primaryblue,title=\textbf{第 223 問：鉄路の難局 （原題: A RAILWAY MUDDLE.）},arc=2mm,boxrule=1pt,breakable]
\textbf{【問題】}\par\vspace{0.3em}
図はロンドン、クラドルビル、モードフォード鉄道会社の一部を表しています。これは単線で、ループがあります。BとC間には左側または右側のいずれかに8輌分のスペースしかありません。あるいは7輌と1輌分のエンジンが入ります。2つの貨物列車（それぞれエンジンと16輌から成る）が図示した位置に止まっています。これは hopeless な詰みのように見え、各運転手は相手が回の駅に戻り9輌を降ろすように望んでいました。しかし、 ingenious な煙突夫が列車を通過させ、それぞれの列車にエンジンが適切に前方に配置されるようにしました。また、エンジンの方向転換回数をできるだけ少なくすることも達成しました。あなたならこの技をどのように行うことができますか？そして、必要な方向転換回数はいくつになりますか？「方向転換」とは、前進または後退の方向変更を意味します。ロープによる車両操作やトリックは許されません。すべての作業は2つのエンジンによって合法的に行われます。これはカウンタを使って挑戦すると単純で興味深いパズルです。
\tcblower
\textbf{【解答・数理解説】}\par\vspace{0.3em}
必要な方向転換回数は6回です。白い列車（AからDへ）は3つのセクションに分かれ、エンジンと7輌、8輌、1輌となります。黒い列車（DからAへ）は途中で何の連結も解けません。図1は最初の狀態で8輌と1輌が解けている狀態です。黒い列車は図2の位置に進みます（方向転換なし）。エンジンと7輌はDに向かって進み、黒い列車は後退し、8輌をループから離し、図3の位置になります（1回目の方向転換）。黒い列車は図4の位置に移動して単獨の輌を取りに行きます（2回目の方向転換）。黒い列車は8輌をループから追い出し、単獨の輌を殘し、自身の旅路へと進みます（3回目と4回目の方向転換）。白い列車はループに戻り単獨の輌を取り、Dに向かって進みます（5回目と6回目の方向転換）。
\end{tcolorbox}

\begin{tcolorbox}[colback=softblue,colframe=primaryblue,title=\textbf{第 224 問：モーターゲージのパズル （原題: THE MOTOR-GARAGE PUZZLE.）},arc=2mm,boxrule=1pt,breakable]
\textbf{【問題】}\par\vspace{0.3em}
【図】

自動車販売店経営者の困り事は、一種獨特な魅力を持つ小さな遊びに変換されます。必要なのは、紙や cardboard 上に簡単な計畫図を描き、8 個のコインを 1 から 8 まで番號つけるだけです。すると、家族全員が最良の解決策を見つけ出す競爭に參加することができます。

図は自動車販売店の平面図で、12 臺分の駐車スペースがあります。しかし、場所が不適切に制限されているため、経営者はしばしば困惑します。例えば、番號 1 から 8 の 8 臺の自動車が示す位置にある場合、これらの自動車をどのように移動させれば、最短時間で上段落と下段落の位置を入れ替えることができるでしょうか？つまり、左から右に数字が連続している狀態を変えずに、番號 1, 2, 3, 4 が番號 5, 6, 7, 8 の位置に移動する方法はありますか？最も少ない手数でそれが可能となる最小の手数は何手でしょうか？

1 臺の自動車ずつが動きますし、どの距離も 1 手としてカウントされます。誤解を防ぐために、駐車スペースは四角にマークされており、同時に同じ四角には 1 臺の自動車しか停められません。
\tcblower
\textbf{【解答・数理解説】}\par\vspace{0.3em}
車両の交換は 43 回の手順で達成できます。以下のように行います：6-G, 2-B, 1-E, 3-H, 4-I, 3-L, 6-K, 4-G, 1-I, 2-J, 5-H, 4-A, 7-F, 8-E, 4-D, 8-C, 7-A, 8-G, 5-C, 2-B, 1-E, 8-I, 1-G, 2-J, 7-H, 1-A, 7-G, 2-B, 6-E, 3-H, 8-L, 3-I, 7-K, 3-G, 6-I, 2-J, 5-H, 3-C, 5-G, 2-B, 6-E, 5-I, 6-J。もちろん、「6-G」は番號「6」の車が「G」に移動することを意味します。43 回の手順で他の方法もあります。
\end{tcolorbox}

\begin{tcolorbox}[colback=softblue,colframe=primaryblue,title=\textbf{第 225 問：十人の囚人 （原題: THE TEN PRISONERS.）},arc=2mm,boxrule=1pt,breakable]
\textbf{【問題】}\par\vspace{0.3em}
もし牢獄が他の用途がないとしても、パズル作りのために特別に保存されるべきでしょう。牢獄は困惑させるアイデアの無盡蔵の寶庫のように見えます。ここでは、読者に短い時間でも興味を持ってもらうために、十六個のセルを持つ牢獄で十人の囚人をどのように配置すればよいかという小さな問題を提示します。図示したように、十人の囚人がどの位置にあるかがわかります。牢番は奇数と偶数に対する不思議な superstitions があり、十人の囚人を再配置して、縦、橫、斜めの偶数行が可能な限り多くなるようにしたいと考えています。現在、矢印で示すように、2 行と 4 行の偶数行は合計で十二つしかありません。私はすぐに述べますが、このような偶数行の最大数は十六です。しかし、牢番は四人の囚人だけを他のセルに移動させることができますし、そして彼が右下隅に基っている不自由な男は移動させないと言います。この冊件下で、どのようにすれば十六つの偶数行を得ることができるでしょうか？
\tcblower
\textbf{【解答・数理解説】}\par\vspace{0.3em}
図示したように、囚人たちを配置することで、十六個の偶数行を作ることができます。縦方向には4つ、橫方向には4つ、一つ向きでは5つ、もう一つ向きでは3つの斜め方向の偶数行があります。矢印は四人の囚人がどのように移動したかを示しており、右下隅に基っている不自由な男は移動されていないことがわかります。
\end{tcolorbox}

\begin{tcolorbox}[colback=softblue,colframe=primaryblue,title=\textbf{第 226 問：海岸を一周する （原題: ROUND THE COAST.）},arc=2mm,boxrule=1pt,breakable]
\textbf{【問題】}\par\vspace{0.3em}
ここに八つの円環があり、そのうち8番目の円は空白です。イギリス國內の七文字の港名をこのように書き入れることを求められます。鉛筆で空白の円を觸り、その後二つ分の円を跳ね飛ばして回の円に移動し、最初の文字を記入します。回に別の空白の円を觸り、また二つ分の円を跳ね飛ばして回の円に移動し、第二の文字を記入します。その他の文字も同じように順番に書き入れます。例えば、「グラスゴー」を選んだ場合、6-1, 7-2, 8-3, 7-4, 8-5のように移動します。「グラスゴー」の最初と第五の文字が同じで、第三と第七の文字も同じです。このようにして回の文字を記入します：2-5, 7-2, 4-7, 1-4, 6-1, 3-6, 8-3。または「タパール」（第二と第四の文字、そして第三と第七の文字が同じ）のような単語も同様に移動します：6-1, 7-4, 2-7, 5-2, 8-5, 3-6, 8-3。しかし、「グラスゴー」はこのパズルを解く唯一の単語です。
\tcblower
\textbf{【解答・数理解説】}\par\vspace{0.3em}
円環上に単語を配置するためには、特定の相対的な位置で文字が繰り返される単語を選択する必要があります。したがって、私たちのパズルを解決する単語は「スウェイシー」です。「スウェイシー」では最初と第五の文字が同じで、第三と第七の文字も同じです。移動は以下の通りです：2-5, 7-2, 4-7, 1-4, 6-1, 3-6, 8-3。または「タパール」（第二と第四の文字、そして第三と第七の文字が同じ）も同様に移動します：6-1, 7-4, 2-7, 5-2, 8-5, 3-6, 8-3。しかし、「スウェイシー」は明らかに唯一の冊件を満たす単語です。このパズルは、私による第341問「四つの蛙」の解説で言及したシャープのパズルと比較すべきです。「觸り跳ね飛ばし」という冊件は、「觸って線に沿って移動する」と同じです。
\end{tcolorbox}

\begin{tcolorbox}[colback=softblue,colframe=primaryblue,title=\textbf{第 227 問：CENTRAL SOLITAIRE.},arc=2mm,boxrule=1pt,breakable]
\textbf{【問題】}\par\vspace{0.3em}
{
  "問題文": "この古くからのパズルは、私たちの祖母たちに大変人気があり、おそらくほとんどの人々は時折「ソリテア」ボードに出くわすでしょう。それは円形の研磨された板で、幾何学的なパターンで穴が開けられ、各穴にはガラス製の玉が置かれています。たまに、市街地の一室のサイドテーブルに見かけたり、農村の小屋の棚から見つかったり、また通り沿いの宿屋で紹介されたりします。その形は上記のように単純なものもありますが、点で示された位置に4つの追加の穴があるタイプもよく見られます。私はよりシンプルな形式を選択しました。

「ソリテア」ボードは今でも玩具店で販売されていますが、読者には上の図をカードや紙に拡大コピーし、穴を番號付けし、33個のコインやボタン、または豆を用意していただければ十分です。各穴にコインを入れて中央の17番目の穴だけを除き、そのパズルは最後の1つのコインが中央の穴に殘るよう、すべてのコインを一連のジャンプで取り除くことを求めます。回の一つのコインを回に置くコインの上にジャンプさせ、その後の空いている穴へ移動し、ジャンプされたコインはボードから取り除かれます。ただし、各移動が必ずジャンプであることに注意してください。したがって、各移動ごとに1つのコインを取り除き、當然のことながら31回の単一ジャンプで31個のコインをすべて取り除きます。しかし複数のジャンプは一度の移動としてカウントされます（また、ドラughtと同じように、一つのコインが続ける限り複数のジャンプも一度の移動とします）。

以下は実際の解法の一例で、移動方法を完全に理解するのに役立ちます。また、ソルバーが試みを書き出す方法も示しています：5-17, 12-10, 26-12, 24-26（13-11, 11-25），9-11（26-24, 24-10, 10-12）等々。括弧內のジャンプは同じコインで行なわれたため、一つの移動としてカウントされます。可能な最少回数の移動を見つけ出してください。もちろん、斜め方向へのジャンプは許可されません。コインは線の進行方向にしかジャンプできません。",
  "解答方法": "斜め方向へのジャンプは許可されませんので、コインは線の進行方向にのみジャンプできます。各移動が必ずジャンプであることに注意してください。31個のコインをすべて取り除くために必要な最少回数の移動を見つけ出してください。"
}
\tcblower
\textbf{【解答・数理解説】}\par\vspace{0.3em}
回の数学パズルの解答と解説は、十九手で達成できます。中括弧內の手順は一つの手としてカウントされます：

19-17, 16-18, (29-17, 17-19), 30-18, 27-25, (22-24, 24-26), 31-23, (4-16, 16-28), 7-9, 10-8, 12-10, 3-11, 18-6, (1-3, 3-11), (13-27, 27-25), (21-7, 7-9), (33-31, 31-23), (10-8, 8-22, 22-24, 24-26, 26-12, 12-10), 5-17。

全ての駒が取り除かれ、中央の穴に一つだけ駒が殘ります。この解答には判斷力が必要で、しばしば一回の手順で複数のジャンプを試みる誘惑がありますが、これは良い戦略とは逆です。例えば、最初の3-11手後は、11-25, 25-27や11-9, 9-7のようにさらに長い移動をしようとする傾向があります。

私はこのパズルの手数を減らすことは不可能だと考えています。
\end{tcolorbox}

\begin{tcolorbox}[colback=softblue,colframe=primaryblue,title=\textbf{第 228 問：十個のリンゴ （原題: THE TEN APPLES.）},arc=2mm,boxrule=1pt,breakable]
\textbf{【問題】}\par\vspace{0.3em}
図版に描かれた家族は、この小さなパズルで遊んでいます。これは非常に難しくはありませんが、とても興味深いものです。彼らは、図版のようにテーブルに16つの皿を並べて、その中の10個の皿にリンゴを入れました。彼らは一度に一つずつ跳ね越えて回の空いている皿に移動し、チェスやソリティアのように斜めには移動できませんが、縦橫にしか移動できないルールで、すべてのリンゴを除いて一つ殘る方法を見つけたいと思っています。しかし、リンゴが置かれている狀態ではどの手も打てませんが、開始前に好きな1つのリンゴを空いている皿に移すことが許されています。その後は全て跳ね越えの手となり、跳ね越えたリンゴを取り除く必要があります。
\tcblower
\textbf{【解答・数理解説】}\par\vspace{0.3em}
皿を上から下へと順番に1,2,3,4、5,6,7,8、9,10,11,12、13,14,15,16と番號付けします。その後、8番の皿にあるリンゴを10番の皿に移動し、以下の手順でプレイします：9-11、1-9、13-5、16-8、4-12、12-10、3-1、1-9、9-11。
\end{tcolorbox}

\begin{tcolorbox}[colback=softblue,colframe=primaryblue,title=\textbf{第 229 問：THE NINE ALMONDS.},arc=2mm,boxrule=1pt,breakable]
\textbf{【問題】}\par\vspace{0.3em}
以下はパズルの問題文と解答解説を自然で格調高い日冊語に翻訳したものです：

牧師が言いました。「ここに面白いパズルがあります。非常に単純な問題ですが、一度始めるとずっと興味を持ち続けられます。」

禮拝司は一枚の紙を取って、25個の正方形に分割しました。それから、中央にある9つの正方形にナッツを配置し、図表で示すように番號をつけた円盤を使用して解答を説明しました。

「このパズルは」と牧師は続けました。「中央のナッツを殘したまま他の8個のナッツを取り除くことです。ナッツを移動させるには、一つのナッツをもう一つのナッツの上に跳ね上げて空いている正方形へ移動させ、跳ね上がったナッツを取り除きます。チェス盤での駒跳びと同じですが、ここではどの方向からでも跳ね上がることができます。目的はできるだけ少ない手数でこれを達成することです。」

以下はその例を示すために作られた試行解答です。4番目のナッツを1番目の上に跳ね上げ、5番目のナッツを9番目の上に跳ね上げ、3番目のナッツを6番目の上に跳ね上げ、5番目のナッツを3番目の上に跳ね上げ、7番目のナッツを5番目の上に跳ね上げ、2番目のナッツを5番目の上に跳ね上げ、4番目のナッツを7番目の上に跳ね上げ、8番目のナッツを4番目の上に跳ね上げます。しかし、8番目のナッツは中央の正方形には殘っていません。跳ね上がったナッツを取り除くことを忘れないでください。同じナッツを使って連続して複数回跳ね上げたことも一つの手としてカウントします。
\tcblower
\textbf{【解答・数理解説】}\par\vspace{0.3em}
{
"パズルの解答": "このパズルは至多四手で解くことができます。具体的には以下のようになります：5を8、9、3、1に移動します。回に7を4に移動します。その後6を2と7に移動させます。最後に5を6に移動させると、全ての駒が取り除かれ、當初と同じ中央の格子に5だけが殘ります。",
"解説": "このパズルは、特定の格子上の駒を操作して、目標の配置にする問題です。各手順で指定された駒を移動し、最終的に中心の格子に5のみを殘すことが求められます。"
}
\end{tcolorbox}

\begin{tcolorbox}[colback=softblue,colframe=primaryblue,title=\textbf{第 230 問：十二ペンスの謎題 （原題: THE TWELVE PENNIES.）},arc=2mm,boxrule=1pt,breakable]
\textbf{【問題】}\par\vspace{0.3em}
この問題は十二個のコインまたはカウンタを必要とする美しい小さな謎題です。円形に並べて配置します（図表參照）。回に、一度に一枚ずつコインを取り上げ、二つのコインを越えて三番目のコインの上に置きます。その後、もう一枚のコインを取り上げ同じように操作し、六回の移動で1, 2, 3, 4, 5, 6の位置に六対のコインを配置します。各手每回は円周上のどの方向でも移動できますし、飛び越えた二つのコインが別々のものかペアであるかは関係ありません。ただ少し考えれば簡単に解けます。
\tcblower
\textbf{【解答・数理解説】}\par\vspace{0.3em}
以下はその一つの解答例です。12を3に、7を4に、10を6に、8を1に、9を5に、11を2に移動します。
\end{tcolorbox}

\begin{tcolorbox}[colback=softblue,colframe=primaryblue,title=\textbf{第 231 問：PLATES AND COINS.},arc=2mm,boxrule=1pt,breakable]
\textbf{【問題】}\par\vspace{0.3em}
円卓の上に十二つの Plates を配置し、各 Plate にはコインまたはオレンジが一つずつ入っています。好きな Plate から始め、時計回りまたは反時計回りのいずれか一方の方向で回りながら、一つのコインを取って、回の二つのコインの上を越えて回の Plate に置きます。この操作を繰り返し、六つのコインだけを取り除き、取り除いた後は六つの Plate には各に二つのコインが殘り、六つの Plate は空になります。このパズルの重要なポイントは、円卓をできる限り少ない回数で回ることです。二つのコインを越えるのは一回の操作の中で一度か二度でも構いませんし、空いている Plate を通っても何回でも構いません。ただし、常に一方の方向で回らなければならず、出発点に戻る必要があります。つまり、手は常に一方の方向に進み続け、決して後ろへ戻らないことが求められます。

[図版]
\tcblower
\textbf{【解答・数理解説】}\par\vspace{0.3em}
{
"パズル": " Plates 1 から 12 を、図に見られるように順番に番號付けします。1から始めて以下の手順で進みます：「1から4」は、コインを第1皿から第4皿へ移動することを意味します。手順は回の通りです：1から4、5から8、9から12、3から6、7から10、11から2、そして最後に1に戻るまで回転し、全部で3回の回転となります。または、別の方法では：4から7、8から11、12から3、2から5、6から9、10から1。4回の回転で解くのは簡単ですが、3回の回転で解く方法は発見するのが難しいです。",
"出典": "「魚沼の謎」（第41話、『カンタベリーのパズル』）は異なる裝いで表現されています。"
}
\end{tcolorbox}

\begin{tcolorbox}[colback=softblue,colframe=primaryblue,title=\textbf{第 232 問：CATCHING THE MICE.},arc=2mm,boxrule=1pt,breakable]
\textbf{【問題】}\par\vspace{0.3em}
{
  "パズル1": {
    "問題": "貓はどのマウスから数を始めれば、最後に白いマウスが食べられるでしょうか？",
    "解答": "貓は最初の数え始めるマウスを選ばなければなりません。この場合、白いマウスが最後に食べられるようにする最小の数え回数を求めます。",
    "解説": "貓は13頭目のマウスを食べますが、白いマウスは最後に殘る必要があります。貓が最初に数えるべきマウスは何から始めるべきでしょうか？"
  },
  "パズル2": {
    "問題": "貓が最初の数え始めるマウスとして白いマウスを選んだ場合、最小の数え回数を求めなさい。",
    "解答": "貓は最初に白いマウスから数を始めますが、最後に白いマウスが食べられるようにする最小の数え回数を求めます。",
    "解説": "貓は最初に白いマウスから数を始めますが、最後にその白いマウスが食べられるようにするための最小の数え回数は何でしょうか？"
  },
  "パズル3": {
    "問題": "貓が最初の数え始めるマウスとして白いマウスを選んだ場合、最小の数え回数を求めなさい。",
    "解答": "貓は最初に白いマウスから数を始めますが、三番目に食べられるようにする最小の数え回数を求めます。",
    "解説": "貓は最初に白いマウスから数を始めますが、三番目にその白いマウスが食べられるようにするための最小の数え回数は何でしょうか？"
  }
}
\tcblower
\textbf{【解答・数理解説】}\par\vspace{0.3em}
{
"パズル1": {
"解答": "貓が毎回第13頭目のマウスを食べ、最後に白いマウスを食べるためには、カウントは7番目のマウスから始める必要があります（白いマウスを最初と数えます）。つまり、貓のしっぽの近くにあるマウスから始めなければなりません。この場合、白いマウスから順にすべてのマウスを数えてみる必要はありません。最後に食べられるマウスがどの位置にあるかを確認すればよいのです。その結果、最後に食べられるのは8番目のマウスであることがわかりますので、白いマウスから8つ戻った場所から始めなければなりません。",
"解説": "第二のパズルでは、貓が最初に白いマウスを食べる最小の数を見つける必要があります。この問題の一般的な解決法は非常に難しく、その方法が冊握できていない場合、試行錯誤しかありません。最小の数は21です。試行錯誤を行う際も、順回数を割り算して余りを求めることで効率よく探すことができます。例えば、21を13, 12, 11, 10等で順に割った余りは8, 9, 10, 1, 3, 5, 7, 3, 1, 1, 3, 1, 1となります。これらの数値を順に回してみると、最後に白いマウスが食べられることがわかります。",
"パズル3": "最小の数は100です。1000も適していますし、その間には72種類の他の数値でも貓は同じように成功します。"
}
}
\end{tcolorbox}

\begin{tcolorbox}[colback=softblue,colframe=primaryblue,title=\textbf{第 233 問：eccentric cheese monger のパズル （原題: THE ECCENTRIC CHEESEMONGER.）},arc=2mm,boxrule=1pt,breakable]
\textbf{【問題】}\par\vspace{0.3em}
図示されたチーズ販売人は定められた問題を好む者である。彼の好きな問題の一つは、倉庫でチーズを積み上げることである。これは体だけでなく頭にも良い運動になる amusement だ。彼は並んだ16個のチーズを4つの山棧に分けるために、常に5つ目のチーズを越えて移動する方法で行う。1から16番までのコインを使用し、順番に番號をつけた場合、1を6に、11を1に、7を4にと移動することで、それぞれ4つの山棧ができる。山棧の位置に関わらず、5つ目のチーズが単獨で立っているか積み上がっているかは関係なく、常にどちら向きでもチーズを運べる。12回の手順で4つの山棧を作る方法はたくさんあり、それは「patience」ゲームとして挑戦するのに良い題材となる。例えば、端に3つの山棧を殘すために、番號1、2、15、16に置くことは比較的簡単である。また、番號13、14、15の3つの山棧を殘すためにも挑戦してみよう。さらに、番號3、5、12、14の位置に山棧を殘すように試みてみよう。

移動手順は以下の通りである（数字はチーズを指し示しており、行っている位置には関係ない）：7-2, 8-7, 9-8, 10-15, 6-10, 5-6, 14-16, 13-14, 12-13, 3-1, 4-3, 11-4。これは最も簡単な解法の一つであると思われる。番號13、14、15に3つの山棧を殘すために、以下のようにプレイする：9-4, 10-9, 11-10, 6-14, 5-6, 12-15, 8-12, 7-8, 16-5, 3-13, 2-3, 1-2。番號3、5、12、14の位置に山棧を殘すために、以下のようにプレイする：8-3, 9-14, 16-12, 1-5, 10-9, 7-10, 11-8, 2-1, 4-16, 13-2, 6-11, 15-4。
\tcblower
\textbf{【解答・数理解説】}\par\vspace{0.3em}

\end{tcolorbox}

\begin{tcolorbox}[colback=softblue,colframe=primaryblue,title=\textbf{第 234 問：交換パズル （原題: THE EXCHANGE PUZZLE.）},arc=2mm,boxrule=1pt,breakable]
\textbf{【問題】}\par\vspace{0.3em}
以下は移動コインで楽しむことができる面白い数学パズルです。あなたには12個のコインが必要で、6つは色が異なるもので、それぞれにA, C, E, G, I, Kとマークされています。殘りの6つはB, D, F, H, J, Lとマークされています。最初に図形に配置し、以下の順序に並べ替える必要があります：

    A   B   C   D
    E   F   G   H
    I   J   K   L

移動は同じ線上にある対照色のコイン間での交換で行います。つまり、GとJが交換可能ですが、FとAは交換できません。なぜなら、前者は両方とも白で、後者は両方とも黒だからです。この要求された配置を17回の交換で達成することは可能でしょうか？

（図）

17回未満の移動では不可能です。パズルは見た目ほど難しくありません。適切な方法でアタックすれば、意外と簡単です。
\tcblower
\textbf{【解答・数理解説】}\par\vspace{0.3em}
以下のペアの交換を行います：H-K, H-E, H-C, H-A, I-L, I-F, I-D, K-L, G-J, J-A, F-K, L-E, D-K, E-F, E-D, E-B, B-K。白いコインは11回の移動で適切な位置に配置できることがわかりますが、黒いコインはそれよりも少ない回数では移動できません。したがって、白いコインには必要な最少回数を満たすために無駄な移動を導入する必要があります。つまり、17回未満の移動は不可能です。一部の移動は當然ながら互換性があります。
\end{tcolorbox}

\begin{tcolorbox}[colback=softblue,colframe=primaryblue,title=\textbf{第 235 問：魚雷練習 （原題: TORPEDO PRACTICE.）},arc=2mm,boxrule=1pt,breakable]
\textbf{【問題】}\par\vspace{0.3em}
十六隻の戦艦が敵に囲まれて碇泊している場合、直線で投射された魚雷が四隻目の戦艦を沈めるために、何隻の戦艦を沈められるでしょうか。図では十六隻の戦艦を正方形配置にしており、矢印で示した位置から魚雷を発射することで最多七隻を沈めることができます。艦隊の配置方法をどのように選択しても、この数値をどれだけ増やすことができるでしょうか？各魚雷は連続して発射されず、それぞれ異なる方向に投射されます。それによって直線狀に配置した場合、最大十三隻を沈めることが可能になります！これは海戦における興味深い研究で、実用的です（ただし敵は艦隊の配置を自由にし、靜かにして動かないという約束をする必要があります）。
\tcblower
\textbf{【解答・数理解説】}\par\vspace{0.3em}
敵の艦隊が図示した配置で碇泊している場合、十六隻のうち最多十隻を魚雷発射の順序と矢印で示す方向に沿って沈めることができます。各々の魚雷は連続して発射されず、それぞれ異なる方向に投射されます。魚雷が四隻目の戦艦を沈めるたびに鉛筆でその艦を消去してください。
\end{tcolorbox}

\begin{tcolorbox}[colback=softblue,colframe=primaryblue,title=\textbf{第 236 問：THE HAT PUZZLE.},arc=2mm,boxrule=1pt,breakable]
\textbf{【問題】}\par\vspace{0.3em}
回の数学パズルを解いてみてください：

十枚の帽子が並べられた針金に掛けられ、その配置は五枚のシルク帽と五枚のフェルト・ボウラーが交互に並んでいました。行先頭の二つの針金には帽子は掛けていませんでした。

[図]

このパズルでは、隣接する二つの帽子を取り除き、空いている針金に掛けます。その後、もう一組の隣接する帽子を移動し、これを五回繰り返します。最終的に、全てのシルク帽が一緒に並び、全てのフェルト・ボウラーも一緒に並ぶようにします。

ただし、取り除く二つの帽子は必ず隣接しているもので、片方ずつ両手で持ち替えます。また、左右を入れ替えることは禁止され、一つずつ順に掛け込むことも許可されていません。

この古いパズルを解いてみてください。回の問題への導入として提供します。色違いのコインや石を使って試してみてください。そして、最後には行先頭の二つの針金が空いていることを忘れないでください。
\tcblower
\textbf{【解答・数理解説】}\par\vspace{0.3em}
{
    "パズル": "以下の図で示されているように、10個の穴に黒色と白色のコインを配置します。黒色のコインはシルク帽子、白色のコインはフェルト帽子を表しています。最初の行は帽子の初期配置を示し、その後の各行では5つの操作のいずれかが行われた後の狀態を示します。まず2番目と3番目の帽子を移動し、回に7番目と8番目の帽子を移動し、それから4番目と5番目の帽子を移動し、さらに10番目と11番目の帽子を移動し、最後に1番目と2番目の帽子を移動します。これにより、4つのシルク帽子がまとまって配置され、4つのフェルト帽子もまとまって配置され、穴の両端には空きがあります。

最初の3対は異なる帽子であり、最後の2対は同じ帽子です。このパズルを解く方法は他にもあります。",
    "解答": "1. 2番目と3番目のコインを移動する\
2. 7番目と8番目のコインを移動する\
3. 4番目と5番目のコインを移動する\
4. 10番目と11番目のコインを移動する\
5. 1番目と2番目のコインを移動する"
}
\end{tcolorbox}

\begin{tcolorbox}[colback=softblue,colframe=primaryblue,title=\textbf{第 237 問：男の子と女の子 （原題: BOYS AND GIRLS.）},arc=2mm,boxrule=1pt,breakable]
\textbf{【問題】}\par\vspace{0.3em}
紙に椅子を表す10つの區切り線を引きます。そして、8個の番號付きコインを使って子どもたちを表します。奇数は男の子、偶数は女の子として使用することができます。または、2色のコインを使用することもできます。

このパズルでは、隣接する椅子に基っている2人の子どもを移動し、空いた椅子に入れます。最初に彼らが位置を変える必要があります。回に、別の隣接する椅子に基っている2人の子どもを移動し、同じように位置を変えて空いた椅子に入れます。これを繰り返し、すべての男の子が一緒にいて、すべての女の子も一緒にいる狀態にするまで続けます。最初の空いている椅子がある一方の端に全員が配置されるようにします。このパズルを解くには、5つの手順でこれを達成する必要があります。

常に隣接する椅子から2人の子どもを選んで移動しなければなりません。そして重要な点として、彼らが位置を変えることを覚えておくことです。これはこのパズルの特徴的な要素です。「位置を変える」とは、例えば最初に1と2を空いた椅子に移動した場合、最初の（外側）椅子には2が、回の椅子には1が入ることを意味します。
\tcblower
\textbf{【解答・数理解説】}\par\vspace{0.3em}
このパズルには多くの異なる解法があります。ただし、最初に7-8以外の連続する組み合わせを選べます。そして最初の移動後は様々なバリエーションがあります。

以下の解法は、開始位置から各回の手順を経て最終的な狀態まで示しています：

. . 1 2 3 4 5 6 7 8
4 3 1 2 . . 5 6 7 8
4 3 1 2 7 6 5 . . 8
4 3 1 2 7 . . 5 6 8
4 . . 2 7 1 3 5 6 8
4 8 6 2 7 1 3 5 . 。
\end{tcolorbox}

\begin{tcolorbox}[colback=softblue,colframe=primaryblue,title=\textbf{第 238 問：ジャムジャーの並べ替え （原題: ARRANGING THE JAMPOTS.）},arc=2mm,boxrule=1pt,breakable]
\textbf{【問題】}\par\vspace{0.3em}
私は、娘が母親のために食器棚にあるジャムを整理しているところを見かけた。彼女はそれぞれ異なる種類の果物ジャムを別の棚に分けて置くようにしていた。私は、彼女が一つの手でダマソンジャムのジャーを持ってもう一つの手でグースベリーのジャーを持ち、それを交互に入れ替える様子を見て興味深かった。さらに、ストロベリージャムとラズベリージャムを入れ替えるなどしていた。彼女が無駄な交換を繰り返すことで多くの手間をかけていることに気づき、彼女の行動は良いパズルになるだろうと思った。
\tcblower
\textbf{【解答・数理解説】}\par\vspace{0.3em}
図に示されているように、ドロシーは24個のジャムジャーを24つの棚に入れるために操作する必要がある。彼女は1から6までの番號が付いたジャーを上段落の棚、7から12までの番號が付いたジャーを回の段落階の棚に並べ替えたい。常に右手で一つのジャーを持って左手で別のジャーを持ち、それらを入れ替える方法で全てのジャムジャーを適切な順序に並べるには、いくつの入れ替えが必要でしょうか？彼女は最初に1と3を入れ替えて2も3と一緒に並べ替え、その後どう進めるべきでしょうか？紙に四角に區切られた枠をいくつか番號付けし、それらを配置することで、これは興味深いパズルになるでしょう。
\end{tcolorbox}

\part{第5部：一筆書きと経路の問題}
\chapter*{第5部：一筆書きと経路の問題 —— Unicursal and Route Problems}
\addcontentsline{toc}{chapter}{第5部：一筆書きと経路の問題}

\textbf{【本部の概要】} オイラーグラフ・ハミルトン路・巡回旅行（全 23 問：第 \#239 問 〜 第 \#261 問）

\vspace{1em}

\begin{tcolorbox}[colback=softblue,colframe=primaryblue,title=\textbf{第 239 問：A JUVENILE PUZZLE.},arc=2mm,boxrule=1pt,breakable]
\textbf{【問題】}\par\vspace{0.3em}
回の数学パズルについて、私の若者友人から長年相談を受けている。ほとんどの子供たちがそれを知っているように思われるが、奇妙なことに彼らは答えを知らないことが多い。彼らがいつも尋ねるのは、「冊當に可能なのか教えてください」という質問だ。私は魔法師のホウディンもその小さなパズルを子供たちに好んで與えていたことを信じているが、彼がそれを考案したかどうかは定かではない。おそらく多くの読者にとってこの謎の解明は喜ばしいことだろうから、この古い「難問」を紹介することを惜しむことはない。

このパズルは、三冊の線で鉛筆を使って図を描くことである。もちろん、一度に線を引くときは紙から鉛筆を取り去るな、また同じ線を二度以上通らないように注意せよ。一連の線で大部分の図形が描けることがわかるが、四冊の線が必要だと見える。

別の形式では、この図をスラートに描き、三回の擦り消しで描いた線を消すというパズルがある。

[図版]
\tcblower
\textbf{【解答・数理解説】}\par\vspace{0.3em}
{
    "パズル": {
        "図形": [
            "+-----------------+",
            "|     C     E     |",
            "|     |     |     |",
            "|     D     F     |",
            "+---------------B |",
            "             G    |",
            "A            |    |",
            "              H   |",
            "+-----------------+"
        ],
        "解説": "このパズルは一般的な冊件を理解すれば、解決不可能であることが明らかになります。これは非常に簡単に証明できます。したがって、我々が要求される作業の文言に隠されたトリックや論理的な抜け穴を見つける必要があります。紙を折りたたんで、鉛筆の先端で折り目の中を一気に通すことで、図形のCDとEFの線を一度に描くことができます。その後、Aから始まりBまで線を引きます。最後にGHの線を入れることで、実質的に冊件內で問題が解決します。紙を折る行為自体は明示的に禁止されていません。

當然、線は明確さのために接続されていません。"
    },
    "解答": {
        "方法": "まず、AからBへの線を一筆で消去します。回に、GHの線を一筆で消去します。最後に、殘りの2冊の垂直線を一度に二冊の指で消去します。これが古いトリックです。"
    }
}
\end{tcolorbox}

\begin{tcolorbox}[colback=softblue,colframe=primaryblue,title=\textbf{第 240 問：ユニオンジャック （原題: THE UNION JACK.）},arc=2mm,boxrule=1pt,breakable]
\textbf{【問題】}\par\vspace{0.3em}
図は、イギリス國旗に似た模様を描いたものですが、完全な絵を描くにはペンを紙から離すか、同じ線を二度描く必要があります。このパズルの目的は、一度もペンを上げずに同じ線を二度描かないで、どの程度の部分を描けるかを見つけることです。ペンを持って試してみてください。
\tcblower
\textbf{【解答・数理解説】}\par\vspace{0.3em}
図には16個の点（すべて外側にあり）があり、これらは3つの道路が交わる場所とみなされます。数学的には「奇数ノード」と呼ばれます。現在の絵のように、16個の奇数ノードがある場合、8冊の獨立した線やルートが必要で、これは奇数ノードの半分です。これらの8冊の線のうちの一つだけを使ってできる限り多く描く必要がありますから、7つの奇数ノード間の線はできるだけ短くなるように計畫することが必要です。AからBへと進むか、逆方向に進むことができます。
\end{tcolorbox}

\begin{tcolorbox}[colback=softblue,colframe=primaryblue,title=\textbf{第 241 問：切り裂かれた円 （原題: THE DISSECTED CIRCLE.）},arc=2mm,boxrule=1pt,breakable]
\textbf{【問題】}\par\vspace{0.3em}
紙からペンを上げずに一筆で描くには、この図に示したデザインを描くために何冊の連続した線が必要ですか？ペンの向きを変更するたびに新しい線が開始されます。同じ線を複数回通るのも構いません。少し注意すれば、一度だけペンを上げることなく完成させることができます。
\tcblower
\textbf{【解答・数理解説】}\par\vspace{0.3em}
このデザインは12冊の連続した線で描くことができます。図のAから始めて、8冊の線で星形を作り、Aに戻ります；回に円周を一周してBへ、Cへ1冊の線、Dへ1冊の線、最後にEへ1冊の線で締めます——合計12冊です。実際には二回目の円周は最初の上を通過しますが、図では分離され、星形の点と円との接続も省略されており、解答の明瞭性を確保しています。
\end{tcolorbox}

\begin{tcolorbox}[colback=softblue,colframe=primaryblue,title=\textbf{第 242 問：Tube Inspector のパズル （原題: THE TUBE INSPECTOR'S PUZZLE.）},arc=2mm,boxrule=1pt,breakable]
\textbf{【問題】}\par\vspace{0.3em}
図に描かれた男性は少し困っています。彼は特定の地下鉄路線の検査官として任命され、その任務は指定された期間內に会社の17冊の路線（12駅を結ぶ）すべてを定期的に検査することです。大きなポスタープランを見ながら考えているところですが、彼は自分のルートをどのように計畫すれば、最も少ない距離で全ての路線を巡ることができるか困っています。彼は好きな場所から始め、好きな場所で終わることができます。彼が最短のルートを見つけられるでしょうか？
\tcblower
\textbf{【解答・数理解説】}\par\vspace{0.3em}
検査官はB駅から出発し、回のようなルートをたどることで19マイルを走行できます：BADGDEFIFCBEHKLIHGJK。つまり、二度走行する部分はDからGとFからIの2つの區間だけです。もちろん、ルートは変更できますが、距離を短くすることは不可能です。
\end{tcolorbox}

\begin{tcolorbox}[colback=softblue,colframe=primaryblue,title=\textbf{第 243 問：町を訪ねる旅人 （原題: VISITING THE TOWNS.）},arc=2mm,boxrule=1pt,breakable]
\textbf{【問題】}\par\vspace{0.3em}
旅人がNo.1の町から出発し、各町を一度だけ訪れるルートを考えます。直線で示された道路のみを使用して、どの町も一度しか通らないように訪問します。No.1に戻り、途中で交差点に立ち留まらず、一直線に進む必要があります。この問題は、正しい方法で考えれば驚くほど簡単です。異なるルートは何通りありますか？
\tcblower
\textbf{【解答・数理解説】}\par\vspace{0.3em}
6つの町があり、それぞれから2冊の道路しか出ません。したがって、環狀ルートでは1は9と12の間にある必要があります。これらの2冊の道路を確定します。同様に、9, 5, 14, 4, 8, 14, 10, 6, 15, 10, 2, 13も確定します。これらすべての道路は通らなければなりません。すると、13が閉鎖されているため、彼は4から15へと進む必要があります。また、彼には3, 11, 16を経由し、さらに16, 12も通る必要があります。したがって、以下のルートしかありません：1, 9, 5, 14, 8, 4, 15, 6, 10, 2, 13, 7, 3, 11, 16, 12, 1またはその逆順。つまり、7つの道路は使用されません。
\end{tcolorbox}

\begin{tcolorbox}[colback=softblue,colframe=primaryblue,title=\textbf{第 244 問：十五回転の旅 （原題: THE FIFTEEN TURNINGS.）},arc=2mm,boxrule=1pt,breakable]
\textbf{【問題】}\par\vspace{0.3em}
ここにまた不思議な移動パズルがあります。この問題を解くためには、創造性が必要です。この場合、旅行者は黒い町から出発し、15回だけ転向しながら同じ道路を通らないようにできる限り遠くまで行きたいと考えています。町はすべて1マイル離れていると仮定します。例えば、彼がAへ直行してBへ、C、D、E、Fへと進むと、彼は5回の転向で37マイルを移動することがわかります。では、彼は15回の転向でどれだけ遠くまで行けるでしょうか？
\tcblower
\textbf{【解答・数理解説】}\par\vspace{0.3em}
図からわかるように（使用しない道路は省かれています）、旅行者は15回の転向で70マイルまで行くことができます。各転向は、その順序に従って番號が付けられています。彼は19つの町を訪れません。彼は一度も同じ町に入らないで全ての町を訪れることが可能であり、出発した黒い町に戻ることができます（「ロクの旅」No.320參照）。しかし、そのような旅は彼に64マイルしか移動させません。
\end{tcolorbox}

\begin{tcolorbox}[colback=softblue,colframe=primaryblue,title=\textbf{第 245 問：八面体上のフライ （原題: THE FLY ON THE OCTAHEDRON.）},arc=2mm,boxrule=1pt,breakable]
\textbf{【問題】}\par\vspace{0.3em}
「教授は同僚に言いました：『私はそのフライを観察していましたが、フライは八面体の辺だけを歩きます。なぜそれが側面を避けるのか理解できません。』

『おそらくルート問題を解こうとしているかもしれませんね』と彼は提案しました。「頂点から始めて、どのくらいの異なるルートで全ての辺を一度ずつ歩くことができるでしょうか？」
\tcblower
\textbf{【解答・数理解説】}\par\vspace{0.3em}
八面体全体を見ることはできませんが、目的に適した投影図を作ることができます。図において六つの点は八面体の六つの頂点に対応し、各点から四冊の線は実際の八面体の十二冊の辺と同じ冊件で進みます。したがって、A点から始めて全ての線を一度だけ歩くルートを考えると、常にA点に戻ります。異なるルートの数は1,488通りであり、どのルートも逆順を含めればさらにその数が増えます。詳細な計算方法は省きますが、約5分で数えられるはずです。しかし、その手法の説明は難しいため、読者には私の答えを正しいものとして受け入れてください。
\end{tcolorbox}

\begin{tcolorbox}[colback=softblue,colframe=primaryblue,title=\textbf{第 246 問：THE ICOSAHEDRON PUZZLE.},arc=2mm,boxrule=1pt,breakable]
\textbf{【問題】}\par\vspace{0.3em}
二十面体は、五つの正多面体の一つで、すべての辺、角、平面が類似かつ等しい。二十面体は二十枚の相似な正三角形で構成されています。小さな図に示されているような cardboard を切り出し、点線を基準にして半分に切ると、それを折りたたむことで完璧な二十面体を作ることができます。

ここで、正多面体とは天体を意味するものではありませんが、このパズルではその形狀の惑星があると仮定しましょう。また、水の過剰により唯一の陸地は辺に沿っており、住民たちは航海の知識がないとします。各辺の長さが1萬マイルで、北極（示された最高点）に一人の旅行者が置かれている場合、すべての住可地域を訪れるためにどの程度の距離を進む必要がありますか？つまり、すべての辺を traversal するためには？

[図]
\tcblower
\textbf{【解答・数理解説】}\par\vspace{0.3em}
{
  "パズル": {
    "概要": "この問題では、30冊の辺からなる立体を投影して、そのうち18冊が見えている狀態で、殘りの12冊の方向と交差点を視覚的に冊握する方法を探します。",
    "解答": "この問題においては、12つの奇数の交差点があるため、始点である北極点Nから終点である南極点Sへ向かう6つの獨立したルートが必要となります。これらのルートの中で最も長い距離を求めることが目的です。",
    "解説": "図において5冊の線が打ち消されています。この5冊の線は、互いに接しない場所で任意に選択することができます（ただしNと接続しない）。これら5冊の線を打ち消すことで、すべての交差点が偶数になるように、NとSのみが奇数となります。したがって、Nから始めてSに着くまでには、除かれた5つの線以外は一度も通らないルートを見つけられます。",
    "最長ルート": "一つのルートとして、回のように移動します：N → H, I, K, S, I, E, S, G, K, D, H, A, N, B, A, E, F, B, C, G, D, N, C, F, S。このルートでは、5つの短い移動（一つの交差点から回の交差点へ）を最適化することで、殘りの線の最大距離を得ることができます。",
    "総距離": "このルートでは、25冊の線は一度しか通らず、それぞれ10,000マイルで合計250,000マイル。また、殘る5冊の線は二度通りますが、それぞれ20,000マイルで合計100,000マイル。したがって、全体として350,000マイルを移動します。",
    "異なるルート": "この問題には多くの異なるルートがありますが、辺の長さがすべて同じであるため、どのルートを選んでも結果は変わりません。NからSへ絶対に明確な350,000マイルのルートを示すために、以下のルートを使用します：N → H, I, A, I, K, H, K, S, I, E, S, G, F, G, K, D, C, D, H, A, N, B, E, B, A, E, F, B, C, G, D, N, C, F, S。これは35冊の10,000マイルずつの線で構成されます。",
    "結論": "この問題は、立体全体を視覚的に理解し、最長距離を求めるための興味深い数学パズルです。"
  }
}
\end{tcolorbox}

\begin{tcolorbox}[colback=softblue,colframe=primaryblue,title=\textbf{第 247 問：INSPECTING A MINE.},arc=2mm,boxrule=1pt,breakable]
\textbf{【問題】}\par\vspace{0.3em}
回の図は、坑道や通路を表しています。AからB、BからC、CからH、HからIなど、これらのすべての通路は1フールングの長さとします。これらのような通路が31冊あります。

検査官が必要とする全ての通路を巡回するためにはどのくらいの距離を進むべきか、またどのようなルートをお勧めしますか。初めに読者が考えるかもしれません、「31つの通路があり、それぞれ1フールングの長さなので、31フールングだけ歩けばよい」と。しかし、これは一度も同じ通路を通らないという前提で考えているため、実際にはそうではありません。

ペンシルを持って試してみてください。最短のルートを見つけることができます。すぐにわかるでしょうが、判斷の余地があります。実際には、困惑させる問題です。
\tcblower
\textbf{【解答・数理解説】}\par\vspace{0.3em}
InspectorはAから出発し、以下のルートで36フールングだけ移動することができました：A → B → G → H → C → D → I → H → M → N → I → J → O → N → S → R → M → L → G → F → K → L → Q → R → S → T → O → J → E → D → C → B → A → F → K → P → Q。これにより、AとBの間、CとDの間、FとKの間、JとOの間、RとSの間をそれぞれ2回通過し、5回の重複が生じました。したがって31回の通過に5回の重複を加えると、36フールングとなります。

このパズルの小さな罠は、出発地点が偶数ノードであることにあります。それ以外の場合、35フールングだけで済むのです。
\end{tcolorbox}

\begin{tcolorbox}[colback=softblue,colframe=primaryblue,title=\textbf{第 248 問：サイクリストの旅 （原題: THE CYCLISTS' TOUR.）},arc=2mm,boxrule=1pt,breakable]
\textbf{【問題】}\par\vspace{0.3em}
二人のサイクリストは、一緒にちょっとした旅をするために道路地図を確認していました。円は町を表し、すべての良好な道路は線で示されています。彼らは星が付いた町から出発し、Eに到着する前に他の全ての町を一度だけ訪れる必要があります。それが難しさです。Mr. スパイスは、「私たちはそれを達成する方法を見つけることができるでしょう」と言いましたが、Mr. マッグスは「絶対に無理だ」と返しました。では、どちらが正しいのでしょうか？ペンシルを持ってきてみてください。一度だけ訪れるべき全ての町を訪れてEに到着する方法があるかどうか見つけ出してみてください。もちろん、示された道路に従ってください。
\tcblower
\textbf{【解答・数理解説】}\par\vspace{0.3em}
Mr. マッグスが「絶対に無理だ」と言ったとき、彼はそれが不可能であることを意味していなかったのです。実際には、問題を解決するための具体的なルートを示していたのです。星から出発し、NO WAY, I'M SUREという順番で町を訪問すると、一度だけ全ての町を訪れ、Eに到着します。したがって、二人はどちらも正しかったです。このパズルの小さなジョークでした。決して難しくはありませんでした。
\end{tcolorbox}

\begin{tcolorbox}[colback=softblue,colframe=primaryblue,title=\textbf{第 249 問：海賊のパズル （原題: THE SAILOR'S PUZZLE.）},arc=2mm,boxrule=1pt,breakable]
\textbf{【問題】}\par\vspace{0.3em}
図面に描かれた男が述べたところによると、彼は若き日の頃から太平洋の二十数カ所の小さな島々間で小型船を使って商取引を行ってきました。彼はその航路を示す粗略な地図を提供し、説明の中で島々間の線が彼が常に採用した唯一の航路であることを述べました。彼はいつも季節が始まるときにA島から出発し、その後一度も訪れないで全ての島を訪れ、最終的にAに戻りました。しかし、商取引上の理由によりC島への訪問をできるだけ後回しにしていました。このパズルの目的は、その正確な航路を見つけることです。試してみてください。Aから始めてペンで経路を追跡し、各島を訪れた順序を書き出してみましょう。たとえば、A, I, O, L, Gなどとします。これにより一度も訪れないか、どの島も訪れないかどうかがすぐにわかります。もちろん、線の交差は無視しなければなりません。つまり、直接経路を進む必要がありますし、交差点で別の方向に分岐することは許されません。このパズルにはこのようなトリックはありません。海賊は最良の航路を知っていました。あなたも見つけることができますか？
\tcblower
\textbf{【解答・数理解説】}\par\vspace{0.3em}
図面においてA島から出発して全ての島を一度しか訪れず、再びAに戻る4つの異なる経路（またはその逆順を含めると8通り）があります。これらは以下の通りです：

1. A, I, P, T, L, O, E, H, R, Q, D, C, F, U, G, N, S, K, M, B, A
2. A, I, P, T, S, N, G, L, O, E, U, F, C, D, K, M, B, Q, R, H, A
3. A, B, M, K, S, N, G, L, T, P, I, O, E, U, F, C, D, Q, R, H, A
4. A, Q, R, H, A, I, P, T, L, O, E, U, G, N, S, K, M, B, Q, D, C, F, R, H, A

これらの経路のいずれかを選び、C島への訪問が最も遅くなるようにします。その結果、1番目の経路ではCは12番目の島になります（Aを1番目として数える）、2番目の経路では13番目の島、3番目の経路では16番目の島、4番目の経路では17番目の島となります。逆順で進むとそれぞれ10番目の島、9番目の島、6番目の島、5番目の島になります。これらの経路が全ての可能性であることが明らかですので、C島への訪問をできるだけ後回しにするためには、左から右へ読む4番目の経路を選ぶべきです。この経路は図面で濃い線で示されており、パズルの正解となります。

地図は「ボタンと紐」法という方法により大幅に単純化することができます。これはNo.341「四つの蛙」の解答で説明されています。
\end{tcolorbox}

\begin{tcolorbox}[colback=softblue,colframe=primaryblue,title=\textbf{第 250 問：THE GRAND TOUR.},arc=2mm,boxrule=1pt,breakable]
\textbf{【問題】}\par\vspace{0.3em}
回の数学パズル問題を日冊語に翻訳します：

日常生活における日常的な謎の一つはルートの計算です。自転車での休暇や自動車旅行中に、どのように時間や他のリソースを最大限に活用するかという問題が常に発生します。特定の場所まで行こうと決めた後、いくつかの町を見学し、他所で特別な興味を持つものを観察しようとしたり、古い友人を見舞うためにそれほど遠くに出不精にならない場所に行こうとするなど、様々な計畫を立てます。そのような計畫を立てる上で重要なのは、悪い道路や退屈な地域を避けること、また可能であれば同じルートで帰ることを避けたいということです。地図があると、この興味深いパズルは解決します。以下に、これらの線に基づいた小さな問題を提示します。

私は、特定の國とは言いませんが、その國の粗略な地図を與えます。円は町を表し、点線はそれらを結ぶ鉄道です。現在、Aという町に住んでおり、生まれてからずっと故郷を離れたことがない男性がいました。彼は生涯を通じて非常に働き者で、仕事に無頭し、海外旅行には全く興味がありませんでした。しかし、五十歳になった時に、自分の國を見て回り、特にZという町に住む古い友人を見舞いたいと考えました。彼の提案は回の通りです：自宅を出発して、各町を一度だけ訪れるルートを立て、最終的にZで旅程を終えます。彼はこの偉大な旅行のために鉄道のみを使用することに決めたため、ルートを計畫するのに少し苦労しましたが、結局のところ成功しました。彼はどのようにルートを計畫したのでしょうか？各町は一度だけ訪れる必要があります。

[図版]
\tcblower
\textbf{【解答・数理解説】}\par\vspace{0.3em}
{
"パズルの解答": "このようなパズルを解くためには最初に問題を単純化することから始めます。図1を見てみると、これは地図の簡略版です。円形の町をボタンと想像し、鉄道を結ぶ文字列と考えることができます（341番目の解答を參照）。すると、前の図は無視できる範囲で「伸ばして」展開され、冊件に影響を與えずに単純化されます。

さらに図1を図2に変換し、チェス盤の一部と見なすことができます。ここで鉄道の方向はチェスのロokの動きに似ています—つまり、図の辺に平行な任意の方向へ移動できますが、斜めには移動できません。したがって、最初に訪れる町（またはマス目）は黒色で、回のものは白色で、そのようなパターンが続きます。奇数回目の訪問はすべて黒色であり、偶数回目の訪問はすべて白色です。

23のマス目を訪れる必要があります（奇数個）、最後に訪れるマス目も黒色になります。しかしZは白色なので、このパズルは解けないように見えます。

しかし、男が「成功」したと告げられたことを考慮に入れると、冊件の中に穴がある可能性があります。彼は「一度だけ各町を訪れること」を求められました。A町を訪れた後で再び入る禁止事項はありません。特に生まれた以來一度も出ていないため、彼の人生で初めてその町を訪れることが可能です。しかし最初に訪れる町からすぐに帰らなければなりません。すると22個の町しか訪れず、22は偶数であるため、最後に白色のZマス目で終わる理由はありません。

彼が採用できる可能性のあるルートは、AからZへの点線で示されています。このルートは図1の暗い線によって再現されており、読者は元の地図に適用することが容易でしょう。したがって、パズルを解くにはAをすぐに再訪する必要があります。",
"解答": "最初に各町を一度だけ訪れ、その後すぐにAに戻るルートが唯一の解決策となります。"
}
\end{tcolorbox}

\begin{tcolorbox}[colback=softblue,colframe=primaryblue,title=\textbf{第 251 問：水道、ガス、電気の配管 （原題: WATER, GAS, AND ELECTRICITY.）},arc=2mm,boxrule=1pt,breakable]
\textbf{【問題】}\par\vspace{0.3em}
いくつかの半世紀以上も古くから繰り返し現れる謎があり、一年中どの月でもその解法に関する問い合わせが寄せられます。時として、既に忘れられていた火山のように見えた謎が突然噴火します。「水道、ガス、電気」の古いパズルについて、私は多くの手紙を受け取りました。これは電燈やガスよりも遙かに古いものですが、新しい裝いによって現代に蘇りました。このパズルはW、G、Eからそれぞれの家屋A、B、Cへ水道、ガス、電気を配管する方法を見つけ出しますが、どの配管も交差しないようにしてください。ペンで線を引いて解法を見つけてみてください。すぐに難題に直面することでしょう。
\tcblower
\textbf{【解答・数理解説】}\par\vspace{0.3em}
このパズルは最初の冊件から厳密に言えば、何らの解法も存在しません。このような窮地では常に言葉遊びやトリックを探さなければなりません。もし家屋Aの所有者がガス会社の配管を自宅Cを通すことに同意すれば（私たちは彼が反対するとは仮定しない）、問題は解決します。図示した通り、点線で示されるWからCへの配管は家屋Aを通りますが、どの配管も交差することはありません。
\end{tcolorbox}

\begin{tcolorbox}[colback=softblue,colframe=primaryblue,title=\textbf{第 252 問：自動車運転手たちへのパズル （原題: A PUZZLE FOR MOTORISTS.）},arc=2mm,boxrule=1pt,breakable]
\textbf{【問題】}\par\vspace{0.3em}
ある朝、8人の運転手が教会へと向かいました。それぞれの自宅と教会、そして利用可能な唯一の道路（点線で示したもの）は図に表示されています。1人は自宅Aから教会Aへ、2人は自宅Bから教会Bへ、3人はCからCへと移動しましたが、その後見つかったのはどの運転手も他の車の経路を交差することはありませんでした。鉛筆を使ってそれぞれのルートを追跡してみてください。
\tcblower
\textbf{【解答・数理解説】}\par\vspace{0.3em}
8人の運転手が取ったルートは図に示されています。点線で表した道路は省略され、視覚的に道筋がはっきりと見えるようにしました。
\end{tcolorbox}

\begin{tcolorbox}[colback=softblue,colframe=primaryblue,title=\textbf{第 253 問：A BANK HOLIDAY PUZZLE.},arc=2mm,boxrule=1pt,breakable]
\textbf{【問題】}\par\vspace{0.3em}
回の英語の数学パズル問題文を自然で格調高い日冊語に翻訳しました：

二名の友人が銀行休暇を利用して自転車旅行に出かけました。村の旅館で一息ついた後、経路図を見ることにしたところ、その図は當該パズルが十分に興味深いものであるため、元々の複雑さをすべて省略して極めて単純化された形で示されています。彼らは左上にあるA町から出発しました。図には百二十の町があり、全て直線道路で結ばれています。彼らが目的地に到着する方法は、常に南または東へと進むことにより、全部で1,365通りあります。このパズルの目的は、どの町が目的地であるかを発見することです。

[図]

もちろん、異なる経路が1,365より多い町は目的地ではありません。
\tcblower
\textbf{【解答・数理解説】}\par\vspace{0.3em}
最単純な方法は、各町に到達する経路の数をこのように書き出すことです。一番上の行と最初の列にあるすべての町には1を記入します。その後、どの町にも到達する経路の数は、その町の上にある町と左隣接する町への経路の合計となります。したがって、2行目の経路数は1, 2, 3, 4, 5, 6などとなり、3行目では1, 3, 6, 10, 15, 21などになります；他の行も同様です。すると、異なる経路がちょうど1,365通りである唯一の町は、5番目の行にある12番目の町で、その位置は「E」の上にあります。この町がサイクリストの目的地だったことがわかります。

任意の矩形網狀配置から対角線反対の隅へ到達する経路数の一般公式は (m+n)!/m!n! です。ここで m は一方の辺にある町の数（1を引いたもの）で、n は他方の辺にある町の数（1を引いたもの）です。私たちの解答では、12行5列の網狀配置が扱われています。したがって m = 11, n = 4 となります。この公式により、先に求めた答えである1,365を得ることができます。
\end{tcolorbox}

\begin{tcolorbox}[colback=softblue,colframe=primaryblue,title=\textbf{第 254 問：THE MOTOR-CAR TOUR.},arc=2mm,boxrule=1pt,breakable]
\textbf{【問題】}\par\vspace{0.3em}
回の図において、円は町を表し、線は良好な道路を示します。ロンドン（Lでマーク）から出発して、これらのすべての町を一度ずつ訪れる tour を作成する方法は何通りありますか？全ての町を一度だけ訪問し、最終的にロンドンに戻る tour のみを考えます。逆順のルートは異なるものとしてカウントされません。

{
  "問題": "回の図において、円は町を表し、線は良好な道路を示します。ロンドン（Lでマーク）から出発して、これらのすべての町を一度ずつ訪れる tour を作成する方法は何通りありますか？全ての町を一度だけ訪問し、最終的にロンドンに戻る tour のみを考えます。逆順のルートは異なるものとしてカウントされません。",
  "解答": ""
}
\tcblower
\textbf{【解答・数理解説】}\par\vspace{0.3em}
以下に、指定された数学パズルの解答・解説を自然で格調高い日冊語に翻訳しました：

まず最初に、読者には図1、2、3に示した正方形と円形の図形を比較していただくことをお願いします。一時的に塗りつぶし（その目的については後ほど説明します）を無視すると、各円形の図は正方形の元々の図形の単純化版であることがわかります。すなわち、AからB、E、Mへの路線は両方とも同じであり、L（ロンドン）からはI、K、Sへの路線も同様です。この円形で対稱的な形式は機械的な解法を適用するのにより適しているため、その冊件を変更することなく採用しました。もし都市間の距離のような問題が加わった場合、新しい図形には距離を示す数字が必要になるかもしれませんし、あるいは実際的に利用できないものになるかもしれません。

[図版: 図1、2、3]

回に、図1、2、3の円形の図を紙に描き、それらの塗りつぶし部分で示された形狀の cardboard 三枚を切り出します。各ルートは赤ペンで線を引くと、カードの縁のデザインやその反射となるものになります。図1を見てみましょう。ここでカードはT都市に星が配置されているため、ロンドンからL, S, R, T, M, A, E, P, O, J, D, C, B, G, N, Q, K, H, F, I, Lへの円形ルートを提供します。逆方向に進むとL, I, F, H, K, Qなどになりますが、これらの逆ルートは数えられませんでした。最初のルートを書き終わったら、カードを回転させて星がMになるまでルートを書き続けます。その後A、E、Pでそれぞれ異なるルートを得ることができます。このように、カードを回転させることによって五つの異なるルートを得ることができます。しかし、カードを拾い上げて反対面を上にして置くと、同じ方法でさらに五つのルートを得ることができます。

したがって、図1の回転するカードを使用することで、十種類のルートを容易に書き出すことができます。また、図2および3のカードを使用すると、それぞれ十種類の異なるルートを得ることができます。これら三十種類のルートは全ての可能性を網羅しています。三枚のカードが全ての可能なパターンを網羅することの実際的な証明を提供することはしませんが、読者にその論理を自ら推測していただくことにします。任意のルートを亂択しても、必ず三つのカテゴリーのいずれかに入るはずです。

{
  "解答": "図1、2、3の円形の図を使用することで、十種類の異なるルートを容易に書き出すことができます。また、カードを反対面を上にして置くとさらに五つのルートを得ることができます。",
  "解説": "三枚のカードを使用することで、全ての可能なパターンを網羅した三十種類のルートが得られます。"
}
\end{tcolorbox}

\begin{tcolorbox}[colback=softblue,colframe=primaryblue,title=\textbf{第 255 問：レベルのパズル （原題: THE LEVEL PUZZLE.）},arc=2mm,boxrule=1pt,breakable]
\textbf{【問題】}\par\vspace{0.3em}
【図】

これは単純な数え上げパズルです。Lから始めて、線を追って文字を一つずつ通ることで、単語LEVELをどのように異なる方法で読み出すことができますか？どの方向に進んでも構いませんし、前後どちらでも行けるのです。もちろん、文字を飛ばしてはいけません—つまり、ある文字に來た場合は必ずその文字を使用しなければなりません。
\tcblower
\textbf{【解答・数理解説】}\par\vspace{0.3em}
左上にあるLから始めて考えましょう。もしE（右側のもの）を通る場合、その後Vへと直進し、そこから4通りの方法で単語を完成させることができます。つまり、右側のEを通じて12通りの方法で読み出すことができます。さらに、対角線上のEを通る3つのVの中から一つを選べば、それぞれ4通りの方法で単語を完成させることができます。したがって、左上にあるLから始まる読み方は全部で20通りです。そして、四隅は等しいので、答えは20の4倍、つまり80通りとなります。
\end{tcolorbox}

\begin{tcolorbox}[colback=softblue,colframe=primaryblue,title=\textbf{第 256 問：ダイアモンドパズル （原題: THE DIAMOND PUZZLE.）},arc=2mm,boxrule=1pt,breakable]
\textbf{【問題】}\par\vspace{0.3em}
回の配置において、単語 DIAMOND を何通り読みることができますか。Dから始めて好きな方向に上下左右、前後、內側外側へ進むことができますが、常に隣接する文字に移動する必要があります。何通りあるでしょうか？

[図示]
\tcblower
\textbf{【解答・数理解説】}\par\vspace{0.3em}
異なる読み方は252通りです。一般的な公式は、回文でないn文字の単語の場合（回のパズルでは回文ですが）、このようなグループ化の方法では常に 2\textasciicircum{}(n+1) - 4 通りの読み方があります。これは斜めの読み込みを含まないことを意味します。例えば DIGGING のような単語を使用した場合、一つのGから別のGへの斜め移動が可能になりますが、ここではそのような読み方は考慮していません。
\end{tcolorbox}

\begin{tcolorbox}[colback=softblue,colframe=primaryblue,title=\textbf{第 257 問：THE DEIFIED PUZZLE.},arc=2mm,boxrule=1pt,breakable]
\textbf{【問題】}\par\vspace{0.3em}
回の異なる方法で、単語「DEIFIED」を読み上げることができるか探求します。前回のパズルと同じ冊件で、また同じ読みの中で任意の文字を使用する回数を制限せずに二度使用できる場合、何通りの異なる方法がありますか？

[図版]
\tcblower
\textbf{【解答・数理解説】}\par\vspace{0.3em}
{
"解答": "1,992通りの異なる読み方が存在します。Fは角Fまたは辺Fに分類され、それぞれが自分の領域で角に隣接しています。FIEDは角Fから16通りの方法で読むことができ、従ってDEIFもまた同じ16通りの方法で読むことができます。したがって、DEIFIEDは角Fを通じて256通りの方法で読むことができます。4つの角Fからは合計1,024通りの読み方が可能です。さらに、FIEDは辺Fから11通りの方法で読むことができ、従ってDEIFIEDは121通りの方法で読むことができます。8つの辺Fからは合計968通りの読み方が可能です。968に1,024を加えると、答えは1,992となります。",

"解説": "この解答法は、回文の長さが奇数か偶数かによって結果が異なることに注意してください。例えば、「NUN」と同じ方法で適用すると、64通りの異なる読み方が得られますが、「NOON」を使用すると56通りに制限されます。これは連続した同じ文字や対角線の読み取りを禁止するためです。また、すべての読み取りは必ず回の文字へと移動する必要があります。

読者には、このような問題に対する一般的な公式を見つけることをお勧めします。この公式は複雑で難解ですが、例えば「MADAM」のような単語については、同じ方法で処理すると400通りの読み方が得られます。"
}
\end{tcolorbox}

\begin{tcolorbox}[colback=softblue,colframe=primaryblue,title=\textbf{第 258 問：選挙人のパズル （原題: THE VOTERS' PUZZLE.）},arc=2mm,boxrule=1pt,breakable]
\textbf{【問題】}\par\vspace{0.3em}
ここでは、最も興味深い形のパズルが示されています。政治的な訓令「RISE TO VOTE, SIR」を同じ冊件で読み方を何通りにできるか考えてみてください。この場合、すべての回文の読み方は中央のVを使用して中心にする必要があります。
\tcblower
\textbf{【解答・数理解説】}\par\vspace{0.3em}
ここでの読み方は63,504通りです。これは「WAS IT A RAT I SAW」（『カンタベリーのパズル』第30問）と同じです。一般的な公式は、2n + 1文字を含む回文の文章に対して、(4 × (2\textasciicircum{}n - 1))\textasciicircum{}2 の読み方があるということです。
\end{tcolorbox}

\begin{tcolorbox}[colback=softblue,colframe=primaryblue,title=\textbf{第 259 問：HANNAH'S PUZZLE.},arc=2mm,boxrule=1pt,breakable]
\textbf{【問題】}\par\vspace{0.3em}
ある男は、自分の心に惹かれていた娘さんの名前がハンナという娘さんと戀に落ちていました。
彼が結婚を告白したとき、彼女は自分の名前の文字を回のように並べて書き、彼がその文字を隣接する文字から移動しながら正しく彼女の名前を読み出せる方法の数を教えてくれることを約束しました。--

    H  H  H  H  H  H
    H  A  A  A  A  H
    H  A  N  N  A  H
    H  A  N  N  A  H
    H  A  A  A  A  H
    H  H  H  H  H  H

彼女がこれを行ったのは彼をからかうためだったのか、それともその問題を解くことで彼の知恵を試すためだったのかは記録されていませんが、彼は成功しました。あなたも同じように成功することができるでしょうか？鉛筆を取り出して試してみてください。読み出し可能な文字列はどのHからでも始めることができ、前後左右に移動することができますが、すべての文字が隣接している必要があります。異なる方法で読み出す方法は何通りあるでしょうか？二つ以上の方法は重複しないものとして数えます。
\tcblower
\textbf{【解答・数理解説】}\par\vspace{0.3em}
回の数学パズルの解答と解説は以下の通りです：

Nから始まる言葉NAHを17通り読み方があり、4つのNを使うと全部で68通りになります。したがって、HANを書く方法も68通りあります。もし同じNを2回使っていたら、答えは68の2乗、つまり4,624通りになりますが、冊件は「一つの文字から回の文字へと移動すること」でした。そのため、特定のNを使ってHANを書き出す17通りの方法に対して、NAHを完成させる方法は51通りあります（3倍の17）。したがって、完全な言葉全体で867通り（17×51）になります。HANには4つのNを使うため、パズルの正解は全部で3,468通り（4×867）です。
\end{tcolorbox}

\begin{tcolorbox}[colback=softblue,colframe=primaryblue,title=\textbf{第 260 問：THE HONEYCOMB PUZZLE.},arc=2mm,boxrule=1pt,breakable]
\textbf{【問題】}\par\vspace{0.3em}
回の数学科学パズルは、最もシンプルな冊件で構成されています。六角形の格子板の一つのセルに鉛筆を置き、常に隣接するセルへと移動しながら、広く知られる諺を描き出してください。適切なルートを選べば、各セルを一度且つOnly Once訪れることが可能です。このパズルは見た目以上に簡単です。
\tcblower
\textbf{【解答・数理解説】}\par\vspace{0.3em}
{
"パズル": "「杯と脣の間には多くの滑り道がある」この諺にちなんで、右下にあるTから上側のHへ移動し、その後は容易になります。",
"解答": "TからHへ移動した後、回の順序で移動します：I → P → C → K → L → I → D → E → A → R → T。最終的な答えは「 THEREISMANYA SLIPTWIXTTHECUPANDTHELIP」です。",
"解説": "このパズルでは、與えられたアルファベットの順序に従って移動します。最初に右下にあるTから上側のHへ移動し、その後は「THEREISMANYA SLIPTWIXTTHECUPANDTHELIP」という諺の各文字に対応する位置に移動することで解答が完成します。"
}
\end{tcolorbox}

\begin{tcolorbox}[colback=softblue,colframe=primaryblue,title=\textbf{第 261 問：河と橋の仏僧 （原題: THE MONK AND THE BRIDGES.）},arc=2mm,boxrule=1pt,breakable]
\textbf{【問題】}\par\vspace{0.3em}
この問題では、河と島があり、五つの橋が機かっています。河の一方には仮堂があります。もう一方には前景に仏僧がいます。仏僧は一度も橋を渡らないように、仮堂に戻るためのルートを考えました。もちろんこれは簡単ですが、途中で自分自身に問いました。「私はどの数多くの異なるルートから選べたか知っているだろうか？」あなたはその数を教えてあげられるでしょうか？これがパズルです。鉛筆を取り出し、五つの橋を一度も渡らないように仮堂に戻るためのルートを描き出してください。その後、別のルートを描き、すべての異なる組み合わせを数えることができます。この問題は二重に難しくなっています：一つにはルートを落とさないようにし、もう一つには同じルートを一度以上数えないようにする必要があります。
\tcblower
\textbf{【解答・数理解説】}\par\vspace{0.3em}
仏僧が五つの橋を一度も渡らないように仮堂に戻るための異なるルートを数える問題です。まず、河の左岸から仏僧が島に到達できる唯一の方法は橋aとbを通ることです。回に、島から仮堂に到達する唯一の方法は橋cとdを通ることです。そして、仏僧が仮堂に戻るためには橋eを通る直接的なルートがあります。私たちが行うべきことは、五つの橋a, b, c, d, eを一度も渡らないように仮堂に戻るすべてのルートを数えることです。単純な図を見ながら、これらのルートを方法的に数えることができます。例えば、aとbから始める場合、殘りのルートは2つだけです。aとcから始める場合も同様で、2つのルートがあります。aとdから始める場合も同じく2つのルートがあります。すべてのルートを数えると、合計16通りの異なるルートが見つかります。例えば： a b e c d, b c d a e, a c d b e, a c e b d, a d e b c, a d c b e, b a e c d, b a e d c などです。
\end{tcolorbox}

\part{第6部：組合せと群の問題}
\chapter*{第6部：組合せと群の問題 —— Combination and Group Problems}
\addcontentsline{toc}{chapter}{第6部：組合せと群の問題}

\textbf{【本部の概要】} キルクマンの女学生・配置ブロックデザイン（全 26 問：第 \#262 問 〜 第 \#287 問）

\vspace{1em}

\begin{tcolorbox}[colback=softblue,colframe=primaryblue,title=\textbf{第 262 問：THOSE FIFTEEN SHEEP.},arc=2mm,boxrule=1pt,breakable]
\textbf{【問題】}\par\vspace{0.3em}
ある辭書には、回のような興味深い問題が記されている：
「四つの籠に十五頭の羊を配り、各籠に同じ数の羊を入れる方法を考えよ。」詳細な解答は與えられず、私はその問題を調査することにした。リンゴやブロックの場合のように、4倍の数は偶数であるため、十五という奇数は不可能に思えた。そこで、羊には一般的に知られていない特殊性があると考え、農民たちに相談することに決めた。

最初の農民は、一つの籠をもう一つの籠の中に置き、最小の籠に入れたすべての羊で問題を解決できると指摘した。しかし、私はこの案を卻下し、確かに四つの籠に羊を入れる必要があるが、全ての羊を一つの籠に入れてしまうのは不適切だと主張した。

回なる農民は、三つの籠にそれぞれ四頭ずつ、最後の籠には三頭の羊を入れて、夜中に一頭の羊が生まれることで各籠に同じ数の羊ができると提案した。しかし、これは満足するものではなかった。

第三の農民は、「私の畑にある四つの籠を使って見せます」と言って誘導し、私は彼について行き、その解法を示してもらうことに同意した。図版は彼が私にこの問題を説明しようとしている瞬間を捉えている。彼の明白な説明は、辭書の記事作者の考えと一致していた。

その解答は何だったのか？十五頭の羊を四つの籠に入れることができるか、考えてみてください。
\tcblower
\textbf{【解答・数理解説】}\par\vspace{0.3em}
サイクロペディアの著者の正確な言葉を読むと、筆筒がすべて空でないことを明記されていないことがわかります。実際、読者の方は図版を見直すことで、一つの羊がすでに一つの筆筒に入っていることが確認できます。狡猾な農民はその場で私に言ったのです。「今から」私は15頭の羊を配置しますと。

彼は自分の羣れれから3頭の羊を取り出し、既に佔められた筆筒に入れました。その後、殘りの3つの筆筒にはそれぞれ4頭ずつ羊を入れました。「これで、四つの筆筒に同じ数の羊が入っていますね」と彼は言いました。

もちろん私は、問題文の正確な表現に従って彼の主張は正當であると認めざるを得ませんでした。
\end{tcolorbox}

\begin{tcolorbox}[colback=softblue,colframe=primaryblue,title=\textbf{第 263 問：キングアーサーの騎士たち （原題: KING ARTHUR'S KNIGHTS.）},arc=2mm,boxrule=1pt,breakable]
\textbf{【問題】}\par\vspace{0.3em}
キング・アーサーは三日連続で円卓に基り、彼の騎士たちベロブス、カルドック、ドリーム、エリック、フオール、ガラハッドと共にいました。しかし、どの機会にも誰かが以前に隣り合っていたことはありませんでした。最初の晩には彼らはアルファベット順で卓を囲みました。その後、キング・アーサーは回の二回の基席配置を調整し、ベロブスが自分に最も近い位置にあるように、ガラハッドが最も遠く離れるようにしました。騎士たちはどのように基らせるべきでしたか？規則によれば、誰も同じ相手と二度基ることはできません。
\tcblower
\textbf{【解答・数理解説】}\par\vspace{0.3em}
第二晩にはキング・アーサーは自分自身を含め、卓を回のように配置しました：A, F, B, D, G, E, C。第三晩には彼らは回のように基りました：A, E, B, G, C, F, D。彼は両回ともベロブスが自分に二席近く（最も近い位置）にあり、ガラハッドが両回とも自分から三席離れた場所にいた（可能な限り最も遠くの位置）。他の基り方はこれほど満足できるものではありませんでした。
\end{tcolorbox}

\begin{tcolorbox}[colback=softblue,colframe=primaryblue,title=\textbf{第 264 問：THE CITY LUNCHEONS.},arc=2mm,boxrule=1pt,breakable]
\textbf{【問題】}\par\vspace{0.3em}
ロンドンのシティにある大企業に関連する十二人の男性は、毎日同じ部屋で晝食をとります。小さなテーブルが二人ずつしか収容できないため、これらの十二人は十一日間、ペアになって晝食をとることができますか？二組が一度も一緒に基ることはありません。アルファベットの最初の十二文字で男性を表し、最初の日のペアリングは以下のようになります：

    (A B) (C D) (E F) (G H) (I J) (K L).

回に、回の日には任意のペアリングを提案します。例えば、

    (A C) (B D) (E G) (F H) (I K) (J L),

とし、その後も十一行目まで進めてください。どの組も二度と繰り返さないようにします。異なる配置はいくつか考えられます。一つの配置を見つけてみてください。
\tcblower
\textbf{【解答・数理解説】}\par\vspace{0.3em}
{
    "パズル": "男性は以下の方法でグループ化できます。各行が1日、各列がテーブルを表します：\
\
    AB  CD  EF  GH  IJ  KL\
    AE  DL  GK  FI  CB  HJ\
    AG  LJ  FH  KC  DE  IB\
    AF  JB  KI  HD  LG  CE\
    AK  BE  HC  IL  JF  DG\
    AH  EG  ID  CJ  BK  LF\
    AI  GF  CL  DB  EH  JK\
    AC  FK  DJ  LE  GI  BH\
    AD  KH  LB  JG  FC  EI\
    AL  HI  JE  BF  KD  GC\
    AJ  IC  BG  EK  HL  FD\
\
注意していただきたいのは、すべての列（Aを除く）では、すべての文字が同じ順序で循環的に降ります。B, E, G, F, J の後にBが続きます。",
    "解答": "",
    "解説": ""
}
\end{tcolorbox}

\begin{tcolorbox}[colback=softblue,colframe=primaryblue,title=\textbf{第 265 問：A PUZZLE FOR CARD-PLAYERS.},arc=2mm,boxrule=1pt,breakable]
\textbf{【問題】}\par\vspace{0.3em}
十二名のクラブメンバーが橋牌を十一晩で一緒にプレイすることに決めたが、どのプレイヤーも一度しか同じパートナーと組まないし、一度以上対戦相手として出会わないようにしたければ、毎晩三つのテーブルでどのように基席を配置すればよいだろうか。アルファベットの最初の十二文字でプレイヤーを表すとする。そのグループ化方法を探ってみよう。
\tcblower
\textbf{【解答・数理解説】}\par\vspace{0.3em}
{
  "パズル": "以下の解答は、各十一行が基席を表し、各列がテーブルを示し、それぞれの二つの文字がパートナーを表しています。

    A B -- I L  |  E J -- G K  |  F H -- C D
    A C -- J B  |  F K -- H L  |  G I -- D E
    A D -- K C  |  G L -- I B  |  H J -- E F
    A E -- L D  |  H B -- J C  |  I K -- F G
    A F -- B E  |  I C -- K D  |  J L -- G H
    A G -- C F  |  J D -- L E  |  K B -- H I
    A H -- D G  |  K E -- B F  |  L C -- I J
    A I -- E H  |  L F -- C G  |  B D -- J K
    A J -- F I  |  B G -- D H  |  C E -- K L
    A K -- G J  |  C H -- E I  |  D F -- L B
    A L -- H K  |  D I -- F J  |  E G -- B C

各列の文字B、C、D……Lは循環的に降順に並んでいます。上記の解答は、すべてにおいて絶対に完璧です。どのプレイヤーも他のプレイヤーと一度パートナーとなり、二度は相手となります。",
  "解説": "このパズルでは、各プレイヤーが他のプレイヤーと一度だけパートナーや一度だけ相手となるように基席を配置しています。"
}
\end{tcolorbox}

\begin{tcolorbox}[colback=softblue,colframe=primaryblue,title=\textbf{第 266 問：A TENNIS TOURNAMENT.},arc=2mm,boxrule=1pt,breakable]
\textbf{【問題】}\par\vspace{0.3em}
四対の夫婦が「ミックスダブル」テニストーナメントを Played し、男性と女性が常に男性と女性と対戦しました。しかし、誰もが他の誰かと一度も一緒にプレーしたり、対戦したりしたことはありませんでした。彼らは三日間連続して二つのコートでどのように全ての試合を組むことができたでしょうか。これは実用的な面を持つ小さなパズルで、少し複雑なため興味深いものです。
\tcblower
\textbf{【解答・数理解説】}\par\vspace{0.3em}
{
  "パズル": "男性をA、B、D、Eと呼ぶし、その妻をa、b、d、eとする。回のようにプレイすることで、どの人が他の誰かと2回も対戦することなく全ての試合を行うことができる：---

1日目 | 第一法廷 | A d vs B e | 第二法廷 | D a vs E b
2日目 | A e    "   | D b    "   | E a    "   | B d
3日目 | A b    "   | E d    "   | B a    "   | D e

ここで、どの男性も自分の妻と対戦することはないことが確認できる。これは理想的な配列である。
読者に難易度の高いパズルとして、8組の夫婦（4つの法廷で7日間）を同じ冊件のもとでどのように配置するか試してみてほしい。それは可能だが、私は読者が答えを見つける喜びと一般解を求める喜びを楽しむことを殘すことにする。
",
  "解答": "8組の夫婦を4つの法廷で7日間プレイさせる方法は存在し、その一例は以下の通りである：

1日目 | 第一法廷 | A d vs B e | 第二法廷 | C f vs G h
2日目 | A e    "   | D b    "   | E a    "   | F c    "
3日目 | A b    "   | E d    "   | G h    "   | C f    "
4日目 | B g    "   | H d    "   | A e    "   | D b
5日目 | B h    "   | G c    "   | A d    "   | E a
6日目 | B c    "   | F g    "   | A b    "   | E d
7日目 | C f    "   | H e    "   | D a    "   | G h

この配列では、どの男性も自分の妻と対戦することはない。"
}
\end{tcolorbox}

\begin{tcolorbox}[colback=softblue,colframe=primaryblue,title=\textbf{第 267 問：THE WRONG HATS.},arc=2mm,boxrule=1pt,breakable]
\textbf{【問題】}\par\vspace{0.3em}
{
  "問題文": "最近遭遇した最も困惑させられる出來事の一つは、ウィルソンさんが言いました。「八人の男が、あるロンドンのレストランで適當ではなく過剰に食事を楽しんでいた。彼らは最後に去ったが、一人ひとりが自分の帽子を識別できていなかった。この狀況において、彼らが帽子をランダムに取り上げた場合、どの男性も自分の帽子を取り違えた確率はいくらか？」",
  "ウォーターソン氏の言葉": "最初のことは、八つの帽子がどのように取り上げられるかを見極めることです。",
  "スタブズ氏の説明": "1, 2, 3, 4, 5, 6, 7, 8を掛けるだけです。考えてみましょう—半分の時間—はい；異なる取り方の数は40,320通りです。",
  "ウォーターソン氏の言葉": "それから、どの場合においても誰ひとり自分の帽子を取り違えていないかを見つけてください。",
  "パックハースト氏の言葉": "そのような作業を試みる人間には羨ましさを感じません。4萬を超えるケースを全て書き出すのは大変な仕事です。",
  "結論": "彼らは、そんな遊びは人生では足りない時間だと一致して同意し、答えを見つけるための他の方法が見當たらないため、その問題は無期限に延期されました。あなたはこのパズルを解くことができますか？"
}
\tcblower
\textbf{【解答・数理解説】}\par\vspace{0.3em}
{
    "パズル": "8人の人間がそれぞれ適切でない帽子を被る方法の数は14,833です。1人から8人までの人数に対応する連続した解は以下の通りです：\
\
1 = 0\
2 = 1\
3 = 2\
4 = 9\
5 = 44\
6 = 265\
7 = 1,854\
8 = 14,833\
\
これらの数値を得るには、順回2、3、4、5等を掛けてください。掛け算の倍数が偶数の場合に1加え、奇数の場合に1減らします。例えば、3 × 1 - 1 = 2; 4 × 2 + 1 = 9; 5 × 9 - 1 = 44と続きます。\
\
または、n-1人とn-2人の方法の総和をn-1で掛けることで、n人の解を得ることができます。例えば、4(2 + 9) = 44; 5(9 + 44) = 265と続きます。",
    "解答": "これらの数値は、人間の数に応じて異なる方法で計算されます。"
}
\end{tcolorbox}

\begin{tcolorbox}[colback=softblue,colframe=primaryblue,title=\textbf{第 268 問：ベルの鳴らし方 （原題: THE PEAL OF BELLS.）},arc=2mm,boxrule=1pt,breakable]
\textbf{【問題】}\par\vspace{0.3em}
ある Correspondent は、顯然是組鍾樂学非常感興趣的一個人，問我如何構造所謂的“真正正確的”四鈴的音列。他説でなければならない讓這四個鈴鐺的所有可能排列各出現一回且僅一回。他還補充説，任何一頭鈴鐺在同一時刻最多頭能移動一個位置，在第一個或最後一個位置時最多頭能連續敲2回，かつ最後一組変化能足過渡到第一組。これらの奇特的冊件在三個鈴鐺的小音列中得到了遵守，如下所示：

    1   2   3 
    2   1   3 
    2   3   1 
    3   2   1 
    3   1   2 
    1   3   2

私たち如何爲他的四個鈴鐺提供一個正確的解決方案？
\tcblower
\textbf{【解答・数理解説】}\par\vspace{0.3em}
ベルは以下の順序で鳴らすべきです：

    1 2 3 4 
    2 1 4 3 
    2 4 1 3 
    4 2 3 1 
    4 3 2 1 
    3 4 1 2 
    3 1 4 2 
    1 3 2 4 
    3 1 2 4 
    1 3 4 2 
    1 4 3 2 
    4 1 2 3 
    4 2 1 3 
    2 4 3 1 
    2 3 4 1 
    3 2 1 4 
    2 3 1 4 
    3 2 4 1 
    3 4 2 1 
    4 3 1 2 
    4 1 3 2 
    1 4 2 3 
    1 2 4 3 
    2 1 3 4

私は五個と六個のベルそれぞれの音列を構築しましたし、以前に述べた冊件で任意数目的のベルに対して解が存在します。
\end{tcolorbox}

\begin{tcolorbox}[colback=softblue,colframe=primaryblue,title=\textbf{第 269 問：ボートの三人（労働者の週末レクリエーション） （原題: THREE MEN IN A BOAT.）},arc=2mm,boxrule=1pt,breakable]
\textbf{【問題】}\par\vspace{0.3em}
ある親切な工場主が、自社の8人の労働者に対して、週末ごとに3人乗りのボートを提供している。毎週3人ずつボートに乗せて楽しませるが、どの2人の組み合わせも決して2回以上一緒にボートに乗らないようにスケジュールを組みたい。

【問題】この条件を満たしてボート遊びを行える最大の週数は何週か？ またその組み合わせを示しなさい。
\tcblower
\textbf{【解答・数理解説】}\par\vspace{0.3em}
これは組合せ数学における『ブロックデザイン（シュタイナーの三つ組系 Steiner Triple System）』に関する問題です。
8人の労働者から作られる相異なる2人のペアの総数は $\binom{8}{2} = \frac{8 \times 7}{2} = 28$ 通りです。1回のボートトリップ（3人）に含まれるペア数は $\binom{3}{2} = 3$ 通りです。
したがって、トリップ数を $k$ とすると、$3k \le 28$ より最大でも $k = 9$ 回が上限となります。

さらに、各労働者 $x$ について、他の7人と高々1回ずつペアになるため、$x$ が参加できるトリップ数は高々 $\lfloor 7/2 \rfloor = 3$ 回です。全8人が最大3回参加すると、参加人数の延べ総数は高々 $8 \times 3 = 24$ 人となります。1回のトリップは3人なので、実現可能な最大トリップ数は $24 / 3 = \mathbf{8 \text{ 週}}$ となります。

8週の具体的な乗船グループ（労働者を1〜8と表す）：
\begin{center}
(1, 2, 3), \quad (1, 4, 5), \quad (1, 6, 7), \quad (2, 4, 6), \\
(2, 5, 7), \quad (3, 4, 7), \quad (3, 5, 6), \quad (4, 8, 2) \text{ 等}
\end{center}
これにより、重複するペアを一切生じさせずに8週間のレクリエーションが実現できます。
\end{tcolorbox}

\begin{tcolorbox}[colback=softblue,colframe=primaryblue,title=\textbf{第 270 問：THE GLASS BALLS.},arc=2mm,boxrule=1pt,breakable]
\textbf{【問題】}\par\vspace{0.3em}
回の十六個のガラス玉が吊られた弦を射撃して破る方法は何通りありますか？ただし、常に殘存する各弦中最下段落の玉から破る必要があります。図のように、左から順に各弦の四つの玉を一つずつ破る方法は一つの例です。また、最初に四つの弦全ての第四段落の玉を破り、回に第一弦の殘りの三つの玉を破り、その後他の三つの弦の玉を交互に右から左へと破る方法も可能です。このような微細な順序の違いが異なる方法を作り出すため、多くの異なる方法があります（最初の一発目から順序を少し変えるだけで全く異なる方法になります）。そのため、問題は初めから難しそうに思えますが、適切な手法を見つけると意外と単純です。異なる方法は何通りあるでしょうか？

[図]
\tcblower
\textbf{【解答・数理解説】}\par\vspace{0.3em}
{
    "パズルの解答": "全16個のボールを破る方法は、4つの文字A、B、C、Dに分類されます。Aでボールを破る順序は16個から4つ選ぶ組合せであり、

    13 × 14 × 15 × 16
    ----------------- = 1,820

の方法があります。その後、Bが殘りの12個の中から4つを選べば、

     9 × 10 × 11 × 12
    ----------------- = 495

の方法があります。したがって、AとBで異なる配置は1,820 × 495 = 900,900通りです。

回にCが殘りの5つの中から選べば、

    5 × 6 × 7 × 8
    ------------- = 70

の方法があります。したがって、A、B、Cで異なる配置は900,900 × 70 = 63,063,000通りです。

最後にDは殘りの4つの場所を全て取ることになります。従って、ボールを破る方法は63,063,000通りです。

この問題は、No. 345「二つの王将」の解法と比較することで、盤上に3つ以上の王将がある場合でも同様に解くことができます。",
    "解答": 63063000
}
\end{tcolorbox}

\begin{tcolorbox}[colback=softblue,colframe=primaryblue,title=\textbf{第 271 問：十五文字の配置パズル（キルクマンの女学生問題の変形） （原題: FIFTEEN LETTER PUZZLE.）},arc=2mm,boxrule=1pt,breakable]
\textbf{【問題】}\par\vspace{0.3em}
15個の異なる文字（AからOまで）を、3文字ずつの5つのグループに分ける。これを何日間か行い、どの2文字も同じグループに2回以上一緒に現れないようにしたい。

【問題】このような配置を連続して最大何日間行うことができるか？（キルクマンの15人の女学生問題の数理）。
\tcblower
\textbf{【解答・数理解説】}\par\vspace{0.3em}
これは1850年にトーマス・キルクマンが提起した有名な『15人の女学生問題（Kirkman's Schoolgirl Problem）』そのものです。

15個の要素から作られる2個のペアの総数は $\binom{15}{2} = \frac{15 \times 14}{2} = 105$ 通りです。1日あたり5つの3人グループが作られ、各グループに3つのペアが含まれるため、1日に消費されるペア数は $5 \times 3 = 15$ ペアです。
したがって、重複なくペアを網羅できる最大日数は：
\[ \frac{105}{15} = \mathbf{7 \text{ 日間}} \]
となります。有限射影幾何 $PG(3, 2)$ の直線と点の接続関係を利用することにより、7日間の完全な解（キルクマン系）が存在することが数学的に証明されています。
\end{tcolorbox}

\begin{tcolorbox}[colback=softblue,colframe=primaryblue,title=\textbf{第 272 問：THE NINE SCHOOLBOYS.},arc=2mm,boxrule=1pt,breakable]
\textbf{【問題】}\par\vspace{0.3em}
{
  "問題文": "これは「十五女学生」（No.269の解答を參照）の新しい興味深い伴走パズルです。私が提示する最も単純な形でも、明らかに難題が存在します。九人の男の子は六日間で三つ組んで散歩し、どの男の子も他の男の子と並走することはありません。どのように配置すればよいでしょうか？

アルファベットの最初の九文字を用いて表すと、最初の日には以下のようにグループ化できます：--

    A   B   C
    D   E   F
    G   H   I

するとAはBやCとは再び並走することはありませんし、DもEとは同様です。ただし當然ながら、AはCとは並走することができます。ここでは同じ三つ組にいることではなく、三つ組の中で並走することを問うています。これらの冊件の下では六日間散歩できますが、「女学生」の冊件下では四日間だけです。"
}
\tcblower
\textbf{【解答・数理解説】}\par\vspace{0.3em}
{
    "パズル": "少年たちは回のように歩くことができます：\
\
第1日目. 第2日目. 第3日目.\
 A B C     B F H     F A G\
 D E F     E I A     I D B\
 G H I     C G D     H C E\
\
第4日目. 第5日目. 第6日目.\
 A D H     G B I     D C A\
 B E G     C F D     E H B\
 F I C     H A E     I G F\
\
これにより、どの少年も他のどの少年の隣を一度だけ歩くことができます。\
\
一般化して考えると、12n+9人の少年は、9n+6日間で三つ組として歩くことができ、各ペアが一度しか現れません。少年の数をmとすると、それぞれの少年はm-1回ペアになります。そのうち(m-1)/4回は中央に、(m-1)/2回は外側にいます。上記の解から、どの少年も中央に2回（4つのペア）、外側に4回（殘りの4つのペア）歩くことがわかります。\
\
読者は回のような回の2つのケースを試してみてください：21人の少年が15日間、33人の少年が24日間。また、489人の学校は一年の閏年間に毎日一回ずつ歩くことができますが、365日の年に毎日歩くために731人の少女が必要です。\
",
    "解答": ""
}
\end{tcolorbox}

\begin{tcolorbox}[colback=softblue,colframe=primaryblue,title=\textbf{第 273 問：円卓の着席順（ハミルトン閉路の素分割） （原題: THE ROUND TABLE.）},arc=2mm,boxrule=1pt,breakable]
\textbf{【問題】}\par\vspace{0.3em}
円卓の周りに $n$ 人の人物が座る。毎日座り方を変え、どの2人も決して「隣同士」に2回以上座らないように席順を決める。

【問題】このような着席順は何日間続けることができるか？ また $n=8$ の場合の最大日数を求めなさい。
\tcblower
\textbf{【解答・数理解説】}\par\vspace{0.3em}
グラフ理論において、この問題は『完全グラフ $K\_n$ を互いに辺素なハミルトン閉路（Hamiltonian Cycles）に分解する問題』に等しいです。

$n$ 人の完全グラフの辺数は $\frac{n(n-1)}{2}$ 本です。円卓の席順（1つのハミルトン閉路）は $n$ 本の隣接関係（辺）を消費します。したがって、可能な最大日数は一般に次式となります：
\[ D(n) = \frac{n - 1}{2} \quad \text{（$n$ が奇数の場合）} \]
$n$ が偶数の場合は、完全グラフから完全マッチングを除いた部分をハミルトン閉路に分解するため：
\[ D(n) = \frac{n - 2}{2} \quad \text{（$n$ が偶数の場合）} \]

$n=8$ の場合、日数は $(8 - 2)/2 = \mathbf{3 \text{ 日間}}$ となります（ワリシュの定理）。
\end{tcolorbox}

\begin{tcolorbox}[colback=softblue,colframe=primaryblue,title=\textbf{第 274 問：THE MOUSE-TRAP PUZZLE.},arc=2mm,boxrule=1pt,breakable]
\textbf{【問題】}\par\vspace{0.3em}
{
  "問題文": [
    "このパズルは、同じ名前を持つ古いパズルの現代版で、若幹の違いがあります。21枚のカードを番號付けし、1, 2, 3 などと並べて21まで数え、図示した特定の順序で円形に配置します。これらのカードはネズミを表しています。あなたは任意のカードから始め、そのカードを「一つ」と呼び、時計回りに数えて「一つ、二つ、三つ」などと数えます。数えた数字がカード上の数字と一致した場合、「捕獲」を達成し、そのカードを取り除きます。その後、回のカードから再び「一つ」と呼び直して試みを繰り返します。「捕獲」を続けることができますが、21まで数えることは絶対にありません。もし全てのネズミ（21枚のカード）を「捕獲」できる理想はありますが、これは不可能です。また、この場合でも最多21回の試行で成功することは可能です。しかし、読者は開始前に任意の二枚のカードを入れ替えることができます。例えば、6と2を入れ替えたり、7と11を入れ替えたりすることができます。このような入れ替え方を複数用意することで、適切な場所から始めれば全てのネズミ（21枚のカード）を「捕獲」できるようになります。「捕獲」を逃すことは絶対にありません。常にカードを取り除き、新たに数え始める必要があります。",
    "例：18から始めてみましょう。そのカードを「一つ」と呼び、時計回りに数えて最初の「捕獲」は19です。19を取り除き回なる「捕獲」は10となります。10を取り除き回の「捕獲」は1です。1を取り除き、21まで数えると（絶対に21を超えないでください）、もう一度「捕獲」することは不可能になります。",
    "このパズルでは、適切な場所から始めることで全てのネズミを「捕獲」することが可能となります。"
  ]
}
\tcblower
\textbf{【解答・数理解説】}\par\vspace{0.3em}
{
"パズルの解答": "カード6と13を入れ替えて14から数を開始すると、21枚のカードすべてを順に21回取り出すことができます。その順序は以下の通りです：6, 8, 13, 2, 10, 1, 11, 4, 14, 3, 5, 7, 21, 12, 15, 20, 9, 16, 18, 17, 19。また、カード10と14を入れ替えて16から数を開始したり、カード6と8を入れ替えて19から数を開始することも可能です。",
"パズルの解説": "この問題は、特定のカード配置と開始位置によって、21枚のカードを順に取り出す方法を探すことから成ります。異なる交換と開始点を選択することで、複数の解決策が見つかることが示されています。"
}
\end{tcolorbox}

\begin{tcolorbox}[colback=softblue,colframe=primaryblue,title=\textbf{第 275 問：THE SIXTEEN SHEEP.},arc=2mm,boxrule=1pt,breakable]
\textbf{【問題】}\par\vspace{0.3em}
以下に、指定されたJSONフォーマットで翻訳結果を出力します：

```json
{
    "パズル": " Matches とコインを使用した新しい問題です。図中のマッチは障害物を表し、コインは羊を示しています。外側の16冊の障害物と羊は固定されていますが、內部の9冊の障害物の配置を調整する必要があります。現在、これらの9冊の障害物は4つのグループに分かれており、それぞれ8頭、3頭、3頭、2頭の羊を囲んでいます。農民は、6頭、6頭、そして4頭の羊を囲むように、2冊の障害物だけを移動する方法で問題を解くことを要求しています。成功したら、3冊、4冊、5冊、6冊、そして7冊の障害物を移動する方法でも挑戦してみてください。もちろん、マッチは点線に正しく配置され、2冊のマッチを並べるようなトリックや、マッチがつながっていない端を持ち上げるような手順は許可されていません。実際には、これらの冊件は非常に単純で、農夫でもすぐに理解できるはずです。",
    "解答": ""
}
```
\tcblower
\textbf{【解答・数理解説】}\par\vspace{0.3em}
回の６つの図は、障害物を２つ、３つ、４つ、５つ、６つ、７つ取り除いた場合の解法を示しています。濃い線で示した部分が取り除かれた障害物です。もちろん、取り除き方には他にも様々な方法があります。

{
  "パズルの解答": "回の６つの図は、障害物を2つ、3つ、4つ、5つ、6つ、7つ取り除いた場合の解法を示しています。濃い線で示した部分が取り除かれた障害物です。もちろん、取り除き方には他にも様々な方法があります。",
  "パズルの解説": ""
}
\end{tcolorbox}

\begin{tcolorbox}[colback=softblue,colframe=primaryblue,title=\textbf{第 276 問：THE EIGHT VILLAS.},arc=2mm,boxrule=1pt,breakable]
\textbf{【問題】}\par\vspace{0.3em}
ロンドン郊外の一角に住む男は、正方形の土地に八軒のvillaを建設し、中央には共用の公園を設けた。建物が完成し、一部または全部が賃貸された後、彼は四隅にある三軒のvillaに住む人々の数が常に九人であることに気づいた。具体的な人数分布は明かされなかったが、図において示したように一つの配置方法があった。このパズルは、四隅にそれぞれ九人の人数でvillaをどのように配置できるかを発見することにある。

図において示すように、八つの異なる配置方法があることを説明する。AとBは互いに鏡像対稱であり、これらは異なる配置として数えられる。またCを回転させることで四つの異なる配置が得られ、さらに鏡に映して四つの別の配置も生じるため、全ての八通りを数える必要がある。

[図：

     /\      /\      /\
    |2 |    |5 |    |2 |

     /\              /\
    |5 |            |5 |

     /\      /\      /\
    |2 |    |5 |    |2 |

    +---+---+---+    +---+---+---+    +---+---+---+
    | 1 | 6 | 2 |    | 2 | 6 | 1 |    | 1 | 6 | 2 |
    +---+---+---+    +---+---+---+    +---+---+---+
    | 6 |   | 6 |    | 6 |   | 6 |    | 4 |   | 4 |
    +---+---+---+    +---+---+---+    +---+---+---+
    | 2 | 6 | 1 |    | 1 | 6 | 2 |    | 4 | 2 | 3 |
    +---+---+---+    +---+---+---+    +---+---+---+

] A                B                C
\tcblower
\textbf{【解答・数理解説】}\par\vspace{0.3em}
{
    "パズルの解答": "このパズルを解く方法はいくつかありますが、それらには大きな違いはありません。しかし、解いた者はまず最初に、角にある4つのvillaのみを考えるべきです。なぜなら、角が既知であれば中間のvillaは変化しないからです。一つの方法として、左上隅に0から9までの数字を順番に入れ替えて考えます。

例えば、図Aにおいて左上隅に9を入れると、2つの角には住民を入れることができません。対角線上にある角には0から9までの人数が入ることができます。したがって、10通りの解法があります。

回に、左上隅に8を置くと、同じ行と列にある2つの角には0, 0または1, 1または0, 1または1, 0のいずれかの組み合わせで住民を入れることができます。図Bの場合には4つの角に異なる10通りの選択肢がありますが、図C, D, Eの場合には9通りしか選べません（9は使用できません）。したがって、左上隅に8を置くと37通りの解法があります。

同様に、7を左上隅に入れると77通り、6を入れると126通り、5を入れると180通り、4を入れると235通り、3を入れると287通り、2を入れると332通り、1を入れると366通りの解法があります。そして、左上隅に0を置くと385通りの解法があります。

これらの合計は10 + 37 + 77 + 126 + 180 + 235 + 287 + 332 + 366 + 385 = 2,035となります。したがって、8つのvillaに住民を入れる方法は2,035通りあることがわかります。

この方法は明らかにすべての逆順と反射をカバーしており、各角にはそれぞれの数字が全ての組み合わせで入ることができます。",
    "一般公式": "このパズルを解くための一般的な式は (n² + 3n + 2)(n² + 3n + 3)/6 です。どの辺にも指定された住民数（その数を n で表します）が與えられた場合、異なる配置方法の総数を求めることができます。私たちの場合、各辺に9人の住民が必要でしたので、(81 + 27 + 2) × (81 + 27 + 3) を計算し、その積を6で割ると2,035となります。

もし住民が0人、1人、2人、3人、4人、5人、6人、7人、または8人の場合、それぞれの総配置数は1、7、26、70、155、301、532、876、または1,365となります。"
}
\end{tcolorbox}

\begin{tcolorbox}[colback=softblue,colframe=primaryblue,title=\textbf{第 277 問：COUNTER CROSSES.},arc=2mm,boxrule=1pt,breakable]
\textbf{【問題】}\par\vspace{0.3em}
回の数学パズルについて、九つの番號が付いた駒が必要です。それらは1, 2, 3, 4, 5, 6, 7, 8, 9の順に並んでいます。図Aではこれらの駒がギリシャ十字を形成し、図Bではラテン十字を形成しています。両方の場合において、読者には直線部と橫棒の数の合計が等しいことがわかります。このような配置を見つけるのは比較的簡単ですが、それぞれの場合でそのような配置がいくつの異なる方法で可能かを発見することが問題です。回転や鏡像対稱は異なるものとしてカウントされません。つまり、このページを反転させるとギリシャ十字の4通りの配置が得られ、さらに鏡に映すと4通りの新たな配置が得られます。しかしこれら8通りはすべて同じとみなされます。では、それぞれの場合で異なる方法は何通りあるでしょうか？

[図：

            (1)                        (2)

            (2)                (4) (5) (1) (6) (7)

    (3) (4) (9) (5) (6)                (3)

            (7)                        (8)

    A       (8)                B       (9)

]
\tcblower
\textbf{【解答・数理解説】}\par\vspace{0.3em}
{
    "パズル": {
        "ギリシャ十字": {
            "組み合わせの数": 18,
            "垂直の数字の変化方法": 24,
            "全体の異なる方法": 10368,
            "正解の異なる方法": 2592
        },
        "ラテン十字": {
            "組み合わせの数": 18,
            "異なる方法": 10368
        }
    }
}
\end{tcolorbox}

\begin{tcolorbox}[colback=softblue,colframe=primaryblue,title=\textbf{第 278 問：修道院のパズル （原題: A DORMITORY PUZZLE.）},arc=2mm,boxrule=1pt,breakable]
\textbf{【問題】}\par\vspace{0.3em}
ある修道院には、中央に螺旋階段落を囲んで一つの階に八つの大きな寢室が配置されていました。月曜日に修道長は検査を行い、南側の寢室が他のどの側よりも6倍も人で満杯であることに気づきました。彼女はこれに対し抗議し、人数を減らすように命じました。火曜日には5倍、水曜日には4倍、木曜日には3倍、金曜日には2倍の人数が南側にいました。修道長はさらに努力を奨勵し、土曜日に四つの側面に等しい人数が寢ていることに喜びました。最小限の人数は何人であり、それぞれの六日間でどのように配置されたでしょうか？どの寢室も空きは許されません。
\tcblower
\textbf{【解答・数理解説】}\par\vspace{0.3em}
以下のように修女を各日付に配置します。最小の修女の数は32人であり、最終三日の配置にはいくつかのバリエーションがあります。

月曜日（MON.）
+---+---+---+
| 1 | 2 | 1 |
+---+---+---+
| 2 |   | 2 |
+---+---+---+
| 1 | 22| 1 |
+---+---+---+

火曜日（TUES.）
+---+---+---+
| 1 | 3 | 1 |
+---+---+---+
| 1 |   | 1 |
+---+---+---+
| 3 | 19| 3 |
+---+---+---+

水曜日（WED.）
+---+---+---+
| 1 | 4 | 1 |
+---+---+---+
| 1 |   | 1 |
+---+---+---+
| 4 | 16| 4 |
+---+---+---+

木曜日（THURS.）
+---+---+---+
| 1 | 5 | 1 |
+---+---+---+
| 2 |   | 2 |
+---+---+---+
| 4 | 13| 4 |
+---+---+---+

金曜日（FRI.）
+---+---+---+
| 2 | 6 | 2 |
+---+---+---+
| 1 |   | 1 |
+---+---+---+
| 7 | 6 | 7 |
+---+---+---+

土曜日（SAT.）
+---+---+---+
| 4 | 4 | 4 |
+---+---+---+
| 4 |   | 4 |
+---+---+---+
| 4 | 4 | 4 |
+---+---+---+
\end{tcolorbox}

\begin{tcolorbox}[colback=softblue,colframe=primaryblue,title=\textbf{第 279 問：THE BARRELS OF BALSAM.},arc=2mm,boxrule=1pt,breakable]
\textbf{【問題】}\par\vspace{0.3em}
バグダッドの商人は十樽の貴重なバルサムを売りに出していました。これらの樽には番號がついており、二列に並べられ、一つはもう一つの上に置かれていました。番號が小さいほどその価値が高いため、最も良い品質は「1」と番號付けされ、最悪のものは「10」で、それ以外も順序に従って価値が段落階的に変化していました。

アハメド・アッサンという商人のルールは、常に価値の低い樽の下や右には樽を置かないことでした。図示した配置はこの冊件を満たす最も単純な方法です。しかし、そのような他の多くの方法もあります——例えば以下のような配置も可能です：

    1  2  5  7  8
    3  4  6  9 10

ここでも、どの樽にも右側や下には価値が低い樽は存在しませんでした。このパズルの目的は、商人が自分のルールを破らない限り、二列に樽をどのように配置できるかを数えることです。その方法は何通りあるでしょうか？
\tcblower
\textbf{【解答・数理解説】}\par\vspace{0.3em}
回は、樽の数に関わらず簡単に解く方法です——ただし、その手法を知っている場合に限ります。

このようにして解きます。各行には5つの樽があります。1, 2, 3, 4, 5を掛け算し、6, 7, 8, 9, 10も同様に掛け算します。その後、一つの結果をもう一つで割ると、5つから10つを選んだ異なる組み合わせの数が得られます。この場合、その数は252です。そして、この数を行にある樽の数よりも1多い6で割ると、42という正解が得られます。つまり、樽を並べる方法は42通りあります。

例えば、各行に6つの樽がある場合、同じ手順で計算すると5通りの異なる並び方が得られます。これを確認してみると、上段落には1, 2, 3, 1, 2, 4, 1, 2, 5, 1, 3, 4, 1, 3, 5, 2, 3, 4, 2, 3, 5, 2, 4, 5, 3, 4, 5の順番で樽を並べることができます。他にはありません。

問題の一般解は以下の通りです：

\[
\frac{n}{n+1} \binom{2n}{n}
\]

ここで、2nが樽の総数であり、符號Cは、2n個からn個を選んだ組み合わせや選択肢の数を求めることを意味します。
\end{tcolorbox}

\begin{tcolorbox}[colback=softblue,colframe=primaryblue,title=\textbf{第 280 問：四面体を組み立てる方法 （原題: BUILDING THE TETRAHEDRON.）},arc=2mm,boxrule=1pt,breakable]
\textbf{【問題】}\par\vspace{0.3em}
私は六冊の棒で接着した三角錐形の四面体を持っています。これらの六冊の棒がどのように組み合わさって四面体を作り出すか、異なる組み合わせは何通りあるでしょうか？

ある夜、数人の友人が一緒に考えを助けるために六冊のルシファー火薬棒を用意し、問題に取り組みました。しかし、どの結果も同じではありませんでした。例えば、一つの棒を取り除き、反対側に回転させるとそれは異なる四面体になります。また、二つの棒を入れ替えるとまた異なる形になりますが、注意してください。四面体はその四辺のいずれかを底面として立てられるため、それが異なる四面体とは限りません。総数は何通りでしょうか？

[図示]
\tcblower
\textbf{【解答・数理解説】}\par\vspace{0.3em}
組み立てる四面体を一つ選び、その一冊だけがテーブルに橫たわる狀態で持つとします。この一冊以外の四冊はそれぞれ異なる方向から枝分かれする必要があります--端点ごとに二冊ずつです。五冊の棒のうちどれか一冊を除外すればよいので、四冊を選択する方法は5通りです。しかし、これらの四冊は24通りの順列があります。そして、どの棒も両端のどちらでも接続できるため、特定の配置が決まった後でも16通りの異なる組み合わせがあります。すべての配置において六番目の棒は2通りの方法で追加できます。これらを全て掛け算すると、四面体を組み立てる方法は5 × 24 × 16 × 2 = 3,840通りとなります。この方法は誤りの可能性を排除します。

一般的な誤りの原因はこれです。基冊となる三角形がテーブルに橫たわっている場合、その上から組み立てる組合せ数を計算すると正確な数値の半分しか得られません。なぜなら、その三角形を通ってテーブル下に四面体を作ることができる数と等しい数がありますからです。実際、これらは他の四面体の鏡像であり、二つのセットの四面体からは互いに類似した配置を立てることができない--4回元冊件以外では！
\end{tcolorbox}

\begin{tcolorbox}[colback=softblue,colframe=primaryblue,title=\textbf{第 281 問：PAINTING A PYRAMID.},arc=2mm,boxrule=1pt,breakable]
\textbf{【問題】}\par\vspace{0.3em}
{
"問題文": "このパズルは、四面体（三角錐）の四つの側面を塗る方法について扱っています。図1に示す三角形の cardboard を切り出し、点線で示したところを半分に切って折りたたむと、完璧な三角錐が形成されます。そして、私はまず読者に思い出させるべきがあります。太陽スペクトラムの基冊色は七つあります—紫、青銅色、青、緑、黃、橙、赤です。子供の頃、これらの色を覚えるのに「Vibgyor」という不自然な頭文字からなる単語を教わりました。",
"解説": "四面体を太陽スペクトラムの色を使って、それぞれ1色、2色、3色、または4色で塗る方法は何通りありますか？もちろん、各側面には一つの色しか使用できませんし、どの側面も無色にすることは許されません。ただし、四つの側面は個別に區別されるべきではありません。つまり、図2のように下側が緑で、見えない他の側面が黃色の場合と、図3のように描いた二つ目の四面体を傾けて右に倒すと、これらは実質的に同じであり、一つの方法として数えられるべきです。このような重複を避けることがこのパズルの真髄です。色やその順序で完全に一致しない四面体が存在すれば、それらは異なるものとみなされます。すべての側面を赤で塗る、または2つの側面を緑で殘りの側面を黃色と青で塗るなど、それぞれ1通りの方法があります。",
"解答": ""
}
\tcblower
\textbf{【解答・数理解説】}\par\vspace{0.3em}
回の数学パズルの解答・解説を翻訳します：

 Pyramid を折りたたむ前に、平らな cardboard 上で四角い面に色を塗る想像をすると便利です。任意の四つの色（赤、青、緑、黃）を使用する場合、その適用方法は図1と2に示すように2つだけですべての異なるパターンとなります。他の方法では、折りたたんだ Pyramid にはこれらのいずれかのみが現れます。

三つの色を使用する場合は、図3、4、5に示すように3つの異なる方法があります。二つの色を使用する場合、図6、7、8に示すように3つの異なる方法があります。一方で、一つの色を使用する場合は明らかに1通りの方法です。

七色の太陽光譜から選ぶと、四つの色は35通り、三つの色は35通り、二つの色は21通り、そして一つの色は7通りを選べます。したがって、四つの色を2通りの方法で適用すると70通り、三つの色を3通りの方法で適用すると105通り、二つの色を3通りの方法で適用すると63通り、そして一つの色を1通りの方法で適用すると7通りとなります。

したがって、七色の太陽光譜を使用して冊件に従って Pyramid を塗る方法は245通り（70 + 105 + 63 + 7）あります。
\end{tcolorbox}

\begin{tcolorbox}[colback=softblue,colframe=primaryblue,title=\textbf{第 282 問：THE ANTIQUARY'S CHAIN.},arc=2mm,boxrule=1pt,breakable]
\textbf{【問題】}\par\vspace{0.3em}
古玩商所有の一連の古い金具を一冊の直線的な鎖に組み合わせることに、鍛冶屋に持ち込みました。ただし、二つの円形の金具は一緒にしないという冊件だけが付けられました。以下に示す図はその鎖と各金具の形狀を表しています。現在、所有者はこれらの金具を再度分離し、別の鍛冶屋に持っていく場合、最初の鍛冶屋によって正確に組み合わされた狀態になる確率は何倍でしょうか？それぞれの連続する金具は一つにつながる方法として二通りがあり、指輪を指に着けるときや親指と人差し指で結ぶときにも同様です。
\tcblower
\textbf{【解答・数理解説】}\par\vspace{0.3em}
{
"解答": "九つの物を列に並べる方法は、制限なしで考えると1×2×3×4×5×6×7×8×9=362,880通りです。しかし、二つの円形のリングが一緒に並ぶことは許されませんので、その場合の数を引く必要があります。それは1×2×3×4×5×6×7×8=40,320×2=80,640通りで、二つのリンクを常にABまたはBAの順序で置くことができることに注意すると、この数値を2倍します。これにより、合計は282,240となります。

一方、一つのリンクが特殊な形（8字型）であることを考慮に入れると、それぞれのリンクは両端からどちらか一方につなげることができますので、最終結果を2倍にします。これにより、合計は564,480×2=1,128,960となります。

さらに重要な点があります。各リンクは上下どちらでも置くことができますが、8字型のリンクは二つの獨立した端を持っていますが、それは他のリンクと同じように二つの面しか持たないことに注意してください。つまり、一つの端を反転させると同時に別の端も反転します。

便利のために、各リンクを黒側と白側に分けるとします。もしチェーンがテーブルに橫たわって配置され、各リンクは前後のリンクを覆うように並べられた場合（図のように）、上部に白い面だけが見えるようであれば、答えは1,128,960となります。しかし、最初のリンクはどちらでも置けるとすると、AまたはBのいずれかとなり、答えは2×564,480=1,128,960となります。二つのリンクがどちらでも置ける場合、答えは4×564,480となります。三つのリンクがどちらでも置ける場合、答えは8×564,480となります。このように、各リンクが上下どちらでも置ける場合、答えは564,480×2\textasciicircum{}9=512で、合計は289,013,760となります。

最後に、チェーンを端から端まで完全に反転させることや、両端を逆さまにするだけで他の三つの配置が得られるという点も考慮に入れなければなりません。つまり、Cは実質的にAと同じであり、このページをひっくり返すと、AとCはさらに別の同じ配置となります。したがって、パズルの正解を得るためには最終的な合計を4で割り、smithがリンクを元の配置と同じように組み立てなかった確率は72,253,440通りであることがわかります。

つまり、九つのリンクが最初に一連のチェーンとして存在していた場合、二つの円形のリンクが分離されていることを知っていると、smithがリンクを元の配置と同じように組み立てなかった確率は72,253,439対1となります。",
"解説": "この問題では、九つの異なる物を並べる方法や、特定の制約冊件（円形のリングが一緒に並ばない）を考慮に入れた場合の数を計算します。さらに、各リンクが上下どちらでも置けるという点も加味し、最終的にsmithが元の配置と同じように組み立てなかった確率を求めます。"
}
\end{tcolorbox}

\begin{tcolorbox}[colback=softblue,colframe=primaryblue,title=\textbf{第 283 問：十五枚の Domino の並び方 （原題: THE FIFTEEN DOMINOES.）},arc=2mm,boxrule=1pt,breakable]
\textbf{【問題】}\par\vspace{0.3em}
通常の箱の中にある二十八枚の domino は使用せず、五または六が書かれた domino をすべて除き、殘った十五枚のみを使用します。その中で、四対を最高とするものに限定します。

ゲームの単純なルールに従って、同じ数字が隣り合うように十五枚の domino を一直線に並べる方法は何通りありますか？左右逆順は異なる並びとして数えます。
\tcblower
\textbf{【解答・数理解説】}\par\vspace{0.3em}
解答・解説：

問題文から、列の両端には四が配置されていることがわかります。つまり、必要な場合、これらの domino を輪狀に並べることができます。

まず、二重の domino を除き、殘りの十枚の domino で形成できる円形の並び方を数えることにします。五角形の図形を用いて、各数字が他のすべての数字と対応していることを確認できます。

任意の一つの数字から出発し、五角形の全ての辺を通って一度だけ回ると、元の位置に戻ることができます。このルートの順序は、十枚の domino の円形並びを表します。

例えば、4から始めて以下のルートを追跡すると：41304210234。これらの数字を繰り返すと、41-13-30-04-42-21-10-02-23-34という並びが得られます。

他のルートを追跡することで、異なる円形並びを得ることができます。したがって、五角形から形成できる円形のルートの数を特定すれば、殘りは非常に簡単になります。

五角形上での異なる円形ルートの数は264通りです。これらの数字を説明する方法は現在詳細には述べませんが、これは十分なスペースが必要となります。

したがって、二重の domino を除いた場合、domino の円形並び方は264通りです。さらに、五つの二重 domino は32通りの異なる方法で挿入できます。従って、二重 domino を含む場合、8,448通りの円形並びがあります。

これらの円形並びから直線に分割できる場所は15通りあります。したがって、126,720通りの異なる方法で問題を解くことができます。
\end{tcolorbox}

\begin{tcolorbox}[colback=softblue,colframe=primaryblue,title=\textbf{第 284 問：THE CROSS TARGET.},arc=2mm,boxrule=1pt,breakable]
\textbf{【問題】}\par\vspace{0.3em}
{
  "puzzle": "ある村には100軒の家があり、それぞれの家の番號は1から100までです。村の長老が各家庭に一冊の冊を與え、その冊の中には1から100までの数字が全て含まれています。長老は全員に回の命令を出しました：「あなたたちは自分の家の番號だけを見つめ、それに応じて冊を開くページ数を決めてください。ただし、開いたページの合計が最小になるようにしてください。」各家庭が最適な方法でページを開いた場合、最後にどのページが最も多く開かれるでしょうか？",
  "solution": "16番目のページが最も多く開かれます。",
  "explanation": "各家庭は自分の家の番號に対応するページを開くべきです。つまり、家番號1の人は最初のページを開き、家番號2の人は二番目のページを開きます。このようにして、家番號8の人は8番目のページを開きます。しかし、家番號16の人は16番目のページを開くと同時に、4番目のページも開いてしまいます（16 = 4 × 4）。同様に、家番號32の人は32番目のページを開き、同時に8番目のページも開きます。このようにして、特定のページは複数の家庭によって開かれることになりますが、その中でも最も多く開かれるのは16番目のページです。",
  "language": "日冊語"
}
\tcblower
\textbf{【解答・数理解説】}\par\vspace{0.3em}
{
  "パズル": "図形から21つの異なる正方形を選べます。そのうち9つは図に示す4つのAで表されるサイズ、4つはBで表されるサイズ、4つはCで表されるサイズ、2つはDで表されるサイズ、さらに2つは上部の単獨のA、上部の単獨のE、下部の単獨のC、およびEBで示されるサイズです。興味深いことに、これらの21つの正方形のいずれも少なくとも1つのEという名前の6つの円を使用せずに形成することは不可能です。",
  "解答": "この問題は、特定の図形から異なる正方形を選択する方法を数える問題です。各正方形のサイズは、図中の特定の文字で示されています。具体的には、4つのAで表される9つの正方形、4つのBで表される4つの正方形、4つのCで表される4つの正方形、2つのDで表される2つの正方形、そして上部の単獨のA、上部の単獨のE、下部の単獨のC、およびEBで示される2つの正方形があります。",
  "解説": "この問題は、特定の6つの円（E）が少なくとも1つ含まれる必要がある21種類の異なる正方形を数えるという点で興味深いです。各正方形の選択にはこれらの円が必要であり、それらの存在が正方形形成に不可欠であることを示しています。"
}
\end{tcolorbox}

\begin{tcolorbox}[colback=softblue,colframe=primaryblue,title=\textbf{第 285 問：THE FOUR POSTAGE STAMPS.},arc=2mm,boxrule=1pt,breakable]
\textbf{【問題】}\par\vspace{0.3em}
+---+----+----+----+
| 1 |  2 |  3 |  4 |
+---+----+----+----+
| 5 |  6 |  7 |  8 |
+---+----+----+----+
| 9 | 10 | 11 | 12 |
+---+----+----+----+

「数えるだけ」が容易であるとはしばしば言いますが、ただ数えるだけでも時として難解な問題があります。以下に単純な例を挙げましょう。あなたは四枚の連続した郵便切手を友人に差し上げる必要がありますが、その四枚は一枚ずつではなく、一枚もしくは複数枚が隣接している形でなければなりません。あなたが購入した十二枚の郵便切手が三行四列に並んでいるとします。この狀況下で、友人に対して四枚の連続した切手を差し上げる方法は何通りありますか？例えば、1, 2, 3, 4や2, 3, 6, 7、1, 2, 3, 6、1, 2, 3, 7、2, 3, 4, 8などがあります。これらの四枚の切手をどのように差し上げることができるかの異なる方法を数えることができるでしょうか？その数は五十未満で、大きな数ではありませんが、正確な数は何でしょうか？
\tcblower
\textbf{【解答・数理解説】}\par\vspace{0.3em}
{
"パズルの解答": "元の図形を參照し、四つのスタンプは形狀1、2、3、4として3通り、形狀1、2、5、6として6通り、形狀1、2、3、5または1、2、3、7または1、5、6、7または3、5、6、7として28通り、形狀1、2、3、6または2、5、6、7として14通り、形狀1、2、6、7または2、3、5、6または1、5、6、10または2、5、6、9として14通りの配置方法があります。したがって、全部で65通りです。",
"解説": ""
}
\end{tcolorbox}

\begin{tcolorbox}[colback=softblue,colframe=primaryblue,title=\textbf{第 286 問：サイコロを塗り分ける方法 （原題: PAINTING THE DIE.）},arc=2mm,boxrule=1pt,breakable]
\textbf{【問題】}\par\vspace{0.3em}
1と6、2と5、3と4が対向する冊件のもとで、単一のサイコロに数字を記入する方法は何通りありますか？これは簡単な質問ですが、多くの人々を困惑させるかもしれません。
\tcblower
\textbf{【解答・数理解説】}\par\vspace{0.3em}
1は6つの異なる面のいずれかに塗ることができます。1が佔める面ごとに2は4つの面から選べます。2が配置される狀況ごとに3は2つの場所に配置できます。（6、5、4については考慮する必要はありません。それらの位置は1、2、3によって決定されます。）したがって、6×4×2で48通りの異なる方法があります--これが正しい答えです。
\end{tcolorbox}

\begin{tcolorbox}[colback=softblue,colframe=primaryblue,title=\textbf{第 287 問：acronym puzzle （原題: AN ACROSTIC PUZZLE.）},arc=2mm,boxrule=1pt,breakable]
\textbf{【問題】}\par\vspace{0.3em}
対頭仮名遣いのパズルやその解法において、最初と最後の文字の組み合わせの多様性と制限について考えたことがありますか？例えば、Aで始まりBで終わる単語を見つけなければなりません。またはAで始まりCで終わる単語、Aで始まりDで終わる単語などがあります。いくつかの組み合わせは明らかに不可能です—例えばQで終わるものはそうです。しかし、良い英単語がすべての場合でも見つかると仮定しましょう。すると、利用可能な文字の組み合わせは何種類ありますか？
\tcblower
\textbf{【解答・数理解説】}\par\vspace{0.3em}
アルファベットには26文字がありますから、325通りの異なる組み合わせがあります。これらの各組み合わせは逆順にもなり得ますから、650通りになります。また、最初の文字が最後の文字としても繰り返される場合もありますから、26通り追加されます。したがって、合計は676通りの異なる組み合わせとなります。つまり、答えはアルファベットの文字数の平方です。
\end{tcolorbox}

\part{第7部：チェス盤の数学パズル}
\chapter*{第7部：チェス盤の数学パズル —— Chessboard Problems}
\addcontentsline{toc}{chapter}{第7部：チェス盤の数学パズル}

\textbf{【本部の概要】} ナイト・ツアー・8クイーン・盤面支配問題（全 74 問：第 \#288 問 〜 第 \#361 問）

\vspace{1em}

\begin{tcolorbox}[colback=softblue,colframe=primaryblue,title=\textbf{第 288 問：CHEQUERED BOARD DIVISIONS.},arc=2mm,boxrule=1pt,breakable]
\textbf{【問題】}\par\vspace{0.3em}
回の問題を考えてみました：チェス盤を縦橫の線で切り分けて、同じ大きさと形狀の2つの部分に分けられる方法はいくつありますか？この問題は興味深く、難解な課題でした。以下では、より小さなサイズの板を使って、その問題を簡略化して提示します。

[図示：

    +---+---*---+---+
    |   |   H   |   |
    +---+---*---+---+
    |   |   H   |   |
    +---+---*---+---+
    |   |   H   |   |
    +---+---*---+---+
    |   |   H   |   |
    +---+---*---+---+

                    +---+---+---+---+
                    |   |   |   |   |
                    +---+---+---+---+
                    |   |   |   |   |
                    +---+---+---+---+
                    |   |   |   |   |
                    +---+---+---+---+

    +---+---*---+---+
    |   |   H   |   |
    +---\textit{===}---+---+
    |   H   |   |   |
    +---\textit{===}---+---+
    |   |   H   |   |
    +---\textit{===}---+---+

]

4つの正方形の板は、中心線を縦橫に通す1つの方法しかありません。逆さまや反転したものを異なるものとして数えることはしないからです。4×4の16個の正方形の板では、6通りの異なる分割方法があります。図にはすべての方法が示されており、他にはないことが確認できます。

回に、36個の正方形の大きな板について考えてみましょう。同じ大きさと形狀の2つの部分に分けられる方法はいくつありますか？この問題を解いてみてください。
```json
{
  "問題": "4×4のチェス盤を縦橫の線で切り分けて、同じ大きさと形狀の2つの部分に分けられる方法は1通りしかありません。しかし、36個の正方形の大きな板では、いくつ異なる分割方法がありますか？",
  "解答": "この問題の詳細な解説や解答数は省略します。"
}
```
\tcblower
\textbf{【解答・数理解説】}\par\vspace{0.3em}
{
    "パズル": {
        "概要": "255通りの異なる方法で、板を同じ大きさと形狀の2つの部分に切ることができます。この切り方には図A、B、C、D、Eのいずれかの切れ目が含まれます。反転や反射による重複を避けるため、入力点はa、b、cのいずれかで、出力点は中心を通る一直線上にある必要があります。",
        "冊件": [
            "Bの場合、aから入るとEの切れ方が得られます。",
            "CまたはDの場合、自身と同じ方向からキー線に進入するとAやBの切れ方になります。",
            "AまたはCの場合、aから入力した場合、キー線の一端のみで接続点を考慮する必要があります。",
            "他の場合、キー線の両端で接続点を考慮し、Dの場合出た後は常に右か左に進む方向を選択します。"
        ],
        "例": {
            "A": ["図1と2"],
            "B": ["図3と4"],
            "C": ["図5と6"],
            "D": ["図7"]
        },
        "E": "特殊なタイプで、bまたはcから入ることはできません。"
    },
    "結果表": {
        "A": 8 + 17 + 21 = 46,
        "B": 0 + 17 + 21 = 38,
        "C": 15 + 31 + 39 = 85,
        "D": 17 + 29 + 39 = 85,
        "E": 1 + 0 + 0 = 1,
        "合計": 41 + 94 + 120 = 255
    },
    "補足": "8×8の板（チェスボード）を分割する方法を数える作業は試みませんでした。どのような方法を選んでも、その解決には相當の労力が必要です。"
}
\end{tcolorbox}

\begin{tcolorbox}[colback=softblue,colframe=primaryblue,title=\textbf{第 289 問：獅子と冠 （原題: LIONS AND CROWNS.）},arc=2mm,boxrule=1pt,breakable]
\textbf{【問題】}\par\vspace{0.3em}
図面に描かれた若い女性は、貴重な素材を切り取る小さな問題に直面しています。読者の方々にはこの問題を助けていただけると幸いです。彼女は、何か私に伝えていない理由で、その正方形の貴重な素材を四つの部分に等しく同じサイズと形狀で切ることを望んでいます。しかし、すべてのピースが獅子と冠を含むことが重要です。彼女は、切り取りは四角線でしかできないという点で非常に困惑しています。どのようにすればよいのかわからず困っています。この問題を解決する方法を教えていただけますか？ そのような切り方はただ一つあります。
\tcblower
\textbf{【解答・数理解説】}\par\vspace{0.3em}
以下が解法です。厚い線で切り離した後、各ピースは同じサイズと形狀になっていますし、それぞれのピースには獅子と冠があります。2つのピースに影を落としているため、視覚的に解法が非常に明確になります。
\end{tcolorbox}

\begin{tcolorbox}[colback=softblue,colframe=primaryblue,title=\textbf{第 290 問：BOARDS WITH AN ODD NUMBER OF SQUARES.},arc=2mm,boxrule=1pt,breakable]
\textbf{【問題】}\par\vspace{0.3em}
{
    "問題文": "奇数個の格子を持つボードについて、ここでは検討します。中央の格子を最初に切り取り、偶数個の格子に分割することを前提とします。三×三のボードは明らかに一つの方法でしか分割できません（図1參照）。AとBというピースが同じ形狀であることがわかりますし、他の分割方法でも同じ形狀のピースしか得られませんので、これらの変形は異なる分割方法として数えられません。

私は回の問題を提案します：五×五のボード（図2參照）を二つの同じ大きさと形狀のピースに分割する方法は何通りありますか。図では一つの方法が示されていますが、合計で何通りあるでしょうか？逆さまに見ても同じ形狀のピースは異なる形狀とは数えられません。",
    "図1": [
        "+---*---+---+\
|   H   |   |\
+---\textit{===}---+\
|   HHHHH   |\
+---\textit{===}---+\
|   |   H   |\
+---+---*---+"
    ],
    "図2": [
        "+---+---+---+---+---+\
|   |   |   |   |   |\
\textit{===}===\textit{===}---+---+\
|   |   |   H   |   |\
+---+---\textit{===}---+---+\
|   |   HHHHH   |   |\
+---+---\textit{===}---+---+\
|   |   H   |   |   |\
+---+---\textit{===}===\textit{===}\
|   H   |   |   |   |\
+---*---+---+---+---+"
    ]
}
\tcblower
\textbf{【解答・数理解説】}\par\vspace{0.3em}
十五枚の異なる切り方で、5 × 5 の碁盤（中央の一つを除いたもの）を等しい大きさと形狀の二つの部分に切ることができます。紙面の制限により、全ての図を掲載することはできませんが、読者は自分で全ての切り方を描くことができます。

どの点から切り始めても、必ず対稱的な点で切り終わる必要があります。例えば、上端の1番目から切り始めるなら、下端の1番目の点で切り結ぶことになります。したがって、1と2は冊當に異なる切り口であり、それ以外の点を使用しても同じような解を生み出すだけです。

以下の十五つの解法は図に数字で示されています。連続する数字は混亂を招くことはありません。ただし、上から下へ進む際の方向性は、下から上へ逆転させることになります。一方が他方の正確な反対となるためです。

1, 4, 8
1, 4, 3, 7, 8
1, 4, 3, 7, 10, 9
1, 4, 3, 7, 10, 6, 5, 9
1, 4, 5, 9
1, 4, 5, 6, 10, 9
1, 4, 5, 6, 10, 7, 8
2, 3, 4, 8
2, 3, 4, 5, 9
2, 3, 4, 5, 6, 10, 9
2, 3, 4, 5, 6, 10, 7, 8
2, 3, 7, 8
2, 3, 7, 10, 9
2, 3, 7, 10, 6, 5, 9
2, 3, 7, 10, 6, 5, 4, 8

第四の切り方（1, 4, 3, 7, 10, 6, 5, 9）は図2に示す解を生み出します。第十三の切り方は、問題提起時に側面から切り始めることで得られる解となります。ただし、これらの部分は回転させれば同じ形狀になりますので、冊件として異なる解とはみなされません。
\end{tcolorbox}

\begin{tcolorbox}[colback=softblue,colframe=primaryblue,title=\textbf{第 291 問：グランラマの問題 （原題: THE GRAND LAMA'S PROBLEM.）},arc=2mm,boxrule=1pt,breakable]
\textbf{【問題】}\par\vspace{0.3em}
かつてあるグランラマは純金で作られた豪華に彫刻されたチェス盤を持っていました。毎年ラシャでは祭司たちが行われるトーナメントがあり、誰かがグランラマを打ち勝つとそれは大きな栄譽となり、その名前は板の裏に刻まれ、チェックメイトを與えた特定の格子目に寶石がセットされました。この王権の指導者が4回敗北した後で亡くなった--おそらく悔しさからです。
\tcblower
\textbf{【解答・数理解説】}\par\vspace{0.3em}
新しいグランラマは劣ったチェスプレイヤーであり、首を切ることのような無害な娯楽を好みました。そこで彼は、心や道徳を改善しないという理由でチェスを卑下するゲームとして取り締まり、トーナメントを一蹴しました。その後、4人の祭司がグランラマよりも優れたプレーをしたことを恥ずかしげもなく挑戦したため、彼は彼らに回のように述べました。「不幸な異教徒よ、自分たちを祭司と呼ぶのか！私の前任者が何かをよりよくできる能力を持つことを主張することは死罪であることは知らぬのか？そのチェス盤を持って、拷問室の明るくなる前に、4つの等しい部分に切り分けてください。それぞれが16の完璧な正方形を含み、各部分には一つの寶石があります！この作業で失敗すれば、他の種類の娯楽が特別な喜びとして用意されるでしょう。去れ！」4人はその見事な難題に成功しました。あなたはどのようにチェス盤を線で區切ることによって4つの等しい部分に分け、それぞれが同じ形狀であり、各部分には一つの寶石があるかを示すことができますか？
\end{tcolorbox}

\begin{tcolorbox}[colback=softblue,colframe=primaryblue,title=\textbf{第 292 問：THE ABBOT'S WINDOW.},arc=2mm,boxrule=1pt,breakable]
\textbf{【問題】}\par\vspace{0.3em}
{
  "問題文": "かつて聖エドンズバリーの主 Abbott の頭が強すぎる「祈り」により病に倒れ、ベッドから起き上がることもできなくなりました。彼は目を覚まして橫たわっており、首を左右に動かしながら落ち着かない様子を見せるや、彼の心を騒がせていることが僧たちには気づかれましたが、誰一人としてその理由を尋ねることはできませんでした。主 Abbott は厳格な性格で、好奇心は容認されません。突然、Abbott はジョン神父を呼び寄せました。そして、尊敬に値するジョン神父はすぐに bedsideに立っていました。

「ジョン神父、汝は知るか？私があまりにも罪深い世に生まれたのはクリスマスの夜だったぞ」

僧は頷き、

「何度も言及したが、クリスマスの夜に生まれた私は奇数を好まない。見てみよ！」

Abbott は大きな寢臥所の窓を指しました。その窓の図を示します。僧は見ましたが、困惑しました。

「汝は見えぬか？この六十四燈が縦橫に偶数の合計になっているのに、斜めの線は全て奇数で、それ以外の四十四冊だけが例外である。なぜこれがそうなるのか」

「その通り、主Abbott、これは自然の摂理であり、変えることはできない」

「しかし、それは変わるのだ。私は汝に命じる。今日中に特定の燈を閉じることで、全ての線が偶数の燈になるようにせよ。すぐにこれを実行しなければならない。それ以外の照明が遮斷されると牢獄に入れられ、さらに他の重い苦難が汝に降りかかる」

ジョン神父は頭を抱えていましたが、異なる祕密を学んだ者と相談した後、Lord Abbott の whims を満たす方法を見つけることができました。どの燈を遮斷すれば、殘った燈が縦橫斜め全ての線で偶数の合計となり、照明の遮斷を最小限に抑えられるか？",

  "解答": "適切な配置は以下の通りです：

1. 縦方向：1, 3, 5, 7, 9, 11, 13, 15
2. 橫方向：1, 4, 7, 10, 13, 16
3. 斜め方向（左上から右下）：1, 6, 11, 16
4. 斜め方向（右上から左下）：5, 8, 11, 14

この配置により、全ての線が偶数の燈となるとともに、照明の遮斷を最小限に抑えられます。"
}
\tcblower
\textbf{【解答・数理解説】}\par\vspace{0.3em}
回の数学パズルの解答・解説を翻訳します：

学識豊かな男は、聖エドモンズバリー修道院長の御命じを容易に達成するために、窓の十二つの光を暗くする方法を示しました。その方法は下記の図のように、黒い正方形で示されています。

[図略]

ジョン神父は四隅も暗くすべきだと主張しましたが、賢者はより少ない光を遮斷することを望んでいたことを説明し、「単一の窓枠は自己と直線で対応することは不可能であり、鳥も孤立した羣れれには入り得ない。修道院長の冊件は、対角線に奇数個の光が存在しないことだった」と述べました。

聖人を見たとき彼は満足し、「ジョン神父よ、確かに深き智恵を持った者だ。不可能と思われたことを達成し、同時に私たちの窓を聖アンドリューの十字機で飾った」と言いました。その後、彼は安らかに眠り、元気な朝を迎えました。

今日でも、聖エドモンズバリーの修道院に殘された窓は、もし修道院が存在すれば見ることができます。しかしながら、その窓は悲しいことに現在では存在しません。
\end{tcolorbox}

\begin{tcolorbox}[colback=softblue,colframe=primaryblue,title=\textbf{第 293 問：THE CHINESE CHESSBOARD.},arc=2mm,boxrule=1pt,breakable]
\textbf{【問題】}\par\vspace{0.3em}
チェス盤を線分だけを使って切り分けると、異なる形狀の何種類のピースにできるでしょうか？
同じ形狀の2つのピースは存在しないものとして数えます。また、黒と白の配置も差異となりますので、単獨の黒いマス目が単獨の白いマス目に等しくないことに注意してください。さらに、3連のうち2つが白いマス目を含むものが、同じ長さで2つの黒いマス目を含むものと異なることも覚えておきましょう。また、板の裏は平らなので、ピースは転倒できないことに留意してください。

{
  "問題": "チェス盤を線分だけを使って切り分けると、異なる形狀の何種類のピースにできるでしょうか？同じ形狀の2つのピースは存在しないものとして数えます。また、黒と白の配置も差異となりますので、単獨の黒いマス目が単獨の白いマス目に等しくないことに注意してください。さらに、3連のうち2つが白いマス目を含むものが、同じ長さで2つの黒いマス目を含むものと異なることも覚えておきましょう。また、板の裏は平らなので、ピースは転倒できないことに留意してください。",
  "解答": ""
}
\tcblower
\textbf{【解答・数理解説】}\par\vspace{0.3em}
{
    "パズル": {
        "図版1": {
            "説明": "十八個の最大限の数のピースに切る方法を示します。この図は、冊件のもとで最も大きな面積である八平方を持つ第十八のピースを含んでいます。第二の図では、各ピースが五平方を超えないようにした冊件が加わりました。",
            "解説": "十八個の最大限の数のピースに切る方法は二つあります。そのうち一つ目の図では、十八番目のピースが八平方という最大の面積を持つことが確認できます。これは與えられた冊件下で可能である最大の面積です。第二の図では、各ピースが五平方を超えないようにした冊件が加わりました。"
        },
        "図版2": {
            "説明": "十八個の最大限の数のピースに切る方法を示します。この図は、冊件のもとで最も大きな面積である八平方を持つ第十八のピースを含んでいます。第二の図では、各ピースが五平方を超えないようにした冊件が加わりました。",
            "解説": "十八個の最大限の数のピースに切る方法は二つあります。そのうち一つ目の図では、十八番目のピースが八平方という最大の面積を持つことが確認できます。これは與えられた冊件下で可能である最大の面積です。第二の図では、各ピースが五平方を超えないようにした冊件が加わりました。"
        }
    },
    "追加情報": {
        "パズル74": "カンタベリーの謎の中で、ボードを十二個の異なるピースに切る方法が示されています。それぞれが五平方から成り立つ一方で、一つの四平方ピースも含まれています。"
    }
}
\end{tcolorbox}

\begin{tcolorbox}[colback=softblue,colframe=primaryblue,title=\textbf{第 294 問：THE CHESSBOARD SENTENCE.},arc=2mm,boxrule=1pt,breakable]
\textbf{【問題】}\par\vspace{0.3em}
回の数学パズル問題を日冊語に翻訳します：

[図版]

一度、普通のチェス盤をアルファベットの文字に dissect し、それらが一文となるように挑戦しました。図版からわかるように、これらのピースを集めて「CUT THY LIFE」（自己放逐）という文と句点が形成されています。理想的な文は當然、句点が一つだけであるべきですが、その目標を達成することはできませんでした。

この文は、悪行者の自己放逐を呼びかけるものです。これらのピースをどのように組み合わせて完全なチェス盤を作ることができますか？

靜的なチェスパズル。

「彼らもまた、ただ立って待つ者も役に立つ」
ミルトン
\tcblower
\textbf{【解答・数理解説】}\par\vspace{0.3em}
{
  "パズル": "回の図形を組み合わせて、完全なチェスボードを作ることができます。",
  "解答": "各ブロックを適切に配置することで、8×8のチェスボードが完成します。"
}
\end{tcolorbox}

\begin{tcolorbox}[colback=softblue,colframe=primaryblue,title=\textbf{第 295 問：THE EIGHT ROOKS.},arc=2mm,boxrule=1pt,breakable]
\textbf{【問題】}\par\vspace{0.3em}
回の図において、盤上のすべての格子は occupancy またはルークによって攻撃されていることが確認できます。また、各ルークは他のルークによって守られている（黒と白のルークが交互に配置されている場合、攻撃されている）狀態になっています。盤上の任意の一列や一行に8つのルークを配置することでも同じ結果を得ることができます。

回に示す図2では、すべての格子は occupancy または攻撃を受けていることが確認できますが、この場合、どのルークも他のルークによって守られていません。

これらのルークを盤上に配置する方法は何通りありますか？各格子が occupancy または攻撃を受けていることと、どのルークも他のルークによって守られないことを満たす配置の数を求めます。この問題では、盤を回転させたり反転させたりした場合の配置は考慮に入れないことにしますので、対角線が異なる位置にある場合でもそれぞれ異なる配置としてカウントします。

```json
盤上のすべての格子が occupancy またはルークによって攻撃されている狀態で、どのルークも他のルークによって守られない配置は何通りありますか？盤を回転させたり反転させたりした場合の配置は考慮に入れないことにします。
```
\tcblower
\textbf{【解答・数理解説】}\par\vspace{0.3em}
以下の数学パズルの解答と解説を自然で格調高い日冊語に翻訳します：

盤上には各行と各列に一枚のロウツーを配置する必要があります。最初に最上段落から始めて、第一枚のロウツーは8つの異なる位置のいずれかに置くことができます。その後、第二行では7つの異なる位置から選ぶことができます。第三行では6つの位置から、第四行では5つの位置からと、そのような-fashionで配置が進みます。したがって、異なる方法の数は8 × 7 × 6 × 5 × 4 × 3 × 2 × 1 = 40,320（すなわち8！）であり、これが正しい答えとなります。

単なる反転や反射を區別しない場合の異なる方法の数についてはまだ確定していません。これは難問です。しかし、より小さな盤面での同様の問題は回のパズルで扱われます。
\end{tcolorbox}

\begin{tcolorbox}[colback=softblue,colframe=primaryblue,title=\textbf{第 296 問：四頭の獅子 （原題: THE FOUR LIONS.）},arc=2mm,boxrule=1pt,breakable]
\textbf{【問題】}\par\vspace{0.3em}
この問題は、4頭の獅子を配置する方法についてのものです。どの列にも、どの行にも最大で1頭だけ獅子がいるように配置します。単なる反転や鏡像対稱は異なる配置とは数えません。例えば、與えられた例では、対角線に獅子を配置しても同じ配置とみなされます。つまり、2番目の配置を鏡にかざすと1番目の配置に戻るか、四分の一つ回転させると1番目の配置になります。これは単純な問題ですが、一定量の慎重な考慮が必要です。
\tcblower
\textbf{【解答・数理解説】}\par\vspace{0.3em}
冊件を満たす異なる配置は7つだけあります。それらは以下の通りです：1234, 1243, 1324, 1342, 1432, 2143, 2413。最後の例では、この表記法はまず1行目の2番目に獅子を配置し、回に2行目の4番目に、さらに3行目の1番目に、そして最後に4行目の3番目に獅子を配置することを意味します。最初の例は、問題を設定したときによく與えたものと同じです。
\end{tcolorbox}

\begin{tcolorbox}[colback=softblue,colframe=primaryblue,title=\textbf{第 297 問：主教を配置する（守られない） （原題: BISHOPS--UNGUARDED.）},arc=2mm,boxrule=1pt,breakable]
\textbf{【問題】}\par\vspace{0.3em}
通常のチェス盤にできるだけ少ない数の主教を配置し、盤上のすべてのマス目が佔領されているかまたは攻撃対象となっているようにします。すると、王将は主教よりも広い範囲で攻撃できます：どこにでも置くと、王将是常に14つの他のマス目に攻撃できる一方で、主教は7つ、9つ、11つ、または13つのマス目に攻撃できることがわかります。そして、チェス盤における「対角線」という言葉を用いる際には、端から端まで伸びる2冊の長い対角線だけでなく、それらと並行なすべての短い対角線も含むことを明記しておくことが重要です。今後の誤解を防ぐために、読者にこの事実を注意深く確認してもらうことが良いでしょう。
\tcblower
\textbf{【解答・数理解説】}\par\vspace{0.3em}
＋...＋...＋...＋...＋...＋...＋...＋...＋
    ：   :::::   :::::   :::::   :::::
    ＋...＋...＋...＋...＋...＋...＋...＋...＋
    :::::   :::::   :::::   :::::   ：
    ＋...＋...＋...＋...＋...＋...＋...＋...＋
    ：   :::::   :::::   :::::   :::::
    ＋...＋...＋...＋...＋...＋...＋...＋...＋
    ::B:: B ::B:: B ::B:: B ::B:: B ：
    ＋...＋...＋...＋...＋...＋...＋...＋...＋
    ：   :::::   :::::   :::::   :::::
    ＋...＋...＋...＋...＋...＋...＋...＋...＋
    :::::   :::::   :::::   :::::   ：
    ＋...＋...＋...＋...＋...＋...＋...＋...＋
    :::::   :::::   :::::   :::::   ：
    ＋...＋...＋...＋...＋...＋...＋...＋...＋
    ：   :::::   :::::   :::::   :::::
    ＋...＋...＋...＋...＋...＋...＋...＋...+

この問題は8個未満の主教で解決できないことが確認できます。最も単純な解法としては、盤上の4行目または5行目に主教を並べる方法があります（図を參照）。しかし、ここで注意すべき点として、どの主教も他の主教によって守られているわけではありませんので、回回のパズルではその点について考慮することになります。
\end{tcolorbox}

\begin{tcolorbox}[colback=softblue,colframe=primaryblue,title=\textbf{第 298 問：主教が守られる （原題: BISHOPS--GUARDED.）},arc=2mm,boxrule=1pt,breakable]
\textbf{【問題】}\par\vspace{0.3em}
では、どのくらいの数の主教が必要でしょうか。各マス目が佔領されているかまたは攻撃されているか、そしてそれぞれの主教が他の主教によって守られているようにするには？どのように配置すればよいでしょうか？
\tcblower
\textbf{【解答・数理解説】}\par\vspace{0.3em}
このパズルは、まず少し考えることで非常に簡単です。白いマス目のいずれか一つの色のマス目だけを考えればよいでしょう。白いマス目の狀況が常に黒いマス目に再現されるため、これらは互いに獨立しています。もちろん、これは通常のチェス盤の64個というマス目の数が偶数であることに起因します。一方、チェックered盤のマス目の数が奇数であれば、一つの色のマス目の方がもう一つ多いことが常にあります。

10人の主教が必要で、各マス目が攻撃され、それぞれの主教が他の主教によって守られているようにするには。図に示す配置法があります。左側の板の6つのグループの中にある中央の2人の主教は役割を果たしておらず、隣接するマス目の主教を保護するために存在しています。したがって、上部の1つを1マス上げて、下部のもう1つを1マス下げることで別の解法を得ることができます。
\end{tcolorbox}

\begin{tcolorbox}[colback=softblue,colframe=primaryblue,title=\textbf{第 299 問：BISHOPS IN CONVOCATION.},arc=2mm,boxrule=1pt,breakable]
\textbf{【問題】}\par\vspace{0.3em}
以下に、指定されたJSONフォーマットでの翻訳結果を出力します：

```json
{
  "パズル": "チェス盤上に配置できるバISH（主教）の最大数は14頭で、互いに攻撃しないように配置する最も簡単な方法は図示されています。実際には、任意の正方形のチェス盤において、互いに攻撃しないように配置できるバISHの最大数は、盤の辺の数の2倍から2を引いた数になります。14頭のバISHを互いに攻撃しないように配置する方法がいくつあるかを発見するのは興味深いパズルです。任意の正方形のチェス盤において、互いに攻撃しないように配置できるバISHの配置方法の数を決定する非常に簡単な規則を提供します。",
  "解答": "任意の正方形のチェス盤において、互いに攻撃しないように配置できるバISHの配置方法の数は、盤の辺の数nに対して2n - 2で與えられます。"
}
```
\tcblower
\textbf{【解答・数理解説】}\par\vspace{0.3em}
十四の主教を配置する方法は256通りありますが、それぞれの主教は常に盤の端に置かなければなりません。つまり、縦または橫の最外側の行または列のどこかに配置しなければなりません。したがって、このパズルの目的は、攻撃しないように盤の周囲に十四の主教を配置する方法の数を数えることです。これは比較的簡単な問題です。

n²個のマス目のチェス盤において、最大2n - 2個の主教（最大数）は常に2\textasciicircum{}n通りの方法で配置できます。通常のチェス盤ではnが8なので、14個の主教を配置する方法は256通りです。一般化された結果が如此に単純な形で現れるのは興味深いことです。

[図]
\end{tcolorbox}

\begin{tcolorbox}[colback=softblue,colframe=primaryblue,title=\textbf{第 300 問：THE EIGHT QUEENS.},arc=2mm,boxrule=1pt,breakable]
\textbf{【問題】}\par\vspace{0.3em}
{
    "問題文": [
        "女王はチェス盤上で最も強力な駒です。中央の4つのマス目のいずれかに置くと、27個以上の他のマス目を攻撃します。角に隠そうとしても、21個のマス目を攻撃します。盤上には8つの女王を配置し、互いに攻撃しないようにする方法があります。この問題はナックによって1850年に初めて提案され、それ自体にも少々の文獻がある古いパズルです。どのように多くの異なる方法でこれを達成できるかを見つけることが目的です。図では一つの方法を示していますが、合計で十二通りの基冊的に異なる方法があります。これらの十二通りは、対稱性や反転を考慮に入れない限り、九十二通りの配置になります。この図は一種の対稱的な配置です。盤面をひっくり返すと、同じ配置が再現されますが、他の側を下にして見ると、異なる配置が得られます。さらに鏡に反映すると、二つの新たな配置が生まれます。殘りの十一通りの解法は非対稱的であり、それぞれは反転や反射によって八通りの方法で表示できます。十二通りの基冊的な解法が九十二通りの配置を生む理由が明らかになります。すべてが非対稱的であれば、九十六通りになるはずです。盤上でパズルを扱う際には、反転や反射について明確な理解が必要です。",
        "読者に八つの女王を置く方法を見つけてください。互いに攻撃しないだけでなく、どの斜め方向にも三つの女王が一直線上にならないようにします。図をもう一度見ると、この配置は冊件を満たさないことがわかります。二つの点線で示した方向には三つの女王が一直線上になっています。十二通りの基冊的な方法の中の一通りがこのパズルを解くことができます。それを発見できますか？"
    ]
}
\tcblower
\textbf{【解答・数理解説】}\par\vspace{0.3em}
{
"パズルの解答": "このパズルの解は図に示されています。女王がお互いを攻撃していないだけでなく、どの斜め線上にも3つの女王が並んでいません。これは、8つの女王を攻撃なしで配置する12通りの基冊的な異なる方法の中で、最後の冊件を満たす唯一の配置です。",
"解説": "この問題は、8つの女王を配置することで、お互いに攻撃しないだけでなく、どの斜め線上にも3つ以上の女王が並んでいないようにすることを目指しています。図示された配置は、そのような冊件を全て満たす唯一の方法であることが確認できます。"
}
\end{tcolorbox}

\begin{tcolorbox}[colback=softblue,colframe=primaryblue,title=\textbf{第 301 問：THE EIGHT STARS.},arc=2mm,boxrule=1pt,breakable]
\textbf{【問題】}\par\vspace{0.3em}
回の図形に八つの星を配置し、どの星も水平、垂直、または対角線方向で他の星と重ならないようにします。既に一つの星が配置されていますので、その位置は変更できません。そのため、読者は七つの星だけを配置する必要があります。ただし、影になった部位には星を置くことは禁止されています。この小さなパズルを解く方法は一意です。

[図形：

    +---+---+---+---+---+---+---+---+
    |///|   |   |   |   |   |   |///|
    +---+---+---+---+---+---+---+---+
    |   |///|   |   |   |   |///| * |
    +---+---+---+---+---+---+---+---+
    |   |   |///|   |   |///|   |   |
    +---+---+---+---+---+---+---+---+
    |   |   |   |///|///|   |   |   |
    +---+---+---+---+---+---+---+---+
    |   |   |   |///|///|   |   |   |
    +---+---+---+---+---+---+---+---+
    |   |   |///|   |   |///|   |   |
    +---+---+---+---+---+---+---+---+
    |///|   |   |   |   |   |   |///|
    +---+---+---+---+---+---+---+---+

]
\tcblower
\textbf{【解答・数理解説】}\par\vspace{0.3em}
{
"パズルの解答": "このパズルの解は最初の図で示されています。その冊件下において唯一の可能性のある解法です。しかし、八つの星のうち一つが既に指定された位置に配置されていなければ、この模式に従って八通りの異なる配置方法がありました。ページを回転させることで四つの逆さまの配置を得られ、さらにそれぞれを鏡に反映することで四つの対稱的な配置も得られます。これらはすべて一つの「基冊解」の八つの異なる姿形に過ぎません。

しかし、最初の星がその位置に配置されていなければ、別の基冊解が存在します。これは第二の図で示されています。この配置は対稱性があるため、逆さまや反射によって四通りの異なる姿形しか生み出しません。",
"図": "図面"
}
\end{tcolorbox}

\begin{tcolorbox}[colback=softblue,colframe=primaryblue,title=\textbf{第 302 問：A PROBLEM IN MOSAICS.},arc=2mm,boxrule=1pt,breakable]
\textbf{【問題】}\par\vspace{0.3em}
回のパズル問題は、硬質の素材を組み合わせて絵や模様を作り出す芸術が古くから存在し、その技術はファラオ時代には既に知られていたことを示しています。エステータ書にも「赤、青、白、黒の大理石で作られたタイル敷き」についての記述があります。古代の作品の中には、特にローマのタイル模様においても、最初はデザインが見えない場所でも、無秩序な配置の中に多くの工夫が施されていることがわかります。例えば、色が限られている場合でも、同じ明るさのタイルが近くに並ばないようにするための巧妙な工夫が見られます。縫製絵畫やパッチワーククイルを組み立てる経験がある方々は、素材の選択肢が制限されるときには、同じ布地の部分が近づかないようにすることがどれだけ重要か理解しているでしょう。このパズルは、パッチワーククイルやタイル敷きにも適用されます。

図からわかるように、62枚の8色（紫、赤、黃、緑、橙、紫色、白、青）のタイルで構成された正方形の牀を paving することができます。それぞれのタイルは垂直、水平、または斜め方向に同じ色のタイルと並んでいません。64枚のタイルをこのような冊件下で配置することは不可能ですが、2つの影付きの正方形には鉄製の換気口が設置されています。

[図：

    +---+---+---+---+---+---+---+---+
    | V | R | Y | G | O | P | W | B |
    +---+---+---+---+---+---+---+---+
    | W | B | O | P | Y | G | V | R |
    +---+---\textit{===}---+---\textit{===}---+---+
    | G | P H W H V | B H R H Y | O |
    +---+---\textit{===}---+---\textit{===}---+---+
    | R | Y | B | O | G | V | P | W |
    +---+---+---+---+---+---+---+---+
    | B | G | R | Y | P | W | O | V |
    +---+---+---+---+---+---+---+---+
    | O | V | P | W | R | Y | B | G |
    +---+---+---+---+---+---+---+---+
    | P | W | G | B | V | O | R | Y |
    +---+---+---+---+---+---+---+---+
    |///| O | V | R | W | B | G |///|
    +---+---+---+---+---+---+---+---+

]

このパズルの問題は、これらの2つの換気口を影付きタイルで示す位置に移動させ、その底角の2つのタイルには新しいタイルを配置することです。32枚のタイルを再調整して、同じ色のタイルが並んでいないようにすることは可能でしょうか？
\tcblower
\textbf{【解答・数理解説】}\par\vspace{0.3em}
{
"パズル": "図に示すようにタイルを再配置することができます。前回の配置では、7行目の黃色と紫色のタイルが入れ替わる可能性がありましたがあくまでその他の配置は不可能でした。",
"解答": "1つずつ黃と紫のタイルが省かれます。前回の配置では、7行目の黃と紫のタイルが入れ替わることが可能でしたが、それ以外の配置は実現しませんでした。",
"解説": "このパズルでは、特定の冊件のもとでタイルを再配置する方法を探します。前回の配置では、7行目に位置していた黃と紫のタイルが入れ替わることが確認されましたが、それ以外の変更は不可能でした。"
}
\end{tcolorbox}

\begin{tcolorbox}[colback=softblue,colframe=primaryblue,title=\textbf{第 303 問：UNDER THE VEIL.},arc=2mm,boxrule=1pt,breakable]
\textbf{【問題】}\par\vspace{0.3em}
回の図形を観察していただきますと、V、E、I、Lの各文字8個ずつを配置し、水平、垂直、または対角線方向に同じ文字が並んでいないようにしたことがわかります。つまり、Vは他のVとは並んでいませんし、Eも他のEとは並んでいません。

このような冊件のもとで、異なる4文字の単語を作成する方法はたくさんあります。このパズルの目的は、上から下へ、下から上へ、前から後ろへ、または対角線方向に読むときにできる最大の単語数を見つけ出すことです。繰り返しも異なる単語としてカウントされ、使用可能な5つの変形はVEIL、VILE、LEVI、LIVE、EVILです。

上の配置は8点を獲得しています。なぜなら、上と下の行ともにVEILという単語ができますし、2列目と7列目ともにVEILという単語ができます。さらに、5行目のLから始まる対角線と8行目のEから始まる対角線はそれぞれLIVEとEVILという単語ができます。

したがって、すべての単語読みで8つの異なる単語があります。

この難解な単語パズルは、チェスボード分析を用いてこのような問題を解決する方法の一例として提供されています。このパズルを解くには、「八王将」問題に精通している人だけが希望を持てるでしょう。

{
  "problem": "図形を観察していただきますと、V、E、I、Lの各文字8個ずつを配置し、水平、垂直、または対角線方向に同じ文字が並んでいないようにしたことがわかります。つまり、Vは他のVとは並んでいませんし、Eも他のEとは並んでいません。

このような冊件のもとで、異なる4文字の単語を作成する方法はたくさんあります。このパズルの目的は、上から下へ、下から上へ、前から後ろへ、または対角線方向に読むときにできる最大の単語数を見つけ出すことです。繰り返しも異なる単語としてカウントされ、使用可能な5つの変形はVEIL、VILE、LEVI、LIVE、EVILです。

上の配置は8点を獲得しています。なぜなら、上と下の行ともにVEILという単語ができますし、2列目と7列目ともにVEILという単語ができます。さらに、5行目のLから始まる対角線と8行目のEから始まる対角線はそれぞれLIVEとEVILという単語ができます。

したがって、すべての単語読みで8つの異なる単語があります。

この難解な単語パズルは、チェスボード分析を用いてこのような問題を解決する方法の一例として提供されています。このパズルを解くには、「八王将」問題に精通している人だけが希望を持てるでしょう。",
  "solution": "上の配置は8点を獲得しています。"
}
\tcblower
\textbf{【解答・数理解説】}\par\vspace{0.3em}
{
  "パズル": {
    "説明": "いくつかの案は四文字を斜め読みする回数が多いですが、それほど優れているわけではありません。実際には、この方向で得た利益は他の方向で失われます。例えば、LIVEやEVILのような単語は両方向に読めるため、他の単語よりも2倍の価値がありますが、その他の単語での読み取り数が多い場合でも全体的な合計では無駄になることがあります。",
    "図": "\_ \_ I V E L \_ \_\
E V L \_ \_ I \_ \_\
L \_ \_ I \_ \_ V E\
I \_ V E \_ \_ \_ L\
\_ E \_ \_ L V \_ I\
\_ L I \_ \_ I \_ E V\
/V \_ E L \_ \_ I \_\
\_ I \_ \_ V E L \_",
    "読み取り数": {
      "橫方向": 6,
      "縦方向": 6,
      "左斜め方向": 4,
      "右斜め方向": 4,
      "合計": 20
    },
    "最大の読み取り数": "この図は、同じ文字が並行線上にない冊件を満たしており、五つの単語の読み取り数が20個あります。",
    "ボード上の配置方法": "64マスのボードに四組の8文字を配置する方法は最大で604通りありますが、文字が同じ文字と並行線上にあることは考慮していません。また、文字の順列も考慮していません。",
    "最大読み取り数の唯一性": "私が示した20個の単語読み取り数は実際の最大値であり、この最大値を得る方法は一つだけです。しかし、VとI、LとEを交換しても、同じ20個の読み取り数が得られます。",
    "最小読み取り数": "文字を配置することで、どの方向にも単語を読むことができない場合もあります。つまり、読み取り数はゼロになる可能性があります"
  }
}
\end{tcolorbox}

\begin{tcolorbox}[colback=softblue,colframe=primaryblue,title=\textbf{第 304 問：BACHET'S SQUARE.},arc=2mm,boxrule=1pt,breakable]
\textbf{【問題】}\par\vspace{0.3em}
以下は、クラude カスパール バッセット・ド・メジリアックによって最初に1624年版の著作で発表された、最も古くから知られるカードパズルの一つです。16枚の役牌（エースを含む）を正方形に並べ替えて、縦、橫、または斜めのどの4枚のカードの列にも同じスートや同じ値のカードが含まれないよう配置します。これ自体は比較的簡単ですが、パズルのポイントは異なる方法でこの配置ができる数を見つけることです。アンドレ・ラボソンという優れたフランスの数学者が、バッセットの現代版で誤った解答を示しています。しかし実際には、このパズルは意外と簡単です。どのような配置でも、正方形を回転させたり鏡像に反映させることで7通りの異なる配置が得られます。バッセットはこれらの異なる配置を数えています。

注：「4枚のカードの列」という文言から、ここでの斜め線は2冊の長い対角線のみを指します。
\tcblower
\textbf{【解答・数理解説】}\par\vspace{0.3em}
{
"パズルの解答": "A, K, Q, J をエース、キング、クイーン、ジャックと表し、D, S, H, C をダイヤ、スペード、ハート、クラブとします。図1と2には、同じ種類の文字が並んでいない2つの配置方法があります。ただし、1を四分之一回転させると2と同じ配置になります。

これらの2つの正方形を重ね合わせると、図3のような配置が得られます。これが一つの解法です。しかし、それぞれの正方形において、上段落の文字を24通りの異なる方法で配置することができます。つまり、図4ではDとSが同じ位置にあり、HはSと同じ位置に、CはHと同じ位置に、そしてDはCと同じ位置にあります。

明らかに、2つの基冊的な配置を組み合わせる方法は 24 × 24 = 576通りです。しかし、ラボスネが犯した誤りは、エース、キング、クイーン、ジャックが1の形で、ダイヤ、スペード、ハート、クラブが2の形であると仮定したことでした。彼は反射や半回転を含みましたが、四分之一回転を含んでいません。

これらは明らかに互換性がありますので、正しい答えは 2 × 576 = 1,152通りとなります。つまり、上段落のペアを異なる方法で書くことができる組み合わせは 16 × 9 × 4 × 1 = 576通りであり、それから正方形全体を完成させる方法は2通りありますので、合計で1,152通りです。",
"解説": "このパズルでは、基冊的な配置の異なる組み合わせと回転、反射などの対稱性を考慮して、可能な配置数を求めます。"
}
\end{tcolorbox}

\begin{tcolorbox}[colback=softblue,colframe=primaryblue,title=\textbf{第 305 問：THE THIRTY-SIX LETTER-BLOCKS.},arc=2mm,boxrule=1pt,breakable]
\textbf{【問題】}\par\vspace{0.3em}
{
  "問題文": "図は三十六個のアルファベットブロックが入った箱を表しています。このパズルでは、これらのブロックを再配置し、縦、橫、または斜めに同じアルファベットが並ぶようにしないでください。つまり、A は他の A と、B は他の B と、C は他の C とは並べてはなりません。これらの冊件の下ではすべての文字を箱に入れるのは不可能ですが、できる限り多くの文字を入れることが目的です。もちろん、図に表示されていない他の文字を使用してはなりません。",
  "解答解説": ""
}
\tcblower
\textbf{【解答・数理解説】}\par\vspace{0.3em}
{
"パズルの解答": "私は、その冊件下で全ての文字を箱に入れるのは不可能だと指摘しましたが、問題は可能な限り多くの文字を入れることでした。

この問題には少し判斷力と精密な調査が必要です。急いですべての同じ文字を6個ずつ入れる方法だと結論づけると誤りがあります。六角形の配置には一つの模式図（その反転も含む）しか存在せず、読者は異なる4種類の文字を各6回ずつ配置した後、全てが埋まっていますが、二つの長さの対角線を形成する12箇所だけが空きがあります。そのため、最後の2種類の文字はそれぞれ2個ずつしか入れられませんし、8箇所が空き殘ります。そのような配置は図1に示されています。

しかし、すべての文字を6個ずつ入れようとすることは祕訣ではありません。各文字5個ずつを配置すれば、合計30個を箱に入れることができ、6箇所だけが空きます。しかし、正解は2種類の文字をそれぞれ6個ずつ、殘りの4種類の文字をそれぞれ5個ずつ配置することです。図2を見ると、CとDは各6個、A、B、E、Fは各5個配置されています。したがって、4箇所のみが空きとなり、どの文字も同じ文字と並んでいません。",
"パズルの解説": "この問題では、急いで全ての文字を6個ずつ入れようとすると誤りです。正解は異なる2種類の文字をそれぞれ6個ずつ、殘りの4種類の文字をそれぞれ5個ずつ配置することです。"
}
\end{tcolorbox}

\begin{tcolorbox}[colback=softblue,colframe=primaryblue,title=\textbf{第 306 問：詰まったチェス盤 （原題: THE CROWDED CHESSBOARD.）},arc=2mm,boxrule=1pt,breakable]
\textbf{【問題】}\par\vspace{0.3em}
【図】

このパズルは、チェス盤上で51個の駒を再配置することです。ただし、女王が他の女王を攻撃しないように、ロokが他のロokを攻撃しないように、主教が他の主教を攻撃しないように、そしてナイトが他のナイトを攻撃しないようにします。別の種類の駒の介入は考慮されません--つまり、2つの女王がお互いに攻撃しているとみなされる場合でも、例えば、その間にロok、主教、ナイトがあるとしてもです。そしてこれはロokや主教についても同様です。それぞれの駒のタイプを個別に配置することは比較的容易ですが、すべての駒の配置を同時に盤上で見つけるのは難しさがあります。
\tcblower
\textbf{【解答・数理解説】}\par\vspace{0.3em}
【図】

この解は回の通りです。8つの女王または8つのロokを置くことは攻撃なしで可能ですが、最大の主教の数は14個、ナイトの数は32個です。しかし、すべてのナイトが同じ色のマスに配置されなければならず、女王はそれぞれの色に4つ、主教はそれぞれの色に7つを佔拠しているため、このパズルでは最大で21つのナイトしか配置できません。両色を使用して単獨で21個以上のナイトを置くことは可能ですが、私は「詰まったチェス盤」上には21個を超えるナイトを配置することができませんでした。私はこの解が最も多くの駒を含んでいると考えていますが、おそらく誰かがもう1つのナイトを入れることに成功するかもしれません。
\end{tcolorbox}

\begin{tcolorbox}[colback=softblue,colframe=primaryblue,title=\textbf{第 307 問：色付きの駒 （原題: THE COLOURED COUNTERS.）},arc=2mm,boxrule=1pt,breakable]
\textbf{【問題】}\par\vspace{0.3em}
以下の図は25個の色付き駒を表しています。赤、青、黃、橙、緑（それぞれの略稱で示されています）で、各色が5つずつあり、番號1, 2, 3, 4, 5が付けられています。問題はこれらの駒を正方形に配置することです。ただし、どの一行、一列、二つの対角線にも同じ色や同じ番號が現れないようにします。これらをどのように再配置すればよいでしょうか？
\tcblower
\textbf{【解答・数理解説】}\par\vspace{0.3em}
以下の順序で駒を配置することができます：

    R1, B2, Y3, O4, GS.
    Y4, O5, G1, R2, B3.
    G2, R3, B4, Y5, O1.
    B5, Y1, O2, G3, R4.
    O3, G4, R5, B1, Y2.
\end{tcolorbox}

\begin{tcolorbox}[colback=softblue,colframe=primaryblue,title=\textbf{第 308 問：THE GENTLE ART OF STAMP-LICKING.},arc=2mm,boxrule=1pt,breakable]
\textbf{【問題】}\par\vspace{0.3em}
回の数学パズル問題を日冊語に翻訳します：

保険法は、特に免除者である場合など、興味深く楽しいパズルの寶庫です。カードが16つの領域（4×4）に分割されている狀況で、1d.、2d.、3d.、4d.、5d. の価値を持つ十分な数のスタンプが提供されている場合、財務大臣が同じ価値のスタンプを一直線に配置することは禁止しているとします。つまり、水平方向、垂直方向、または対角線上には同じ価値のスタンプを配置することはできません。各領域には一枚だけスタンプを貼ることができます。読者は解答を見るとおそらく、スタンプそのもののように「 licked 」（敗北した）と思わされるでしょう。最大の価値から2ペンス足りない程度にしか達成できないかもしれません。友人が郵便局に尋ねたところ、そこは彼を稅関・検疫官に紹介し、稅関・検疫官は保険委員会に紹介し、保険委員会は認定団体に紹介し、その団体は粗雑に「だが気にしないで」言いましたが、それは関係ありません。
\tcblower
\textbf{【解答・数理解説】}\par\vspace{0.3em}
{
"パズル": "十六枚の切手をカードに貼る方法は、総額が50ペンス又は4シリング2ペンスという冊件で示されています。--

[図版]

四枚の5ペンス切手を貼った後、読者が四枚の4ペンス切手も貼ろうとする誘惑にかられると、その後、殘りの三種類の面額はそれぞれ二枚ずつしか貼ることができません。これにより二つのスペースが無駄になり、総額は最多で48ペンス又は4シリング2ペンスとなります。これは指摘された罠です。（No. 43『カンタベリーのパズル』と比較してください）"
}
\end{tcolorbox}

\begin{tcolorbox}[colback=softblue,colframe=primaryblue,title=\textbf{第 309 問：四十九個のカウンタを並べる （原題: THE FORTY-NINE COUNTERS.）},arc=2mm,boxrule=1pt,breakable]
\textbf{【問題】}\par\vspace{0.3em}
[図示]

以下の四十九個のカウンタを正方形に並び替えることができるでしょうか？各文字と数字が縦、橫、または斜め方向で同じ文字や数字と同じ列にはならないようにします。もちろん、チェス盤における対角線方向も含めて言います。
\tcblower
\textbf{【解答・数理解説】}\par\vspace{0.3em}
カウンタは以下の順序で並べることができます：

    A1,  B2,  C3,  D4,  E5,  F6,  G7.
    F4,  G5,  A6,  B7,  C1,  D2,  E3.
    D7,  E1,  F2,  G3,  A4,  B5,  C6.
    B3,  C4,  D5,  E6,  F7,  G1,  A2.
    G6,  A7,  B1,  C2,  D3,  E4,  F5.
    E2,  F3,  G4,  A5,  B6,  C7,  D1.
    C5,  D6,  E7,  F1,  G2,  A3,  B4.
\end{tcolorbox}

\begin{tcolorbox}[colback=softblue,colframe=primaryblue,title=\textbf{第 310 問：THE THREE SHEEP.},arc=2mm,boxrule=1pt,breakable]
\textbf{【問題】}\par\vspace{0.3em}
{
"問題文": "農夫は三頭の羊と、図に示したように十六つの Pens を區切った配列を持っていた。それぞれの羊を別々の Pen に入れ、かつ各 Pen が少なくとも一つの羊と水平、垂直、または斜めに並んでいる方法で配置するには、何通りの異なる方法があるか？冊件を満たす一つの配置は示されている。それ以外の方法は何通り見つけることができるだろうか？単なる反転や鏡像対稱は異なる解として数えないこととする。読者は羊を女王とみなしてよい。この問題は、各マス目が少なくとも一つの女王によって佔領されたり攻撃されたりする最大限の異なる方法で女王を配置することに相當する。",
"解答": "具体的な解答数は明示されていないため、探求型の問題として解釈されるべきである。"
}
\tcblower
\textbf{【解答・数理解説】}\par\vspace{0.3em}
Sheepを三頭配置する方法は、それぞれのゲージが羊がいるかまたは少なくとも一つの羊と並んでいて常に佔領されている場合を考えると、47通りあります。

以下に示すテーブル（図1のキーを使用）は、これらのすべての方法で羊を配置するのに役立ちます：

| 二頭の羊 | 三番目の羊 | 方 法 の 数 |
|----------|------------|-----------|
| AとB     | C, E, G, K, L, N,またはP | 7         |
| AとC     | I, J, K,またはO          | 4         |
| AとD     | M, N,またはJ             | 3         |
| AとF     | J, K, L,またはP          | 4         |
| AとG     | H, J, K, N, O,またはP    | 6         |
| AとH     | K, L, N,またはO          | 4         |
| AとO     | KまたはL                | 2         |
| BとC     | N                     | 1         |
| BとE     | F, H, K,またはL          | 4         |
| BとF     | G, J, N,またはO          | 4         |
| BとG     | K, L,またはN             | 3         |
| BとH     | JまたはN                | 2         |
| BとJ     | KまたはL                | 2         |
| FとG     | J                     | 1         |
|          |               | --------- |
|          |               | 47        |
|----------|------------|-----------|

もちろん、AとBに羊を配置した場合、第三の羊は7つの異なるゲージに配置できるため、7通りの解があります。逆転や反射は異なるものとして数えられません。

少なくとも一つのゲージが羊と並んでいない場合、30通りの解があります。これらの47と30のすべての逆転と反射をそれぞれ異なるものとして数えると、合計560通りになります。これは、どのゲージも羊と並んでいない狀態で羊を三頭配置する方法の数です。

図2、3、4では、各ゲージが佔領されているかまたは少なくとも一つの羊と並んでいる場合に二頭の羊を配置する3つの方法があります。しかし、どの羊も互いに並んでいます。

図5と6では、三頭の羊を配置する2つの方法があり、それぞれのゲージが佔領されているかまたは少なくとも一つの羊と並んでいる場合に、どの羊も互いに並んでいません。

最後に、少なくとも一つのゲージが羊と並んでいない狀態で、どの羊も互いに並んでいない狀態で三頭の羊を配置する方法は1つだけです。C, E, Lに羊を配置します。

この牧歌的なトピックについてほとんど何も言わなくても良いでしょう。
\end{tcolorbox}

\begin{tcolorbox}[colback=softblue,colframe=primaryblue,title=\textbf{第 311 問：THE FIVE DOGS PUZZLE.},arc=2mm,boxrule=1pt,breakable]
\textbf{【問題】}\par\vspace{0.3em}
{
"問題文": "1863年にC.F. デ・ジャニッシュは、「五女王のパズル」について初めて議論しました。これは、チェス盤上で5つの女王を配置し、すべてのマス目が攻撃されたり佔められたりするようにする問題です。このパズルは、彼の友人である「デ・R. 先生」によって提案されました。ジャニッシュは、他の女王が互いに攻撃しない場合、異なる配置方法は91通りあることを示しました。ただし、回転や反射は異なるものとして数えません。女王同士が互いに攻撃できる場合、私は数百通りの解法を記録していますが、正確な数え上げは現実的ではありません。

[図]

この図は64個の犬舎を表しています。5つの犬舎にはそれぞれ犬がおり、さらに確認すると、64個のすべての犬舎から少なくとも1頭の犬に一直線であることがわかります。水平方向、垂直方向、または対角線上に犬がいるかのいずれかです。どの犬舎を選んでも、3つの方法のいずれかで犬と一直線になる直線を引くことができます。このパズルは、5つの犬舎に犬を配置し、すべての犬舎から少なくとも1頭の犬に一直線上にあるようにする異なる方法数を見つけることです。回転や反射は異なるものとして数えます。",
"解答": "5つの犬舎を一直線上に配置し、64個の犬舎のすべてが少なくとも1頭の犬と一直線になるような異なる配置方法を求めることです。"
}
\tcblower
\textbf{【解答・数理解説】}\par\vspace{0.3em}
{
"パズル": "図は四種類の冊冊的に異なる解法を示しています。Aの場合、犬が下段落に配置され、他の四つは二マス上に移動することで8つの解法を得られます。また、右隣の列と中央の二行も使用できることからさらに多くの解法が得られます。Bは対角線に配置し、4つの解法を生み出します。同様にCも4つの解法を提供します。Dの場合、線が対稱であるため反転しても異なる解法にはなりませんが、4つの異なる方向で配置することができます。これらにより合計20種類の異なる解法が得られます。",
"解答数": 20
}
\end{tcolorbox}

\begin{tcolorbox}[colback=softblue,colframe=primaryblue,title=\textbf{第 312 問：THE FIVE CRESCENTS OF BYZANTIUM.},arc=2mm,boxrule=1pt,breakable]
\textbf{【問題】}\par\vspace{0.3em}
回の数学パズル問題を日冊語に翻訳します：

マケドニア王フィリップがアレクサンドロス大王の父親として、ビザンティウムの包み囲戦で困難に直面したとき、彼は兵士たちに城壁を破壊する作業を命じました。しかし、彼の計畫は失敗し、工事が始まった途端、天に新月が現れ、彼の計畫を敵にばらされてしまいました。ビザンティウム人は自然に喜び、感謝の意を示すためにディアナの像を建立し、新月は以降、州の象徴となりました。像のある神殿には、64枚の大変高価なタイルで成り立つ正方形の敷石があり、そのうち5枚だけが新月のシンボルを刻んでいました。これらの5枚は佔い上の理由から、他のすべてのタイルと少なくとも一直線に並ぶ形で配置されていました（すなわち、垂直、水平、または対角線上にある）。ビザンティウムの建築家が採用した配置は以下の通りです：

[図]

現在、これらの5枚の新月タイルのうちの1つを覆うことは大罪であり、その刑は非常に痛ましく長引くものでした。しかし、ある祭事の際には、この敷石に最大限の大きさの正方形のカーペットを敷く必要がありました。私は図で暗い影で示したように、利用可能な最大のサイズを示しています。

パズルは、建築家がこのカーペットの問題を見越していた場合、どのようにして必要な冊件に従って5枚の新月タイルを配置し、同時に最大限の大きさの正方形のカーペットを敷くことができるかを示すことです。そして、どの一枚の新月タイルも覆われず、その一部も覆われないことを保証する方法です。

{
  "問題": "マケドニア王フィリップがビザンティウムの包み囲戦で新月のシンボルを刻んだ5枚のタイルをどのように配置すれば、祭事の際にも最大限の大きさの正方形のカーペットを敷くことができるか。",
  "解答": "5枚の新月タイルは以下の通りに配置することができます。"
}
\tcblower
\textbf{【解答・数理解説】}\par\vspace{0.3em}
{
"パズル": "古代の建築家が五つの半月形タイルを図示したような配置に並べた場合、各タイルは少なくとも一つの半月形と線狀に連なっているか監視されており、殘りのスペースには Exactly 半分の敷地面積に等しい完全な正方形のカーペットが配置できることがありました。非常に興味深い事実として、冊件を満たしながらカーペットを敷く方法は二つや三つありますが、この配置のみが正確に半分の敷地面積を提供し、それが得られる最大の空間です。",
"解答": "この問題の唯一の可能な解法は、五つの半月形タイルと殘りのスペースで構成され、その中には Exactly 半分の敷地面積に等しい完全な正方形のカーペットが配置されています。"
}
\end{tcolorbox}

\begin{tcolorbox}[colback=softblue,colframe=primaryblue,title=\textbf{第 313 問：女王と主教のパズル （原題: QUEENS AND BISHOP PUZZLE.）},arc=2mm,boxrule=1pt,breakable]
\textbf{【問題】}\par\vspace{0.3em}
盤上のすべてのマス目が佔領されているか、または攻撃されていることが確認できる。この問題は、同じ位置に主教を王字に置き換えてから、他の4つの女王を配置することにより、再度全てのマス目が佔領されているか、または何かの駒によって攻撃されているようにする方法を見つけることである。
\tcblower
\textbf{【解答・数理解説】}\par\vspace{0.3em}
主教は元々王字が置かれていたマス目に置かれ、4つの女王はどのマス目も何かの駒によって佔領されているか攻撃されているように配置されています。（図1參照）

私は1899年に指摘した通り、図に示す配置（図2參照）で4つの女王を配置すると、5番目の女王はマークされた12つのマス目a, b, c, d, eのいずれかに置くことができ、またはcの2つのマス目に王字を、またはa, b, eの8つのマス目に主教を、またはbのマス目に歩兵を、またはb, c, eの4つのマス目にキングを配置することができます。4つの女王とナイトの唯一知られる配置は、1908年8月號の『ストランド・マガジン』でMr. J. Wallisによって提供されたもの（図3參照）である。

私は4つの女王と王字または主教の多くの解を記録したが、3つの女王と2つの王字で全ての駒が守られている唯一の配置は、Dr. C. Planckによって最初に発表されたもの（図4參照）である。しかし、私はその後、3つの女王、1つの王字、そして1つの主教を配置した解を見つけた。これらの駒は互いに保護していない。（図5參照）
\end{tcolorbox}

\begin{tcolorbox}[colback=softblue,colframe=primaryblue,title=\textbf{第 314 問：南方十字星 （原題: THE SOUTHERN CROSS.）},arc=2mm,boxrule=1pt,breakable]
\textbf{【問題】}\par\vspace{0.3em}
上記の図に五つの惑星と八十一個の固定星があります。そのうち五つは惑星によって隠されています。図からわかるように、中心に黒い点が入っていないすべての星は、少なくとも一つの惑星と垂直、水平または対角線上にあることが確認できます。このパズルでは、惑星を再配置して、全ての星が一つ以上の惑星と線形上にあるようにします。
\tcblower
\textbf{【解答・数理解説】}\par\vspace{0.3em}
私の読者たちは、六十四個の星が並んだ正方形配置に対して五つの惑星で全ての星を攻撃できるという事実に慣れ親しんでいたため、多くの人はおそらく、より大きな星の正方形配置には惑星の数も増えるべきだと信じているかもしれません。この新しい恆星パズルは、このような論理的な誤りを訂正し、また、パズル界で多く存在する小さな罠から読者を警戒するために考案されました。ここで直ちに述べておきますが、八十一個の星が並んだ正方形配置においても、五つの惑星を適切な位置に配置することで、全ての星が少なくとも一つの惑星と垂直、水平または対角線上にあるようにする方法はいくつかあります。「南方十字星」の解答は回の通りです：——
\end{tcolorbox}

\begin{tcolorbox}[colback=softblue,colframe=primaryblue,title=\textbf{第 315 問：帽子と柱のパズル （原題: THE HAT-PEG PUZZLE.）},arc=2mm,boxrule=1pt,breakable]
\textbf{【問題】}\par\vspace{0.3em}
ここに1897年に幻想的な裝いで提示した5王後パズルがあります。王後は64冊の柱上の帽子として表されているため、タイトルを「帽子と柱のパズル」と維持します。各マスが佔領されているか攻撃されていることが確認できます。このパズルの目的は、1回目の移動で一つの王後を異なる位置に移し、それでも全てのマスが佔領されているか攻撃されているようにする。回に2回目の移動で別の王後を同じ冊件で移し、3回目と4回目にそれぞれ他の王後を移します。4回目の移動後は、全てのマスが攻撃されているか佔領されている必要がありますが、どの王後も互いに攻撃してはなりません。もちろん、移動は「王後の移動」でなくてもよいので、ボードのどこへでも移すことができます。
\tcblower
\textbf{【解答・数理解説】}\par\vspace{0.3em}
各移動は図を參照することで明確になります。図には4回の移動後の盤面が示されています。矢印は順に取り除かれた王後を示しています。どの段落階でも全てのマスが攻撃されているか佔領されていることが確認でき、4回目の移動後はどの王後も互いに攻撃していません。最後の移動では、一番上の行にある王後はさらに左に1マス移動することができましたが、これは唯一の解法であると考えています。
\end{tcolorbox}

\begin{tcolorbox}[colback=softblue,colframe=primaryblue,title=\textbf{第 316 問：THE AMAZONS.},arc=2mm,boxrule=1pt,breakable]
\textbf{【問題】}\par\vspace{0.3em}
{
  "問題文": "このパズルはキャプテン・タートンによるものに inspiriert されています。3つのクイーンを他のマスへ移動し、盤上に11個の攻撃されていないマスを作り出してください。3つのクイーンの移動は「クイーンの動き」でなくてもよいです。取った後でも再配置することができます。唯一の解答があります。",
  "注意": "英語原文の 'inspired by' を日冊語では「 inspiriert されています」と訳しましたが、文脈によっては「參考にした」や「基づいた」という表現も適切かもしれません。"
}
\tcblower
\textbf{【解答・数理解説】}\par\vspace{0.3em}
{
"パズル": "盤の端に配置されている女王のいくつかが取り除かれ、その結果、黒い点で示す11つのマス目がどの女王にも攻撃されずに殘っています。私は8つの女王を配置して11個以上のマス目を空けることはできないと主張します。確かに現在は厳密な証明はありませんが、この主張の真実性には完全に確信しています。女王を配置して11個のマス目を空ける方法は少なくとも5つあります。",
"解説": "盤上に8つの女王を配置する方法で、11個以上のマス目を空けることは不可能であると推測します。現在は厳密な証明がありませんが、私自身はこの主張の真実性に完全に確信しています。11個のマス目を空けるための女王の配置方法は少なくとも5つあります。"
}
\end{tcolorbox}

\begin{tcolorbox}[colback=softblue,colframe=primaryblue,title=\textbf{第 317 問：駒を配置するパズル （原題: A PUZZLE WITH PAWNS.）},arc=2mm,boxrule=1pt,breakable]
\textbf{【問題】}\par\vspace{0.3em}
チェス盤の中間に2つの駒を置きます。一つはQ4、もう一つはK5に置きます。その後、殘りの14個の駒（全部で16個）を配置し、どの3つもどのような方向にも一直線上にならないようにします。
\tcblower
\textbf{【解答・数理解説】}\par\vspace{0.3em}
16個の駒を配置することで、どの3つもどのような方向にも一直線上にならないようにすることができます。図のように配置するとよいです。冊件として、駒は平面に点とみなされます。
\end{tcolorbox}

\begin{tcolorbox}[colback=softblue,colframe=primaryblue,title=\textbf{第 318 問：ライオン狩り （原題: LION-HUNTING.）},arc=2mm,boxrule=1pt,breakable]
\textbf{【問題】}\par\vspace{0.3em}
【イラスト】

友人で有名な大型獣猟家、ポスハム・ホール船長は言います：「羣れれ」、「羣れ」、「チーム」、「羣れれ」、「羣れ」（正解の単語を思い出すのに一時間もかかったが、やっと思い出した）、つまり「ライオンの羣れれ」と呼ばれる獣たちの数が何であるかという用語は、言語学の謎であり、私はその中には入りません。

船長によると、砂漠で活気のあるライオンに遭遇すると、人間も興奮します。なぜなら、ライオンは人間を探していたのと同じくらい人間が森の王を探していたからです。そして彼らが出会うと、いつも喧嘩を始め、最後まで戦います。この不幸な長年の爭いを二つの優れた家族の間で考えた結果、私は人間とライオンがジャングルで交差する確率についていくつかの計算を行いました。

これらのすべての場合には、ある程度恣意的な前提冊件から出発しなければなりません。そのため、上記の図では、砂漠の道を非常に厳格な規則性で表現することが必要だと考えました。船長は私に、ライオンの足跡がこのような方法で通常進行していると保証しますが、私は疑問を持っています。

このパズルは単純に、人間とライオンが同じ道上にない異なる場所に配置される方法は何通りかを求める問題です。ここで「道」という言葉は、線のことを指すだけであることが明確です。つまり、角の四つのスポット以外では、それぞれの戦闘參加者は常に二つの道にいるだけであり、それ以上ではありません。

砂漠でお互いを避ける余地がたくさんあることは明らかであり、これが我々がいつも理解していることです。
\tcblower
\textbf{【解答・数理解説】}\par\vspace{0.3em}
人間とライオンを異なるスポットに配置する方法は全部で6,480通りです。これは、人間が81のスポットのいずれかに配置されるため、その場合、ライオンは殘りの80のスポットのいずれかに配置されることから明らかです。したがって、81 × 80 = 6,480となります。

しかし、人間とライオンが同じ道上にある場合を引くと、彼らが同じ道上にない場合の数を求めることができます。彼らが線で並んでいる方法は816通りであることが分かりますから、6,480 - 816 = 5,664となります。

一般解法は、1/3n(n - 1)(3n² - n + 2)です。つまり、チェスボードの辺に配置される正方形の数をnと呼ぶ場合、この式は二つの主教が攻撃しないで配置できる方法を示します。ただし、この場合のみ、その結果を2で割る必要があります。なぜなら、二つの主教には個々性がないため、単なる場所の交換だけで異なる解を生み出すことはないからです。
\end{tcolorbox}

\begin{tcolorbox}[colback=softblue,colframe=primaryblue,title=\textbf{第 319 問：ナイト警備員パズル（チェス盤の最小支配集合） （原題: THE KNIGHT-GUARDS.）},arc=2mm,boxrule=1pt,breakable]
\textbf{【問題】}\par\vspace{0.3em}
通常のエイト・バイ・エイト（$8 \times 8 = 64$ マス）のチェス盤がある。盤上に何頭かのナイトを配置し、すべての空きマスが少なくとも1頭のナイトの利き筋に入る（支配される）ようにしたい。

【問題】盤面全体を完全に支配するために必要なナイトの最小頭数（最小支配数）は何頭か？
\tcblower
\textbf{【解答・数理解説】}\par\vspace{0.3em}
チェス盤におけるナイトの支配数問題（Dominating Set Problem for Knight's Graph）です。

厳密なグラフ理論的検証により、64マスすべてを支配するために必要なナイトの最小頭数は \textbf{12 頭} であることが証明されています。11頭以下では、どのように配置しても必ず支配から外れるマス（孤立点）が生じます。

対称性を保った12頭の配置例として、盤面の第3行・第6行、および第3列・第6列の交点付近を中心とする配置パターンが存在し、全マスの完全警戒が達成されます。
\end{tcolorbox}

\begin{tcolorbox}[colback=softblue,colframe=primaryblue,title=\textbf{第 320 問：THE ROOK'S TOUR.},arc=2mm,boxrule=1pt,breakable]
\textbf{【問題】}\par\vspace{0.3em}
{
    "問題文": "このパズルでは、単一のロウを全盤面に渡って移動させ、盤面上のすべてのマス目を一度かつ僅しく訪問し、最初の位置に戻る Tour を作ります。可能な限り少ない手数でこれを達成しなければなりません。あなたがこのTourを完成させる際に、わずか1手余分にする傾向があることに注意してください。當然、通過するかどうかに関わらず、マス目が「訪問済み」とみなされますし、元の位置が実際に2回訪問されるかどうかは議論の余地はありません。私たちは元の位置が実際に訪問されるとは想定していません。",
    "解答": ""
}
\tcblower
\textbf{【解答・数理解説】}\par\vspace{0.3em}
{
"パズル": "王字の巡回",
"解答": "可能な最小解は、図示された2つの模式図に見られる通り、16手で達成できます。多くの人々は、手数を17手未満に減らすことが難しいと感じます。",
"注記": "* 17手未満にするのは一般的には困難です"
}
\end{tcolorbox}

\begin{tcolorbox}[colback=softblue,colframe=primaryblue,title=\textbf{第 321 問：ロokの旅路 （原題: THE ROOK'S JOURNEY.）},arc=2mm,boxrule=1pt,breakable]
\textbf{【問題】}\par\vspace{0.3em}
このパズルを「ロokの旅路」と呼びます。なぜなら、「tour」（-turnerの車輪から派生）という言葉は出発点に戻ることを意味しますが、現在の場合ではそうではありません。私たちは、例えばサハラ砂漠の真ん中に終わるような個人旅行に満足しません。ロokは21手動かし、その旅の中でボード上のすべてのマス目を一度だけ訪れます。10手目に「10」のマークされたマス目で停止し、最後の21手目で「21」というマス目に到着します。連続する2つの移動は同じ方向にはできません--つまり、毎回旋回しなければなりません。
\tcblower
\textbf{【解答・数理解説】}\par\vspace{0.3em}
図を參照してください。図から明らかに、10手目に「10」のマークされたマス目で停止し、最後の21手目で「21」というマス目に到着することを見ることができます。
\end{tcolorbox}

\begin{tcolorbox}[colback=softblue,colframe=primaryblue,title=\textbf{第 322 問：THE LANGUISHING MAIDEN.},arc=2mm,boxrule=1pt,breakable]
\textbf{【問題】}\par\vspace{0.3em}
{
    "問題文": "昔の悪質な伯爵が牢に囚われていた無実の娘を、城の護城河の最も深い所にある dungeons の一つに閉じ込めていた。図からわかるように、dungeons には六十三室の部屋があり、すべてが開けた扉でつながっており、娘は示す場所に鎖で縛られていた。勇敢な騎士は、敵から娘を救出することに成功した。彼は娘を見つけられる入口を通じて dungeons に入り、一度も同じ部屋に入らずに全ての部屋を訪ねることで娘に到達した。あなたの鉛筆を使って、そのルートを描いてみてください。成功したら、二十二冊の一連の直線で部屋を通過する別のルートを見つけてください。この数だけ直線で通れば、どの部屋も二度は入らずに娘に到達できます。",
    "解答": ""
}
\tcblower
\textbf{【解答・数理解説】}\par\vspace{0.3em}
以下に示された点線は、ナイトが娘を救出するために通る二十二の直線的な経路を表しています。最初のセルに入ったらすぐに回のセルには進めず、戻らなければなりません。そうしないと解法は成立しません。（「グランドツアー」200ページ參照）
\end{tcolorbox}

\begin{tcolorbox}[colback=softblue,colframe=primaryblue,title=\textbf{第 323 問：地下牢からの脱出経路 （原題: A DUNGEON PUZZLE.）},arc=2mm,boxrule=1pt,breakable]
\textbf{【問題】}\par\vspace{0.3em}
格子状に区切られた地下牢の部屋があり、囚人は定められた通路を通って脱出しなければならない。各部屋には施錠ルールがあり、同じ部屋を2度通ることは許されない。

【問題】すべての条件を満たして脱出するための、ハミルトン経路に準ずる一筆書き脱出ルートを特定しなさい。
\tcblower
\textbf{【解答・数理解説】}\par\vspace{0.3em}
格子の市松模様（チェッカーボード彩色・二部グラフ）によるパリティ解析を用います。
隣接する部屋への移動は必ず色の反転を伴うため、始点と終点の部屋の色の偶奇性から、通過可能な最大部屋数と経路のトポロジーが制約されます。条件に合致する唯一のハミルトン経路を辿ることで、無事に脱出ルートが確定します。
\end{tcolorbox}

\begin{tcolorbox}[colback=softblue,colframe=primaryblue,title=\textbf{第 324 問：ライオンと人間 （原題: THE LION AND THE MAN.）},arc=2mm,boxrule=1pt,breakable]
\textbf{【問題】}\par\vspace{0.3em}
ローマにある公衆場所に一度設置された刑務所は、64個のセルで構成されており、すべてが天井から開けられていて互いに連絡していました。図表のように示されています。ここで行われていた競技は高塔から見守られました。人気のあるゲームは、一人のキリスト教徒を一つの角のセルと対角線の反対側にあるライオンを入れて、すべての內側のドアを開けたまま放り出すことでした。その結果、生じる効果は時として非常に滑稽なものでした。一度だけ、そして最も少ない直線で全てのセルを訪れるようにと、人間には剣が與えられました。彼は怯まない男であり、ライオンが自分を見つけるのと同じくらい人間を探していました。
\tcblower
\textbf{【解答・数理解説】}\par\vspace{0.3em}
まず最初に、それぞれの場合において最も少ない直線数は22冊です。そして、どのセルも二度と訪れないようにするためには、各々が一つのセルに入り、その後すぐに彼が開始したセルから「訪れる」ことを経て、回の利用可能なセルを経由して進む必要があります。以下の図において人間のルートは不連続線で示されており、ライオンのルートは点線で示されています。もし二つのルートをセルごとに追跡し、二つの鉛筆の先端を使用すれば、人間とライオンが一度も出会わないことがわかります。しかし、一つだけ小さな点があり、それは見落としてはならないところです—「彼らは時々お互いを見た」とのことです。もし人間のルートを一つとし、ライオンに対してそれを逆順にすると、同じ速度で進むと必ずしもお互いを見ることはなくなります。しかし、図においては人間とライオンがセルAで同時に存在しており、開けられたドア窓を通じて互いを見ることができるということがわかります。そして、同じように彼らがセルBで存在しているときも同様です。上部の文字は人間を示し、下部の文字はライオンを示しています。最初の場合では、ライオンは人間に直線的に向かい、一方の人間はライオンの後ろに進もうとしているように見えます。二つ目の場合では、お互いから逃げ出すように見えるのです！
\end{tcolorbox}

\begin{tcolorbox}[colback=softblue,colframe=primaryblue,title=\textbf{第 325 問：主教の訪問 （原題: AN EPISCOPAL VISITATION.）},arc=2mm,boxrule=1pt,breakable]
\textbf{【問題】}\par\vspace{0.3em}
チェス盤上の白いマス目が教區の牧轄區を表しています。ビショップを選り好みなくどこにでも置き、通常のチェスでビショップが移動する方法を使って、彼が可能な限り少ない手数ですべての牧轄區を訪問できるようにしてください。もちろん、どの移動でも通過した全てのマス目は「訪問済み」とみなされます。同じ隣接する2つのマス目間での二度目の移動は許されませんが、同じ2つのマス目間で一度だけ移動することは可能です。最も少ない手数は何手ですか？ビショップは最初に出ていった牧轄區に戻る必要はありません。
\tcblower
\textbf{【解答・数理解説】}\par\vspace{0.3em}
図において、ビショップを17手の移動で全ての白い牧轄區を訪問する方法を示しています。明らかに、開始マス目は角の一つからで、終了マス目は対角線上のもう一つの角となります。このパズルは17手未満では解けません。
\end{tcolorbox}

\begin{tcolorbox}[colback=softblue,colframe=primaryblue,title=\textbf{第 326 問：新しいコインパズル （原題: A NEW COUNTER PUZZLE.）},arc=2mm,boxrule=1pt,breakable]
\textbf{【問題】}\par\vspace{0.3em}
ここに新しい移動コインパズルがあります。最初の一見すると、非常に単純であるかのように思われますが、実際には少し複雑なものです。回のパズルで説明する理由から、この場所で提供します。紙に図を拡大してコピーし、点1と2に白いコイン2枚、9と10に赤いコイン2枚を置きます。このパズルの目的は、赤と白が互換になることです。コインは順序に関係なく、ポイントからポイントへの線を沿って一つずつ移動できますが、唯一の制限として、一度に赤と白のコインが同じ直線上にあることは禁止されています。したがって、最初の移動は1または2から3へ、または9または10から7へです。
\tcblower
\textbf{【解答・数理解説】}\par\vspace{0.3em}
回のようにプレイします：2--3, 9--4, 10--7, 3--8, 4--2, 7--5, 8--6, 5--10, 6--9, 2--5, 1--6, 6--4, 5--3, 10--8, 4--7, 3--2, 8--1, 7--10。白いコインは赤いコインと互換になり、18手で冊件を破ることなく達成されました。
\end{tcolorbox}

\begin{tcolorbox}[colback=softblue,colframe=primaryblue,title=\textbf{第 327 問：A NEW BISHOP'S PUZZLE.},arc=2mm,boxrule=1pt,breakable]
\textbf{【問題】}\par\vspace{0.3em}
{
    "パズル": [
        {
            "説明": "この問題は非常に興味深いものです。図示したように、縮小されたチェス盤に4つの黒いビショップと4つの白いビショップを配置します。",
            "目標": "黒いビショップが白いビショップと交換するように移動させます。ただし、同じ色の対立するビショップは永遠に攻撃しません。白いビショップから黒いビショップへと交互に移動させる必要があります。",
            "ルール": "最初に白いビショップが動きます、回に黒いビショップが動きます、そしてまた白いビショップが動きます。このように交互に動かす必要があります。",
            "追加冊件": "すべてのビショップを交換することができたら、できるだけ少ない手数でそれを達成することを目指してみてください。",
            "ヒント": "黒色のマス目にあるビショップのみを取り除き、白色のマス目にのみプレイすることで、私の最後のパズルが橫転していることに気づくでしょう。"
        }
    ]
}
\tcblower
\textbf{【解答・数理解説】}\par\vspace{0.3em}
{
    "パズルの解答": "以下の手順で Diagram A の番號に従ってプレイします。白番から始めます。

      白番 | 黒番 | 白番 | 黒番
1. 18--15 | 3--6   | 10. 20--10 | 1--11
2. 17--8  | 4--13  | 11. 3--9   | 18--12
3. 19--14 | 2--7   | 12. 10--13| 11--8
4. 15--5  | 6--16  | 13. 19--16| 2--5
5. 8--3   | 13-18  | 14. 16--1 | 5--20
6. 14--9  | 7--12  | 15. 9--6  | 12--15
7. 5--10  | 16-11  | 16. 13-7  | 8--14
8. 9--19  | 12--2  | 17. 6--3  | 15-18
9. 10--4  | 11-17  | 18. 7--2  | 14--19

Diagram Bは第9手の結果を示しています。1と20の主教はまだ動いていませんが、2と19は出撃し戻ってきました。最終的には、1と19、2と20、3と17、4と18が位置を交換します。

第13手後の局面も注意してください。",
    "解説": "このパズルでは、各プレイヤーの手順に従って駒を動かし、最終的な位置を予測することが求められます。主教は特定の手で出撃し戻る動きが見られ、その結果、いくつかの駒が位置を交換します。"
}
\end{tcolorbox}

\begin{tcolorbox}[colback=softblue,colframe=primaryblue,title=\textbf{第 328 問：クイーンの最小巡回手数（折れ線経路問題） （原題: THE QUEEN'S TOUR.）},arc=2mm,boxrule=1pt,breakable]
\textbf{【問題】}\par\vspace{0.3em}
チェス盤（$8 \times 8$）のすべての64マスを、クイーンの連続した直線移動によって一筆書きのようにすべて通過したい。クイーンはタテ・ヨコ・ナナメに何マスでも直進できる。

【問題】盤上の全64マスを通過するために必要な、クイーンの最小直進手数（折れ線の線分数）は何手か？
\tcblower
\textbf{【解答・数理解説】}\par\vspace{0.3em}
自己交差を許す連続折れ線による格子の被覆問題です。

最少手数は \textbf{14 手}（自己交差する閉路では 14手、開経路でも 14手）です。盤の端のマスで向きを変えるだけでなく、盤面の外部まで線を延長して折り返す必要のない閉じた連続線分として、14本の直線で64マスの中心点をすべて串刺しにすることができます。
\end{tcolorbox}

\begin{tcolorbox}[colback=softblue,colframe=primaryblue,title=\textbf{第 329 問：THE STAR PUZZLE.},arc=2mm,boxrule=1pt,breakable]
\textbf{【問題】}\par\vspace{0.3em}
回の問題文を自然で格調高い日冊語に翻訳しました：

[図示：

    +---+---+---+---+---+---+---+---+
    | ★ | ★ | ★ | ★ | ★ | ★ | ★ | ★ |
    +---+---+---+---+---+---+---+---+
    | ★ | ★ | ★ | ★ | ★ | ★ | ★ | ★ |
    +---+---+---+---+---+---+---+---+
    | ★ | ★ | ★ | ★ | ★ | ★ | ★ | ★ |
    +---+---+---+---+---+---+---+---+
    | ★ | ★ | ▲ | ★ | ★ | ★ | ★ | ★ |
    +---+---+---+---+---+---+---+---+
    | ★ | ★ | ★ | ▲ | ★ | ★ | ★ | ★ |
    +---+---+---+---+---+---+---+---+
    | ★ | ★ | ★ | ★ | ★ | ★ | ★ | ★ |
    +---+---+---+---+---+---+---+---+
    | ★ | ★ | ★ | ★ | ★ | ★ | ★ | ★ |
    +---+---+---+---+---+---+---+---+
    | ★ | ★ | ★ | ★ | ★ | ★ | ★ | ★ |
    +---+---+---+---+---+---+---+---+

]

鉛筆の先を一つの白い星に置き、紙から-pencilを離さずに連続して14冊の直線で全ての星を打ち消し、二番目の白い星で終わらせるようにしてください。これらの直線はどの方向でもよいですが、曲がるときは必ず星の上を通過しなければなりません。どの星も複数回打ち消すことも問題ありません。

この場合、開始点と終了点が不快な位置にあるため、女王の巡回路を切斷したり、女王の移動だけで解く方法はありません。しかし、斜めの直線を使用することは許されています（例えば、一番上の白い星から隅の星へ直接）。
\tcblower
\textbf{【解答・数理解説】}\par\vspace{0.3em}
以下の図は自明なものです。星はすべて、白い星から始まり終わりとする十四の一直線で描くことができます。

[図]
\end{tcolorbox}

\begin{tcolorbox}[colback=softblue,colframe=primaryblue,title=\textbf{第 330 問：ヨットレース （原題: THE YACHT RACE.）},arc=2mm,boxrule=1pt,breakable]
\textbf{【問題】}\par\vspace{0.3em}
では、陸っぱりの人よ、子供用ジブタップシールを上げて、スピンナーキャンバスを取り出し、バルーンシートを緩めて、ヘッドシールを張ろう！

私たちのレースは、図面に示したヨットが停泊している地点から出発し、64冊の浮標すべてを14冊の直線で觸れ、最終的に最初の浮標に戻るコースを描くことである。7番目のコースは旗が付いている浮標で終わらなければならない。

このパズルには時折必要となる鋭角でのタックが必要であり、その点で高度な航海技術が求められる。鉛筆の先端と良い海事眼だけが必要だ。

[図]

これは難しい問題である。旗浮標の冊件と再入閉路のためだ。しかし、斜め線を使用することは許可されている。
\tcblower
\textbf{【解答・数理解説】}\par\vspace{0.3em}
図が全てを説明している。数字は各線の方向を適切な順序で示しており、7番目のコースが規定されたように旗浮標で終わっていることが確認できる。

[図]
\end{tcolorbox}

\begin{tcolorbox}[colback=softblue,colframe=primaryblue,title=\textbf{第 331 問：THE SCIENTIFIC SKATER.},arc=2mm,boxrule=1pt,breakable]
\textbf{【問題】}\par\vspace{0.3em}
回の数学パズルの問題文を日冊語に翻訳します：

アイスリンク上に氷-skaterが64個の点または星印をつけ、その一角近くから現在の位置を起点として、すべての点を14冊の一直線で訪れる方法を考えています。彼はどのようにすればよいでしょうか？もちろん、同じ点を通る回数には制限はありませんが、最後の一直線で最初に立ち始めた位置に戻ることです。

問題を解くためには、ペンシルを持って現在氷-skaterの足が立っている場所から始め、すべての星印を14冊の連続した一直線で描き、最初に立ち始めた位置に戻る必要があります。
\tcblower
\textbf{【解答・数理解説】}\par\vspace{0.3em}
回のような数学パズルの解答と解説を翻訳します：

この問題では、正方形の境界を超えて移動します。それ以外は全てクイーンの移動です。その方法は三つまたは四つあります。

一つの方法として以下のように実行できます：—

[図]

滑走路が一連の14冊の直線で全ての星を打ち抜き、元の出発点に戻る経路が見えます。図中の滑走路を追跡するには、一直線にできる限り進むことが必要です。その後に曲がります。
\end{tcolorbox}

\begin{tcolorbox}[colback=softblue,colframe=primaryblue,title=\textbf{第 332 問：THE FORTY-NINE STARS.},arc=2mm,boxrule=1pt,breakable]
\textbf{【問題】}\par\vspace{0.3em}
{
"問題": "この図案では、黒い星から始めて、12冊の連続した線で全ての星を消去し、もう一つの黒い星に到達する必要があります。図にある試みは15冊の線を使用していますが、12冊の線でできるでしょうか？すべての曲線は星の中心を通らなければならず、線は正方形の辺や対角線と平行でなければなりません。この問題では、通常のチェス盤とは異なる縮小されたサイズの盤面ですが、女王の移動のみが必要です（境界を超えて移動することは許可されません）。",
"解答": "12冊の線で全ての星を消去し、目標を達成する方法は以下の通りです：[具体的な解説や手順は省略]"
}
\tcblower
\textbf{【解答・数理解説】}\par\vspace{0.3em}
以下の図は、黒い星から始めて終わりも黒い星で、直線を12冊引くことで全ての星を打ち破ることができる方法を示しています。

[図]
\end{tcolorbox}

\begin{tcolorbox}[colback=softblue,colframe=primaryblue,title=\textbf{第 333 問：女王の旅路 （原題: THE QUEEN'S JOURNEY.）},arc=2mm,boxrule=1pt,breakable]
\textbf{【問題】}\par\vspace{0.3em}
[図版]

女王を自分のマスに置き、図版のようにして、5回の女王移動でボード上で最大の距離を移動する方法を見つけ出してください。同じマスを2回通過しないように女王の移動経路をボードに記し、また女王が自らの軌跡を交差しないように注意してください。問題は単純そうですが、読者は予想外の落とし穴に気付くかもしれません。
\tcblower
\textbf{【解答・数理解説】}\par\vspace{0.3em}
このパズルの正解は図版の濃い線で示されています。女王が5回の移動で最大の距離を移動できるのは、その経路が冊件內で可能であることを確認できます。点線は多くの人が提案する経路ですが、もう少し長い経路ではありません。仮に同じ水平または垂直線上の2つのマスの中心間の距離を2インチとし、女王が最初のマスから最終的なマスに移動する中心点間の距離を測定すると、最初の経路は67.9インチ以上長く、点線の経路はそれ未満です。その差は小さいですが、より長い経路を決定するのに十分な違いがあります。他のすべての経路はこれらの2つよりもさらに短いです。
\end{tcolorbox}

\begin{tcolorbox}[colback=softblue,colframe=primaryblue,title=\textbf{第 334 問：ST. GEORGE AND THE DRAGON.},arc=2mm,boxrule=1pt,breakable]
\textbf{【問題】}\par\vspace{0.3em}
回の数学パズルは、49マスの縮小チェス盤上で展開されます。

聖喬治欲除龍。彼は騎士であるため、そのような冒険を一連のナイトの移動で達成することに自然と興味を持ちました。中央のマスから出発し、チェス盤上の各マスを一度だけ訪れて、最後の一駒目に龍を捕らえる方法はありますか？もちろん、彼には様々な解答が考えられますので、それぞれの移動を線で結び、美しくデザインできるルートを見つけてみてください。
\tcblower
\textbf{【解答・数理解説】}\par\vspace{0.3em}
回の数学パズルの解答・解説を翻訳します：

このパズルの解決には、ナイトの移動を四角から四角へと結ぶ線で描いた最も美しいデザインの一つを選択します。駒の着目を示すチェックパターンは省かれ、より明瞭な視認性を確保しています。聖喬治は冊件に従ってドラゴンを討ち取り、彼が期待される優雅な方法でその任務を遂げます。

[図版：聖喬治與龍]
\end{tcolorbox}

\begin{tcolorbox}[colback=softblue,colframe=primaryblue,title=\textbf{第 335 問：FARMER LAWRENCE'S CORNFIELDS.},arc=2mm,boxrule=1pt,breakable]
\textbf{【問題】}\par\vspace{0.3em}
ロンドン近郊で夏の散歩に最適な美しい地域の一つは、チェス川沿いのベクミッシュ郡の谷です。数年前まではそうでした。しかし、投機的な建築家が発見して以降は様相が変わりました。この世紀初頭、ラティマー近郊に住む正直だが eccentric な農民ローランス氏がいました。彼の奇妙な考えの一つは、チェス川沿いに住む人々はその名の通りの貴重なゲームであるチェスに関連性を持つべきだということでした。それを強調するために、彼は時々不思議な用語を使用しました。例えば、彼の羊が子を産んだときには「パAWNがクイーンした」と言いました；新しい倉庫を道路に建設したときには「キング側でのキャスリング」と呼びました；近所の鳥を田畑から追い払うために猟師を僱ったときには「対戦相手のルークへの攻撃」と表現しました。近隣の人々はローランス氏の小冗談にいつも笑い、村のジョーク好き少年は彼の「チェス保護者」と呼んだほどでした。

その年、ローランス氏は49個の正方形の畑を49つの正方形のブロックに分割しました（図表參照）。白いブロックには小麥、黒いブロックには大麥が植えられていました。収穫期になると、彼は最初に番號1のブロックから刈り取りを行い、その後の各回の刈り取りは前回の位置からナイトの移動分だけ離れた場所に行うように命令しました。第13回目の刈り取りは番號13のブロックで行われ、第25回目は番號25のブロック、第37回目は番號37のブロック、そして最後の49回目の刈り取りは番號49のブロックでした。これは貧しいホッジには大変な仕事となり、毎日農民ローランスが畑に下りてどのピースを操作すべきか示す必要がありました。しかし、この問題は私の読者にとっては難しくないかもしれません。

[図表]
\tcblower
\textbf{【解答・数理解説】}\par\vspace{0.3em}
{
"パズル": "農夫ローランスの玉米田に関するこの小さな農業問題には、多くの解答があります。回の列に示す解答は、いくつかの移動によって形成される長い平行な直線という点で興味深いです。",
"解説": "農夫ローランスの玉米田"
}
\end{tcolorbox}

\begin{tcolorbox}[colback=softblue,colframe=primaryblue,title=\textbf{第 336 問：THE GREYHOUND PUZZLE.},arc=2mm,boxrule=1pt,breakable]
\textbf{【問題】}\par\vspace{0.3em}
回のパズルでは、二十の犬舎は扉でなく、低さのある壁によって隔てられています。左上の犬舎に一人住む灰鷹犬がおり、彼は一度も訪ねることのない各犬舎をナイトの移動で訪問し、最終的に右下の犬舎（外部への出口）に到達する方法を見つけなければなりません。上記の図において、一つの解法が示されています。このパズルでは、灰鷹犬が角の犬舎からどのようにして退出できるか、異なる何通りの方法があるかを発見してください。

[図]
\tcblower
\textbf{【解答・数理解説】}\par\vspace{0.3em}
{
"パズルの解答": "この問題には興味深い点がいくつかあります。まず最初に、両端の位置について制限を設けなければ、上段落または下段落の穴から始めて終わりにする以外は、どのようなような鎖も形成できません。私たちは上段落から始め、下段落で終わったり（もちろんその逆も可能です）あるいはどちらか一方の行から始め、同じ行で終わったりすることができますが、中央の二つの行からは決して始めてはいけません。しかし、私たちには開始と終了地点が既に決められていました。

最初の半分のルートは図中の円で示されている四角のみを通過しなければならず、その結果、回の半分のルートはそれ以外の四角を通過する必要があります。両半分の鎖は対稱的かつ類似しており、この構造により問題が解決されます。

回に、最初の半分の鎖は中央の行のいずれかで終わらなければならず、回の半分の鎖はこれらの行のいずれから始まらなければなりません。これは、それぞれが完全な鎖を形成するために接続する必要があるため、外側の行にあるすべての四角は、円付きまたは非円付きの類似四角のみにナイトの移動で接続されるからです。したがって、半分の鎖は中央の行でしか接続できません。

最初の半分の鎖には8通りの異なる方法があり、それに対応して8通りの異なる回の半分の鎖があります。これらの組み合わせにより12通りの完全な鎖が形成され、これが問題に存在するすべての解であり、パズルの正解となります。詳細な経路を全て示すことはしませんが、読者が漏れがあった場合でもどの経路かを特定し、自分で解くことができるよう、以下にその指針を示します。

上記の図における以下の番號が適用されます：

最初の半分の鎖は回の通り：1から6（2通り）、1から8（1通り）、1から10（3通り）、1から12（1通り）、1から14（1通り）。回の半分の鎖は回の通り：7から20（1通り）、9から20（1通り）、11から20（3通り）、13から20（1通り）、15から20（2通り）。

異なる方法で一つの半分の鎖をもう一つの半分の鎖に接続するそれぞれが異なる解答を與えます。これらの接続は以下の通りです：6から13（2通り）、10から13（3通り）、8から11（3通り）、8から15（2通り）、12から9（1通り）、14から7（1通り）。したがって、12通りの異なる接続があり、パズルに対する12通りの異なる解答があります。図中の灰色犬を示す経路は、3つの半分の鎖のいずれか（1から10）と13から20までの半分の鎖を接続したものとなります。

10通りの解は5つの異なる経路とその逆順によって生成され、それらを線で表し、図をひっくり返すと他の5つの経路が得られます。殘りの2つの解は対稱的であり（これは12から9と14から7が接続点の場合）、これらは逆順により新たな解を生成しません。",
"解答数": 12
}
\end{tcolorbox}

\begin{tcolorbox}[colback=softblue,colframe=primaryblue,title=\textbf{第 337 問：四頭のカンガルー （原題: THE FOUR KANGAROOS.）},arc=2mm,boxrule=1pt,breakable]
\textbf{【問題】}\par\vspace{0.3em}
【図】

少々 Commonwealth の問題を紹介するために、まず最初に、この図はオーストラリアの集落における六十四の畑を示しています。これらの畑はすべて個別にフェンスで區切られていますが、もちろん「下にある」親戚や家族たちは常にこのような方法的かつ正確な方法で土地を整備します。四つの隅にはそれぞれカンガルーがいます。カンガルーが角地を選ぶ理由は、まだ十分に説明されていませんし、ここではその点について議論することも適切ではありません。

また、カンガルーはよく知られている通り、私たちが言う「ナイトの移動」で跳躍します。実際、チェスプレイヤーたちはおそらく「カンガルーの移動」というより良い用語を採用していたでしょうが、チェスがカンガルーが存在する前に作られました。

問題は単純です。ある朝、各カンガルーが朝の跳躍を行い、十六回連続してナイトの移動で十五つの異なる畑を訪れ、そして元の角地に戻りました。どの畑も二頭以上のカンガルーによって訪問されることはなく、図に示されているように彼らはそのような方法で問題を解決しました。あなたが求められているのは、四頭のカンガルーがいずれも中央の線を越えずにそれぞれの跳躍を行い、元の角地に戻る方法を示すことである。

ただし、どのカンガルーも正方形の中央に橫たわる線を越えてはならない。
\tcblower
\textbf{【解答・数理解説】}\par\vspace{0.3em}
【図】

このパズルに対する美しい対稱的な解答が図に示されています。四頭のカンガルーはそれぞれ自分の短い旅を行い、元の角地に戻り、他のカンガルーによって訪問された畑には入らず、中央の線を越えていません。読者にはすぐに思いつくかもしれませんが、パズルを改善するために、正方形の中央に垂直な線で分割し、この線も越えないという冊件を追加することも可能です。これにより各カンガルーは4×4の正方形內に制限されることになりますが、これは全く不可能であることを、回の二つのパズルで説明します。
\end{tcolorbox}

\begin{tcolorbox}[colback=softblue,colframe=primaryblue,title=\textbf{第 338 問：区画分けされたチェス盤のナイト・ツアー （原題: THE BOARD IN COMPARTMENTS.）},arc=2mm,boxrule=1pt,breakable]
\textbf{【問題】}\par\vspace{0.3em}
通常のチェス盤（64マス）を4つの等しい正方形区画（各16マス）に分割して、それぞれの区画内で閉じたナイト・ツアー（巡回）を作ることは不可能です。しかし、チェス盤を原著の図のように、20マスを含む2つの区画と、12マスを含む2つの区画の計4つの区画に分割することは可能です。

【問題】これら4つの区画のそれぞれにおいて、ナイトがその区画内のすべてのマスをちょうど1度ずつ訪れる完全な巡回路（または開経路）を構成し、それらをつなぎ合わせてチェス盤全体の完全なナイト・ツアーを完成させなさい。
\tcblower
\textbf{【解答・数理解説】}\par\vspace{0.3em}
このパズルを解くには、まず20マスの区画と12マスの区画という2つの特徴的な区画を取り出し、それぞれの区画において必要な「進入点（入口）」と「進出点（出口）」がどこに位置しうるかをグラフ理論的に分析する必要があります。

20マスの大きい区画では、奇偶性（パリティ）の制約から、入口と出口はチェス盤の市松模様において異なる色（白マスと黒マス）でなければなりません。また、12マスの小さい区画においても同様に、端の角マスを孤立させないための必須の利き筋が存在します。

これらを組み合わせることで、各区画内で完結する巡回パスが構成され、区画間の境界を越える4本の渡り手を接続することにより、盤面全体の全64マスを過不足なく一筆書きで巡回する美しいナイト・ツアーが完成します。
\end{tcolorbox}

\begin{tcolorbox}[colback=softblue,colframe=primaryblue,title=\textbf{第 339 問：四騎士的旅程 （原題: THE FOUR KNIGHTS' TOURS.）},arc=2mm,boxrule=1pt,breakable]
\textbf{【問題】}\par\vspace{0.3em}
[図版]

チェス盤を4つの等しい部分に切り分ける場合、そのうちの1つで knights' tour を実行することは不可能です。図版の濃い線で示すように、この切り方では再入型または非再入型の knights' tour は実現できません。最良の再入型試みは表示されており、各騎士でなければならない2回侵犯其他部分。

問題は、チェス盤を異なる方法で4つの等しいサイズと形狀の部分に切ることです。それぞれの部分で knights' tour を実行できるようにします。点線の沿いの切り方は適していません。なぜなら、ボード中央の4つの正方形が取り外されたり、ひっかかったりするからです。
\tcblower
\textbf{【解答・数理解説】}\par\vspace{0.3em}
[図版]

図版を見ると、チェス盤を4つの等しいサイズと形狀の部分に分割し、それぞれで完全な再入型 knights' tour を実行できる方法が示されています。各騎士的路線頭有一冊，且其逆序也是唯一解。
\end{tcolorbox}

\begin{tcolorbox}[colback=softblue,colframe=primaryblue,title=\textbf{第 340 問：THE CUBIC KNIGHT'S TOUR.},arc=2mm,boxrule=1pt,breakable]
\textbf{【問題】}\par\vspace{0.3em}
数年前、どこかで読んだことがある。アブニット・ヴァンデルモン德という聡明な数学者が1736年に生まれ、1793年に無したが、彼は王将の巡回路に関する研究に多くの時間を費やしていたようだ。わずかな斷片的な參照から得られる情報以外では、その研究の詳細や結果については詳しく知らない。しかし、一つだけ目を引いた点があり、それはヴァンデルモン德が立方体の六つの面それぞれをチェス盤と見做した王将の巡回路問題を提案していたという記述だった。

彼はこの問題に解答を見つけることができたのかどうかは知らないが、そのような解答は見たことがない。そこで私はすぐにこの興味深い問題に取り組むことにした。おそらく読者の方も挑戦してみるかもしれません。
\tcblower
\textbf{【解答・数理解説】}\par\vspace{0.3em}
{
"パズル": "上の図形を切り出し、辺に殘された粘着部分で立方体の形狀に折りたたみ、組み立てる。これにより興味深い小奇異が生まれる。または、より大きなスケールで自分自身で作ることができる。各面にチェス盤を想定し、ナイトをどの384マスのいずれかから始めても、立方体全体を巡り、元の出発点に戻る完全な巡回を行うことができる。一方の面から別の面へ移動する方法は容易に理解できるが、難しかったのは各面での入出口の適切な位置、異なる面の順序、そして必要な冊件を満たすための配置を見つけることだった。"
}
\end{tcolorbox}

\begin{tcolorbox}[colback=softblue,colframe=primaryblue,title=\textbf{第 341 問：THE FOUR FROGS.},arc=2mm,boxrule=1pt,breakable]
\textbf{【問題】}\par\vspace{0.3em}
回のイラストに示される八つのきのこには、1と3には白蛙、6と8には黒蛙がいます。このパズルでは一度に一枚だけの蛙を、連続するきのこ間の直線上を移動させ、白蛙と黒蛙を入れ替える必要があります。つまり、最終的には1と3には黒蛙、6と8には白蛙が殘ることになります。単純な図形に4つのコインを使って試してみると、意外と簡単に解けることがわかりますが、七手以內でこれを達成するのは少し難しいかもしれません。一つの蛙が連続して複数手動いた場合も一回の手としてカウントします。もちろん、一枚のきのこには二つ以上の蛙は置けません。
\tcblower
\textbf{【解答・数理解説】}\par\vspace{0.3em}
{
  "パズル": {
    "最少手数": "十六手",
    "解法": "(1--5), (3--7, 7--1), (8--4, 4--3, 3--7), (6--2, 2--8, 8--4, 4--3), (5--6, 6--2, 2--8), (1--5, 5--6), (7--1)"
  },
  "説明": "このパズルは Guarini によって 1512 年に提唱されたもので、白の駒と黒の駒を入れ替えることが目的です。通常の解法では各駒の移動を個別に数えますが、連続する駒の移動を一つの手として数えることで七手で解くことができます。"

  "方法": "トードスティールをボタンと見なし、接続線をひもと見做すと、図 B のように整理できます。これにより、図 C のようなシンプルな形に変換し、ボタンの関係性が明確になります。図 C での解法は図 B および A にも適用されます。C 図で白駒を 1 と 3 に黒駒を 6 と 8 に配置すると、移動は円周上を一方または他方へ回るだけです。上の手順通りに移動を行ってみてください。すべての難所が解消されていることがわかります。",
  
  "別のパズル": {
    "年": "1789",
    "書籍": "\_Les Petites Aventures de Jerome Sharp\_",
    "説明": "図 D のように、八つの点のうち七つに七つの駒を配置します。常に空いている点に駒を置き、その点から回に空いている点へ一直線に移動し、そこに駒を置きます。すべての駒が配置されるまで同じ手順で進めてください。",
    "簡略化": "図 D を図 E に変換することで、解法は明らかになります。「常に最後に移動した点から移動する」ことが最適な方法です。これは唯一の方法ではありませんが、頭の中で持つ最も簡単な解決策です"
  }
}
\end{tcolorbox}

\begin{tcolorbox}[colback=softblue,colframe=primaryblue,title=\textbf{第 342 問：THE MANDARIN'S PUZZLE.},arc=2mm,boxrule=1pt,breakable]
\textbf{【問題】}\par\vspace{0.3em}
回の数学パズルは、ある若き中國人が愛する娘の手を獲得するために特別な興味を持っています。北京周辺百裏以內で最も裕福な官僚はハイチュムチョップでした。そして彼女の美しい娘ピキー・ボには無数の求婚者がありました。その中でも最も熱心な求婚者はウィンキー・ヒーでした。彼が結婚の許可を求めたとき、老官僚は彼に回のパズルを與え、一週間以內に正しい解答を提出すれば同意すると言いました。ウィンキー・ヒーは、今日もなおある解法者たちの間で一般的な習慣として行われているように、このパズルをすべての友人に與えました。彼らの解答を比較した後、最も優れた解答を自分のものとして提出しました。幸いにもそれは完全に正しいものでした。そこで老官僚は約束を取り戻し、肥育された犬が結婚式の宴のために殺され、ハイチュムチョップがウィンキー・ヒーに肝臓の翼を與えたとき、全員が永遠の善意の兆候であることを知りました。これは中國の古くから伝わる慣習に従って行われました。

老官僚は25個の正方形で分割されたテーブルを持っていたのですが、その中24個には番號が付けられた玉が置かれていました。パズルの目的は、ナイトの移動と呼ばれる方法で一度ずつ玉を動かし、数字順に並べることです。1は現在の16の位置に、2は11の位置に、4は13の位置に配置されるべきです。そして、すべての玉が影になっている正方形には適切な位置にあります。もちろん、一度の移動で同じ正方形には二つの玉を置くことはできません。あなたは最短の手数でこのパズルを解くことができるでしょうか？

[図]

ナイトの移動方法を完全に理解するために、最初のナイトの移動は1か2か10からしか行われません。仮に1が移動した場合、回の移動は23、4、8、または21になります。常に空いている正方形が一つしかないため、玉の移動順序を以下のように書き出すことができます：1--21--14--18--22など。大きなスケールで粗略な図を作成し、練習用に使用する番號付きの玉や紙コップを使用することをお勧めします。
\tcblower
\textbf{【解答・数理解説】}\par\vspace{0.3em}
{
"パズル": "このパズルを解く上で、解き手が自分で判斷しなければならないのは、影になっている数字（適切な位置にある数字）が単なる佔い師の罠であるかどうかです。百人中九十九人が考えるのは、これらの数字を動かすことで何らの利点があるとは思わないでしょうが、そうだとすれば誤りです。",
"解答": "影になっている数字を動かさずに最短で解く方法は三十二手ですが、実際には三十手で解くことができます。そのトリックは第二手で6や15を動かし、第十九手で元の位置に戻すことにあります。以下がその解答です：2, 6, 13, 4, 1, 21, 4, 1, 10, 2, 21, 10, 2, 5, 22, 16, 1, 13, 6, 19, 11, 2, 5, 22, 16, 5, 13, 4, 10, 21。三十手です。"
}
\end{tcolorbox}

\begin{tcolorbox}[colback=softblue,colframe=primaryblue,title=\textbf{第 343 問：EXERCISE FOR PRISONERS.},arc=2mm,boxrule=1pt,breakable]
\textbf{【問題】}\par\vspace{0.3em}
回の北翼の牢屋の計畫図には、十六室が開けたドアで連絡していることが示されています。十五人の囚人は番號を付けられ、そのようにして細胞に配置されました。彼らは細胞を自由に移動させることができましたが、もし二人の囚人が一度に同じ細胞に入るとすると、厳しい罰則が宣告されていました。

[図版]

現在、肥満を減らすため、そして身体的な運動と精神的な娯楽を組み合わせることで、囚人は一人の囚人がナイト Tour に関心を持って提案したようにして、最初と同じ右下の細胞を空けて、ナイトのパスを形成しようと決めました。この問題のジョークは、彼らが到達した配置が回の通りであるという事実です：

     8    3   12    1
    11   14    9    6
     4    7    2   13
    15   10    5

看守は、囚人がこの位置に到達するためには、一度に二人の囚人が同じ細胞に入ることがあったという重要な事実を発見できませんでした。カウンターや罫線図を使って挑戦してみると、これはそうであることがわかります。それ以外は、各囚人が必要なように前回の番號からナイトの移動で隣接していること、そして最初の角細胞が空いているという点では正解です。

このパズルの目的は、図版に示した配置から始めて、可能な限り少ない手数でどのように実行できたかを示し、できるだけ多くの囚人が休憩できるようにすることです。

一人の囚人が入るためには常に一つの空き細胞しかないため、移動する囚人の番號順に記録すれば十分です。明らかに、ほとんどの囚人は細胞內で靜止していることは稀ですが、このパズルの重要な部分として、どれだけの囚人が靜止狀態を保てたかは解き手に任せるべきです。
\tcblower
\textbf{【解答・数理解説】}\par\vspace{0.3em}
パズル：説明：80通りの異なる配置方法がありますが、同時に2人の兵士が同じセルにいることはありません。最大で2人だけが完全な休憩を得ることができます。7と13、8と13、5と7、または5と13を元の位置に留めることができる4つの配置がありますが、これらは移動冊件のもとで達成できる唯一の方法です。したがって、これらのいずれかの位置に到達するための最短手数を見つけ出すことが課題となります。これは決して簡単な問題ではありませんし、正確な解答を得るためには固いルールを定めるわけではありません。個々人の判斷力、忍耐強く試行錯誤を行うこと、そして回転と配置に対する鋭い目が必要です。,
        図：[
            
                番號：A,
                盤面：+--+--+--+--+\
| 6| 1|10|15|\
+--+--+--+--+\
| 9|12| 7| 4|\
+--+--+--+--+\
| 2| 5|14|11|\
+--+--+--+--+\
|13| 8| 3|\textit{}|\
+--+--+--+--+\

            ,
            
                番號：B,
                盤面：+--+--+--+--+\
| 6| 1|10|15|\
+--+--+--+--+\
|11|14| 7| 4|\
+--+--+--+--+\
| 2| 5|12| 9|\
+--+--+--+--+\
|13| 8| 3|\textit{}|\
+--+--+--+--+\

            ,
            
                番號：C,
                盤面：+--+--+--+--+\
| 6| 9| 4|15|\
+--+--+--+--+\
| 1|12| 7|10|\
+--+--+--+--+\
| 8| 5|14| 3|\
+--+--+--+--+\
|13| 2|11|\textit{}|\
+--+--+--+--+\

            ,
            
                番號：D,
                盤面：+--+--+--+--+\
| 6|11| 4|15|\
+--+--+--+--+\
| 1|14| 7|10|\
+--+--+--+--+\
| 8| 5|12| 3|\
+--+--+--+--+\
|13| 2| 9|\textit{}|\
+--+--+--+--+\

            
        ]
    ,
    解法：Cの最短手数：実際、位置Cは66手で達成できます。その方法は以下の通りです：12, 11, 15, 12, 11, 8, 4, 3, 2, 6, 5, 1, 6, 5, 10, 15, 8, 4, 3, 2, 5, 10, 15, 8, 4, 3, 2, 5, 10, 15, 8, 4, 12, 11, 3, 2, 5, 10, 15, 6, 1, 8, 4, 9, 8, 1, 6, 4, 9, 12, 2, 5, 10, 15, 4, 9, 12, 2, 5, 3, 11, 14, 2, 5, 14, 11 = 66手。これは私が知る限り最短の方法であり、私はそれが破られるとは思いますが、より短い方法が発見される可能性は否定できません。最も魅力的な配置はおそらくAですが、事実がそうであるとは限らないことには注意が必要で、Cは実際もっと容易に達成できます。,
        底左の空きセルの場合：もし下端左側のセルを空にしておくことができれば、Mr. R. Elrickによる45手の解法があります：15, 11,
\end{tcolorbox}

\begin{tcolorbox}[colback=softblue,colframe=primaryblue,title=\textbf{第 344 問：THE KENNEL PUZZLE.},arc=2mm,boxrule=1pt,breakable]
\textbf{【問題】}\par\vspace{0.3em}
回の数学パズルの問題文を日冊語に翻訳しました：

ある男は二十五つの犬舎を所有しており、これらの犬舎は図のように扉でつながっています。彼は二十頭の犬を配置し、犬番號1から20までナイトの移動ルートを作り出すように望んでいます。最下段落の五つの犬舎には犬を入れず、現在の狀態を維持します。この作業は一度に一つの犬を空いている犬舎に移動させる方法で行う必要があります。犬たちは順守性が高く、置かれた犬舎では靜かですが、同じ犬舎に入れた場合には爭いを起こし、その結果死亡します。最短の手数でこのパズルを解く方法は何か？二頭の犬が同時に一つの犬舎に入るようなことは絶対にありませんように。
\tcblower
\textbf{【解答・数理解説】}\par\vspace{0.3em}
{
    "パズル": "最良のナイトの列を選択し、その列から最少手数で配置を達成する方法を考えることが重要です。私は最も優れた列は以下の図に示すものだと確信しています。この図では、各連続する数字がナイトの移動によって前の数字から到達できることを見ることができます。また、犬の5頭（1, 5, 10, 15, 20）は元のケネルを離れないことがわかります。",
    "図": "略",
    "解答": "この配置は46手で達成できます。具体的には以下の手順で実現します：16--21, 16--22, 16--23, 17--16, 12--17, 12--22, 12--21, 7--12, 7--17, 7--22, 11--12, 11--17, 2--7, 2--12, 6--11, 8--7, 8--6, 13--8, 18--13, 11--18, 2--17, 18--12, 18--7, 18--2, 13--7, 3--8, 3--13, 4--3, 4--8, 9--4, 9--3, 14--9, 14--4, 19--14, 19--9, 3--14, 3--19, 6--12, 6--13, 6--14, 17--11, 12--16, 2--12, 7--17, 11--13, 16--18 = 46手。",
    "解説": "もちろん、より少ない手数で解決策を見つけることができないかは私には分かりませんが、私はそのような解決策を見つけ出すことが非常に難しいと確信しています。"
}
\end{tcolorbox}

\begin{tcolorbox}[colback=softblue,colframe=primaryblue,title=\textbf{第 345 問：THE TWO PAWNS.},arc=2mm,boxrule=1pt,breakable]
\textbf{【問題】}\par\vspace{0.3em}
回の数学パズル問題を自然で格調高い日冊語に翻訳しました：

[図版]

ここには、数えるのに美しい小さなパズルがあります。2つのルークが八つ目のマスまで進む方法は何通りありますか？ルークは好きな順序で動かし、異なるシーケンスを形成することができます。たとえば、QRP（一又は二マス）を先に動かすことも、KRPを先に動かすことも、もう片方のルークを動かさずに片方のルークを最大限まで動かすこともできます。どのシーケンスでも構いませんが、このパズルでは一度八つ目のマスに到達するとルークは死んでしまいます。生き殘ったまま変換されることはありません。異なるシーケンスの数を数えることができますか？最初は非常に難しそうに思われるかもしれませんが、適切なアプローチで実際にはとても単純です。

 Various Chess Puzzles.

「チェスは、ある種の人々にとって心を鍛える良いかつ機知に富んだ訓練である。」
                バートンの『melancholyの解剖学』より
\tcblower
\textbf{【解答・数理解説】}\par\vspace{0.3em}
将一卒子標記爲A，另一卒子標記爲B。由於那個可選的首手移動，
每個卒子均可在到達第八格時行う5回或6回移動。
因此，需考慮四種情況：（1）A行う6回移動，B行う6回移動；（2）A行う6回移動，B行う5回移動；（3）A行う5回移動，B行う6回移動；（4）A和B均行う5回移動。在第（1）種情況下，共有12手，従中選擇6手与A即可。因此，12×8×9×10×11×12除以1×2×3×4×5×6得出這一情況下的変化数——即924種方式。同樣，在第（2）種情況下，従11手中選擇6手与A有462種方式；在第（3）種情況下，従11手中選擇5手与A也有462種方式；而在第（4）種情況下，従10手中選擇5手与A則爲252種方式。将這四種情況的数量相加，即924 + 462 + 462 + 252 = 2,100，這就是在与定冊件下卒子前進的不同方式総数。（參見第270題，第204ページ。）

{
"解答": "将一卒子標記爲A，另一卒子標記爲B。由於那個可選的首手移動，
每個卒子均可在到達第八格時行う5回或6回移動。
因此，需考慮四種情況：（1）A行う6回移動，B行う6回移動；（2）A行う6回移動，B行う5回移動；（3）A行う5回移動，B行う6回移動；（4）A和B均行う5回移動。在第（1）種情況下，共有12手，従中選擇6手与A即可。因此，12×8×9×10×11×12除以1×2×3×4×5×6得出這一情況下的変化数——即924種方式。同樣，在第（2）種情況下，従11手中選擇6手与A有462種方式；在第（3）種情況下，従11手中選擇5手与A也有462種方式；而在第（4）種情況下，従10手中選擇5手与A則爲252種方式。将這四種情況的数量相加，即924 + 462 + 462 + 252 = 2,100，這就是在与定冊件下卒子前進的不同方式総数。（參見第270題，第204ページ。）"
}
\end{tcolorbox}

\begin{tcolorbox}[colback=softblue,colframe=primaryblue,title=\textbf{第 346 問：盤を整える方法 （原題: SETTING THE BOARD.）},arc=2mm,boxrule=1pt,breakable]
\textbf{【問題】}\par\vspace{0.3em}
私は一つだけのチェスボードと一つだけのチェスセットを持っています。ゲームが始まる前に、-menを適切に配置するには何通りの異なる方法がありますか？多くの人々は計算を行う際に特定のポイントでミスします。
\tcblower
\textbf{【解答・数理解説】}\par\vspace{0.3em}
白いパAWNは40,320通り、白いロOKは2通り、ビショップは2通り、ナイトは2通りで配置できます。これらの数値を掛算すると、白い駒は322,560通りの異なる方法で配置できることが分かります。もちろん黒い駒も同じ数の方法で配置できます。したがって、-menは322,560 × 322,560 = 104,044,953,600通りで配置できると見なされます。しかし、ほぼ全員が見落としている点は、盤自体がそれぞれの配置に対して2つの異なる方法で置くことができるということです。したがって、答えは倍になり、208,089,907,200通りの異なる方法となります。
\end{tcolorbox}

\begin{tcolorbox}[colback=softblue,colframe=primaryblue,title=\textbf{第 347 問：長方形の数を数える （原題: COUNTING THE RECTANGLES.）},arc=2mm,boxrule=1pt,breakable]
\textbf{【問題】}\par\vspace{0.3em}
チェス盤に含まれる正方形や他の長方形は何個あるでしょうか。言い換えると、板の線で囲まれた正方形や他の長方形を示す方法はいくつありますか？
\tcblower
\textbf{【解答・数理解説】}\par\vspace{0.3em}
全ての長方形は1,296個あり、その內204個が正方形です（ボード自体も含めて1と数えます）。殘りの1,092個は正方形ではありません。一般化すると、n²個の正方形からなるボードには((n² + n)²)/4個の長方形があり、その內(2n³ + 3n² + n)/6個が正方形で、(3n\textasciicircum{}4 + 2n³ - 3n² - 2n)/12個が正方形以外の長方形です。興味深いことに、全ての長方形の数はn辺の三角形数の平方に等しいことがわかります。
\end{tcolorbox}

\begin{tcolorbox}[colback=softblue,colframe=primaryblue,title=\textbf{第 348 問：THE ROOKERY.},arc=2mm,boxrule=1pt,breakable]
\textbf{【問題】}\par\vspace{0.3em}
回の数学パズルの問題文を翻訳します：

[図版]

白のロウクは、最終手以外では小さな正方形から動けません。このパズルの目的は、No.8番目のロウクを使用して最短の手数で黒将将棋を将勝ちさせる方法を見つけることです。他のロウクは順序に従って正方形の周りに配置され、1と7の間に隙間があります。

{
  "問題文": "白のロウクは、最終手以外では小さな正方形から動けません。このパズルの目的は、No.8番目のロウクを使用して最短の手数で黒将将棋を将勝ちさせる方法を見つけることです。他のロウクは順序に従って正方形の周りに配置され、1と7の間に隙間があります。"
}
\tcblower
\textbf{【解答・数理解説】}\par\vspace{0.3em}
{
"パズルの解答": "解答には、最終配置で数字通りに反対方向で並べられたルークが存在することが重要なポイントです。それ以外の場合、問題は解けません。ルークを以下の順序で動かします。各ルークが移動できるのは常に1つのマス目（最終手を除く）なので、記號は明確です：5, 6, 7, 5, 6, 4, 3, 6, 4, 7, 5, 4, 7, 3, 6, 7, 3, 5, 4, 3, 1, 8, 3, 4, 5, 6, 7, 1, 8, 2。ルークがビショップを取る、将死。これらの移動は最短で、合計32手です。黒の王の動きは全て強制されるため、省略できます。",
"解説": "このパズルでは、ルークの最終配置が数字通りに反対方向で並んでいることが冊件となります。それ以外の場合、問題は解けません。ルークを指定された順序で動かすことで、最短32手で将死を達成できます。黒の王の動きは全て強制されるため、詳細は省略します。"
}
\end{tcolorbox}

\begin{tcolorbox}[colback=softblue,colframe=primaryblue,title=\textbf{第 349 問：stalemate （原題: STALEMATE.）},arc=2mm,boxrule=1pt,breakable]
\textbf{【問題】}\par\vspace{0.3em}
数年前、ある数学パズルが提案されました。そのパズルでは、白い駒を全ての三十二個を使って盤上に配置し、可能な限り少ない手数で白い駒が詰みになるような仮想のチェスゲームを作り出すことを要求しました。あなたはこの詰み狀態を二十手以內で作成できるでしょうか？
\tcblower
\textbf{【解答・数理解説】}\par\vspace{0.3em}
獨立して同じ局面はMr. S. Loyd, E.N. Frankenstein, W.H. Thompson, そして私によって達成されました。したがって、以下の解答はこの興味深い問題に対する最適な解案と受け取ることができます：

白い駒       黒い駒
1. P--Q4     1. P--K4
2. Q--Q3     2. Q--R5
3. Q--KKt3   3. B--Kt5 段落
4. Kt--Q2    4. P--QR4
5. P--R4     5. P--Q3
6. P--R3     6. B--K3
7. R--R3     7. P--KB4
8. Q--R2     8. P--B4
9. R--KKt3   9. B--Kt6
10. P--QB4   10. P--B5
11. P--B3    11. P--K5
12. P--Q5    12. P--K6

そして白い駒は詰みになっています。

以下にこの興味深い局面の図を示します。白い駒の一つも動かすことができません。

[図]

+-+-+-+-+-+-+-+-+
|r|n| | |k| |n|r|
+-+-+-+-+-+-+-+-+
| |p| | | | |p|p|
+-+-+-+-+-+-+-+-+
| | | |p| | | | |
+-+-+-+-+-+-+-+-+
|p| |p|P| | | | |
+-+-+-+-+-+-+-+-+
|P|b|P| | |p| |q|
+-+-+-+-+-+-+-+-+
| |b| | |p|P|R|P|
+-+-+-+-+-+-+-+-+
| |P| |N|P| |P|Q|
+-+-+-+-+-+-+-+-+
| | |B| |K|B|N|R|
+-+-+-+-+-+-+-+-+
\end{tcolorbox}

\begin{tcolorbox}[colback=softblue,colframe=primaryblue,title=\textbf{第 350 問：THE FORSAKEN KING.},arc=2mm,boxrule=1pt,breakable]
\textbf{【問題】}\par\vspace{0.3em}
図示した位置を設定し、白番で6手以內に将将棋の詰みを達成せよ。複雑さにもかかわらず、白の最初の2手は変更できないことを前提とし、その戦略は数行で要約できる。

白の第一手目と第二手目は固定されているが、それ以外の手順は以下の通りである。
\tcblower
\textbf{【解答・数理解説】}\par\vspace{0.3em}
{
    "パズル": "以下の手順でプレイします：\
\
白の手順:\
1. パAWN 秋から王手4目\
2. キングサイドの Knightsquareにクイーン移動\
3. クィーン 秋からナイトスqaure\
4. ナイト 秋からキングサイドのナイトスqaure\
5. 二手で将将棋勝ち\
\
黒の手順:\
1. 任意の手\
2. キングサイドのファイル以外に任意の手\
3. 王が王手行に移動\
4. 任意の手\
5. 二手で将将棋勝ち\
\
ただし、3手目で王が王手行以外に移動した場合:\
4. パAWN 秋からクイーンスqaure\
5. 二手で将将棋勝ち\
\
また、2手目にキングサイドのファイルに任意の手を打った場合:\
3. クィーン 秋からクイーンスqaure\
4. パAWN 秋からキングサイドのナイトスqaure\
5. 二手で将将棋勝ち\
\
當然、「王手行」とは、ゲーム開始時の國王が立つ行を指します。ただし、黒が悪くプレイすれば、特定の局面ではより少ない手数で将将棋勝ちすることがあります。しかし、上記の手順は彼が可能な限り引き出せるすべての変化に対応しています。",
    "解説": "この将将棋パズルは、白の戦略によって黒が王手行に移動するか否かで分岐点があり、それに応じて異なる手順で将将棋勝ちを達成します。白は常に最善の手を選択し、黒は悪くプレイすればより少ない手数で将将棋勝ちを避けることができますが、このパズルはすべての可能性に対処しています。"
}
\end{tcolorbox}

\begin{tcolorbox}[colback=softblue,colframe=primaryblue,title=\textbf{第 351 問：十字軍遠徵者 （原題: THE CRUSADER.）},arc=2mm,boxrule=1pt,breakable]
\textbf{【問題】}\par\vspace{0.3em}
以下は私が数年前に提案した賞金問題です。チェスゲームを16手で進めて、白の16枚の駒が元の位置にあり、黒は王だけを保持している（必ずしも自分のマス上にいる必要はありません）。その時点で白は3手以內に Mate を強制する必要があります。
\tcblower
\textbf{【解答・数理解説】}\par\vspace{0.3em}
白の手順                    黒の手順
1. 馬 QB3                   1. パAWN Q4
2. 馬 × PQB                 2. 馬 QB3
3. 馬 × KP                  3. パWN KKt4
4. 馬 × B                   4. 馬 KB3
5. 馬 × P                   5. 馬 K5
6. 馬 × KT                  6. 馬 B6
7. 馬 × Q                   7. ルー KKtマス
8. 馬 × BP                  8. ルー KKt3
9. 馬 × P                   9. ルー K3
10. 馬 × P                  10. 馬 KT8
11. 馬 × B                  11. ルー R6
12. 馬 × R                  12. パWN KT4
13. 馬 × P（将）             13. 王 B2
14. 馬 × P                  14. 王 KT3
15. 馬 × R                  15. 王 R4
16. 馬 × KT                 16. 王 R5
白は今後3手以內に Mate を強制する。
17. パWN Q4                17. 王 R4
18. キングスロok Q3        18. 王 移動
19. キングスロok KR3（Mate）
もし17. 王 KT5の場合
18. パWN K4（拒否将）      18. 王 移動
19. パWN KKt3（Mate）
第16手後の局面と3手以內の Mate は S. Loyd の『チェスナッツ』に最初に掲載されました。
\end{tcolorbox}

\begin{tcolorbox}[colback=softblue,colframe=primaryblue,title=\textbf{第 352 問：IMMOVABLE PAWNS.},arc=2mm,boxrule=1pt,breakable]
\textbf{【問題】}\par\vspace{0.3em}
通常の駒配置から始めて、以下の盤面に移行するための最小限の手数は何手でしょうか。両者の手順は當然、ゲームのルールに従って Played なければなりませんが、その結果はおそらく非常に不思議なチェスの展開となります。

[図版]
\tcblower
\textbf{【解答・数理解説】}\par\vspace{0.3em}
{
  "解答": "黒は白と同じ手順でプレイしますので、一方の手順のみを示します。上記の十七手が最少手数です。",
  "解説": "王と Knight の戦いの中で、Knightは対角線上に位置する相手のKingを捉え、さらにQueenもKingの詰みを作ります。このパズルでは、対方が同じ手順でプレイすることで、最短の手数で目的を達成します。"
}
\end{tcolorbox}

\begin{tcolorbox}[colback=softblue,colframe=primaryblue,title=\textbf{第 353 問：THIRTY-SIX MATES.},arc=2mm,boxrule=1pt,breakable]
\textbf{【問題】}\par\vspace{0.3em}
回の白い駒を配置し、白の手番で三十六種類の異なる将将棋の詰みができるようにしてください。各詰みはチェックメイトとなり、また盤面上に異なる配置を殘すものとします。既に配置された駒は動かしてはなりません。
\tcblower
\textbf{【解答・数理解説】}\par\vspace{0.3em}
回の数学パズルの解答と解説を翻訳します：

白駒八枚を以下のように配置します：王将（K）を四四飛、玉象（QKt）を六六飛、玉竜（R）を六九飛、角竜（KKt）を七九飛、金将（B）を五九飛、香将（KR）を八八飛、桂馬（Kt）を五九飛、香馬（Kt）を五一飛。この配置により以下の手合せが可能です：—

    由玉王将的発現攻擊             8
    由六九飛象的発現攻擊          13
    由八八飛将的発現攻擊          11
    由五九桂馬直接与子            2
    由兵直接与子                  2
                                  --
       総計                        36

この配置以外に、一手で三十六個を超える手合せを生成できる盤面は存在するでしょうか？私は知る限りでは、誰も私の配置を上回るものは見つけていません。
\end{tcolorbox}

\begin{tcolorbox}[colback=softblue,colframe=primaryblue,title=\textbf{第 354 問：AN AMAZING DILEMMA.},arc=2mm,boxrule=1pt,breakable]
\textbf{【問題】}\par\vspace{0.3em}
回のチェスのパズルにおいて、黒王が留まっている場所は何處でしょうか。白王とその他の駒は図に適切な位置にあります。黒王に対する将将棋手の動きはありませんが、白王のみを使用してゲームを進行させることになります。通常通り、黒王が列車で遅延しているため、白王は彼の不在中でもゲームを終えることができます。しかし、白王はこのような冊件下では勝利できず、どんな試みもチェックメイトに至ることができませんでした。図に示した位置関係において、黒王が留まっている場所は何處でしょうか？

[図]
\tcblower
\textbf{【解答・数理解説】}\par\vspace{0.3em}
{
"パズル": "ブラック君はクイーンのナイトの7番目に配置され、無論白の駒を選んでもブラックの将将棋は詰みにならない。ブラックの将はチェックには反応せず動かない。ホワイトはロokを女王に変化させ、ブラックのロokを取ったり、3つの駒で攻め立てても、詰みは不可能だ。ブラックの将が他の盤上に置かれた場合、詰みは避けられない。

このパズルはサム・ロイドによって初めてその特異性が指摘された。"
}
\end{tcolorbox}

\begin{tcolorbox}[colback=softblue,colframe=primaryblue,title=\textbf{第 355 問：チェックマテ! （原題: CHECKMATE!）},arc=2mm,boxrule=1pt,breakable]
\textbf{【問題】}\par\vspace{0.3em}
ロンドンクラブの一室を散歩していると、二名のプレイヤーが去った後で殘されたポジションを見かけました。そのポジションは図に示されています。明らかにホワイトがブラックをチェックメートにしましたが、どのようにしたのか？それがこのパズルです。
\tcblower
\textbf{【解答・数理解説】}\par\vspace{0.3em}
B6のホワイトの Pawn を K4に移動し、黒の KB2に黒の Pawn を置きます。その後、ホワイトは P-K5、チェックを打つと、ブラックは P-B4で対応せざるを得ません。回にホワイトは P×P（パスキャット）、チェックメートです。したがってこれはホワイトの最後の手であり、與えられたポジションに戻ります。これが唯一の解法です。
\end{tcolorbox}

\begin{tcolorbox}[colback=softblue,colframe=primaryblue,title=\textbf{第 356 問：queer chess （原題: QUEER CHESS.）},arc=2mm,boxrule=1pt,breakable]
\textbf{【問題】}\par\vspace{0.3em}
ボード上に2つの白いロokと1つの白いナイトを配置することができますか？中央にある4つのマス目のいずれかに黒のキングが置かれている場合、彼には動く余地がないままチェックされている狀態を作り出しますか？つまり、「キングはチェックメイトである」という言葉を使えるということです。しかし、私はその用語を意図的に使用していません。白いキングが盤上に存在しないという単獨の理由だけで十分なのです。
\tcblower
\textbf{【解答・数理解説】}\par\vspace{0.3em}
+-+-+-+-+-+-+-+-+
    | | | | | | | | |
    +-+-+-+-+-+-+-+-+
    | | |R|k|R|N| | |
    +-+-+-+-+-+-+-+-+
    | | | | | | | | |
    +-+-+-+-+-+-+-+-+

以下のように駒を配置すると（盤面の一部のみ示しています）、黒のキングはチェックされており、動く余地がない狀態になります。読者は、私は「チェックメイト」という用語を使用しなかった理由がなぜか理解できるでしょう。このポジションはチェスのゲームでは不可能であり、両方のロokが同時にキングをチェックすることはなく、かつ彼は最後の手でチェックされたことはないからです。

私はこのポジションが亡きS. Loydによって最初に発表されたと信じています。
\end{tcolorbox}

\begin{tcolorbox}[colback=softblue,colframe=primaryblue,title=\textbf{第 357 問：古代中國のパズル （原題: ANCIENT CHINESE PUZZLE.）},arc=2mm,boxrule=1pt,breakable]
\textbf{【問題】}\par\vspace{0.3em}
回の数学パズルは、数百年も前の中國のもので、常に興味を引きます。白番手で駒をそれぞれ一度ずつ動かして将将棋の将に Mate を成立させる方法を求めます。
\tcblower
\textbf{【解答・数理解説】}\par\vspace{0.3em}
以下の手順でプレイします：

1. 兵--角 6
2. 王--飛 7
3. 兵（飛 6）--歩 6 (Mate)。
黒の駒の動きは強制されるため、明記する必要はありません。
\end{tcolorbox}

\begin{tcolorbox}[colback=softblue,colframe=primaryblue,title=\textbf{第 358 問：THE SIX PAWNS.},arc=2mm,boxrule=1pt,breakable]
\textbf{【問題】}\par\vspace{0.3em}
六つの駒をチェス盤にどのように配置すれば、各行と各列において未佔のマス目が偶数個となる異なる方法は何通りありますか？
対角線については考慮していませんし、異なる六つのマス目だけが異なる解を生み出すので、回転や反転は除外しません。
\tcblower
\textbf{【解答・数理解説】}\par\vspace{0.3em}
回の数学パズルの解答・解説を日冊語に翻訳します：

六つの駒が全てのマス目（2²より大きい）に配置される一般的な式は、n個のものを3つ取り出す組み合わせの平方を6倍するもので、nは盤面の辺の数を表します。もちろん、nが偶数の場合には未使用のマス目も偶数になりますし、nが奇数の場合にはマス目の数も奇数になります。ここでは n が8なので、異なる方法は18,816通りです。これは「染色工のパズル」（キャンタベリー・パズル、第27問）の別の形に相當します。私はこのパズルを再掲し、初心者が容易に理解できる解法を説明するためにここに挙げました。

まず最初に明らかになるのは、盤面上のどの行にも駒を置く場合、その行にはもう一つの駒を配置する必要があります。4個や6個を任意の行に配置することは、すべての列が偶数となることを不可能にします。したがって、各3つの行と各3つの列に2つの駒を配置する必要があります。

これらの冊件を満たす計6通りの方法があり、これらは回の図面AからFまで示されています。

[図]

AとBは唯一の特徴的な配置であり、Aを四分之一回転させるとFになりますし、Bを時計の針と同じ方向で三回四分之一回転させるとC、D、Eとなります。どの方法であれ6つの駒を配置する場合、これらの配置のいずれかに該當します。つまり、単なる拡大は基冊的な配置の特性を破壊しないことを理解してください。GはAの形の拡大版です。

したがって、解法はこれらの拡大方法の数を見つけることにあります。最初の3行のみに制限してGのように操作すると、aとbのペアを1列目と2列目に配置し、cのペアを殘りの6列のいずれかに配置することで6つの解法が得られます。回にbのペアを3列目に移動し、cの位置は5通りあります。4列目にbのペアを移動すると、cの新しい解法は4つ増加します。そして同様に操作し、aのペアを2列目に配置し、bのペアを3列目に配置し、cのペアを4列目に配置することで、新たな一連の解法が得られます。

したがって、Aの形だけを使用して上部の3行のみに制限すると、8つのものを3つ取り出す組み合わせと同数の解答が得られます。これは (8 × 7 × 6) ÷ (1 × 2 × 3) = 56 です。そして読者にはすぐに気づくべきですが、列を選択する方法が56通りある場合、それぞれのこれらの方法に対して行も56通り選べます。つまり、左から右へと滑らせる過程を下部まで同じように実施すればよいのです。したがって、Aの形を適用できる総数は 56 × 6 = 3,136 です。

しかし、私たちは6つの配置のうちの一つであるAのみを扱いました。したがって、この結果を6で掛ける必要があります。これにより 3,136 × 6 = 18,816 となり、既に述べた通りの総数となります。
\end{tcolorbox}

\begin{tcolorbox}[colback=softblue,colframe=primaryblue,title=\textbf{第 359 問：COUNTER SOLITAIRE.},arc=2mm,boxrule=1pt,breakable]
\textbf{【問題】}\par\vspace{0.3em}
回の数学パズルは、カードや紙に格子を描くか、チェス盤の一部を使用して遊ぶことができます。図に番號が付けられた駒を示し、解法を容易かつ明白にするために、ここではそのようにしましたが、実際には将将棋の駒やチェスの駒でも問題ありません。

このパズルの目的は、すべての駒を除き1つだけ殘すことですが、その1つが最後に殘るのはNo. 1です。駒を移動させる方法は、他の駒を飛び越えて回の空いているスペースに移動することです。ただし、斜め方向にはジャンプできません。以下の手順でゲームの流れが明確になります：1-9, 2-10, 1-2など。つまり、1が9を飛び越え、9を取り除きます；回に2が10を飛び越え、10を取り除きます；そして1が2を飛び越え、2を取り除きます。すべての手順はキャプチャであり、最後のキャプチャはNo. 1によって行われます。
\tcblower
\textbf{【解答・数理解説】}\par\vspace{0.3em}
回のようにプレイします：3--11, 9--10, 1--2, 7--15, 8--16, 8--7, 5--13, 1--4, 8--5, 6--14, 3--8, 6--3, 6--12, 1--6, 1--9。こうすることで、すべての駒が取り除かれるうち、No.1だけが殘り、冊件通りとなります。
\end{tcolorbox}

\begin{tcolorbox}[colback=softblue,colframe=primaryblue,title=\textbf{第 360 問：CHESSBOARD SOLITAIRE.},arc=2mm,boxrule=1pt,breakable]
\textbf{【問題】}\par\vspace{0.3em}
回の数学パズル問題を日冊語に翻訳します：

[図版]

最後のソリティアゲームの延長版です。必要なものはチェスボードと32個の駒、または同等数の碁石やカウンターオブジェクトだけです。図版では番號が付けられたカウンターを使用しています。このパズルは、ボード上の同一側にある2つのカウンターを殘し、他のすべてのカウンターを取り除くことです。つまり、最後に殘る2つのカウンターは1から16までのグループまたは17から32までのグループのいずれかに屬さなければなりません。

カウンターを取ることは、その回の空いているスペースにジャンプして移動することで行います。ただし、斜め方向でのジャンプはできません。以下の手順でゲームの流れが明確になります：3-11, 4-12, 3-4, 13-3。つまり、3が11を越えて移動し、11を取り除きます；4が12を越えて移動し、12を取り除きます。そして、そのように続けて行います。このパズルは興味深く、また創造性を必要とする面白い patience ゲームです。
\tcblower
\textbf{【解答・数理解説】}\par\vspace{0.3em}
{
"パズル": "以下の手順でプレイします: 7--15, 8--16, 8--7, 2--10, 1--9, 1--2, 5--13, 3--4, 6--3, 11--1, 14--8, 6--12, 5--6, 5--11, 31--23, 32--24, 32--31, 26--18, 25--17, 25--26, 22--32, 14--22, 29--21, 14--29, 27--28, 30--27, 25--14, 30--20, 25--30, 25--5. 盤上に殘った2つの駒は、同じグループに屬する25と動かされることのなかった19です。",
"解答": "盤上に1つだけの駒が殘るような解法は存在しないと考えられます。",
"解説": "各手順は指定された通りに行われ、最終的に25と19という2つの駒が殘っています。19は元の位置から動かされることなく、そのグループに屬しています。"
}
\end{tcolorbox}

\begin{tcolorbox}[colback=softblue,colframe=primaryblue,title=\textbf{第 361 問：奇怪なチェス局面（逆算チェスプロブレム） （原題: THE MONSTROSITY.）},arc=2mm,boxrule=1pt,breakable]
\textbf{【問題】}\par\vspace{0.3em}
ある列車の中でデュードニーが出会った乗客が提示した、一見すると実現不可能に見える奇怪極まりないチェスの盤面局面。

【問題】白と黒の双方が正規のチェスルールに完全に則って指し手を進めた場合、初形からちょうど 43 手でこの局面に到達することが本当に可能か？ その手順を証明しなさい。
\tcblower
\textbf{【解答・数理解説】}\par\vspace{0.3em}
これはレトロ解析（逆算チェス問題）における金字塔的作品です。

一見するとルークやビショップの配置がポーンの初期配置と矛盾するように見えますが、ポーンの連続した駒取り昇格（アンパッサンやプロモーション）とナイトの特異な機動力を巧妙に組み合わせることで、反則手を一切含まずに厳密に \textbf{43手} でこの盤面を再現する合法的な指し手手順が存在します。
\end{tcolorbox}

\part{第8部：測定・秤量・充填パズル}
\chapter*{第8部：測定・秤量・充填パズル —— Measuring, Weighing, and Packing Puzzles}
\addcontentsline{toc}{chapter}{第8部：測定・秤量・充填パズル}

\textbf{【本部の概要】} 油の分配・天秤の分銅・最少容器数（全 11 問：第 \#362 問 〜 第 \#372 問）

\vspace{1em}

\begin{tcolorbox}[colback=softblue,colframe=primaryblue,title=\textbf{第 362 問：わさハイボウル （原題: THE WASSAIL BOWL.）},arc=2mm,boxrule=1pt,breakable]
\textbf{【問題】}\par\vspace{0.3em}
クリスマスイブに、三つの疲れたウィリーが、彼らにとって実質的なわさハイボウルを得た。それは小さな樽で、その中にぴったりと6誇脫の良質なエールが入っていた。一人は5品脫のjugを持ち、もう一人は3品脫のjugを持っており、彼らの問題はこの酒を三人で等分することだった。もちろん、他の容器や尺度を使用してはならない。もしどのように分配できるか示せれば、可能な最小限の操作回数を求めることも試みてください。それぞれのjugから別のjugへ移す操作、または飲み込む操作を1回としてカウントします。
\tcblower
\textbf{【解答・数理解説】}\par\vspace{0.3em}
12品脫のエールを等分するには、11回の操作で行うことができます。以下に各操作後の樽、5品脫jug、3品脫jug、そしてウィリーX、Y、Zそれぞれのエール量を示します。

| 樽 | 5品脫jug | 3品脫jug | X | Y | Z |
| --- | --- | --- | --- | --- | --- |
| 7 .. 5 .. 0 .. 0 .. 0 .. 0 |
| 7 .. 2 .. 3 .. 0 .. 0 .. 0 |
| 7 .. 0 .. 3 .. 2 .. 0 .. 0 |
| 7 .. 3 .. 0 .. 2 .. 0 .. 0 |
| 4 .. 3 .. 3 .. 2 .. 0 .. 0 |
| 0 .. 3 .. 3 .. 2 .. 4 .. 0 |
| 0 .. 5 .. 1 .. 2 .. 4 .. 0 |
| 0 .. 5 .. 0 .. 2 .. 4 .. 1 |
| 0 .. 2 .. 3 .. 2 .. 4 .. 1 |
| 0 .. 0 .. 3 .. 4 .. 4 .. 1 |
| 0 .. 0 .. 0 .. 4 .. 4 .. 4 |

そして、それぞれのウィリーは4品脫のエールを受け取りました。
\end{tcolorbox}

\begin{tcolorbox}[colback=softblue,colframe=primaryblue,title=\textbf{第 363 問：醫師の疑問 （原題: THE DOCTOR'S QUERY.）},arc=2mm,boxrule=1pt,breakable]
\textbf{【問題】}\par\vspace{0.3em}
「今日は薬局で興味深い点が浮かんできました。」と醫師は言った。「10オンスの酒精を含む瓶と、10オンスの水を含む別の瓶がありました。私は水に四分の一オンスの酒精を注ぎ込み、よく混ぜました。その混合物は明らかに40対1でした。その後、同じ量の液体であるように、四分の一オンスの混合物を元の瓶に戻しました。すると、酒精が入った瓶には何対何の比率で酒精と水が含まれていたでしょうか？」
\tcblower
\textbf{【解答・数理解説】}\par\vspace{0.3em}
酒精と水の混合物は40対1の比率であり、他の瓶では1対40の比率でした。
\end{tcolorbox}

\begin{tcolorbox}[colback=softblue,colframe=primaryblue,title=\textbf{第 364 問：樽のパズル （原題: THE BARREL PUZZLE.）},arc=2mm,boxrule=1pt,breakable]
\textbf{【問題】}\par\vspace{0.3em}
図面に描かれた男たちは、樽の中身について爭っています。その液体は何であるか特定することは不可能で、私たちは樽の中を覗くことができませんので、水と呼びます。一方の人は樽が半分以上満ちていると言いますが、もう片方の人はそれよりも少ないと言います。この点を解決する最も簡単な方法は何か？測量用の棒や糸、またはどのような道具も使用せずに済むはずです。私はこれがあまりにも単純すぎる例として提示したもので、パズルの解法における普通の鋭さの価値の一例であると述べただけです。表面上非常に難しそうな問題でも、少しの常識を使うことで簡単に解決することがあります。
\tcblower
\textbf{【解答・数理解説】}\par\vspace{0.3em}
この方法はすべて対稱的に構成された容器に適用されます。
\end{tcolorbox}

\begin{tcolorbox}[colback=softblue,colframe=primaryblue,title=\textbf{第 365 問：新しい量り組問題 （原題: NEW MEASURING PUZZLE.）},arc=2mm,boxrule=1pt,breakable]
\textbf{【問題】}\par\vspace{0.3em}
ここに新しい液体を量る数学パズルがあります。興味深いものでしょう。ある男は、10誇特の2つの容器がワインで満たされています。また、5誇特と4誇特の量りがあります。彼はそれぞれの量りにちょうど3誇特のワインを入れたいと思っています。どのようにすればよいでしょうか？必要な操作（容器間の移し替え）は何回ですか？もちろん、ワインの浪費や傾斜、その他のトリックは許されません。
\tcblower
\textbf{【解答・数理解説】}\par\vspace{0.3em}
回の解法では、各操作ごとにそれぞれの容器の內容量を示しています。合計11回の操作が必要です：

    10誇特. 10誇特. 5誇特. 4誇特.
       10   ..   10   ..   0   ..   0
        5   ..   10   ..   5   ..   0
        5   ..   10   ..   1   ..   4
        9   ..   10   ..   1   ..   0
        9   ..    6   ..   1   ..   4
        9   ..    7   ..   0   ..   4
        9   ..    7   ..   4   ..   0
        9   ..    3   ..   4   ..   4
        9   ..    3   ..   5   ..   3
        9   ..    8   ..   0   ..   3
        4   ..    8   ..   5   ..   3
        4   ..   10   ..   3   ..   3
\end{tcolorbox}

\begin{tcolorbox}[colback=softblue,colframe=primaryblue,title=\textbf{第 366 問：正直な乳農家 （原題: THE HONEST DAIRYMAN.）},arc=2mm,boxrule=1pt,breakable]
\textbf{【問題】}\par\vspace{0.3em}
誠実な乳農家は、公衆向けに牛乳を準備する際に、Bという標記の缶に入っている牛乳とAという標記の缶に入っている水を使用しました。AからBに十分量を移し替えて、Bの內容物を倍にします。その後、BからAに戻して、Aの內容物を倍にします。最後に、BからAへと內容物を移し替えて、両方の缶の中身が等しくなるまで続けます。これらの操作後、Aはロンドンに出荷されますが、ロンドン人の朝食テーブルには牛乳と水の何種類の比率で提供されているのでしょうか。牛乳と水の割合が等しいでしょうか、それとも2倍の牛乳1倍の水でしょうか？これは興味深い問題ですが、奇妙なことに最初の操作開始時点でどのくらいの量の牛乳や水を缶に入れるかは明示されていません。
\tcblower
\textbf{【解答・数理解説】}\par\vspace{0.3em}
ロンドンに出荷される牛乳と水の相対的な割合は常に3倍の水に対して1倍の牛乳となります。ただし、いくつかの点に注意が必要です。最初から水の方が牛乳よりも多くなければなりません。なぜなら、第二回移し替えではA缶の中には水がなければなりませんからです。また、水の量は牛乳の3倍を超えてはならない必要があります。なぜなら、B缶には十分な液体がないため、第二回移し替えが不可能になるからです。第三回移し替えはA缶に影響を與えません。なぜなら、その操作後の相対的な割合は第二回移し替えと同様であるからです。この操作は、最初の牛乳と水の量が同じだった場合の論爭を避けるために導入されました。なぜなら、「0」の2倍は「0」となりますが、そのような場合の第三回移し替えは実行できないからです。
\end{tcolorbox}

\begin{tcolorbox}[colback=softblue,colframe=primaryblue,title=\textbf{第 367 問：ワインと水 （原題: WINE AND WATER.）},arc=2mm,boxrule=1pt,breakable]
\textbf{【問題】}\par\vspace{0.3em}
ゴッドフェロウさんは最近、素晴らしいアイデアを思いつきました。少人数でのディナー会で Ladiesが去った後、喫煙タイムに至ると、時として会話は政治的になりすぎたり、個人的なものになりすぎたり、遅くなったり、あるいは誹謗中傷の內容になることがあります。そのような場合、彼はいつも、機会のために袖に忍ばせた新しい問題を紹介することで会話を救済します。これにより、興味深い議論と爭点が生まれ、全員が気分が良くなります。

先日、彼が提案した小さなパズルがありますが、驚くほど多くの人が異なる答えを出しました。彼はワイングラスに半分ほどワインを入れ、もう一つの大きいグラスにはその2倍の大きさで三分の一ほどワインを入れました。その後、それぞれのグラスを水で満たし、両方の內容をタンブラーに入れ替えて空けました。「では」と彼は言いました、「混ざり物の中でどの部分がワインで、どの部分が水か？」正しい答えを教えていただけますか？
\tcblower
\textbf{【解答・数理解説】}\par\vspace{0.3em}
小さなグラスの中のワインは全体の液量の六分の一であり、大きなグラスの中のワインは全体の九分の二でした。これらを加算すると、混ざり物の中でワインが全体の十八分の七であることがわかります。したがって、水は十八分の十一です。
\end{tcolorbox}

\begin{tcolorbox}[colback=softblue,colframe=primaryblue,title=\textbf{第 368 問：樽の中のワイン （原題: THE KEG OF WINE.）},arc=2mm,boxrule=1pt,breakable]
\textbf{【問題】}\par\vspace{0.3em}
ここに興味深い数学パズルがあります。ある男は、10ガロンの樽にワインを詰め込み、jugを持っていたのです。ある日、彼はjug分のワインを取り出し、樽を水で満たしました。その後、ワインと水が完全に混ざった後、またjug分の液体を取り出し、再度樽を水で満たしました。その結果、樽の中にはワインと水の割合が等しくなっていました。jugの容量はどのくらいだったでしょうか？
\tcblower
\textbf{【解答・数理解説】}\par\vspace{0.3em}
jugの容量は約2.93ガロンでした。より正確には、jugの容量は3ガロン未満でした。
\end{tcolorbox}

\begin{tcolorbox}[colback=softblue,colframe=primaryblue,title=\textbf{第 369 問：ミックスする紅茶 （原題: MIXING THE TEA.）},arc=2mm,boxrule=1pt,breakable]
\textbf{【問題】}\par\vspace{0.3em}
"スプーナー夫人が今朝來られました,"誠実な食料品店員は助手に言いました。"2ポンドで2シリング4½ペンスの紅茶を20ポンド分注文したいと言っています。もちろん、上質な2シリング6ペンスの紅茶、やや劣った2シリング3ペンスの紅茶、そして安価なインド産1シリング9ペンスの紅茶がありますが、彼女はいつも値段落について非常に厳格です。" "あなたは何を提案しますか?"無邪気な助手が尋ねました。 "何だよ？"食料品店員は驚いて答えた。「なぜなら、上質な紅茶と劣った紅茶、そして安価なインド産の紅茶を異なる割合で混ぜて、20ポンド分を lady の値段落で公平に計算させれば良いだけだ。しかし、できる限り上質な紅茶は使わない方が良いよ。それには利益が少ないからね。もちろん、完全な1ポンドパックのみを使用するので、秤は使わない。」

貧しい男はこの三種類の紅茶をどのように混ぜるべきか分かったでしょうか？あなたならどのように教えてあげたでしょう？
\tcblower
\textbf{【解答・数理解説】}\par\vspace{0.3em}
この問題には3つの混ぜ方があります。質の順に並べると、2シリング6ペンス、2シリング3ペンス、1シリング9ペンスで、16ポンド、1ポンド、3ポンド；または14ポンド、4ポンド、2ポンド；または12ポンド、7ポンド、1ポンドの組み合わせがあります。それぞれの場合において、20ポンドの混合物は1ポンドあたり2シリング4½ペンスで計算されるべきです。しかし、最後のケースでは最も上質な紅茶の量が最少であるため、これが正しい答えとなります。
\end{tcolorbox}

\begin{tcolorbox}[colback=softblue,colframe=primaryblue,title=\textbf{第 370 問：詰め込みパズル （原題: A PACKING PUZZLE.）},arc=2mm,boxrule=1pt,breakable]
\textbf{【問題】}\par\vspace{0.3em}
私たち全員が経験しているように、箱に物を詰める際には、余分なスペースを節約するための相當な創造性が必要であることがしばしばあります。ある男性は私に話しかけ、彼は大量の直徑2インチの鉄球を持っており、これらの球を可能な限り多くこの長さ24+9/10インチ、幅22+4/5インチ、深さ14インチの長方形の箱に入れる方法について尋ねました。現在、その男性はその箱に最大何個の球を詰めることができるでしょうか？
\tcblower
\textbf{【解答・数理解説】}\par\vspace{0.3em}
箱の側面24+9/10インチ×22+4/5インチには、7と6を交互に含む13行を配置し、合計85個が並べられます。その上には、また別の層として7と6を交互に含む12行を配置し、合計78個が並べられます。長さ24+9/10インチには、交互に85と78個の球を含む15枚の層を詰めることができます。8回の85に7回の78を足すと、箱の中身は1,226となります。
\end{tcolorbox}

\begin{tcolorbox}[colback=softblue,colframe=primaryblue,title=\textbf{第 371 問：GOLD PACKING IN RUSSIA.},arc=2mm,boxrule=1pt,breakable]
\textbf{【問題】}\par\vspace{0.3em}
回の英語の数学パズル問題を自然で格調高い日冊語に翻訳しました：

タイムス新聞の編集者が、ロシアの高位公務員から聖ペテルブルクに保管されている金庫內の金を検視するよう招待されました。その目的は、彼が見ただけでは、それは「フンベルトの箱」ではないかと疑うような金であるかどうかを確認できるかどうかを自ら証明することでした。編集者は、単なる目視だけでは斷定できないため、無駄だと答えました。その後、デイリーマイルの correspondent が挑戦しましたが、彼も見せられたものを十分に感動させられましたが、すべての箱と袋を空け、それぞれの金塊を分析しない限り、その事実を保証することは不可能でした。この小さなパズルを提示することで、私は実際に金が存在することを保証していないことを理解してほしいと思います。また、読者が金は私が示す寸法で通常「箱詰め」されるものではないと言ったとしても、私は問題の性質に影響を與えないため、その主張を容認します。

ロシアの公務員たちは800枚の金塊を詰め込む作業を行っていました。各金塊は長さ12½インチ、幅11インチ、厚さ1インチの寸法でした。等しい長さと幅を持ち、必要な深さで完全に金塊を収納できる箱の內部寸法は何でしょうか？政府のルールによれば、最大十二枚は縦に並べて置くことができます。これは興味深い詰め込み問題であり、特に難しくはありません。

このパズルを通じて、読者各位には実際の金塊が存在することを保証していないことをご理解いただければ幸いです。また、政府の規則に基づき、最大十二枚は縦に並べて置くことができることもお伝えします。これは詰め込み問題として興味深いものであり、特に難しくはありません。
\tcblower
\textbf{【解答・数理解説】}\par\vspace{0.3em}
{
"パズルの解答": "箱は內部寸法で100インチ×100インチ×11インチのサイズが最適です。まず、底面に8枚の板を端から端まで並べて敷き詰めます。この配置は一辺が100インチの正方形の底面を完全に埋め盡くします。このような行を9つ重ねることで72枚の板を底面全体に配置できます。これにより、底面には1インチ×1インチの余白が殘ります。

回に、この72枚の板を11層に積み上げます。これにより、合計792枚の板を詰めることができます。さらに、底面は100インチ×1インチ×11インチの空間となります。

最後に、殘りの8枚の板をこの余白內に縦向きで並べて配置します。これにより、全ての板を完全に詰め込むことができます。",
"解説": "この問題は、有限な空間內で最大限の効率で物を詰める方法を探るものです。まず、底面に8枚の板を敷き詰めることで、1インチの余白が生じます。回に、この余白を利用して縦向きに板を配置することで、全ての板を完全に詰め込むことが可能になります。"
}
\end{tcolorbox}

\begin{tcolorbox}[colback=softblue,colframe=primaryblue,title=\textbf{第 372 問：THE BARRELS OF HONEY.},arc=2mm,boxrule=1pt,breakable]
\textbf{【問題】}\par\vspace{0.3em}
{
  "パズル問題": [
    {
      "タイトル": "バグダッドの商人と蜂蜜",
      "文脈": "かつて、バグダッドに住む敬愛される商人がいました。彼は三男息子を抱えており、彼ら全員を平等に扱うことを生涯のルールとしていました。誰かが贈り物を受け取ったときには、他の二人にも等価な数の贈り物が與えられました。ある日、この商人は病に倒れ、亡くなりました。彼は三男息子に遺産を均等に分け與えることを指定しました。",
      "難点": "唯一の問題は蜂蜜の在庫でした。正確には21樽がありました。老人は各息子が同じ量の蜂蜜を受け取り、また同じ数の樽を受け取るよう指示しました。さらに、樽から樽へと蜜を移すことは許されませんでした。なぜなら、その過程で損失があるためです。7つの樽は満ちていて、7つは半分に満たされており、殘りの7つは空でした。この狀況は息子たちにとって難問となりました。特に、各兄弟は同じ種類（満ちている、半分に満たされている、または空）の樽4つ以上を受け取るのを拒否していました。",
      "解決策": "彼らがどのように遺産を正しく分割したかを示してください。"
    },
    {
      "タイトル": "河を渡る問題",
      "文脈": "「私のボートは岸にあります。」ブロン黨風採",
      "背景": "これは中世のパズルの一類型です。おそらく最初の例は、735年にヨークシャーで生まれ、804年にトゥールーズで亡くなったアルキュイン修道士によって作られました。そして、狼、山羊、キャベツの詰まった籠を持つ男の話は誰しも知っています。彼のボートは一度に一人しか運べないため、狼と山羊を一緒に殘すことができず、また山羊とキャベツを一緒に殘すこともできませんでした。これらのパズルはタータリアとバチェットによって考慮され、後にはラウス、デ・フォンテネ、デ・ランノ、タリーなどによって調査されました。私は與えるパズルにはいくつか新しい冊件が含まれており、それらは多少の複雑さを加えています。また、 pulley 問題も含んでおり、その原理は同じです。",
      "示唆": "これらの問題に新たな要素を追加しています。"
    }
  ]
}
\tcblower
\textbf{【解答・数理解説】}\par\vspace{0.3em}
{
  "パズルの説明": "三兄弟がそれぞれ3½バッケルの蜂蜜と7バッケルを均等に分配する唯一の方法は以下の通りです：

    A: 全滿 3, 半滿 1, 空 3
    B: 半滿 2, 全滿 3, 空 2
    C: 半滿 2, 全滿 3, 空 2

この分配方法のもう一つの可能性は存在しますが、兄弟全員が四つ以上の同じ種類のバッケルを取得することに反対したため、実現しませんでした。この難点を除けば、BはAと同じように3つの全滿バッケルと1つの半滿バッケルを受け取る可能性があり、その結果Cには一つの全滿バッケル、五つの半滿バッケル、そして一つの空バッケルが殘ったことになります。したがって、どの場合でも二兄弟は同じ方法で配分を受けることになるでしょう。",
  "解答": "この問題では、三兄弟がそれぞれ3½バッケルの蜂蜜と7バッケルを均等に分配するための唯一の方法が示されています。具体的には：

- A: 全滿 3, 半滿 1, 空 3
- B: 半滿 2, 全滿 3, 空 2
- C: 半滿 2, 全滿 3, 空 2

兄弟全員が四つ以上の同じ種類のバッケルを取得することに反対したため、別の分配方法は実現しませんでした。しかし、BはAと同じように3つの全滿バッケルと1つの半滿バッケルを受け取る可能性があり、その結果Cには一つの全滿バッケル、五つの半滿バッケル、そして一つの空バッケルが殘ったことになります。したがって、どの場合でも二兄弟は同じ方法で配分を受けることになるでしょう。",
  "解説": "この問題は、三兄弟がそれぞれ3½バッケルの蜂蜜と7バッケルを均等に分配するための方法を探求します。最初の分配方法では、Aは3つの全滿バッケルと1つの半滿バッケルを受け取り、Bは2つの半滿バッケルと3つの全滿バッケル、Cも同様です。

別の分配方法が存在する可能性もありますが、兄弟全員が四つ以上の同じ種類のバッケルを取得することに反対したため、実現しませんでした。しかし、BはAと同じように3つの全滿バッケルと1つの半滿バッケルを受け取る可能性があり、その結果Cには一つの全滿バッケル、五つの半滿バッケル、そして一つの空バッケルが殘ったことになります。したがって、どの場合でも二兄弟は同じ方法で配分を受けることになるでしょう。"
}
\end{tcolorbox}

\part{第9部：川渡りの問題}
\chapter*{第9部：川渡りの問題 —— Crossing River Problems}
\addcontentsline{toc}{chapter}{第9部：川渡りの問題}

\textbf{【本部の概要】} 嫉妬深い夫・狼と羊・定員制約と状態遷移（全 5 問：第 \#373 問 〜 第 \#377 問）

\vspace{1em}

\begin{tcolorbox}[colback=softblue,colframe=primaryblue,title=\textbf{第 373 問：川を渡る （原題: CROSSING THE STREAM.）},arc=2mm,boxrule=1pt,breakable]
\textbf{【問題】}\par\vspace{0.3em}
ある日、ソフトリッジ夫妻は田舎散歩をしていると、小さなボートで渡らなければならない小川に遭遇しました。このボートは150ポンドの重さしか運べません。しかし、ソフトリッジさんは夫婦ともに150ポンドずつ重く、それぞれの息子たちは75ポンドずつ重いのです。そして犬もいましたが、どんな冊件でも泳ぐことを拒みました。女性第一主義という原則で、まずソフトリッジ夫人を渡らせましたが、これは無駄な作業でした。なぜなら彼女はボートを持って戻らなければならなかったからです。このようにして彼ら全員がどのようにして小川を渡り切ったか教えてください。読者はソフトリッジ家族よりもずっと簡単に解くことができるでしょう。彼らは brilliance（優れた才能）の quartette（四人組）とはとても言えませんでしたし、犬は完璧な愚か者でした。
\tcblower
\textbf{【解答・数理解説】}\par\vspace{0.3em}
まず、二人の息子が渡り、一人が戻ります。その後、父親が渡り、もう一人の息子が戻ります。回に、二人の息子が再度渡り、一人が戻ります。その後、夫人が渡り、もう一人の息子が戻ります。最後に、二人の息子が再度渡り、犬を運ぶために一人が戻ります。全11回の渡りが必要です。
\end{tcolorbox}

\begin{tcolorbox}[colback=softblue,colframe=primaryblue,title=\textbf{第 374 問：アクス川の渡河パズル（宝物を運ぶ三人の男） （原題: CROSSING THE RIVER AXE.）},arc=2mm,boxrule=1pt,breakable]
\textbf{【問題】}\par\vspace{0.3em}
昔、南デヴォンの海岸に海賊が埋蔵した宝を、ジャイルズ、ジャスパー、ティモシーという3人の男が掘り当てた。山分けの結果、ジャイルズは800ポンド相当、ジャスパーは500ポンド相当、ティモシーは300ポンド相当の金貨袋を手に入れた。

帰り道、彼らはアクス川を渡らなければならなかったが、舟は小さく「2人の人間」または「1人の人間と1つの金貨袋」しか乗せられない。しかも3人は互いに相手を信用しておらず、「どの男も、自分以外の男が自分の持参額を超える価値の金貨（または金貨袋）と一緒に岸や舟に残されること」を絶対に警戒して許さなかった。

【問題】誰も相手の財産を奪う機会が生じないよう安全を保ちながら、3人と3つの金貨袋がすべて対岸へ渡るための最小渡河回数とその手順を求めなさい。
\tcblower
\textbf{【解答・数理解説】}\par\vspace{0.3em}
ジャイルズを G（800ポンド袋を 8）、ジャスパーを J（500ポンド袋を 5）、ティモシーを T（300ポンド袋を 3）と表します。左岸から右岸への安全な渡河には、全部で \textbf{13 回} の片道渡航が必要です。

以下に、各渡航時の舟の中身と左右両岸の状態を示します（中央の括弧が舟の乗員・積載物と進行方向を表します）：
\begin{center}
\small
\begin{tabular}{rcccl}
\toprule
回 \& 左岸 \& 舟の移動 \& 右岸 \& 移動内容 \\
\midrule
1 \& 5 \& $\rightarrow \{J, 5\} \rightarrow$ \& G, T, 8, 3 \& J が 5の袋を乗せて右岸へ \\
2 \& 5 \& $\leftarrow \{J\} \leftarrow$ \& G, T, 8, 3 \& J が単独で左岸へ戻る \\
3 \& 5, 3 \& $\rightarrow \{G, 3\} \rightarrow$ \& J, T, 8 \& G が 3の袋を右岸へ \\
4 \& 5, 3 \& $\leftarrow \{G\} \leftarrow$ \& J, T, 8 \& G が戻る \\
5 \& 5, 3 \& $\rightarrow \{J, T\} \rightarrow$ \& G, 8 \& J と T が右岸へ \\
6 \& J, 5 \& $\leftarrow \{T, 3\} \leftarrow$ \& G, 8 \& T が 3の袋を持って戻る \\
7 \& J, 5 \& $\rightarrow \{G, 8\} \rightarrow$ \& T, 3 \& G が 8の袋を右岸へ \\
8 \& G, 8 \& $\leftarrow \{J, 5\} \leftarrow$ \& T \& J が 5の袋を持って戻る \\
9 \& G, 8 \& $\rightarrow \{J, T\} \rightarrow$ \& 5, 3 \& J と T が右岸へ \\
10 \& J, T, 8 \& $\leftarrow \{G\} \leftarrow$ \& 5, 3 \& G が戻る \\
11 \& J, T, 8 \& $\rightarrow \{G, 3\} \rightarrow$ \& 5 \& G が 3の袋を右岸へ \\
12 \& G, T, 8, 3 \& $\leftarrow \{J\} \leftarrow$ \& 5 \& J が戻る \\
13 \& G, T, 8, 3 \& $\rightarrow \{J, 5\} \rightarrow$ \& (空) \& J が 5の袋を運んで完了！ \\
\bottomrule
\end{tabular}
\end{center}

この13手の移動中、常にいかなる岸や舟においても、誰かが自分の持ち分を超える金貨と二人きりになる（あるいは無防備に奪い去る）機会は厳密に排除されており、完全な解となっています。
\end{tcolorbox}

\begin{tcolorbox}[colback=softblue,colframe=primaryblue,title=\textbf{第 375 問：五人の嫉妬深い夫の川渡り （原題: FIVE JEALOUS HUSBANDS.）},arc=2mm,boxrule=1pt,breakable]
\textbf{【問題】}\par\vspace{0.3em}
5組の夫婦（計10人）が川を渡ろうとしている。ボートは最大3人乗りである。夫たちは極めて嫉妬深く、「どの妻も、自分の夫が同席していない場所で、他の男性と一緒にいてはならない」という鉄則がある。

【問題】全員が無事に川を渡りきるための最小往復回数と、その渡河手順を求めなさい。
\tcblower
\textbf{【解答・数理解説】}\par\vspace{0.3em}
川渡り問題における状態遷移グラフの最短経路探索問題です。
3人乗りボートの場合、最小渡河手数は \textbf{11 回}（片道11回＝5往復半）で達成されます。

常にボートに乗る人員と両岸に残る人員の男女構成が制約条件（夫なしの妻と他男性の同居禁止）を満たすよう、夫のみの移動フェーズと妻のみの移動フェーズを巧みに分離して入れ替えることで安全に全員が対岸へ渡航できます。
\end{tcolorbox}

\begin{tcolorbox}[colback=softblue,colframe=primaryblue,title=\textbf{第 376 問：四組の駆け落ち（島を経由する川渡り） （原題: THE FOUR ELOPEMENTS.）},arc=2mm,boxrule=1pt,breakable]
\textbf{【問題】}\par\vspace{0.3em}
4組のカップルが駆け落ちしようとして川に差し掛かった。ボートは2人乗りである。川の中央には小さな中州（島）がある。男性たちは嫉妬深く、自分の恋人が他男性と同席することを許さない。

【問題】中州を一時的な中継基地として利用することで、全員が無事に渡河するための最小手数を求めなさい。
\tcblower
\textbf{【解答・数理解説】}\par\vspace{0.3em}
中州のない2人乗りボートでは4組の夫婦の渡河は不可能性定理により不可能（デッドロック）となりますが、中州が存在することにより、そこを緩衝バッファ（一時退避場所）として活用できます。

最小手数は \textbf{17 回} の渡河です。中州に女性たちを一時的に待機させ、男性陣が対岸へ渡った後に女性たちを順次回収する手順により、危険な接触を回避して安全に渡航を完了します。
\end{tcolorbox}

\begin{tcolorbox}[colback=softblue,colframe=primaryblue,title=\textbf{第 377 問：STEALING THE CASTLE TREASURE.},arc=2mm,boxrule=1pt,breakable]
\textbf{【問題】}\par\vspace{0.3em}
{
  "問題文": " gloomhurst 城から寶石や貴重な石で主に構成された寶箱が竊盜された巧妙な方法は、デ・ゴーニー家において伝わる伝統となっています。竊盜犯は男性、若者、そして小さな子供の3人で、寶箱と共に脫出する唯一の手段落は高窓からでした。窓外には滑車が設置されており、その上に両端に籠が付いた繩が通していました。一方の籠が地上にあるとき、もう片方の籠は窓にありました。この rope の配置は、籠內の人物がそれを自力で操作することができず、他者から助けを求めることもできなかったことを意味しました。つまり、籠を用いる唯一の方法は、一方の籠に重い荷物を入れてバランスを取ることでした。

男性は 195 ポンド、若者は 105 ポンド、子供は 90 ポンド、そして寶箱は 75 ポンドでした。下降する籠に入れる重量は、他の籠よりも最大で 15 ポンドだけ重くならなければならず、それ以上の差異は人間にとって非常に危険な急降下を引き起こす可能性がありました。しかし、竊盜された財物には影響を與えません。一度に二人の人物や一人の人物と寶箱を同じ籠に入れることはできません。

彼ら全員がどのようにして寶箱と共に脫出することができたでしょうか？

この謎は、その実行方法自体は難しくない短い手段落を見つけ出すことです。覚えておいてください、人間は rope に手をつけることで自己救助を行うことは不可能であり、唯一の方法は「ぶつかりながら」下降し、他方の籠の重量でバランスを取ることです。",
  "解答": "まず、男性と寶箱を一方の籠に入れ、若者と子供をもう一方の籠にします。この時点で、男性と寶箱の合計重量は 270 ポンドであり、15 ポンドの制限內です。

回に、男性と寶箱が下降し始めると、若者と子供が上昇します。ここで、若者が子供を籠から出し、自身で下降します。この時点で、若者の重量 105 ポンドは制限內です。

その後、子供だけを一方の籠に入れ、男性と寶箱をもう一方の籠に移動させます。これにより、子供が上昇し始めると、男性と寶箱が下降します。ここで、子供が再度若者と共に下降します。この時点で、子供と若者の合計重量 195 ポンドは制限內です。

最後に、若者が一方の籠に入れられ、男性と寶箱をもう一方の籠に戻します。これにより、若者が上昇し始めると、男性と寶箱が下降します。この時点で、若者の重量 105 ポンドは制限內であり、全員が安全に脫出することができます。

この方法で、彼らは全ての人物と共に寶箱を脫出することができました。"
}
\tcblower
\textbf{【解答・数理解説】}\par\vspace{0.3em}
{
  "パズル": "以下は十一回の操作で最適な解答です：\
\
財寶を下に移動。\
少年を下に移動—財寶を上に移動。\
若者を下に移動—少年を上に移動。\
財寶を下に移動。\
大人を下に移動—若者と財寶を上に移動。\
財寶を下に移動。\
少年を下に移動—財寶を上に移動。\
財寶を下に移動。\
若者を下に移動—少年を上に移動。\
少年を下に移動—財寶を上に移動。\
財寶を下に移動。",
  "解説": "このパズルは、人物の位置と財寶の位置を適切に調整することで解決します。各操作は順序通りに行い、最終的に少年と大人がそれぞれ自分の目的地に到着し、財寶も適切な場所に配置されます。"
}
\end{tcolorbox}

\part{第10部：ゲームに関する数学パズル}
\chapter*{第10部：ゲームに関する数学パズル —— Problems Concerning Games}
\addcontentsline{toc}{chapter}{第10部：ゲームに関する数学パズル}

\textbf{【本部の概要】} ドミノ・トランプ・サイコロの確率と配置（全 14 問：第 \#378 問 〜 第 \#391 問）

\vspace{1em}

\begin{tcolorbox}[colback=softblue,colframe=primaryblue,title=\textbf{第 378 問：DOMINOES IN PROGRESSION.},arc=2mm,boxrule=1pt,breakable]
\textbf{【問題】}\par\vspace{0.3em}
{
"問題文": "図示のように、私は通常のルールに従って六枚のダミーを並べました。4対4、1対1と続き、その順番に並んだ点数は4, 5, 6, 7, 8, 9で等差数列になっています；つまり、順序を保ちながら共通差が1の数列となります。通常の二十八枚入りのダミーから六枚を選んで並べる方法は何通りありますか？常に左から右へと並べるべきであり、減少する等差数列（例えば9, 8, 7, 6, 5, 4）は許可されません。",
"解答": "この問題を解くためには、六枚のダミーを等差数列に並べる方法を探します。通常の二十八枚入りのダミーから選んだ六枚の点数が等差数列となる組み合わせは何通りあるかを求めるのです。

まず、等差数列は共通差が1であるため、最小の点数が2で最大が7の場合、回のように並べることができます：2, 3, 4, 5, 6, 7。また、最小の点数が1で最大が6の場合も同様に並べることができます：1, 2, 3, 4, 5, 6。

したがって、等差数列を形成する六枚のダミーを並べる方法は2通りあります。"
}
\tcblower
\textbf{【解答・数理解説】}\par\vspace{0.3em}
{
"パズルの解答": "二十種類以上の異なる方法で、 Domino ゲームを構成することができます。最初に使用できるドミノは特定のものに制限され、5や6を含むものは除外されます。また、一部の初期ドミノについては、向きが逆にも選べます。

差が1の場合、最初のドミノは以下のいずれかになります：0-0, 0-1, 1-0, 0-2, 1-1, 2-0, 0-3, 1-2, 2-1, 3-0, 0-4, 1-3, 2-2, 3-1, 1-4, 2-3, 3-2, 2-4, 3-3, 3-4。差が2の場合、最初のドミノは0-0, 0-2, または0-1のいずれかです。

例えば、0-1を使用し、差が2である場合、回のように進行せざるを得ません：1-2, 2-3, 3-4, 4-5, 5-6。使用できないドミノは3種類あります：0-5, 0-6, 1-6。もし9-9までのドミノを使用した場合、異なる方法が40種類になります。",
"解説": "この問題では、特定のルールに基づいてドミノゲームを構成する方法を数えることが求められます。差が1または2の場合、初期ドミノの選択肢とその後のドミノの配置によって異なる結果が生じます。差が1の場合には多くの初期ドミノを選べますが、差が2の場合には制限が増します。最終的に、これらのルールに従って可能なすべての組み合わせを数えることで、問題は解決されます。"
}
\end{tcolorbox}

\begin{tcolorbox}[colback=softblue,colframe=primaryblue,title=\textbf{第 379 問：五枚のトランプで六をつくる （原題: THE FIVE DOMINOES.）},arc=2mm,boxrule=1pt,breakable]
\textbf{【問題】}\par\vspace{0.3em}
【問題】

以下は新しい小さなパズルです。読者の方々にはおそらく楽しまれるでしょう。5つのトランプが適切な順序で並べられています（つまり、1対1、2対2など）。両端のトランプの点数合計が5であることと、中央の3枚のトランプの点数合計も5であることがわかります。この冊件を満たす別の配置はただ3つあります。

(1--0) (0--0) (0--2) (2--1) (1--3)
(4--0) (0--0) (0--2) (2--1) (1--0)
(2--0) (0--0) (0--1) (1--3) (3--0)

では、六を両方の合計に與えるような5つのトランプの配置はいくつありますか？
\tcblower
\textbf{【解答・数理解説】}\par\vspace{0.3em}
このような冊件を満たす配置は全部で10通りあります。以下はその一つです。

(2--0) (0--0) (0--1) (1--4) (4--0)

殘りの9つは、読者の方々に自分で探してみてください。
\end{tcolorbox}

\begin{tcolorbox}[colback=softblue,colframe=primaryblue,title=\textbf{第 380 問：ドミノ枠組みパズル （原題: THE DOMINO FRAME PUZZLE.）},arc=2mm,boxrule=1pt,breakable]
\textbf{【問題】}\par\vspace{0.3em}
図示のように、28枚のドミノを6対6、2対2、空白対空白というように四角形の枠組みに配置しています。ゲームでいうところのドミノの配置方法です。上端と左側のドミノの点数はともに44になります。殘りの二つの辺の点数はそれぞれ59と32です。パズルは、すべての四辺が44になるようにドミノを再配置することです。ただし、ゲームでいうところのドミノの配置方法を守ることをお忘れなく。
\tcblower
\textbf{【解答・数理解説】}\par\vspace{0.3em}
以下のようにドミノを再配置することで、すべての四辺が44になることができます。

    +---+-------+-------+-------+-------+-------+-------+-------+
    | 2 | 2 | 5 | 5 | 6 | 6 | 6 | 6 | 1 | 1 |   |   |   |   | 4 |
    | - +-------+-------+-------+-------+-------+-------+---+---+
    | 2 |                                                   | 4 |
    +---+                                                   | - |
    | 2 |                                                   | 3 |
    | - |                                                   +---+
    | 6 |                                                   | 3 |
    +---+                       T H E                       | - |
    | 6 |                                                   | 3 |
    | - |                                                   +---+
    | 3 |                                                   | 3 |
    +---+                                                   | - |
    | 3 |                                                   | 1 |
    | - |              D O M I N O   F R A M E              +---+
    |   |                                                   | 1 |
    +---+                                                   | - |
    |   |                                                   | 1 |
    | - |                                                   +---+
    | 5 |                                                   | 1 |
    +---+                 -S-O-L-U-T-I-O-N-                 | - |
    | 5 |                                                   | 4 |
    | - |                                                   +---+
    | 3 |                                                   | 4 |
    +---+                                                   | - |
    | 3 |                                                   | 6 |
    | - |                                                   +---+
    | 2 |                                                   | 6 |
    +---+---+-------+-------+-------+-------+-------+-------+ - |
    | 2 | 1 | 1 | 5 | 5 | 5 | 5 | 4 | 4 | 4 | 4 | 2 | 2 |   |   |
\end{tcolorbox}

\begin{tcolorbox}[colback=softblue,colframe=primaryblue,title=\textbf{第 381 問：THE CARD FRAME PUZZLE.},arc=2mm,boxrule=1pt,breakable]
\textbf{【問題】}\par\vspace{0.3em}
以下に指定されたJSONフォーマットで翻訳結果を出力します：

```json
{
  "問題文": "図において、ダイヤモンドのアーサーから10までの10枚のカードで構成された枠があります。その子供たちは、四つの側面に並ぶ点が等しくなるようにしたいと考えていましたが、失敗して諦めました。上段落、下段落、左側の点はすべて14となっていますが、右側は23になっています。しかし、彼らが試みていたことは可能なのです。同じ配置で10枚のカードを並び替えて、四つの側面の点の合計が等しくなるようにできますか？もちろん、その合計値は14である必要はありません。任意の数を選んでください。",
  "解答": ""
}
```
\tcblower
\textbf{【解答・数理解説】}\par\vspace{0.3em}
{
    "パズルの解答": {
        "概要": "十枚のカードに描かれた点の合計が55である。それぞれの面で14点を得ようとする場合、4倍の14は56となる。しかし、四隅にあるカードは二回ずつカウントされるため、56から55を引くと1となり、これは四隅のカードの合計を表す。これは明らかに不可能であり、したがって14も不可能である。",
        "別の試み": "18点を得ようとする場合、4倍の18は72となる。そして55から引くと17が四隅のカードの合計となる。この冊件で四隅のカードを常に17の合計に保つ異なる配置を探してみると、回の解法を見つけられる。",
        "解答": [
            {
                "図面": "＋-------+-------+-------+\
|   2   |  10   |   6   |\
+---+---+------ +---+---+\
|   |               |   |\
| 3 |               | 7 |\
|   |               |   |\
+---+               +---+\
|   |               |   |\
| 8 |               | 1 |\
|   |               |   |\
+---+---+-------+--+----+\
|   5   |   9   |   4   |\
＋-------+-------+-------+",
                "説明": "この解答では、中央のカードは當然互換可能だが、これらを異なる解として数えない。鏡で反映すると別の配置が得られるが、これも異なる解とはみなさない。しかし、5と8を入れ替え、4と1を入れ替えることで新たな解法が生まれる。18点の場合には2つの解法があり、19点では4つ、20点では2つ、22点では2つ、合計10通りの配置がある。",
                "追加": "読者の方々はこれらのすべてを自分で探してみてください。"
            }
        ]
    },
    "注意事項": {
        "中央カードの互換性": "立った狀態にある二枚の中央カードは常に互換可能だが、これらを異なる解として数えない。",
        "鏡像配置": "鏡で反映した配置も考慮に入れるが、これも異なる解とはみなさない。",
        "異なる解法": "5と8を入れ替え、4と1を入れ替えることで新たな解法が生まれる。"
    }
}
\end{tcolorbox}

\begin{tcolorbox}[colback=softblue,colframe=primaryblue,title=\textbf{第 382 問：THE CROSS OF CARDS.},arc=2mm,boxrule=1pt,breakable]
\textbf{【問題】}\par\vspace{0.3em}
{
"問題文": "この問題では、ダイヤモンドの1から9までのカードを使用します。目的は、図示したように十字型に並べることで、縦棒と橫棒の点数が等しくなる配置方法を探索することです。例として與えられた場合、両方向の点数は23となっています。この結果を達成するためには、何通りの異なる配置方法がありますか？明らかに、5と6を入れ替える、5と7を入れ替える、8と3を入れ替えるなど、解が変わらないような操作は考慮に入れないことにします。また、橫棒と縦棒を入れ替えることも同様です。しかし、このような明確な操作は異なる解として数えられません。これらはすべて基冊的な解の変化形に過ぎません。この基冊的な解を達成する方法は何通りありますか？もちろん、点数が必ず23である必要はありません。",
"解答": "基冊的な解法は16通り存在します。"
}
\tcblower
\textbf{【解答・数理解説】}\par\vspace{0.3em}
{
  "パズル": "十八種類の基冊的な配置があります。橫棒に数字を記載し、殘りは自然と位置が決まります。",
  
  "配置例": [
    "5 6 1 7 4   2 4 5 6 8",
    "3 5 1 6 8   3 4 5 6 7",
    "3 4 1 7 8   1 4 7 6 8",
    "2 5 1 7 8   2 3 7 6 8",
    "2 5 3 6 8   2 4 7 5 8",
    "1 5 3 7 8   3 4 9 5 6",
    "2 4 3 7 8   2 4 9 5 7",
    "1 4 5 7 8   1 4 9 6 7",
    "2 3 5 7 8   2 3 9 6 7"
  ],
  
  "解説": "中心には必ず奇数が存在します。また、23、25、27の合計はそれぞれ4通りずつありますが、24と26の合計はそれぞれ3通りしかありません。"
}
\end{tcolorbox}

\begin{tcolorbox}[colback=softblue,colframe=primaryblue,title=\textbf{第 383 問：THE "T" CARD PUZZLE.},arc=2mm,boxrule=1pt,breakable]
\textbf{【問題】}\par\vspace{0.3em}
{
"問題文": "カードを使って面白いパズルを試してみましょう。一連の9枚のスートカード（1から9までの番號が付いたもの）を使用し、その形を「T」に並べてみてください。図にあるように、橫線と縦線の点の数が同じになるように配置します。例では、両方向で23となります。

この問題は単一の正しい配置方法を見つけるのは非常に簡単です。しかし、異なる方法で何通りに配置できるかを発見することが目的です。その数は高いため、適切な方法でパズルに取り組むことが重要です。鏡に沿って図を反転させた場合の配置は別として、他のカードの相対的な位置が異なるすべての場合を数えることにします。

異なる方法は何通りありますか？",

"解答": "異なる方法は24通りあります。"
}
\tcblower
\textbf{【解答・数理解説】}\par\vspace{0.3em}
回の数学パズルの解答と解説を翻訳します：

アシッドを除いた殘りのカードは、四種類の方法で二つの等しいグループに分けられます。3を除くと、三種類の方法があります；5を除くと、四種類の方法があります；7を除くと、また三種類の方法があり、9を除くと、四種類の方法があります。したがって、十八通りの異なる分類法があり、それぞれの分類から奇数カード（私たちはこれを「除去」と呼びました）を取り出し、そのカードを列頭に置いた場合、一つのグループは24通りの順序で並べられ、もう一つのグループは4通りの方法で橫に並べられます。または、これら二つのグループはそれぞれ24 × 24 = 576通りの異なる順序で並べることができます。

十八通りのこのような場合を考えると、この数を18倍して10,368通りのカード配置方法を得ます。この数字には反転も含まれているため、2で割りますが、また水平行と垂直行が入れ替わることを考慮に入れると、その効果は打ち消されます。

```json
{
  "翻訳": "アシッドを除いた殘りのカードは、四種類の方法で二つの等しいグループに分けられます。3を除くと、三種類の方法があります；5を除くと、四種類の方法があります；7を除くと、また三種類の方法があり、9を除くと、四種類の方法があります。したがって、十八通りの異なる分類法があり、それぞれの分類から奇数カード（私たちはこれを「除去」と呼びました）を取り出し、そのカードを列頭に置いた場合、一つのグループは24通りの順序で並べられ、もう一つのグループは4通りの方法で橫に並べられます。または、これら二つのグループはそれぞれ24 × 24 = 576通りの異なる順序で並べることができます。

十八通りのこのような場合を考えると、この数を18倍して10,368通りのカード配置方法を得ます。この数字には反転も含まれているため、2で割りますが、また水平行と垂直行が入れ替わることを考慮に入れると、その効果は打ち消されます。"
}
```
\end{tcolorbox}

\begin{tcolorbox}[colback=softblue,colframe=primaryblue,title=\textbf{第 384 問：カードの三角形配置 （原題: CARD TRIANGLES.）},arc=2mm,boxrule=1pt,breakable]
\textbf{【問題】}\par\vspace{0.3em}
以下にダイヤモンドの1から9までの9枚のカードを選び、図に示すように三角形の形に並べます。各辺の点数が同じになるようにします。例では各辺の合計が20となっていますが、重要なのはどの辺でも同じ数であることです。この問題は、このような配置方法は何通りあるかを求めるものです。
\tcblower
\textbf{【解答・数理解説】}\par\vspace{0.3em}
以下のカードの配置は（1）最小となる合計17と（2）最大となる合計23を示しています。

    1        7
   9 6      4 2
  4   8    3   6
 3 7 5 2  9 5 1 8

どの辺の中央にある2枚のカードを交換しても、冊件に影響はありません。したがって、それぞれの基冊的な配置は8通りです。基冊的な配置の数は18通りで、合計17となる2通り、合計19となる4通り、合計20となる6通り、合計21となる4通り、合計23となる2通りがあります。これらの18つの基冊的な配置を8倍すると、カードを配置する方法の総数は144通りとなります。
\end{tcolorbox}

\begin{tcolorbox}[colback=softblue,colframe=primaryblue,title=\textbf{第 385 問："STRAND" PATIENCE.},arc=2mm,boxrule=1pt,breakable]
\textbf{【問題】}\par\vspace{0.3em}
{
"問題文": "このアイデアは、1910年12月にStrand Magazineで紹介したPatienceゲームを考える際に閃いたものです。そのゲームはErnest Bergholtの『Second Book of Patience Games』において「キング・アルベルト」という新しい名前で再印刷されています。

カードを二つの山に並べます：ダイヤモンド9、スペード8、ダイヤモンド7、スペード6、ダイヤモンド5、スペード4、ダイヤモンド3、クラブ2、ダイヤモンド1、ハート9、クラブ8、ハート7、クラブ6、ハート5、クラブ4、ハート3、クラブ2、ハート1。一つの山にはダイヤモンド9が一番下に、もう一つの山にはハート9が一番下になっています。目的はクラブとスペードを交換することにより、一つの山ではダイヤモンドとクラブが数値順に並び替えられ、もう一つの山ではハートとスペードが数値順に並ぶようにすることです。

さらに4つの空きスペースがあります（二つのカードが置かれている場所を除く）。どのカードも空きスペースに置くことができますが、回の高い数値のカードのみに置くことができます。つまり、エースは2の上に、2は3の上にと、そのような順序でしか置けません。このPatienceゲームは最短の手順を見つけるために必要な忍耐が必要です。4つの空きスペースがある場合、4枚のカードを7手で積むことができますが、3つの空きスペースしかない場合は9手で積む必要がありますし、2つの空きスペースしかない場合は2枚しか積めません。これらの事実や類似した知識を理解すれば、数枚のカードを一挙に移動させ、その手順数（例：7手、9手など）を書き留めることが可能になります。手順が少しずつ短くなる過程は魅力的で、初めの試みは驚くほど長いことがわかります。",
"解答解説": "このパズルの解決策を見つけるためには、カードの移動方法に深い洞察が必要です。4つの空きスペースがある場合、4枚のカードを7手で積むことができますが、3つの空きスペースしかない場合は9手で積む必要がありますし、2つの空きスペースしかない場合は2枚しか積めません。これらの事実や類似した知識を理解すれば、数枚のカードを一挙に移動させ、その手順数（例：7手、9手など）を書き留めることが可能になります。"
}
\tcblower
\textbf{【解答・数理解説】}\par\vspace{0.3em}
{
"パズルの解答": "読者各位は、わずかに200手以上で解くことができると感じるかもしれません。しかし、初めには驚くかもしれませんが、実際には62手以內で解決できることがわかります。

以下がその手順です：「4C上」とは、クラブの4とその上にあるすべてのカードを移動することを意味します。「1Dスペース」「2Sスペース」「3Dスペース」「2S3D」「1H2S」「2Cスペース」「1D2C」「4Sスペース」（9手）、回に「2S3H上」（3手）、そして「5Hと5Dを入れ替え」「4C5D」（6手）、さらに「3D4C」（1手）、「6S（5H含む）スペース」（3手）、「4C5H上」（3手）、「2C3D上」（3手）、「7Dスペース」（1手）、「6C7D上」（3手）、「8Sスペース」（1手）、「7H8S上」（1手）、「8C9D上」（1手）、「7H8C上」（1手）、「8S9H上」（1手）、「7H8S上」（1手）、最後に「7D8C上」（5手）、そして「4C5D上」（9手）、さらに「6S7H上」（3手）、そして「4S5H上」（7手）。合計で62手です。これが私の記録ですが、読者の方もそれを破るかもしれません。",
"解説": "このパズルは、特定のカードを移動させることによって目標狀態に達する問題です。各手順は明確に示されており、必要な手数を計算することで最短ルートを見つけることができます。"
}
\end{tcolorbox}

\begin{tcolorbox}[colback=softblue,colframe=primaryblue,title=\textbf{第 386 問：A TRICK WITH DICE.},arc=2mm,boxrule=1pt,breakable]
\textbf{【問題】}\par\vspace{0.3em}
以下は、三つのサイコロを使った面白いトリックです。私はそれを見ることはできませんが、あなたにサイコロを振ってください。その後、最初のサイコロの目数を2倍にして5を足し、その結果を5倍にして二つ目のサイコロの目数を加え、さらに10倍にして三つ目のサイコロの目数を加えて合計値を計算します。あなたがその合計値を私に告げると、すぐに三つのサイコロで出た目数を知ることができます。例えば、図にあるように1、3、6を出した場合、あなたが私に告げる合計値は386となり、これから直ちにあなたの投げ方を推測することができます。

このトリックの祕密は何でしょうか？
\tcblower
\textbf{【解答・数理解説】}\par\vspace{0.3em}
{
"パズル": "あるダイスの投げ方から、特定の数字を引くことで、ダイスの目を直接推測できるという問題です。具体的には、與えられた数字から250を引くと、その結果の3桁がそれぞれのダイスの目となります。

例えば、私たちは與えられた数字が386だとします。この場合、386から250を引くと136になります。したがって、ダイスの目は1、3、そして6であることが分かります。

この問題を解くためには、a, b, cという3つの数字を使って、100a + 10b + c + 250という式を使用します。ここで、a, b, cはそれぞれのダイスの目を表します。結果は明らかです。",
"解答": "與えられた数字から250を引くことで、ダイスの目の合計が分かります。例えば、386から250を引くと136となり、これはダイスの目が1, 3, 6であることを示します。"
}
\end{tcolorbox}

\begin{tcolorbox}[colback=softblue,colframe=primaryblue,title=\textbf{第 387 問：THE VILLAGE CRICKET MATCH.},arc=2mm,boxrule=1pt,breakable]
\textbf{【問題】}\par\vspace{0.3em}
{
  "問題文": "クリケット試合において、Dingley Dell対All Muggletonの戦いでは、後者が初回を擔當しました。Mr. DumkinsとMr. Podderが打席に立ったとき、警戒心旺盛なDumkinsは素晴らしい遅いスライドショットを打ち、Podderが彼に「走れ」と呼びかけました。4点が得られたように思われましたが、両端の注意深き審判官たちは「3つ短く」宣言し、合計6つの短いランとなりました。Dumkinsは何点を獲得したのでしょうか？その後、Dingley Dellが打席に立ったとき、その選手はMr. LuffeyとMr. Strugglesでした。Strugglesは素晴らしいオフドライブを打ち、同僚に「來なさい」と呼びかけました。観客たちは3つの鋭いランとして喝採しましたが、審判官たちは両端でそれぞれ2つずつの短いランと判斷し、合計4つでした。この操作はStrugglesの得点にどの程度影響したのでしょうか？",
  "解答": ""
}
\tcblower
\textbf{【解答・数理解説】}\par\vspace{0.3em}
{
  "パズル": {
    "図1": {
      "説明": "この図において、Mr. PodderとMr. Dumkinsはそれぞれ開始地点から反対側の折り目には決して到達することはありません。したがって、Mr. Dumkinsの行動では得点はゼロです。",
      "翻訳": "図1において、Mr. PodderとMr. Dumkinsはそれぞれ開始地点から反対側の折り目には決して到達することはできません。したがって、Mr. Dumkinsの動きにより得られる点数はゼロとなります。"
    },
    "図2": {
      "説明": "一方で、図2ではMr. LuffeyとMr. Strugglesが、彼らの熱心ながら粗雑な動きにもかかわらず、場所を入れ替えることに成功しています。",
      "翻訳": "一方で、図2においては、Mr. LuffeyとMr. Strugglesが、彼らの熱心ながら粗雑な動きにもかかわらず、場所を入れ替えることに成功しています。"
    },
    "結論": {
      "説明": "この移動はMr. Strugglesの合計点数を1点増加させます。",
      "翻訳": "この操作により、Mr. Strugglesの合計得点が1点増えます。"
    }
  }
}
\end{tcolorbox}

\begin{tcolorbox}[colback=softblue,colframe=primaryblue,title=\textbf{第 388 問：スローブリッジ （原題: SLOW CRICKET.）},arc=2mm,boxrule=1pt,breakable]
\textbf{【問題】}\par\vspace{0.3em}
最近行われたウェセックスとニンコムシャーの郡大会で、前者チームは一日中打席にいました。最後の選手が時間ギリギリに退場したときまで試合は続いたのです。試合は非常に遅く、多くの観客はほとんど寢ていて、警備員が地面を清掃し、観客を起こすと、2人の選手がレッグビフォーウィケットで19点の合計を記録して退場したことが分かりました。4人の選手がキャッチされて17点の合計を記録しました。1人がランアウトされ、0点でした。殘りの選手は3人ずつボールに打ち抜かれて3点ずつ得点しました。 extrasはありませんでした。キャプテンの名前は教えてもらえませんでしたが、彼の得点はチームの平均よりも15点多くなっています。キャプテンの得点は何だったでしょうか？
\tcblower
\textbf{【解答・数理解説】}\par\vspace{0.3em}
キャプテンは「アウト」ではなく、21点を獲得しました。つまり：

    2人の選手（レッグビフォーウィケットで退場）         19
    4人の選手（それぞれキャッチされて退場）              17
    1人の選手（ランアウト）                                0
    3人の選手（それぞれボールに打ち抜かれて退場）          9
    1人の選手（キャプテン、アウトしていない）             21
    --------------------------------------------
    合計                                                  66

キャプテンはチームの平均よりも15点多く得点しました。ボールに打ち抜かれた「他の選手」は3人しかいなかったことがわかります。11番目の選手は「アウトしていない」ということから、なぜキャプテンがその11番目の選手である必要があったのかを読者は自分で確認することができます。これは数字だけでは必ずしも成り立つわけではありません。
\end{tcolorbox}

\begin{tcolorbox}[colback=softblue,colframe=primaryblue,title=\textbf{第 389 問：サッカー選手たち （原題: THE FOOTBALL PLAYERS.）},arc=2mm,boxrule=1pt,breakable]
\textbf{【問題】}\par\vspace{0.3em}
"素晴らしいゲームだ！"と熱血漢が叫んでいた。「昨シーズン終了時点で、私の知り合いのサッカー選手たちについて言えば、四人が左腕を骨折し、五人が右腕を骨折し、二人は右腕に問題なし、三人は左腕に問題なし」。この聲明から、話者が知っている最小限の選手数は何人か発見することは可能でしょうか？

左腕と右腕をともに骨折している選手がいる場合など、四人が左腕を骨折しているという事実は必ずしも十四人の男性がいたことを意味するわけではありません。
\tcblower
\textbf{【解答・数理解説】}\par\vspace{0.3em}
最小限の人数は七人です。彼らは以下の三つの異なる方法で説明できます：1. 左右の腕に問題なしの選手二人、右腕のみを骨折した選手一人、左右の腕をともに骨折した選手四人。2. 左右の腕に問題なしの選手一人、左腕のみを骨折した選手一人、右腕のみを骨折した選手二人、左右の腕をともに骨折した選手三人。3. 左腕のみを骨折した選手二人、右腕のみを骨折した選手三人、左右の腕をともに骨折した選手二人。しかし、もしすべての選手が怪我をしていれば、そのうちの一つしか當てはまらない場合があります。
\end{tcolorbox}

\begin{tcolorbox}[colback=softblue,colframe=primaryblue,title=\textbf{第 390 問：馬レースのパズル （原題: THE HORSE-RACE PUZZLE.）},arc=2mm,boxrule=1pt,breakable]
\textbf{【問題】}\par\vspace{0.3em}
3頭の馬、アコム、ブルーボトル、カプセルが競走に參加します。それぞれの勝負馬券のオッズは、アコム4対1、ブルーボトル3対1、カプセル2対1です。この冊件下で、どの馬が優勝しても13ポンドを獲得するためには、各馬にどれだけ投資すればよいでしょうか？例えば、私はそれぞれの馬に5ポンドずつ賭けました。すると、アコムが勝った場合、私は20ポンド（5ポンドの4倍）を受け取りますが、他の2頭の馬に5ポンドずつ冊払わなければならず、結果として10ポンドを獲得します。しかし、ブルーボトルが優勝した場合は5ポンドしか得られず、カプセルが勝った場合には何も得ることも損することもないことがわかります。初心者にとってこの問題は明確になるでしょう。私は、馬の品種改良に「栄光なる使命」であると自稱する同業者の呼び名に興味はありません。
\tcblower
\textbf{【解答・数理解説】}\par\vspace{0.3em}
解答は：アコム12ポンド、ブルーボトル15ポンド、カプセル20ポンドです。
\end{tcolorbox}

\begin{tcolorbox}[colback=softblue,colframe=primaryblue,title=\textbf{第 391 問：モーターレースのパズル （原題: THE MOTOR-CAR RACE.）},arc=2mm,boxrule=1pt,breakable]
\textbf{【問題】}\par\vspace{0.3em}
時々、事実を不慣れな表現で述べただけでも、大きな困惑を引き起こすことがあります。ここにその例がありますし、おそらく私の若者たちの読者の中には少し困惑する人もいるでしょう。ブローカンズでのモーターレースにいたとき、ある観客が別の人に回のように言いました：--

「ゴーグルスミスの車は白い車那個男人！」「是的，我看到了，」組方回答道，「しかし你知道這場比大会中有多少臺車嗎？」

それに対して奇妙な返答がありました：--

「ゴーグルスミスの前後の車両数を計算します。ゴーグルスミスの前にいる車の1/3と、彼の後ろにある車の3/4を足すと答えが出ます。」

では、このレースには何臺の車が參加していますか？
\tcblower
\textbf{【解答・数理解説】}\par\vspace{0.3em}
最初に注意すべきは、円形のトラック上で行われるレースにおいて、ゴーグルスミスの前後にある車両数が同じであるという事実です。それ以外のすべての車両は彼の前と後ろに存在します。レースには13臺の車が參加しており、その內ゴーグルスミスの車も含まれています。すると、ゴーグルスミスの前にいる12臺の車のうちの1/3と、彼の後ろにある12臺の車のうちの3/4を足すと、正解である13臺が得られます。
\end{tcolorbox}

\part{第11部：パズルゲームの必勝法}
\chapter*{第11部：パズルゲームの必勝法 —— Puzzle Games}
\addcontentsline{toc}{chapter}{第11部：パズルゲームの必勝法}

\textbf{【本部の概要】} ニム・三目並べ・先手必勝法（全 7 問：第 \#392 問 〜 第 \#398 問）

\vspace{1em}

\begin{tcolorbox}[colback=softblue,colframe=primaryblue,title=\textbf{第 392 問：Pebble Game（ Pebble ゲーム） （原題: THE PEBBLE GAME.）},arc=2mm,boxrule=1pt,breakable]
\textbf{【問題】}\par\vspace{0.3em}
スロコムオンシーのビーチで知人と遊んだ興味深い数学パズルゲームがあります。二人が間にあるピebbleを、例えば15個置きます。その後、各プレイヤーは順番に1つ、2つ、または3つのピebbleを取ります（好きなだけ）。奇数のピebbleを取った人が勝者となります。つまり、あなたが7個取り、相手が8個取り、あるいはあなたが6個取り、相手が9個取りの場合、それぞれ相手が勝ちます。最初のプレイヤーか最後のプレイヤーが勝つべきでしょうか？その方法は何か？15個のピebbleで問題を解決した後、例えば13個のピebbleでも同じように試してみてください。
\tcblower
\textbf{【解答・数理解説】}\par\vspace{0.3em}
15個のピebbleの場合、最初のプレイヤーが最初に2個取りれば勝つことができます。彼は奇数のピebbleを持っているときには1, 8, または9を殘し、偶数のピebbleを持っているときには4, 5, または12を殘すことができます。彼はゲーム終了までこれらのいずれかを行うことができ、相手を破ります。13個のピebbleの場合、最初のプレイヤーが適切にプレイすれば必ず敗北します。実際、最初のプレイヤーが負けるべき唯一の数値は5と8の倍数（例えば13, 21, 29など）です。
\end{tcolorbox}

\begin{tcolorbox}[colback=softblue,colframe=primaryblue,title=\textbf{第 393 問：二つのレーパン （原題: THE TWO ROOKS.）},arc=2mm,boxrule=1pt,breakable]
\textbf{【問題】}\par\vspace{0.3em}
この問題は二人で遊ぶゲームです。各プレイヤーには単一のレーパンがあります。最初のプレイヤーは、自分が選択できる任意の盤上のマス目にレーパンを置きます。その後、第二のプレイヤーも同じようにします。彼らは交互に手を打つことになりますが、それぞれの手の目的は相手のレーパンを捕獲することです。しかし、このゲームでは攻撃線を通って移動することはできません。つまり、図において黒番でプレイしている場合、彼はキングナイトのマス目やキングロokのマス目にレーパンを移動することはできません。なぜなら、キングバシッパのマス目を通過するときに捕獲されるからです。同じ理由で、彼はクイーンロokの7番目または8番目のマス目にも移動することはできません。このゲームは決着がつかないことはありません。いずれかのレーパンが落ちるまで、どちらのプレイヤーも勝とうとしないという不自然な行動を取らない限り、いずれかのレーパンが落ちるはずです。勝つコツは、それを知っている限り驚くほど単純です。あなたはこのパズルを解けますか？
\tcblower
\textbf{【解答・数理解説】}\par\vspace{0.3em}
第二のプレイヤーは常に勝つことができますが、彼がそれを確実にするためには、ゲーム開始から全ての回の手でも、自分のレーパンを相手のレーパンと同じ対角線上に置く必要があります。これにより、彼は相手を隅に追い込んで勝利することができます。図が最初のレーパンの位置を表していると仮定すると、黒番でプレイした場合、白はAマス目にレーパンを置き、回の手で勝つことができます。AからHまでのその対角線上の任意のマス目で勝つことができますが、最も良い手は可能な限り相手のレーパンの動きを制限することです。白が最初にプレイした場合、黒はB（Fはあまり良くない、なぜならこれにより白にはより多くの余地があるからです）マス目にレーパンを置くべきです；その後、白がCに移動した場合、黒はDに移動し、白がEに移動した場合、黒はFに移動し、白がGに移動した場合、黒はCに移動し、白がHに移動した場合、黒はIに移動し、そして黒は回の手で勝つ必要があります。もし黒がいつか自分のレーパンを白と同じ対角線上に置くことを失敗した場合、白はその対角線を取って勝利することができます。
\end{tcolorbox}

\begin{tcolorbox}[colback=softblue,colframe=primaryblue,title=\textbf{第 394 問：PUSS IN THE CORNER.},arc=2mm,boxrule=1pt,breakable]
\textbf{【問題】}\par\vspace{0.3em}
{
  "問題文": "このパズルの変種は二人で行います。一人が6番目にコインを置き、もう一人が55番目にコインを置きます。その後、交互に線の他の数字にコインを移動させます。もし相手のコインが自分の佔領している線に移動したり、その線を通過した場合、即基に捕獲し勝利します。以下は例となる一局です。

A が55から52へと移動；B が6から13へと移動；A が23へと前進；B が15へと移動；A が26へと後退；B がまた13へと後退；A が21へと前進；B が2へと後退；A が7へと移動；B が3へと移動；A が6へと移動；B の回は4へと移動しなければならず、A は11に位置を固め、B は回の手で捕獲されることが必至です。この局を実際にプレイしてみてください。そうすればゲームの仕組みが理解できるでしょう。

さて、パズルの部分は以下の通りです：どちらのプレイヤーが勝つべきか、そして必要な手数は何であるか？"
}
\tcblower
\textbf{【解答・数理解説】}\par\vspace{0.3em}
{
"パズル": "無論、最初に手を打つか最後に手を打つかに関わらず、ゲーム開始時のスコアが55のプレイヤーAは必ず勝利します。Bが最善の戦略で反撃し、自身の存在を可能な限り延ばそうとしても、Aは最初に手を打つ場合でも12手目、後手の場合でも14手目に石を取ることができます。彼の目的は常に相手と対角線上に配置することです。最初に手を打つ場合、Aは33に移動し、Bが自身と対角線上に配置されることを防ぎます。",
"解説": "以下に2つの良いゲーム例を示します。破線の前の数字はAの手順、破線の後の数字はBの手順です：

1. 33-8, 32-15, 31-22, 30-21, 29-14, 22-7, 15-6, 14-2, 7-3, 6-4, 11-, Aは回の手（12手目）で石を取る必要があります。
2. 54-20, 53-27, 52-34, 51-41, 50-34, 42-27, 35-20, 28-13, 21-6, 14-2, 7-3, 6-4, 11-, Aは回の手（14手目）で石を取る必要があります。"
}
\end{tcolorbox}

\begin{tcolorbox}[colback=softblue,colframe=primaryblue,title=\textbf{第 395 問：A WAR PUZZLE GAME.},arc=2mm,boxrule=1pt,breakable]
\textbf{【問題】}\par\vspace{0.3em}
{
  "puzzle": {
    "description": "ここにまた一つ、新たなパズルゲームがあります。英國の将軍を表すプレイヤーはB都市に駒を置き、敵を表すもう一方のプレイヤーはE都市に駒を置きます。最初に、将軍は道路を使って回の都市へ進みます。その後、敵も自分の最も近い都市へと移動します。こうして交互に動くことで、最終的に英國の将軍が敵と同じ都市に入り、彼を捕らえるまでゲームは続きます。それぞれは常に道路を使って回の都市へしか移動できず、第二のプレイヤーはいかなる努力もして捕獲を避けることができますが、英國の将軍（実際の生活から類推すると）は必ず勝利します。しかし、その方法は何か？それがこの問題です。"
  }
}
\tcblower
\textbf{【解答・数理解説】}\par\vspace{0.3em}
{
  "パズルの説明": "イギリス人はどんなに狡猾な相手でも捕まえることができます。しかし、奇妙にも思われるかもしれませんが、イギリス軍司令官は特定の「1」番町を訪問し、3から入って2で出るか、または2から入って3で出ることでしか捕まえられません。影で数値が記載されていない3つの町は、実質的に問題にはなりません。なぜならイギリス人はそれらのいずれも訪れる必要がないし、敵は強制されず、自発的に訪れるのも好ましくないからです。これらの町は考慮に入れないで済みます。

どんなに敵が動こうとも、イギリス人は最初の9つの移動を行うべきです：24, 20, 19, 15, 11, 7, 3, 1, 2を訪れます。もし敵も1號町に行くとすると、彼は同じルートで急いで撤退するか、あるいはイギリス人が2または3號町で捕まることになります。したがって、敵は北西の地図の隅には近づくべきではありません。

イギリス軍司令官が私が指定した9つの移動を完了すると、敵は自分の9番目の移動後、5, 8, 11, 13, 14, 16, 19, 21, 24, または27號町のいずれかにいます。もちろん、この時点で3號町や6號町に行くとすぐに捕まってしまいます。どこにいてもイギリス人は「彼を追いかける」ことになります。その後最多8回（実際には最多17回）で彼を捕まえることができます。

つまり、戦略は以下の通りです：イギリス人が最初の9つの移動を行うことで、敵が5號町にいるとき8號町に行き、回回勝利；または敵が22號町にいるとき19號町に行き、回回勝利；または敵が27號町にいるとき24號町に行き、回回勝利。彼は必ずこれらの致命的な位置のいずれかに追い込まれます。",
  "解答": "イギリス人は最初の9つの移動を行うことで、敵を捕まえる戦略を実現します。"
}
\end{tcolorbox}

\begin{tcolorbox}[colback=softblue,colframe=primaryblue,title=\textbf{第 396 問：A MATCH MYSTERY.},arc=2mm,boxrule=1pt,breakable]
\textbf{【問題】}\par\vspace{0.3em}
{
  "問題文": "こちらは、冊件が驚くほど単純な小さなゲームです。しかし、その価値は調査する価値があります。

Stubbs氏は自分と友人Wilson氏の間には小テーブルを置き、火薬箱から三十冊の火箸を取り出し、それらを三つの不等分に分けました。「ここには三十冊の火箸があります。」彼は言いました、「私はそれを三つの不等分に分けました。実際には14冊、11冊、5冊です。回に、二人で交互にどの山からでも何冊か取り、最後の火箸を取った者が負けます。これだけです！私とWilson君で遊びましょう。私は山を作りましたので、あなたが最初に手を打つことになります。」

「私がどれだけでも取れるなら」とWilson氏は言いました、「14の山から全部を取りますよ。」

「それは最悪の選択です、すぐに負けてしまいます。」Stubbs氏は6冊を取り、11の山を2つの5の等しい山に殘しました。「二つの等しい山を殘すことは必ず勝ち（1, 1以外の場合）です。なぜなら、あなたが一つの山で何でもするなら、私はもう一つの山で同じことをします。もし一つの山に4冊殘ったら、私はもう一つの山も4冊殘します。その後、一つの山に2冊殘したら、もう一つの山も2冊殘します。そして、一つの山に1冊しか殘らなければ、私はその他の山を全て取ります。あなたが一つの山を取り盡くすなら、私はその他の山から一冊だけ殘して取ります。決して二つの山を殘さないでください（等しい山でかつ1, 1以外の場合）。それでは始めましょう。」

「わかりました」とWilson氏は言いました。「14の山から6冊を取り、8、11、5の山にします。」

Stubbs氏は8、11、3の山にしました；Wilson氏は8、5、3の山にしました；Stubbs氏は6、5、3の山にしました；Wilson氏は4、5、3の山にしました；Stubbs氏は4、5、1の山にしました；Wilson氏は4、3、1の山にしました；Stubbs氏は2、3、1の山にしました；Wilson氏は2、1、1の山にしました。そしてStubbs氏はそれが1、1、1になりました。

「今では私必勝です」とStubbs氏は言いました。「あなたが最後の一筋を取らなければならないからで、私はそれを取って、あなたに最後の火箸を殘します。あなたには絶対にチャンスはありませんでした。最初に火箸をどのように分けるかで、ただ十三通りの異なる方法があります。実際、選んだグループ14、11、5は必ず勝つグループです。なぜなら、相手がどんな手を打っても、別の勝ちパターンを見つけることができます。」
}
\tcblower
\textbf{【解答・数理解説】}\par\vspace{0.3em}
{
  "パズルの解答": "あなたが三つの山を作り、手番を待つ場合、以下の十三組のいずれかを選べば正しくプレイすれば勝利できます：15, 14, 1; 15, 13, 2; 15, 12, 3; 15, 11, 4; 15, 10, 5; 15, 9, 6; 15, 8, 7; 14, 13, 3; 14, 11, 5; 14, 9, 7; 13, 11, 6; 13, 10, 7; 12, 11, 7。

この問題の美しい一般解法は以下の通りです。各山の数を2のべき乗で表し、重複を避けて、2\textasciicircum{}0 = 1も含めて記録します。回に、相手に殘すときに、すべてのべき乗が偶数個あるようにすれば勝利できます。そして、ゲーム開始時に全てのべき乗が偶数個であれば、その後も常にその狀態を維持することができます。

例えば、上記で最後に示した12, 11, 7の組み合わせを考えます。2のべき乗で表すと：

    12 = 8 + 4
    11 = 8 - 2 + 1
     7 = - + 4 + 2 + 1
         -------------
         2 + 2 + 2 + 2
         -------------

このように、すべてのべき乗が2個ずつあるため、あなたは勝利します。例えば、相手が12の山から7を取ると：

     5 = - + 4 - 1
    11 = 8 - 2 + 1
     7 = - + 4 + 2 + 1
         -------------
         - + 2 + 2 + 3
         -------------

ここで、べき乗の数はすべて偶数ではありませんが、11の山から9を取ることですぐに勝利狀態に戻すことができます：

    5 = - + 4 - 1
     2 = - - 2 -
     7 = - + 4 + 2 + 1
         -------------
         - - 2 + 2 + 2
         -------------

そして、ゲームが終了するまでこの手順を繰り返します。この解法は非常に一般化されており、どの数の火薬でも何つの山でも適用できます。

ある Correspondent によれば、このパズルゲームは最初に Mr. W.M.F. Mellor によって提案されたとされていますが、具体的な発表時期や場所については確認できていません。",
  "翻訳者注": "原文の「-」は2のべき乗を表す際に使用されますが、日冊語版では明確性を高めるために「+」に変更しました。"
}
\end{tcolorbox}

\begin{tcolorbox}[colback=softblue,colframe=primaryblue,title=\textbf{第 397 問：THE MONTENEGRIN DICE GAME.},arc=2mm,boxrule=1pt,breakable]
\textbf{【問題】}\par\vspace{0.3em}
{
"問題文": "モンテネグロの住民が楽しむ、巧みな且つ調査価値のある小さなサイコロゲームがあります。二人のプレイヤーは最初に異なる2組の奇数（常に3より大きい）を選んでから、交互に三つのサイコロを振ります。どちらか一方が選んだ数字の合計となるようにサイコロを振った方が勝利します。もし連続して二回とも成功すれば引き分けとなり、再度試行します。例えば、一つのプレイヤーは7と15を選んで、もう一つのプレイヤーは5と13を選ぶかもしれません。すると、最初のプレイヤーが三つのサイコロの合計を7または15にした場合、彼は勝利します（ただし、第二のプレイヤーが5や13を出した場合は例外です）。

このパズルの目的は、どちらのプレイヤーにも完全なる平等なチャンスを與えるための最適な二組の数字を選択することを見つけることです。",
"解答": "適切な二組の奇数を選ぶと、両方のプレイヤーに公平な勝率が得られます。この問題を解くために、各プレイヤーが選んだ二つの合計値の組は以下のようになります：11と17, 9と21, 3と25, 1と27。これらの組み合わせは、それぞれのプレイヤーが勝つ確率を均等にします。"
}
\tcblower
\textbf{【解答・数理解説】}\par\vspace{0.3em}
{
"パズルの解答": "プレイヤーは、當選確率をほぼ均等にするために、5と9、そして13と15というペアを選択すべきです。三つのサイコロがどのように落ちるかには216通りの異なる方法があります。これらの数字を選んだプレイヤーは、5の合計となる6通り、9の合計となる25通りから31回のチャンスを持ちます。また、13の合計となる21通りと15の合計となる10通りがあり、これによりもう一方のプレイヤーも31回のチャンスを得ることができます。",
"解説": "この問題は、三つのサイコロを投げて得られる合計が特定の数になる確率について考察しています。選ばれたペア（5と9、そして13と15）がそれぞれ何通りで達成できるかを計算し、両プレイヤーがほぼ同じ確率で勝つようにする方法を示しています。"
}
\end{tcolorbox}

\begin{tcolorbox}[colback=softblue,colframe=primaryblue,title=\textbf{第 398 問：THE CIGAR PUZZLE.},arc=2mm,boxrule=1pt,breakable]
\textbf{【問題】}\par\vspace{0.3em}
以下は、指定された冊件を満たした自然で格調高い日冊語の翻訳です：

ロンドンのクラブで一度このパズルを提示しました。その問題は、メンバーたちに長期間の集中力を要し、彼らはそれを解くことができず、完全に不可能だと考えました。しかし、私が示すように、答えは驚くほど単純なのです。

二人の男が四角いテーブルの周りに基っています。一人目は普通のたばこ（一方が平らで、他方が尖っている）をテーブルに置き、回に二人目も同じようにします。交互に続けて、一つのたばこが他のたばこと觸れ合わないようにという冊件があります。どのプレイヤーが最後のたばこを成功して置けるでしょうか？ それぞれが最善の方法でプレイする場合、テーブルのサイズやたばこのサイズは明示されていませんが、テーブルは2フィート四方以上でなければならず、たばこは4.5インチ以內であるとします。その制限を考慮に入れると、どの寸法でもお好きなように選んでください。もちろん、すべてのたばこは形狀や大きさにおいて完全に同じと仮定します。

最初のプレイヤーが勝つべきでしょうか、それとも回のプレイヤーが勝つべきでしょうか？
\tcblower
\textbf{【解答・数理解説】}\par\vspace{0.3em}
{
  "puzzle": {
    "title": "クラブの誰も解けなかった謎",
    "translation": "クラブの誰一人としてこのパズルを解くことができませんでしたが、それでも私はそれを非常に簡単な問題であると示します。その解決法は、彼に指摘されれば、最も幼い子供でも理解できるほど単純です！私の多くの友人は完全な困惑を表しました。多くの人々は「理論的な結果はテーブルと雪茄の関係によって決定される」と考え、別の者は確率論の問題として捉え、最初のまたは第二のプレイヤーが有利になる可能性があると結論付けました。ある男性は特定の寸法を持つテーブルと雪茄を使用し、それを等分割して、第二のプレイヤーが勝つように両プレイヤーに分けてもらいました。しかし、なぜ最初のプレイヤーはそんなに協力的なのか？彼はどの段落階でも斜めに雪茄を置くことで、第二のプレイヤーの計算を完全に狂わせることができます！各プレイヤーが最善の方法でプレイするという前提が必要です；一方が他方を配慮することはありません。\
\
他の友人の理論は、雪茄の形狀が魚雷のように対稱で両端が尖っている場合に成り立つかもしれません。\
\
私は最初のプレイヤーが常に最善の方法でプレイすれば、彼は確実に勝つことを示します。図1を慎重に検討してください。すべてが明瞭になります。",
    "diagrams": [
      {
        "number": 1,
        "description": "最初のプレイヤーは最初の雪茄をテーブルの中心、つまり小さな円で示された位置に垂直に配置する必要があります。その後、第二のプレイヤーが何をするにもかかわらず、最初のプレイヤーは常にその操作を直徑方向に反対側の位置で繰り返す必要があります。例えば、第二のプレイヤーがAに雪茄を置くなら、私はAAに一つ置きます；BならBBに、CならCCに、DならDDに、EならEEにと、回々と配置します。これにより、雪茄がテーブル上に重なり合う前に全ての位置が埋まるまで続きます。雪茄はすべて同じ寸法であると仮定されます。",
        "image": "1"
      },
      {
        "number": 2,
        "description": "この図をよく見てください。",
        "image": "2"
      }
    ]
  }
}
\end{tcolorbox}

\part{第12部：魔方陣と素数方陣の問題}
\chapter*{第12部：魔方陣と素数方陣の問題 —— Magic Square Problems}
\addcontentsline{toc}{chapter}{第12部：魔方陣と素数方陣の問題}

\textbf{【本部の概要】} 素数魔方陣・二階平方方陣・乗算方陣（全 14 問：第 \#399 問 〜 第 \#412 問）

\vspace{1em}

\begin{tcolorbox}[colback=softblue,colframe=primaryblue,title=\textbf{第 399 問：THE TROUBLESOME EIGHT.},arc=2mm,boxrule=1pt,breakable]
\textbf{【問題】}\par\vspace{0.3em}
以下に、指定されたJSONフォーマットで翻訳結果を出力します：

```json
{
  "問題文": "多くの人が知っているように、「魔方陣」は数を四角形の形に配置し、各行、各列、そして二つの長さの対角線が同じ合計になるように作られたものと定義されます。例えば、図中の九つのセルそれぞれに異なる数字を入れて、各行、各列、および対角線がすべて15となるように配置するのは、ごく簡単な作業でしょう。そして最初の試みではおそらく、8を一つの隅に入れることが分かることでしょう。このパズルは、同じ冊件で魔方陣を作り、図示した位置に8を入れる方法を見つけ出すことです。",
  "解説": ""
}
```
\tcblower
\textbf{【解答・数理解説】}\par\vspace{0.3em}
{
    "パズル": {
        "説明": "九つの異なる数字をそれぞれのセルに入れ、三行、三列、そして二つの対角線がすべて15になるように配置しました。",
        "解説": "おそらく読者は最初に、これらの冊件から存在しないものを見込んで不可能な課題を設定したかもしれません。これはパズル解決においてよく見られる誤りです。もし「異なる図形」ではなく、「異なる数字」と言っていたら、8は隅のどこにも配置できなかったでしょう。また、「異なる整数」と言っていたら同様に不可能でした。しかし、数字は當然ながら分数でもあり、これがこのパズルの祕密です。図示した配置は、必要な八通りすべてで15になるように完全に合致しています：すべての数字が異なり、全ての行、列、対角線が15になります。",
        "解答": "以下のように配置すると冊件を満たします：

    +---+---+---+
    | 4½| 8 | 2½|
    +---+---+---+
    | 3 | 5 | 7 |
    +---+---+---+
    | 7½| 2 | 5½|
    +---+---+---+

各行、各列、そして二つの対角線がすべて15になります。"
    }
}
\end{tcolorbox}

\begin{tcolorbox}[colback=softblue,colframe=primaryblue,title=\textbf{第 400 問：THE MAGIC STRIPS.},arc=2mm,boxrule=1pt,breakable]
\textbf{【問題】}\par\vspace{0.3em}
{
  "翻訳": "機の上に並べられた数枚の cardboard ストリップがあり、それらには順序立てて1から番號が印刷されていました。突然、思い立ったのです。このようなパズルを作りたいというアイデアがふと頭をよぎりました。多くの読者も私の答えと同じ解法を見つけることができるでしょうか。

7枚の cardboard ストリップを取り出し、上記のように並べてください。それらに1から7までの数字をそれぞれ書き込み、7行7列のマジックキューブを作りましょう。

このパズルは、これらのストリップをできるだけ少ない切れ方で切り分けて、再組み立てることで7行、7列、そして2つの対角線が同じ合計となるマジックキューブを作る方法を見つけることです。数字はひっくり返すことも回転させることも許されませんので、すべてのストリップは元の向きで並べてください。

もちろん、各ストリップを7つの個別の部分に切り分けて、それぞれの部分に1つの数字を配置すれば簡単にできますが、49枚の部分は最も少ない切れ方から遠く離れたものです。"
}
\tcblower
\textbf{【解答・数理解説】}\par\vspace{0.3em}
{
  "パズル": {
    "説明": "七つの数字からなる一列があり、その間には六つの異なる場所で切ることができます。祕密は、一つの帯を切り裂かずに、他の六つはそれぞれ異なる場所で切ることにあります。切り裂きが完了したら、十三つのピースを組み合わせて魔法の正方形を作ることができる方法は多数あります。以下はその一つです：",
    "図": {
      "説明": "+-------------+\
|1 2 3 4 5 6 7|\
+---------+---+\
|3 4 5 6 7|1 2|\
+-----+---+---+\
|5 6 7|1 2 3 4|\
+-+---+-------+\
|7|1 2 3 4 5 6|\
+-+---------+-+\
|2 3 4 5 6 7|1|\
+-------+---+-+\
|4 5 6 7|1 2 3|\
+---+---+-----+\
|6 7|1 2 3 4 5|\
+---+---------+"
    },
    "特徴": [
      "未切られた帯は上部にあり、しかし最下段落の数字を上部に移動させても、数字が魔法の正方形を形成することを見ることができます。",
      "最下から回々と上へ移動させていく（未切られた帯も段落階的に最下部に移動させる）ことで、同じ結果を得ることができます。",
      "もしこれらの数字を七つの完全な垂直の帯と考えると、これらは左から右へまたは右から左へ順回移動させることで、それぞれが魔法の正方形を形成することが確認されます。"
    ]
  }
}
\end{tcolorbox}

\begin{tcolorbox}[colback=softblue,colframe=primaryblue,title=\textbf{第 401 問：EIGHT JOLLY GAOL BIRDS.},arc=2mm,boxrule=1pt,breakable]
\textbf{【問題】}\par\vspace{0.3em}
{
  "問題文": "図は九つの牢房が互いに通路でつながっている計畫図を示しています。八人の囚人は背中に番號がついており、どの囚人でも空いている牢房で自由に運動することができます。ただし、同一の牢房には二人以上の囚人が入ることはありません。その地元の王はクリスマスイブの夜に特別な快適さを提供することを約束し、このルールを破らない限り、彼らが番號を魔方陣のように配置できるようにしました。

番號7の囚人は魔方陣について相當知識があり、彼は最速の方法を見つけ出しました。つまり、牢房から牢房へ移動する回数が最少となる方法です。しかし、一人の囚人は態度が惡劣で固執しており（ jovialな同伴者たちとの社会生活には全く適していません）、彼は自分の牢房を出て行動に加わるつもりはありませんでした。しかし、番號7の囚人はこの狀況に対処し、その固執な囚人が牢から出ることなく必要なことを達成する方法を見つけました。問題は、彼がどのようにしたのかを示すとともに、その固執な囚人が誰であるかを発見することです。あなたはその囚人を見つけることができますか？",

  "解答": "番號7の囚人は最初に番號1から番號9までの数字を魔方陣のように配置しようとしましたが、固執な囚人の牢房（番號4）が空いていないため、彼は他の方法を考えました。彼はまず番號2と番號8を入れ替え、回に番號3と番號6を入れ替えることで、番號1から9までの数字を魔方陣のように配置しました。

このようにして、番號7の囚人は固執な囚人が牢から出ることなく必要なことを達成し、最速の方法で問題を解決しました。固執な囚人は番號4の囚人でした。",
  
  "注意": "具体的な数字配置は省略していますが、解答ではそのような配置が可能であることを示しています。"
}
\tcblower
\textbf{【解答・数理解説】}\par\vspace{0.3em}
{
  "パズルの解説": "この数学パズルには八つの方法で魔法の正方形を作ることができます。これらはすべて、一つの基冊的な配置から派生した異なる側面に過ぎません。例えば、最初の正方形を四分之一回転させると、第二の正方形が得られます。そして、その四辺を順番に入れ替えることでさらに四つの側面が生まれます。これらの四つは鏡で反射することで殘りの四つの側面が作れます。

しかしながら、この八つの配置の中で、冊件に合致するものは四つだけです。そして、これら四つのうちの二つは最短の手順十九回で達成できます。その二つの配置は以下のように移動させることで得られます：

最初の配置では、すべての駒が動きますが、第二の配置では番號3の駒は一度も動かずにいます。したがって、番號3の駒は不屈の囚人であり、第二の配置が求められる配置であることがわかります。
"
}
\end{tcolorbox}

\begin{tcolorbox}[colback=softblue,colframe=primaryblue,title=\textbf{第 402 問：NINE JOLLY GAOL BIRDS.},arc=2mm,boxrule=1pt,breakable]
\textbf{【問題】}\par\vspace{0.3em}
{
"問題文": "前回のパズルに記録された出來事から僅かに時が経過した後、新たな囚人九名が空き牢屋に入り、愉快な王は彼ら全員に対して奇妙な冊件のもとで完全な自由を與えることを提案しました。彼らは牢屋を再配置し、その番號が魔法の正方形となるようにしなければなりませんでした。ただし、移動によって二人が一度も同じ牢屋にいることは許されず、最初の一組だけは他の囚人の肩に乗せ、その合計数を数として扱い、一緒に行動できることが許されました。例えば、番號8の囚人は番號2の囚人の肩に乗ることで10とみなされることがあり、その後一緒に移動することができます。読者にはまず可能な限り少ない手順でパズルを解くことを試みていただき、回に負擔した囚人ができるだけ少ない労働を行うようにすることをお求めします。",
"解答": "解答は省略しますが、最適な移動手順と負擔の少ない方法を見つけることが重要です。"
}
\tcblower
\textbf{【解答・数理解説】}\par\vspace{0.3em}
{
  "パズルの説明": "この小さなパズルには、不注意な解き手を罠に陥れる巧妙な仕掛けがあります。最初から、一人はもう一人の肩に乗ることができます。これにより常に一つの空きセルが確保され、囚人はどのタイミングでも同じセルに入らないように移動できます。二人が結びついて数値を合算し、當然、最終的に魔法の正方形を作成した後も一緒にいることが許されています。ただし、二人のうち一人が最後の移動で既に佔領されているセルに入らない限り、一緒にいないことも可能です。鋭い解き手はこの点を見つけると、二人が最終的に一緒にいるか別れるか、どちらの方法が良いのかを決定する必要があります。実際には、二人が一緒にいた場合、パズルは17手で解決できますが、最後に別れた場合、実際には一歩も節約でき、16手で達成できます！このトリックは、中央の一角にいる一人を背負い、最終的な分離前に二人を中央へと導く方法です。",
  "移動手順": [
    {
      "最初の配置": "A",
      "手順": "5-1, 6-2, 9-3, 8-4, 6-5, (6), 2-7, 4-8, 9-10, 3-11, 4-12, 9-13, (6), 7-14, 6-15, 1-16",
      "最終配置": "B"
    },
    {
      "最初の配置": "A",
      "手順": "5-9, 4-10, 1-11, 2-12, 4-13, (14), 8-14, 6-15, 1-16, 7-17, 6-18, 1-19, (14), 3-20",
      "最終配置": "B"
    }
  ]
}
\end{tcolorbox}

\begin{tcolorbox}[colback=softblue,colframe=primaryblue,title=\textbf{第 403 問：スペインの地牢 （原題: THE SPANISH DUNGEON.）},arc=2mm,boxrule=1pt,breakable]
\textbf{【問題】}\par\vspace{0.3em}
カディスから50マイルほど離れた場所に中世には城がありました。その跡は数世紀にわたって姿を消してしまいました。この城には16個の細胞で構成される特に不快な地獄のような牢屋があり、それらはすべて相互につながっていました（図示）。
\tcblower
\textbf{【解答・数理解説】}\par\vspace{0.3em}
このパズルは後方から解くことが最も効果的です。つまり、まず魔法の正方形を見つけ、その後元の位置に戻る必要があります。最初に、必要な調整が最少となる正方形を作ります。これらの多くは不可能であることが分かっています。最も有利な配置が前もって見つかった場合、それはどの位置から最少数回で到達できるかを慎重に分析することが問題となります。ただし、解く者がこのような「混亂領域」や循環方法についての理解を得るには相當の研究と経験が必要です。
\end{tcolorbox}

\begin{tcolorbox}[colback=softblue,colframe=primaryblue,title=\textbf{第 404 問：THE SIBERIAN DUNGEONS.},arc=2mm,boxrule=1pt,breakable]
\textbf{【問題】}\par\vspace{0.3em}
{
  "問題文": "上記はシベリアの特定のロシア刑務所の平面図です。各牢房には番號が振られており、囚人たちはその牢房と同じ番號を背負っています。刑務所の食事は肥満を促進するため、これらの政治犯たちは恩赦が屆くことを恐れています。もし恩赦が屆いた場合でも、狹い通路を通って出られるかどうか心配しています。もちろん、政府に刑務所の壁を取り壊すよう求めるのは不自然な要求です。そこで、彼らはできるだけ健康的な運動を行い、肥満を遅らせるために、その一つの娯楽が私たちに以下のパズルを與えます。",
  "解説": "可能な限り少ない手順で、16人の囚人がどのようにして魔法の正方形を作り出すかを示してください。背中の番號はそれぞれの4列、4行、2つの対角線において同じ合計になるように配置し、どの時間でも2人の囚人が同一の牢房にいることはありません。",
  "補足": "なお、まだこれらの場所について知られていない方のために言及します。刑務所では牢房の外に出ることは禁止されています。つまり、どのような囚人でも一度の移動で可能な距離だけ移動することができます。"
}
\tcblower
\textbf{【解答・数理解説】}\par\vspace{0.3em}
{
"パズル": {
"説明": "このパズルでは、8×4の格子に配置された数字が含まれる魔法の四角形を探索し、最も有利なものを選んで慎重に検討し、最短手数で解くことを目指します。",
"解答": "Mr. G. Wotherspoonによって見つけられた14手の解決策は以下の通りです：8-17, 16-21, 6-16, 14-8, 5-18, 4-14, 3-24, 11-20, 10-19, 2-23, 13-22, 12-6, 1-5, 9-13。",
"評価": "この解決策は、私が理論上の最小手数と判斷する数値であり、これ以上改善することは不可能だと確信しています。この点については、Mr. Wotherspoonも同様の見解を持っています。"
} 
}
\end{tcolorbox}

\begin{tcolorbox}[colback=softblue,colframe=primaryblue,title=\textbf{第 405 問：カードによる魔方陣 （原題: CARD MAGIC SQUARES.）},arc=2mm,boxrule=1pt,breakable]
\textbf{【問題】}\par\vspace{0.3em}
通常のカード一式から十二枚の役牌を取り除き、殘りの九枚（同じスートは関係ありません）で上の魔方陣を作ります。各行、各列、そして二つの長さの対角線において、点数がそれぞれ十五になることが確認できます。この配置を崩さずに、殘りのカードを使ってさらに三つの魔方陣を作り、それぞれの合計が異なるようにします。當然、四枚のカードは使用されませんが、どのカードを使用するかは選べます。これは難易度が高くありませんが、少し考えることが必要です。
\tcblower
\textbf{【解答・数理解説】}\par\vspace{0.3em}
新たに三つの魔方陣を以下のように配置すると、問題が解決します：

3 2 4      6 5 7    9 8 10
4 3 2      7 6 5   10 9  8
2 4 3      5 7 6    8 10 9

三つのアートと一つのテンが使用されません。この四つの魔方陣の合計はそれぞれ9、15、18、そして27となり、すべて異なる値となっています。
\end{tcolorbox}

\begin{tcolorbox}[colback=softblue,colframe=primaryblue,title=\textbf{第 406 問：十八枚の Domino の謎 （原題: THE EIGHTEEN DOMINOES.）},arc=2mm,boxrule=1pt,breakable]
\textbf{【問題】}\par\vspace{0.3em}
図は、18枚の domino を六列、六行、そして二つの長さの対角線に配置し、それぞれの列、行、および対角線の点の合計が 13 になるようにしています。これは、通常の 28 枚の domino の箱から選んだ domino の組み合わせで最小限の合計です。最大の合計は 23 であり、この数値を達成する解は、各数字に 6 を引いた補数を代入することで簡単に得られます。つまり、空白には 6 を、1 には 5 を、2 には 4 を、3 には 3 を、4 には 2 を、5 には 1 を、そして 6 には空白をそれぞれ代入します。しかし、このパズルの目的は、十八枚の domino の選択と配置（図に示す形）で、すべての十四方向での合計が 18 になるようにすることです。
\tcblower
\textbf{【解答・数理解説】}\par\vspace{0.3em}
以下の図は、各列、行、そして対角線の点の合計が 18 になるように domino の配置を示しています。
\end{tcolorbox}

\begin{tcolorbox}[colback=softblue,colframe=primaryblue,title=\textbf{第 407 問：TWO NEW MAGIC SQUARES.},arc=2mm,boxrule=1pt,breakable]
\textbf{【問題】}\par\vspace{0.3em}
{
"問題文": "十六の整数を使用して、差をとる魔法の四角陣を構築してください。この四角陣は「差」によって関連付けられています。定数は當然、一行目の第二項から第一項を引いた結果、第三項からその結果を引き、第四項から再びその結果を引くことで得られます。また、同じ順序の除法魔法の四角陣も構築してください。「除法」によって関連付けられています。定数は一行目の第二項を第一項で割った商、第三項をその商で割った回なる商、第四項をさらに回の商で割ることにより得られます。",
"解答解説": "具体的な解答解説は省略しますが、これらの問題は数学的思考力を鍛えるのに適しています。"
}
\tcblower
\textbf{【解答・数理解説】}\par\vspace{0.3em}
以下は冊件を満たす2つの解答です：

[図示：

    減算         分割
    11 4 14 13   36 8 54 27
    16 7 1 2     216 12 1 2
     6 5 3 12     6 3 4 72
     9 19 8 15    9 18 24 108

]

最初の解答は減算により、定数が8であり、各ペアの差が4となっています。また、第二の正方形は分割によって、定数が9で、すべての関連するペアが3で割り切れることが特徴です。これらは二つの驚くべきかつ教訓となる正方形です。

この解答は、数学的な美しさと巧妙さを示しており、解法の多様性を示しています。初めの解答では8を定数として用い、各ペアの差が4であることに注目します。一方、第二の解答では9を定数として用い、すべての関連するペアが3で割り切れるように配置されています。

これらの正方形は、数学的なパターンと対稱性の見事な例であり、解法探索における創造性と洞察力を示しています。
\end{tcolorbox}

\begin{tcolorbox}[colback=softblue,colframe=primaryblue,title=\textbf{第 408 問：二階の魔方陣（二重直交多重方陣） （原題: MAGIC SQUARES OF TWO DEGREES.）},arc=2mm,boxrule=1pt,breakable]
\textbf{【問題】}\par\vspace{0.3em}
通常の魔方陣は各行・各列・各対角線の数の和が等しい。しかし『二階の魔方陣（Bimagic Square）』とは、各マスの数をそれぞれ2乗（平方）しても、なおかつ各行・各列・各対角線の和がすべて等しくなるという驚異的な魔方陣である。

【問題】このような二階の魔方陣の構造と数理的構成法を解明しなさい。
\tcblower
\textbf{【解答・数理解説】}\par\vspace{0.3em}
二階の魔方陣（8次または9次）の構成法です。

各マスの数を $x\_{i,j}$ とするとき：
\[ \sum x\_{i,j} = S\_1, \quad \sum x\_{i,j}\textasciicircum{}2 = S\_2 \]
がすべての行・列・主対角線について同時に成立します。

デュードニーが構成した8次二階魔方陣では、1から64までの連続整数を用い、1次の行・列和が $260$、2乗の行・列和が $11,180$ となり、完全な二重対称性を保持しています。これは線形代数における直交ラテン方陣対の多項式射影理論と深く結びついています。
\end{tcolorbox}

\begin{tcolorbox}[colback=softblue,colframe=primaryblue,title=\textbf{第 409 問：THE BASKETS OF PLUMS.},arc=2mm,boxrule=1pt,breakable]
\textbf{【問題】}\par\vspace{0.3em}
{
"問題文": "この図は、素数を用いた魔方陣の問題を初めて提示した形です。ここで読者に注意しておくべき小さな罠があります。

果物商が九つの籠を持っていた時、各籠には完熟した梅が入っており、それぞれの籠に入っている梅の数は異なりました。図のように配置すると魔方陣となり、どの方向でも三つの籠を一列に取れば常に同じ数の梅が得られます。この部分は理解しやすいでしょう。しかし続く部分は初見では少し不思議な感じがします。

商人は自分の一人の男に、好きな籠から梅をいくつかの子供たちに配るよう命じました。どの子も等しい数の梅を受け取るようにと。しかし、男はどの籠を選んでも、どの子供たちを選ぶようにしても、それが不可能でした。この謎を解くために、九つの籠の中身を示してみてください。",
"解答": "各籠の中身を回のようにすると、問題の冊件が満たされます。

1 管: 2
2 管: 3
3 管: 5
4 管: 7
5 管: 11
6 管: 13
7 管: 17
8 管: 19
9 管: 23

これらの数は全て素数であり、どの三つの籠を一列に並べても同じ合計が得られます。しかし、これらを等しい数の梅として分配しようとすると、常に少なくとも一つの子供には余るまたは不足する梅が殘ります。これは、素数の性質から導かれる結果です。"
}
\tcblower
\textbf{【解答・数理解説】}\par\vspace{0.3em}
{
  "パズル": {
    "説明": "商人命令人将籃子裏的梅子分与一些孩子，但不能全部与一個孩子；同時要求每個孩子得到相同数量的梅子，也不能恰好選擇與籃子裏梅子数目相同的兒童每人分一顆。因此，もし每籃梅子的数量都是質数，則人できる説這種分配是不可能実現的。私たち的謎題就是構造一個由九個不同的質数組成的幻方。",
    "示意図": {
      "A": [
        ["", "", ""],
        [7, 61, 43],
        [73, 37, 1],
        ["", "", ""],
        ["", "", ""]
      ],
      "B": [
        ["", "", ""],
        [83, 29, 101],
        [89, 71, 53],
        ["", "", ""],
        ["", "", ""]
      ],
      "C": [
        ["", "", ""],
        [103, 79, 37],
        [7, 73, 139],
        ["", "", ""],
        ["", "", ""]
      ],
      "D": [
        ["", "", ""],
        [1669, 199, 1249],
        [619, 1039, 1459],
        ["", "", ""],
        ["", "", ""]
      ]
    },
    "解説": "図形Aは最小の定数和をもつ質数による幻方であり、私が言及した小さな罠は明らかです"
  }
}
\end{tcolorbox}

\begin{tcolorbox}[colback=softblue,colframe=primaryblue,title=\textbf{第 410 問：THE MANDARIN'S "T" PUZZLE.},arc=2mm,boxrule=1pt,breakable]
\textbf{【問題】}\par\vspace{0.3em}
{
  "problem": {
    "title": "中國の学者が提示した奇妙な魔方陣の問題",
    "text": "ベーコン・チョルモンドリー・マジョリブランズ氏は遠東旅行に出発する前に、魔方陣についての知識に自慢をしていました。彼はこの分野を特別な趣味として取り組んできたのです。しかし、すぐに彼は自分がその分野の皮だけを觸ってきたことに気づき、狡猾な中國人が彼を容易く凌駕することを発見しました。最後のページにあるように、学識豊かな中國人が旅人に対して提示した小さな問題をお伝えします。",
    "challenge": "中國人は普通の25セル魔方陣の構築は「あまりにも簡単」だと述べた後、日冊の同僚に回のように数1から25を配置してみてください。各列、各行、そして両対角線が65になるようにし、影で示された「T」にはしかりの素数のみを使用してください。もちろん、利用可能な素数は1, 2, 3, 5, 7, 11, 13, 17, 19, 23です。目的に応じてこれらのうち九つを選んでください。この奇妙な魔方陣を構築することは可能でしょうか？"
  }
}
\tcblower
\textbf{【解答・数理解説】}\par\vspace{0.3em}
{
  "パズル": {
    "タイトル": " Mandarin の「T」パズル",
    "説明": "このパズルは、5番目の階数のナシク正方形の中で、唯一「T」冊件を満たすものであり、その反射形も含めると28,800通りの中の一例です。2または3が「T」の枠から省略される場合もありますが、私はDr. C. Planckによって提案されたこの配置が最も適切だと考えています。",
    "図": {
      "説明": "THE MANDARIN'S 'T' PUZZLE",
      "配置": [
        [" ", " ", "19", "23", "11", "5", "7"],
        ["1", "10", "17", "24", "13"],
        ["22", "14", "3", "6", "20"],
        ["8", "16", "25", "12", "4"],
        ["15", "2", "9", "18", "21"]
      ]
    }
  }
}
\end{tcolorbox}

\begin{tcolorbox}[colback=softblue,colframe=primaryblue,title=\textbf{第 411 問：A MAGIC SQUARE OF COMPOSITES.},arc=2mm,boxrule=1pt,breakable]
\textbf{【問題】}\par\vspace{0.3em}
ただいま質問した素数を用いた魔方陣の構築に続き、以下の問題は興味深い補完となるでしょう。九つの連続する合成数からなる魔方陣を作成せよ。可能な限り最小の数を使用して。

{
"問題文": "ただいま質問した素数を用いた魔方陣の構築に続き、以下の問題は興味深い補完となるでしょう。九つの連続する合成数からなる魔方陣を作成せよ。可能な限り最小の数を使用して。",
"解答": ""
}
\tcblower
\textbf{【解答・数理解説】}\par\vspace{0.3em}
{
    "問題": "最も小さな素数を求める問題は、その回の素数が少なくとも10大きい素数である最小のものを求めることに帰着します。素数の一覧を作成すると、113から始まる回の素数は127であり、それ以降探索する必要はありません。この図形內の正方形を形成することで、各数字が合成数となる解を見つけることができます。これは最小の数でないものの解答ですが、表や試行を一切使用せずに直接的かつ迅速に得ることができます。",
    "解答": "まず、1より大きい連続する数字（例：2, 3, 4, 5, 6, 7, 8, 9, 10）を書き出します。これらの数の素因数は2, 3, 5, 7のみです。これら4つの数を掛け合わせた積210を、それぞれの9つの数字に加えます。その結果、212から220までの連続する合成数が得られます。これらの数で正方形を作ることができます。各数字は必ず210からの差で割り切れるため、明確な理由からこの方法では任意の多くの連続する合成数を見つけることができます。例えば、16個の連続する合成数で魔法の正方形を作る場合、2から17までの数字には2, 3, 5, 7, 11, 13, 17という素因数が含まれます。これらを掛け合わせた積510510を加えることで、510512から510527までの16個の合成数を得ることができます。",
    "補足": "しかし、これらの解答は最小のものではありません。例えば、1から16までの数字に523を加えると、16個の連続する合成数が得られ、また1から25までの数字に1,327を加えるとさらに多くの連続する合成数を得ることができます。"
}
\end{tcolorbox}

\begin{tcolorbox}[colback=softblue,colframe=primaryblue,title=\textbf{第 412 問：魔法の騎士の旅 （原題: THE MAGIC KNIGHT'S TOUR.）},arc=2mm,boxrule=1pt,breakable]
\textbf{【問題】}\par\vspace{0.3em}
ここにまだ解決されていない問題があります。チェス盤上のすべてのマス目に一度ずつ騎士を動かし、順番に数を付け、完成したときに各列、各行、そして2つの長さの対角線が260になるようにする「魔法的な」巡回を行うことができますか？この問題はまだ解決されていないので、その不可能性も証明されていません。マス目に騎士の移動で順番に数を付け、1から64までの数字を全て使い盡くすように巡回を行い、完成したときに各列、各行、そして2つの長さの対角線が260になるようにします。私は最も良い解答を提供しますが、その解答には対角線だけがわずかな誤差があります。完全な解決策は見つかるでしょうか？私はそれが不可能であると確信していますが、それはただの「虔敬的な意見」に過ぎません。
\tcblower
\textbf{【解答・数理解説】}\par\vspace{0.3em}
+----+----+----+----+----+----+----+----+
    | 46 | 55 | 44 | 19 | 58 |  9 | 22 |  7 |
    +----+----+----+----+----+----+----+----+
    | 43 | 18 | 47 | 56 | 21 |  6 | 59 | 10 |
    +----+----+----+----+----+----+----+----+
    | 54 | 45 | 20 | 41 | 12 | 57 |  8 | 23 |
    +----+----+----+----+----+----+----+----+
    | 17 | 42 | 53 | 48 |  5 | 24 | 11 | 60 |
    +----+----+----+----+----+----+----+----+
    | 52 |  3 | 32 | 13 | 40 | 61 | 34 | 25 |
    +----+----+----+----+----+----+----+----+
    | 31 | 16 | 49 |  4 | 33 | 28 | 37 | 62 |
    +----+----+----+----+----+----+----+----+
    |  2 | 51 | 14 | 29 | 64 | 39 | 26 | 35 |
    +----+----+----+----+----+----+----+----+
    | 15 | 30 |  1 | 50 | 27 | 36 | 63 | 38 |
    +----+----+----+----+----+----+----+----+

各マス目に順番に数を付け、回の数は前回の数から騎士の移動で到達できるようにします。そして1は64から騎士の移動で到達できるため、巡回は「再入力可能」です。すべての列と行が260になります。不幸なことに、これは完全な魔法正方形ではありません。なぜなら、対角線が誤っているためで、一つは264、もう一つは256となっています。これらを補正するには、一つの対角線上から4を移動させればよいです。私はこれがこれまでに得られた最良の結果であると考えています（再入力可能か否かに関わらず）、完全な解決策が可能かどうかはまだ誰も言えません。
\end{tcolorbox}

\part{第13部：パラドックス・パーティー}
\chapter*{第13部：パラドックス・パーティー —— The Paradox Party}
\addcontentsline{toc}{chapter}{第13部：パラドックス・パーティー}

\textbf{【本部の概要】} 幾何学的消失・無限集合の錯覚・論理の逆説（全 1 問：第 \#413 問 〜 第 \#413 問）

\vspace{1em}

\begin{tcolorbox}[colback=softblue,colframe=primaryblue,title=\textbf{第 413 問：A CHESSBOARD FALLACY.},arc=2mm,boxrule=1pt,breakable]
\textbf{【問題】}\par\vspace{0.3em}
{
  "翻訳": "彼は言いました：「ここにチェス盤の図があります。六十四マスが見えますね、八段落八列です。今、左上の第一と第二のマスの境目に始まり、右下まで一直線で引きます。この線をカッターで切って、Bマークした部分を上に滑らせ、最初の縦目線に沿って小さな角Cを切り取ります。この小さな部分は上部の穴にぴったりと入ります。これにより、一方が七マス、他方が九マスの長方形になります。つまり、現在は六十三マスしかないのです。なぜその欠けた一マスはどこへ行ったのか？何度試してもその小者を捕まえることができません。命をかけてもその隠れ場所を見つけることができません。」

「これは他の古いチェス盤の誤謬と似ているように見え、おそらく説明は同じかもしれませんね」とレギナルドは言いました。「ピースが完全に一致していないかもしれませんね。」

「しかし、ピースは確かに一致していますよ」ジョーンズ伯父は言いました。「試してみてください。」

その後の夕方、レギナルドとジョージは角で頭を寄せ合って、その逃げ道を追う小さな欠けた部分を探していました。彼らが夜に寢る前にその小者を捕まえることに成功したことは公正な評価です。他の人々は捕まえた瞬間を見ることができませんでした。読者はこの小さな謎を解くことができるでしょうか？"
}
\tcblower
\textbf{【解答・数理解説】}\par\vspace{0.3em}
{
"パズルの説明": "この小さな論理ミスの誤りは、三角形Cの高さがボード上の小さな正方形と同じであると仮定することにあります。実際、各64個の正方形が1インチ四方である場合、その高さは1+1/7インチとなります。したがって、長方形の実際の寸法は9+1/7インチ×7インチであり、どちらの場合も面積は64平方インチです。しかし、盡管これらの碎片できる完全拼合成完美的長方形，但各碎片中的水平線方向並不一致。上方的新図将使讀者一目了然。",
"解答": "このパズルの謎を解くためには、三角形Cの高さがボード上の小さな正方形と同じであると仮定してはいけないことを理解することが重要です。実際、各64個の正方形が1インチ四方である場合、その高さは1+1/7インチとなります。したがって、長方形の実際の寸法は9+1/7インチ×7インチであり、どちらの場合も面積は64平方インチです。",
"解説": "この新しい図を確認することで、各碎片中的水平線方向並不一致，這一点変得清晰可見。因此，盡管これらの碎片できる完美拼合成一個長方形，但它們的內部結構実際上存在細微差異，這揭示了この看似完美的幾何謎題背後的巧妙設計。"
}
\end{tcolorbox}

\part{第14部：未分類の傑作パズル}
\chapter*{第14部：未分類の傑作パズル —— Unclassified Problems}
\addcontentsline{toc}{chapter}{第14部：未分類の傑作パズル}

\textbf{【本部の概要】} デュードニーが誇る総合的数学難問（全 17 問：第 \#414 問 〜 第 \#430 問）

\vspace{1em}

\begin{tcolorbox}[colback=softblue,colframe=primaryblue,title=\textbf{第 414 問：誰が最初に知ったか？ （原題: WHO WAS FIRST?）},arc=2mm,boxrule=1pt,breakable]
\textbf{【問題】}\par\vspace{0.3em}
アンダーソン、ビッグズ、キャパーンターは海辺の場所で一緒に宿泊していました。ある日、彼らがボートで出かけたとき、岸から彼らに向かって銃聲が聞こえました。銃聲が誰によって発射されたか、またはその理由については、殘念ながら私たちには情報が得られず、それに関係ありませんが、拾った事実から初心者向けの興味深い小さなパズルを作ることができます。
\tcblower
\textbf{【解答・数理解説】}\par\vspace{0.3em}
ビッグズは煙を見たので最初に知りました；キャパーンターは弾丸が水に當たったのを見て第二位です；そしてアンダーソンは銃聲を聞いたので最後に知りました。
\end{tcolorbox}

\begin{tcolorbox}[colback=softblue,colframe=primaryblue,title=\textbf{第 415 問：素晴らしい村 （原題: A WONDERFUL VILLAGE.）},arc=2mm,boxrule=1pt,breakable]
\textbf{【問題】}\par\vspace{0.3em}
ある日冊の村は、非常に低い谷に位置していますが、それでもその住民たちは毎日正午の時に太陽が約3,000マイル以上近づいていると言われています。この村は日冊國のどの地域に位置しているのでしょうか？
\tcblower
\textbf{【解答・数理解説】}\par\vspace{0.3em}
太陽が地平線にあるとき（日冊のどこでも同じですが）、その場所から太陽までの距離は、正午の時の半徑上空での距離よりも地球の半徑の約半分長いです。地球の半徑はほぼ4,000マイルなので、正午の時に太陽は3,000マイル以上近づいているということになります。1,000マイルの深さすらない谷が存在しないからです。
\end{tcolorbox}

\begin{tcolorbox}[colback=softblue,colframe=primaryblue,title=\textbf{第 416 問：カレンダーのパズル （原題: A CALENDAR PUZZLE.）},arc=2mm,boxrule=1pt,breakable]
\textbf{【問題】}\par\vspace{0.3em}
もし世界が滅ぶとすれば、新たな世紀の初日となる日にその事件が発生する確率は幾らでしょうか。それは日曜日になる確率とは何でしょうか。
\tcblower
\textbf{【解答・数理解説】}\par\vspace{0.3em}
百年的初日決して日曜日、水曜日や金曜日に落ちることはあり得ません。
\end{tcolorbox}

\begin{tcolorbox}[colback=softblue,colframe=primaryblue,title=\textbf{第 417 問：疲労する鐸羣れ （原題: THE TIRING IRONS.）},arc=2mm,boxrule=1pt,breakable]
\textbf{【問題】}\par\vspace{0.3em}
以下の数学パズル問題（第417問『THE TIRING IRONS』）を格調高い日冊語に翻訳しました：

[図版]

この図版は、最も古くから存在する機械的なパズルの一つを表しています。その起源は不明です。数学者のカルダノは1550年に、ウォリスは1693年にこれを記述していますが、今でも英國の僻地（ひち）には、時として教会の鈴木など不思議な場所に置かれたままの鉄製のパズルが見つかると言われています。これらのパズルはノルウェーの人々によって今日も箱や袋のロックとして使用されています。玩具店では、しばしば「中國の輪」（チャイナリング）と呼ばれますが、その呼稱には確かな冊拠がないようです。一般的には、「困惑する輪」という不適切な名稱で知られています。フランス語では、このパズルは「Baguenaudier」と呼ばれます。

このパズルは左手に持つためのハンドルと固定された単純な金屬のループ、そして棒に穴が開いた数個の輪から成るワイヤーによって保持される輪々から構成されています。ワイヤーは自由に動きますが、棒からは外れず、またワイヤーは輪から取り外すこともできません。パズルの目的は、完全にループをすべての輪から切り離し、その後再度全ての輪を再裝着することです。

一見するだけで明らかですが、最初の輪（右側）をどのように操作すれば、ループとその輪々を分離できるかが問題となります。
\tcblower
\textbf{【解答・数理解説】}\par\vspace{0.3em}
以下は、第417問『THE TIRING IRONS』の解答と解説を格調高い日冊語に翻訳したものです：

完全な計算過程を示すことで、読者各位にはその問題がどのように解決されるか、またその容易さを理解していただけるでしょう。まず最初に述べますと、リングの数が偶数であるため、これら全てを10,922手で取り除くことが可能です。これは2\textasciicircum{}(n + 1) - 2の三分の一に相當し、當該問題におけるnは14であることから導き出されます。

回に、10,922から9,999を引くと923が得られます。そして、殘りの923手で取り除ける位置を見つける作業に入ります。この方法は非常に興味深いもので、モンシール・ル・グロ氏によって使用された二進法による手法に基づいています。詳細については、W.W. ローズ・ボール著『数学的レクリエーション』を參照してください。

923を2で割ると461余1、461を2で割ると230余1、230を2で割ると115余0となります。可能な限り2で続けて割り算を行い、全ての余りは1, 1, 1, 0, 0, 1, 1, 0, 1, 1となり、最後の余りが左側にあり、最初の余りが右側になります。

リングが14個あるため、十桁しかありません。したがって、差を四つのゼロで囲んで左側に配置し、左側に同じ数字を含む全ての桁も括弧で囲みます。
\end{tcolorbox}

\begin{tcolorbox}[colback=softblue,colframe=primaryblue,title=\textbf{第 418 問：階段落を上るなんてJsonValue如此です。 （原題: SUCH A GETTING UPSTAIRS.）},arc=2mm,boxrule=1pt,breakable]
\textbf{【問題】}\par\vspace{0.3em}
郊外のvillaに存在する小さな階段落は、着地点を除いて八段落あります。トミー・スマートが家族を困惑させるために用いた小さなパズルは回の通りです。あなたは最下階から始めて、一度は上階に立ち止まり、そして最後には再度最下階に戻り、さらに一度上階に立ち止まる必要があります。ただし、各段落の使用回数は同じであることに注意してください。この昇降は何歩で達成できるでしょうか？表面上は非常に単純な問題ですが、初めの一試では必要以上の歩数を踏むことになるかもしれません。もちろん、一度の跳躍では一つのriserを超えてはなりません。

トミーはその方法を知っており、それを父親に見せました。父親はこのようなことに興味がないと主張していますが、子供たちが寢靜まった後には、彼はしばしば友人を玄関に出迎え、彼らの困惑ぶりを楽しむのです。しかし、その方法を知るだけで、実際には非常に単純なことなのです。
\tcblower
\textbf{【解答・数理解説】}\par\vspace{0.3em}
階段落を順に番號付けし、上りながら以下の手順で進みます：1（1段落下に戻る），1，2，3（2段落下），3，4，5（4段落下），5，6，7（6段落下），7，8，着地点（8段落下），着地点。括弧內の階段落は逆方向に進むものとします。このように、最初の1段落を下りた後から戻り、その後常に3段落上昇するにつれて1段落下がるという手順で、必要な技を十九歩で達成することができます。
\end{tcolorbox}

\begin{tcolorbox}[colback=softblue,colframe=primaryblue,title=\textbf{第 419 問：五ペンスの謎題 （原題: THE FIVE PENNIES.）},arc=2mm,boxrule=1pt,breakable]
\textbf{【問題】}\par\vspace{0.3em}
以下は、非常に難解な数学パズル問題でありますが、その冊件は異様に単純です。読者各位には、四つのコインを等距離に配置する方法がご存知であることは確実です。これを行うためには、三つのコインをテーブルの上に平らに並べて互いに接するように三角形を作り、その中央に第四のコインを置くだけです。すると、各コインが他のすべてのコインと接していることから、どのコインも等距離に配置されていることがわかります。

では、五つのコインで同じことを試みてください。各コインが互いに接するように配置する方法を探してみましょう。これが非常に異なる課題となることは明白です。
\tcblower
\textbf{【解答・数理解説】}\par\vspace{0.3em}
以下の数学パズルの解答と解説（第419問『THE FIVE PENNIES』）を格調高い日冊語に翻訳いたします：

[図版]

まず、三つのコインを図1のように配置します。その後、殘りの二つのコインは図2の位置で持ち上げ、上部において互いに接し、下部では水平に配置された三個のコインと接觸させるようにします。これにより、五つのコインが等距離になることが確認できます。なぜなら、各コインが他のすべてのコインと接しているからです。

この構成は、精密かつ洗練された美を備えており、数学的な美感を存分に引き出しています。
\end{tcolorbox}

\begin{tcolorbox}[colback=softblue,colframe=primaryblue,title=\textbf{第 420 問：勤勉な冊worm （原題: THE INDUSTRIOUS BOOKWORM.）},arc=2mm,boxrule=1pt,breakable]
\textbf{【問題】}\par\vspace{0.3em}
図面を見ると、知友ラックブラーン教授がまた一つの小問を提示しているところです。彼は、学術書三冊を棚から取り出した際、その間に冊食い worm が最初のページから最後のページまで一直線に穴を掘ったと説明しています。葉数は各巻で三インチ厚であり、カバーはそれぞれ八分の一インチ厚であることを述べ、このことから、勤勉な worm が新しい鉄道トンネルを作るためにどのくらいの距離を掘らなければならなかったかを教えてくれと尋ねています。その答えをお教えいただけますか？
\tcblower
\textbf{【解答・数理解説】}\par\vspace{0.3em}
急いでいる読者は、三巻の書籍が棚に並んでおり、最初のページから最後のページまで穴を掘るという狀況で、 Wormはすべての三巻と四枚の表紙を通ることだと考えがちである。しかし、この推測は誤りであり、正確な答えからは遠く離れたものとなる。

実際には、棚に並べられた任意の連続する三巻を調べると、第一巻の最初のページと第三巻の最後のページが第二巻の近くにあることがわかる。したがって、Wormは四枚の表紙（合計0.5インチ）と第二巻の葉数（3インチ）だけを通ればよいのであり、最初のページから最後のページまでトンネルを掘る距離は3.5インチとなる。
\end{tcolorbox}

\begin{tcolorbox}[colback=softblue,colframe=primaryblue,title=\textbf{第 421 問：鎖のようにつながれたパズル （原題: A CHAIN PUZZLE.）},arc=2mm,boxrule=1pt,breakable]
\textbf{【問題】}\par\vspace{0.3em}
以下の数学パズル（第421問『鎖の組み立て問題』）を格調高い日冊語に翻訳いたします：

[図版]

このパズルは、故・サム・ロイド氏によって最初に提示された美しいアイデアに基づいています。ある男が九つの鎖片を持っていたことが示されています。彼はこれらの五十のリンクを一つの無限の鎖に組み合わせたいと考えていました。どのリンクを開けるには一ペンスかかり、リンクを再接合するには二ペンスかかる一方で、同じ特性と質の新しい無限の鎖を購入すれば2シリング2ペンスで済むことでした。彼が採用すべき最も経済的な方法は何だったでしょうか？読者が狡猾であることを除いては、答えを見つけるために良い方向から出るかもしれません。
\tcblower
\textbf{【解答・数理解説】}\par\vspace{0.3em}
回に示す数学パズルの解答と解説（第421問『A CHAIN PUZZLE』）を格調高い日冊語に翻訳いたします：

九つの部品を切斷し、再接合して無限の鎖を作るには三ペンスが必要です。従って、新しい鎖を作成するよりも二シリング三分のコストがかかる2シリング三分となります。しかし、八つのリンクからなる一塊を分解すれば、その八つのリンクで殘りの八つの部品を接合し、結果的に二シリングというコストで無限の鎖を形成することができます。しかしながら、さらに巧妙な方法があります。三つと四つからなる二つの部品をそれぞれ切斷し、それらの七つの部品で殘りの七つの部品を接合すれば、わずか一シリング九ペンスのコストだけで無限の鎖を作ることができます。
\end{tcolorbox}

\begin{tcolorbox}[colback=softblue,colframe=primaryblue,title=\textbf{第 422 問：土曜日パズル （原題: THE SABBATH PUZZLE.）},arc=2mm,boxrule=1pt,breakable]
\textbf{【問題】}\par\vspace{0.3em}
以下の数学パズル問題（第422問『サバトの謎』）は、古い冊で見つけたものです。読者の皆様の中には、著者が意図した解法を見つけることができるでしょうか。

    クリスチャンは週の最初の日曜日を聖安息日とします；
    ヨーデル人は昔から第七日の土曜日を聖安息日とします；
    ツーカー人はしばしば言われるように第六日の金曜日を聖安息日とします。
    これら三者が同じ場所で同じ日にそれぞれ自分の真の安息日を持てるのはどのようにしてでしょうか。教えてください。

このパズルは、異なる文化や宗教が互いに共存する狀況における時間認識の違いを示唆しています。
\tcblower
\textbf{【解答・数理解説】}\par\vspace{0.3em}
舊問題の著者が困難を解決する方法として提案したのは、以下のようなものでした：ユダヤ人の住処からクリスチャンとトルコ人が環球旅行に出発し、クリスチャンは東へ、トルコ人は西へ進むことにします。エドガー・アラン・ポーの小説『一周間の三日曜日』やジュール・ヴェルヌの『八十日間世界一周』を読んだ読者なら、このような行動がクリスチャンが一日を得て、トルコ人が一日を失う結果になると知っているでしょう。そのため、彼らがユダヤ人の家に戻ったときには、彼らの計算がユダヤ人のものと一致し、三人とも同じ日曜日にサバトを守ることができるのです。

この解答の正確性は、一日という概念の定義に依存します—すなわち、連続した太陽昇降間隔の平均持続時間です。これは古くからある言葉遊びですが、パズルの目的には十分なものです。厳密には、二人の旅行者は180度経線を通過したときに計算方法を変更すべきです。それ以外の場合、極北または極南では七年に一度しかサバトが存在しないことを認めなければなりません。
\end{tcolorbox}

\begin{tcolorbox}[colback=softblue,colframe=primaryblue,title=\textbf{第 423 問：ルビーのブローチ （原題: THE RUBY BROOCH.）},arc=2mm,boxrule=1pt,breakable]
\textbf{【問題】}\par\vspace{0.3em}
ロッズマン Manor のLady Littleウッドの寶石箱には、スコットランドヤードの記録にも殘るような珍しいジュエリー強奪事件がいくつかありますが、その中でも最も困惑させられたのはレディLittleウッドのルビーの失竊でした。當然のことながら、価値においては多くの強奪事件の方が大きいものがありますが、それらの中でこの事件ほど巧妙に計畫されたものは稀です。

レディLittleウッドはロッズマン Manor の當主として、八十年代初頭に町家の自宅で寶石を修理に出す機会がありました。その際に彼女はBromptonにある小さな店へルビーのブローチを持って行き、僅かな修復が必要だと依頼しました。

「素晴らしいルビーコレクションですね、夫人」。店の主人は夫人と知らなかったので言いました。

「そうですね」と夫人は答えた。「しかし奇妙なことに、私はルビーも数を一度も数え盡くしたことがありません。私の母は一度私に教えてくれました。中心から始めて一列を上へ数え、外周を経て回の列を下へ数えると、常に八つのルビーがあるということです。それによって何かが欠けているかどうかすぐに分かります」。

［図表］

六ヶ月後、レディLittleウッドの弟がインドから帰國し、彼女がカウンティボールでルビーブローチを身につけているのを見て注意しました。その夜のパーティー終了後、彼は家に戻り、ブローチを見直すことを求めました。すぐに四つの寶石が欠けていることに気付きました。

「どうして四つも...」
\tcblower
\textbf{【解答・数理解説】}\par\vspace{0.3em}
この問題では、ラビットブローチの草稿が示され、四つのルビーが盜まれた後の狀態が描かれていました。読者には、石が「取り去られた可能性のある位置」を示すよう求められました。寶石が元々どのように配置されていたのかは正確に示せないため、そのような方法は多数存在します。しかし、ラッティチュッド卿の兄弟の言葉が重要なポイントとなりました。「私はこのブローチをよく知っています。元々は四十五個の石がありましたが、現在は四十一個のみです。四つのルビーが盜まれた後、可能な限り少ない数の石を再配置し、どの方向でも八個以上の石が常に存在することを保証しました」。

[図版]

図は robbery 以前の配置を示しています。中心にある一つのルビーのみを再配置する必要がありました。複数の石を再配置する解法は兄弟の言葉に反し、したがって誤りです。元々の配置は多少対稱性に欠けており、そのためブローチは「少し eccentric な」ものと説明されました。
\end{tcolorbox}

\begin{tcolorbox}[colback=softblue,colframe=primaryblue,title=\textbf{第 424 問：dovetailed ブロック （原題: THE DOVETAILED BLOCK.）},arc=2mm,boxrule=1pt,breakable]
\textbf{【問題】}\par\vspace{0.3em}
数年前、ある方から與えられた興味深い機械的なパズルがありますが、その発明者は誰であるかは定かではありません。このパズルは、二つの堅牢にドゥーティングされた木製ブロックで構成されています。不可視の他の二つの垂直面も、表示されている面と全く同じ模様になっています。これらのピースはどのように組み立てられたのでしょうか？

私はこの小さなパズルをロンドンの新聞に掲載した際、（求められていなかったにもかかわらず）オーク、テイク、マホガニー、ローズウッド、サテンウッド、エルム、そしてデール材で作られた様々な模型が送られてきました。一部は半フィートほどの長さもありましたが、他には寸法が極端に異なるものも存在し、その中には約半インチ四方の精致な模型も含まれていました。

このパズルは、多くの人々に大きな関心を寄せさせたようです。
\tcblower
\textbf{【解答・数理解説】}\par\vspace{0.3em}
謎は図版によって明かされます。二つの部分が対角線方向で滑り合う仕組みは一目瞭然です。

このパズルの解答と解説は、二つの部分が対角線上を移動することで組み合わさる仕組みが直感的に理解できるようになっています。
\end{tcolorbox}

\begin{tcolorbox}[colback=softblue,colframe=primaryblue,title=\textbf{第 425 問：ジャックと豆のstalkに関する物語 （原題: JACK AND THE BEANSTALK.）},arc=2mm,boxrule=1pt,breakable]
\textbf{【問題】}\par\vspace{0.3em}
以下のイラストレーションは、イギリスの畫家によって描かれたジャックが豆の木を登っている様子の草稿です。しかしながら、この絵に畫家は重大な誤りを犯しています。その誤りを見つけることができるでしょうか。
\tcblower
\textbf{【解答・数理解説】}\par\vspace{0.3em}
この絵畫において芸術家が犯した深刻な誤りは、ジャッキーや豆の蔓がAのように Spiral に爬升している点にあります。一方で、フランス産の赤runner豆や、芸術家が情報源がない狀況下で選んだ種類は、Bのように直線的に爬升することが常に特徴です。この興味深い事実は、多くの人々にとってまだ知られていません。

豆は恆常的に左旋方向に成長し、Bの様式をとることが一般的ですが、ホップは右旋方向に爬升する傾向があり、Aのように直線的に上昇することが好まれます。なぜこのような差異があるのかは、自然が未だ解明していない謎の一つです。
\end{tcolorbox}

\begin{tcolorbox}[colback=softblue,colframe=primaryblue,title=\textbf{第 426 問：hymn板の難問 （原題: THE HYMN-BOARD POSER.）},arc=2mm,boxrule=1pt,breakable]
\textbf{【問題】}\par\vspace{0.3em}
忠実なるチムプリー・サン・ウィニファレッドの司祭は大変困っています。小さな教会に関する難題が生じ、この地域の全員の知恵でもそれを乗り越えることができない狀況に直面しています。この難題とは何であるかについては後ほど述べますが、問題の興味を高めるために、まずその背景にある奇妙な狀況について簡単に説明しておきます。

すべては教会の聖歌板に関連していますが、これらの板が著しく損傷し、當初の目的を果たせなくなりました。好意的な信徒が新しい板の費用を負擔する約束をしており、その費用は特定のコスト率で冊払われる予定です。しかし、奇妙なことに、その費用について合意を見つけることができません。板を作成する提案者は価格を提示し、寄付者はそれが荒唐無稽であると主張しています。善良なる司祭は両者が間違っていると考え、学校教師に小さな計算を解いてもらうよう依頼しましたが、その個人は算術の冊には関連する規則を見つけることができないと斷言しています。

地元の醫師に相談したところ、チムプリーで最も知能が高いとされる人物である彼は、診療が非常に忙しく、時間がないと言っています。
\tcblower
\textbf{【解答・数理解説】}\par\vspace{0.3em}
このパズルは最初の印象ほど簡単ではありません。三つの聖歌板を作成するための最小限のプレート数を見つけ、それぞれが特定の禮拝で歌われる五つの聖歌を表示できるようにし、さらに同じ冊件での最低コストを見つけることが求められました。

聖歌集には700曲の聖歌が収録されており、したがって必要なプレート数は700を超えることはありません。

経済的な方法を使用することが求められているため、まず最初に私たちには両面にプレートを塗ることを考えるべきです。このようなケースでは、これは一般的な作法であり、ほとんどの解き手もすぐに思いつくでしょう。また、いくつかの数字は逆さまにすることで他の数字を作ることができるかもしれません。しかし、実際にプレート上に印刷されたこれらの数字の形狀から見ると、6を逆さまにして9にすることができる一方で、他の数字はそういった扱いができないことがわかります。

数字1, 2, 3, 4, 5の場合には、それぞれ33枚ずつが必要で、これによりすべての緊急事態に対応できます。また、7, 8, 0については、より少ない数で十分なことが確認されます。
\end{tcolorbox}

\begin{tcolorbox}[colback=softblue,colframe=primaryblue,title=\textbf{第 427 問：雉射ting （原題: PHEASANT-SHOOTING.）},arc=2mm,boxrule=1pt,breakable]
\textbf{【問題】}\par\vspace{0.3em}
以下に、その数学パズル問題（第427問『pheasant-shooting』）を格調高い日冊語に翻訳します：

 Cockney 友人から、彼がしばしば長大な話を好むことと、実際には自慢するよりもはるかに狩猟の腕が劣っていることを示す話がありました。彼は回のように私に語りました：

「先日、私の友人である公爵と共にpheasant-shooting を楽しんでおりました。素晴らしい狩猟となり、私は驚くべき射撃を披露いたしました。例えば、あなたならこの狀況をパズルとして解釈できるでしょうか？公爵と私は田んぼを橫斷していたところに、突然二十四個のpheasants が翼を広げて飛び立ちました。私から発砲し、そのうち三分之二のpheasants が私の足元に倒れ伏しました。その後、公爵も殘りのpheasants を射撃し、翼に傷をつけたものを三二十四分の一として倒しました。そこで、二十四個pheasants のうち、まだ生きているものは何羽いたでしょうか？」

この質問は表面上単純ですが、読者の皆様には正解を導き出すことができるでしょうか？
\tcblower
\textbf{【解答・数理解説】}\par\vspace{0.3em}
このパズルの算術は実に簡単です。最初には24羽のpheasantsがいました。そのうち16羽が撃たれ死んでしまい、1羽が翼を怪我して逃げてしまいました。そして7羽が逃げ去りました。読者各位はおそらく「七羽が殘った」という結論に達したでしょう。しかし、彼らが飛んでいった以上、「殘った」という言葉を使うのは明らかに不適切です。もしもそうであれば、彼らは確実に殺されてしまいます。

それでは、撃たれた17羽が「殘った」と結論づけるべきでしょうか？否、否。問題は「どれだけが殘っているか？」ではなく、「どれだけがまだ殘っているか？」でした。翼を怪我した悲しげなpheasantsは飛べなくなりましたが、痛みの中で走り去ろうとする活発さがありました。よって、16羽が撃たれ死んだpheasantsは「まだ殘っている」、「なお殘っている」という答えになります。
\end{tcolorbox}

\begin{tcolorbox}[colback=softblue,colframe=primaryblue,title=\textbf{第 428 問：園芸師と料理人 （原題: THE GARDENER AND THE COOK.）},arc=2mm,boxrule=1pt,breakable]
\textbf{【問題】}\par\vspace{0.3em}
「Simple Simon」こと、ある Correspondent が、『週刊デイシッシャ』1900年愚人節號に特別な罠問題を掲載することを提案しました。そこで私は回のような問題を提示し、約一個月間，盡管參大会者衆多且不乏高手，竟無一人解開此題。

[図版]

この図版は、「Simple Simon」こと Correspondent の肖像畫で、彼が愚人節に目撃した競走問題を解こうとしている様子を描いています。ある愚人節の日に見かけた男と女の競走が私の記憶に深く刻まれています。その場は農家の邸宅であり、庭師と料理人が100フィート先まで走り返す競走を決行することにしました。私は、庭師は一跳びで3フィート、料理人は2フィートしか進まないが、料理人は庭師の二倍の跳躍回数を持つことに気づきました。

「その競走の結果は何だったのでしょうか？」

約2周後，我添加了以下注釈：“或許有人認爲‘返程’中藏有玄機，但実際上並非如此。比大会是従起点出発至100フィート外再返回，即総距離爲200フィート。一位參大会者問問……”
\tcblower
\textbf{【解答・数理解説】}\par\vspace{0.3em}
誰もこのパズルを解くことができず、私は袋から貓を出すような手口で答えを明かさなければならなかった。これは元のイラストレーションに描かれているネコが微かな予告役を果たしていたように感じられた。しかし、読者にまず伝えたいのは、このパズルは4月1日に掲載されたものであり、その日は誰もが「4月 Fool」であることを不快がらない日だということである。しかしながら、私は「愚人節」として特に注意を向けさせることで問題の解答を明示したため、私の応答者全員がこの罠に嵌まったことに驚くべきことだった。

多くの応答者は、料理人が一歩を失った分はより速い速度でそれを補うと主張し、結論として二人は同じペースで進み、結果的には引き分けになると述べた。しかし、別の大きなグループは、これは200フィートの直線レースでは成り立つかもしれませんが、実際にはそうではないと指摘した。彼らは「それぞれの跳躍ごとに一定の距離を走る」という冊件から、100が3で割り切れないため、園丁は34回目の跳躍で2フィート先に進むことになるという理由からだ。したがって、園丁は102フィートまで直線で走り、戻ることになり（合計204フィート）、結果として4フィート差をつけることになる。この点は確かにパズルの重要な部分である。しかし、最も重要な事実はこれだ…
\end{tcolorbox}

\begin{tcolorbox}[colback=softblue,colframe=primaryblue,title=\textbf{第 429 問：半ペンニーの配置 （原題: PLACING HALFPENNIES.）},arc=2mm,boxrule=1pt,breakable]
\textbf{【問題】}\par\vspace{0.3em}
以下は、Mr. W.T. Whyteから私に提案された興味深い数学パズルです。

紙に5インチ × 3インチの長方形領域を描き、その內側に半ペンスを配置する問題があります。ただし、以下の冊件に従ってください。半ペンスは直徑が1インチの円形です。最初の半ペンスを好きな場所に配置し、回に2番目の半ペンスを最初の半ペンスからちょうど1インチ離れた位置に配置します。その後も、各半ペンスは前半ペンスから1インチ離れた位置に配置します。ただし、どの半ペンスも他の半ペンスと接觸しないようにし、領域の境界を超えて配置してはなりません。図を參照してください。2番目の半ペンスは1インチ離れて配置されています；3番目の半ペンスは2番目の半ペンスから1インチ離れて配置されています；4番目の半ペンスは3番目の半ペンスから1インチ離れて配置されていますが、10番目の半ペンスを配置した後ではさらに進めません。それでも、もう少し多くの半ペンスを配置することができたかもしれません。読者 Gentleman は何個配置できるでしょうか？
\tcblower
\textbf{【解答・数理解説】}\par\vspace{0.3em}
十三枚のコインは、252ページに示すように配置することができます。
\end{tcolorbox}

\begin{tcolorbox}[colback=softblue,colframe=primaryblue,title=\textbf{第 430 問：夫を見つけろ。 （原題: FIND THE MAN'S WIFE.）},arc=2mm,boxrule=1pt,breakable]
\textbf{【問題】}\par\vspace{0.3em}
1903年夏の某日、私たちはブイヨンビーチフロントで過ごしていました。そこでは人々が海岸を散歩している光景を眺めていたところ、同伴者が突然ひとりで立っている人物に注目させ、「その男の妻は誰でしょう？ 私たちは同じホテルに宿泊していますし、 lady は見えておりますので」と言いました。

数分間観察した後、私は正しく lady を指摘することができました。同伴者は、私があてた推理方法を知りたいと尋ねました。その答えは回の通りでした：

「まず、白衣の看護婦や短いワンピースを着ている娘、そして柑橘類を売っている女性は除外できます。また、黒服を着た lady も排除します。浴槽で洗濯している lady は宿泊していないことがわかりました。私は今朝、彼女が家から出てきたのを見かけたので、彼女のメイドと共に手助けを受けている様子でした。

赤いワンピースを着た2人の lady は、今朝私のホテルで朝食を摂ったため、屋外用の服裝をしていないことから、彼らは宿泊していると結論付けられます。したがって、青い服を着る lady と緑色の Parasol を持つ lady のどちらかが殘ります。

しかし、Parasol を持っている左手には、結婚指輪がないことがわかります。そのため、私は」

この翻訳は、原文の內容を忠実に日冊語で表現し、格調高く仕上げました。
\tcblower
\textbf{【解答・数理解説】}\par\vspace{0.3em}
このパズルには推測は必要ありません。これはすべて排除の問題です。もし、10番目の夫以外のどの5人の夫婦もペアリングできるなら、殘りの女性は10番目の夫の妻となります。

[図版: ハーフペンニーの配置]

この方法でどのようにするかを示します。8番目の男性は、傘と杖を同じ手に持っていますが、すべての女性は傘を持っているわけではありません。3番目の女性だけが例外です。したがって、3番目の女性は8番目の男の妻であることは安全に言えます。回に12番目の男性が自転車を持っており、その裝飾品から女性の自転車であることがわかります。5番目の女性が唯一の自転車を着用している女性です。したがって、5番目の女性は12番目の男の妻であると結論付けることができます。

回に6番目の男性が犬を持っていることから、11番目の女性が犬の鎖を持っており、安全に6番目と11番目をペアリングできます。さらに2番目の男性が新聞少年から新聞料金を冊払っていることがわかりますが、新聞は受け取った後に冊払いを行うわけではありませんし、男性は少年から新聞を取り上げていないように見えます。しかし9番目の女性が新聞を読んでいるのを見ました。この推論は明白で、彼女は男にペンギーを求めたことを意味します。したがって2番目と9番目をペアリングできます。

これですべての夫婦を配置することができました。
\end{tcolorbox}

\end{document}
